\documentclass[11pt,oneside,openany]{book}
\usepackage[letterpaper,margin=1in]{geometry}
\usepackage{amsmath,amsthm,mathtools}
\usepackage{fontspec}
\usepackage{unicode-math}
\setmainfont{lmroman10-regular.otf}[
  BoldFont=lmroman10-bold.otf,
  ItalicFont=lmroman10-italic.otf,
  BoldItalicFont=lmroman10-bolditalic.otf
]
\setsansfont{lmsans10-regular.otf}[
  BoldFont=lmsans10-bold.otf,
  ItalicFont=lmsans10-oblique.otf,
  BoldItalicFont=lmsans10-boldoblique.otf
]
\setmonofont{lmmono10-regular.otf}[
  ItalicFont=lmmono10-italic.otf
]
\setmathfont{latinmodern-math.otf}
\usepackage{microtype}
\usepackage{enumitem}
\usepackage{booktabs,longtable,array}
\usepackage{titlesec}
\usepackage{fancyhdr}
\usepackage{hyperref}
\usepackage{bookmark}
\usepackage{xcolor}
\hypersetup{
  unicode=true,
  colorlinks=true,
  linkcolor=black,
  urlcolor=black,
  pdftitle={Projection Craft S Essays: Books 0--1},
  pdfsubject={Ambiguity-resolved LaTeX edition of the stable S0--S1 source corpus}
}
\setlist[itemize]{leftmargin=2em,itemsep=0.25em,topsep=0.4em}
\setlist[enumerate]{leftmargin=2.4em,itemsep=0.25em,topsep=0.4em}
\setlength{\parindent}{0pt}
\setlength{\parskip}{0.55em}
\allowdisplaybreaks
\raggedbottom
\emergencystretch=3em
\pagestyle{fancy}
\fancyhf{}
\fancyhead[L]{\small Projection Craft}
\fancyhead[R]{\small \leftmark}
\fancyfoot[C]{\thepage}
\renewcommand{\headrulewidth}{0.4pt}
\titleformat{\chapter}[display]
  {\normalfont\bfseries\sffamily}
  {\filleft\Large\chaptername\ \thechapter}
  {1ex}
  {\Huge}
\titlespacing*{\chapter}{0pt}{-25pt}{30pt}

\newcommand{\PCchapter}[3]{%
  \cleardoublepage
  \chapter*{\texorpdfstring{#1\\[0.4em]#2}{#1: #2}}%
  \markboth{#1}{#1}%
  \addcontentsline{toc}{chapter}{#1: #2}%
  \begin{center}\small\sffamily #3\end{center}
  \vspace{1em}
}
\newcommand{\PCsection}[1]{%
  \section*{#1}\addcontentsline{toc}{section}{#1}%
}
\newcommand{\PCsubsection}[1]{%
  \subsection*{#1}%
}
\newcommand{\PCsubsubsection}[1]{%
  \subsubsection*{#1}%
}

\begin{document}
\begin{titlepage}
  \centering
  \vspace*{0.14\textheight}
  {\Huge\bfseries\sffamily Projection Craft\par}
  \vspace{1.2em}
  {\LARGE\sffamily S Essays\par}
  \vspace{0.7em}
  {\Large Books 0--1\par}
  \vfill
  {\large Ambiguity-Resolved Source Edition\par}
  \vspace{0.6em}
  {\small Reconciled against the canonical definitions, errata, and scope rules established in the project thread.\par}
  \vspace*{0.08\textheight}
\end{titlepage}
\frontmatter
\chapter*{Editorial Resolution Notice}
\addcontentsline{toc}{chapter}{Editorial Resolution Notice}
This edition preserves the numbering, theorem order, and canonical operator identities of the stable S0--S1 corpus. It resolves only ambiguities already identified in the project thread. No physical interpretation is promoted to theorem status.

\begin{enumerate}
\item \textbf{Carrier convention.} Analytic calculations use the lifted coordinate \(x\in\mathbb R\). Periodic phase statements use the quotient circle \(\mathbb T=\mathbb R/(2\pi\mathbb Z)\). The notation \(x\equiv x+2\pi\) means equality of quotient classes, not literal equality in \(\mathbb R\).
\item \textbf{Dihedral convention.} \(D_4\) denotes the order-eight symmetry group of the four-quadrant phase circle. Some references denote the same group by \(D_8\).
\item \textbf{Boundary terminology.} Closing a single open quadrant to \([b_k,b_{k+1}]\) is a legitimate componentwise compactification. The countable disjoint union \(\widehat{\mathcal D}=\bigsqcup_k\overline Q_k\) is instead called the \emph{split-boundary extension domain}; it is not asserted to be compact as a whole.
\item \textbf{Validity-domain logic.} In S1.II, the declared domain \(\Omega_{\mathcal P}\) is separated from the maximal operational admissibility envelope. Lawfulness requires inclusion in that envelope; equality means the declared domain is maximal. This removes the former circularity caused by using the instance-domain guard to define the domain against which that same guard was tested.
\item \textbf{Naming.} \emph{Projection Craft} is the project title, \emph{S Essays} is the corpus title, and \emph{R-operator} remains the established name of the mathematical operator family. These names do not designate competing document series.
\end{enumerate}

A complete change ledger is included in \texttt{AMBIGUITY\_RESOLUTION.md} in the source package.

\tableofcontents
\mainmatter
\part*{Book 0: Formal Core}
\addcontentsline{toc}{part}{Book 0: Formal Core}
% ===== BEGIN S0_I =====
% Source: S0.i.docx
\PCchapter{S0.I}{The Pythagorean Foundation of the Bounded Phase Operators}{Book 0 - Formal Core}

\PCsubsection{IDENTIFICATION DIVISION}

\PCsubsubsection{PROGRAM-ID}

R-THEORY-PYTHAGOREAN-CORE

\PCsubsubsection{PURPOSE}

To derive the exact algebraic properties of the R-operators from the Pythagorean identity

\[
\sin ^2x+\cos ^2x=1,
\]

together with standard definitions of the reciprocal and quotient trigonometric functions.

\PCsubsubsection{NON-PURPOSE}

This division does not establish:

\begin{itemize}
\item physical conservation of energy, mass, or information;
\item uniqueness of the R-operator family;
\item empirical applicability;
\item replacement of exponential or logarithmic laws;
\item a law of temporal evolution.
\end{itemize}
Those require additional arguments beyond the Pythagorean foundation.


\PCsection{ENVIRONMENT DIVISION}

\PCsubsection{Definition E.1 --- Phase variable}

Analytic calculations use a lifted phase coordinate

\[
x\in\mathbb R
\]

measured in radians. The operator family is \(2\pi\)-periodic. When one revolution is identified, the phase circle is

\[
\mathbb T=\mathbb R/(2\pi\mathbb Z),
\]

and \([x]\in\mathbb T\) denotes the quotient class of \(x\). Thus

\[
x\equiv x+2\pi
\]

means equality in \(\mathbb T\), not literal equality of real numbers. All analytic domains, derivatives, and one-sided limits in Book 0 are taken on the lifted real line unless the quotient circle is named explicitly.

\PCsubsection{Definition E.2 --- Admissible domain}

Define

\[
\mathcal D = \left\{ x\in\mathbb R: \sin x\neq0 \text{ and } \cos x\neq0 \right\}.
\]

Equivalently,

\[
\mathcal D = \mathbb R \setminus \left\{ \frac{k\pi}{2}:k\in\mathbb Z \right\}.
\]

All expressions below are evaluated only for \(x\in\mathcal D\).

\PCsubsection{Definition E.3 --- Unit-circle coordinates}

Let

\[
s=\sin x, \qquad c=\cos x.
\]

The point

\[
P(x)=(c,s)
\]

lies on the unit circle.


\PCsection{AXIOM DIVISION}

\PCsubsection{Inherited Theorem A.1 --- Pythagorean theorem}

For a right triangle with legs \(a,b\) and hypotenuse \(h\),

\[
a^2+b^2=h^2.
\]

\PCsubsection{Corollary A.2 --- Unit-circle Pythagorean identity}

Since the radius of the unit circle is \(1\),

\[
\boxed{s^2+c^2=1}.
\]

Equivalently,

\[
\boxed{\sin ^2x+\cos ^2x=1}.
\]

This identity is the foundational algebraic relation used throughout the remainder of the proof.


\PCsection{DATA DIVISION}

\PCsubsection{Definition D.1 --- Absolute reciprocal functions}

Define

\[
M=|\sec x| =\frac1{|c|},
\]

and

\[
N=|\csc x| =\frac1{|s|}.
\]

Both satisfy

\[
M>0, \qquad N>0.
\]

\PCsubsection{Definition D.2 --- Primitive displayed operators}

Define Sininverex:

\[
\operatorname{srx}(x) = N+\cot x.
\]

Define Sinexpon:

\[
\operatorname{sxp}(x) = \frac1{\operatorname{srx}(x)}.
\]

Define Cosinexpon:

\[
\operatorname{cxp}(x) = M+\tan x.
\]

Define Cosininverex:

\[
\operatorname{crx}(x) = \frac1{\operatorname{cxp}(x)}.
\]

\PCsubsection{Definition D.3 --- UNA operators}

Define Unainverex:

\[
\operatorname{urx}(x) = \operatorname{srx}(x) - \operatorname{crx}(x).
\]

Define Unaexpon:

\[
\operatorname{uxp}(x) = \operatorname{cxp}(x) - \operatorname{sxp}(x).
\]

This sign convention is fixed throughout the calculus.

\PCsubsection{Definition D.4 --- Saw operators}

Define

\[
\operatorname{saw}_{r}(x) = \frac1{\operatorname{urx}(x)},
\]

and

\[
\operatorname{saw}_{x}(x) = \frac1{\operatorname{uxp}(x)}.
\]

\PCsubsection{Definition D.5 --- Flatwave}

Define

\[
\operatorname{Flatwave}(x) = \operatorname{saw}_{r}(x) + \operatorname{saw}_{x}(x).
\]

Thus,

\[
\operatorname{Flatwave}(x) = \frac1{\operatorname{urx}(x)} + \frac1{\operatorname{uxp}(x)}.
\]


\PCsection{PROCEDURE DIVISION}

\PCsection{SECTION 1 --- Reciprocal Trigonometric Identities}

\PCsubsection{Lemma 1.1 --- Secant–tangent Pythagorean identity}

For every \(x\in\mathcal D\),

\[
\boxed{M^2-\tan ^2x=1}.
\]

\begin{proof}
Since

\[
M=|\sec x|,
\]

we have

\[
M^2=\sec ^2x.
\]

Using the definitions of secant and tangent,

\[
\sec ^2x-\tan ^2x = \frac1{c^2}-\frac{s^2}{c^2}.
\]

Therefore,

\[
\sec ^2x-\tan ^2x = \frac{1-s^2}{c^2}.
\]

By the Pythagorean identity,

\[
1-s^2=c^2.
\]

Hence,

\[
\sec ^2x-\tan ^2x = \frac{c^2}{c^2} = 1.
\]

Therefore,

\[
\boxed{M^2-\tan ^2x=1}.
\]

\end{proof}

\PCsubsection{Lemma 1.2 --- Cosecant–cotangent Pythagorean identity}

For every \(x\in\mathcal D\),

\[
\boxed{N^2-\cot ^2x=1}.
\]

\begin{proof}
Since

\[
N=|\csc x|,
\]

we have

\[
N^2=\csc ^2x.
\]

Thus,

\[
\csc ^2x-\cot ^2x = \frac1{s^2}-\frac{c^2}{s^2} = \frac{1-c^2}{s^2}.
\]

By the Pythagorean identity,

\[
1-c^2=s^2.
\]

Therefore,

\[
\csc ^2x-\cot ^2x = \frac{s^2}{s^2} = 1.
\]

Hence,

\[
\boxed{N^2-\cot ^2x=1}.
\]

\end{proof}

\PCsubsection{Corollary 1.3 --- Strict positivity of the conjugate factors}

Because

\[
M^2-\tan ^2x=1>0
\]

and \(M>0\),

\[
M>|\tan x|.
\]

Therefore,

\[
M+\tan x>0, \qquad M-\tan x>0.
\]

Likewise,

\[
N>|\cot x|,
\]

so

\[
N+\cot x>0, \qquad N-\cot x>0.
\]

No denominator introduced below vanishes on \(\mathcal D\).


\PCsection{SECTION 2 --- Exact Reciprocal Conjugates}

\PCsubsection{Theorem 2.1 --- Cosine-family reciprocal conjugate}

For every \(x\in\mathcal D\),

\[
\boxed{ \frac1{M+\tan x} = M-\tan x }.
\]

\begin{proof}
From Lemma 1.1,

\[
M^2-\tan ^2x=1.
\]

Factoring the difference of squares gives

\[
(M+\tan x)(M-\tan x)=1.
\]

Since \(M+\tan x\neq0\), division gives

\[
\frac1{M+\tan x}=M-\tan x.
\]

\end{proof}

\PCsubsection{Theorem 2.2 --- Sine-family reciprocal conjugate}

For every \(x\in\mathcal D\),

\[
\boxed{ \frac1{N+\cot x} = N-\cot x }.
\]

\begin{proof}
From Lemma 1.2,

\[
N^2-\cot ^2x=1.
\]

Therefore,

\[
(N+\cot x)(N-\cot x)=1.
\]

Since \(N+\cot x\neq0\),

\[
\frac1{N+\cot x} = N-\cot x.
\]

\end{proof}

\PCsubsection{Corollary 2.3 --- Explicit forms of the reciprocal operators}

The reciprocal R-operators may be written as

\[
\boxed{ \operatorname{crx}(x) = M-\tan x },
\]

and

\[
\boxed{ \operatorname{sxp}(x) = N-\cot x }.
\]

\PCsubsection{Corollary 2.4 --- Exact reciprocal pairing}

\[
\boxed{ \operatorname{cxp}(x)\operatorname{crx}(x)=1 },
\]

and

\[
\boxed{ \operatorname{srx}(x)\operatorname{sxp}(x)=1 }.
\]

\PCsubsection{Corollary 2.5 --- Positivity of all four displayed operators}

For every \(x\in\mathcal D\),

\[
\boxed{ \operatorname{srx}, \operatorname{sxp}, \operatorname{cxp}, \operatorname{crx}>0 }.
\]

The primitive operators carry magnitude but no negative sign. Signed direction first appears in the UNA differences.


\PCsection{SECTION 3 --- Sign Variables and Branch Forms}

\PCsubsection{Definition 3.1 --- Quadrant signs}

Define

\[
\alpha=\operatorname{sgn}(s), \qquad \beta=\operatorname{sgn}(c).
\]

Since \(x\in\mathcal D\),

\[
\alpha,\beta\in\{-1,+1\}.
\]

Because

\[
|s|=\alpha s, \qquad |c|=\beta c,
\]

we obtain

\[
N=\frac{\alpha}{s}, \qquad M=\frac{\beta}{c}.
\]

\PCsubsection{Lemma 3.2 --- Signed rational forms}

The primitive operators satisfy

\[
\boxed{ \operatorname{srx}(x) = \frac{\alpha+c}{s} },
\]

\[
\boxed{ \operatorname{sxp}(x) = \frac{\alpha-c}{s} },
\]

\[
\boxed{ \operatorname{cxp}(x) = \frac{\beta+s}{c} },
\]

and

\[
\boxed{ \operatorname{crx}(x) = \frac{\beta-s}{c} }.
\]

\begin{proof}
For example,

\[
\operatorname{srx} = N+\cot x = \frac{\alpha}{s}+\frac{c}{s} = \frac{\alpha+c}{s}.
\]

The other formulas follow identically.


\end{proof}

\PCsection{SECTION 4 --- UNA Structure}

\PCsubsection{Definition 4.1 --- Auxiliary branch quantity}

Define

\[
p=\alpha c-\beta s.
\]

\PCsubsection{Lemma 4.2 --- Exact UNA branch forms}

For every \(x\in\mathcal D\),

\[
\boxed{ \operatorname{urx}(x) = \frac{1+p}{sc} },
\]

and

\[
\boxed{ \operatorname{uxp}(x) = \frac{1-p}{sc} }.
\]

\begin{proof}
Using the signed rational forms,

\[
\operatorname{urx} = \operatorname{srx}-\operatorname{crx}.
\]

Therefore,

\[
\operatorname{urx} = \frac{\alpha+c}{s} - \frac{\beta-s}{c}.
\]

Using the common denominator \(sc\),

\[
\operatorname{urx} = \frac{c(\alpha+c)-s(\beta-s)}{sc}.
\]

Expanding,

\[
\operatorname{urx} = \frac{\alpha c+c^2-\beta s+s^2}{sc}.
\]

Since

\[
s^2+c^2=1,
\]

we obtain

\[
\operatorname{urx} = \frac{1+\alpha c-\beta s}{sc}.
\]

Thus,

\[
\operatorname{urx} = \frac{1+p}{sc}.
\]

Similarly,

\[
\operatorname{uxp} = \operatorname{cxp}-\operatorname{sxp},
\]

so

\[
\operatorname{uxp} = \frac{\beta+s}{c} - \frac{\alpha-c}{s}.
\]

Combining denominators gives

\[
\operatorname{uxp} = \frac{\beta s+s^2-\alpha c+c^2}{sc}.
\]

Hence,

\[
\operatorname{uxp} = \frac{1+\beta s-\alpha c}{sc} = \frac{1-p}{sc}.
\]

\end{proof}

\PCsubsection{Lemma 4.3 --- Bound on the auxiliary quantity}

For every \(x\in\mathcal D\),

\[
\boxed{|p|<1}.
\]

\begin{proof}
Squaring \(p\),

\[
p^2 = (\alpha c-\beta s)^2.
\]

Expanding,

\[
p^2 = \alpha^2c^2+\beta^2s^2-2\alpha\beta sc.
\]

Since

\[
\alpha^2=\beta^2=1,
\]

and

\[
s^2+c^2=1,
\]

we have

\[
p^2 = 1-2\alpha\beta sc.
\]

But

\[
\alpha\beta sc=|sc|.
\]

Therefore,

\[
p^2=1-2|sc|.
\]

Since \(x\in\mathcal D\),

\[
|sc|>0.
\]

Thus,

\[
p^2<1,
\]

and therefore

\[
|p|<1.
\]

\end{proof}

\PCsubsection{Corollary 4.4 --- Positivity of the UNA numerators}

Because \(|p|<1\),

\[
1+p>0, \qquad 1-p>0.
\]

\PCsubsection{Theorem 4.5 --- Common UNA sign}

For every \(x\in\mathcal D\),

\[
\boxed{ \operatorname{sgn}(\operatorname{urx}) = \operatorname{sgn}(\operatorname{uxp}) = \operatorname{sgn}(sc) }.
\]

Since

\[
\sin 2x=2sc,
\]

we also have

\[
\boxed{ \operatorname{sgn}(\operatorname{urx}) = \operatorname{sgn}(\operatorname{uxp}) = \operatorname{sgn}(\sin 2x) }.
\]

\begin{proof}
Both UNA numerators \(1+p\) and \(1-p\) are positive. Their common denominator is \(sc\). Therefore, both UNA operators have the sign of \(sc\).


\end{proof}

\PCsection{SECTION 5 --- Double-Angle Closure}

\PCsubsection{Theorem 5.1 --- UNA sum identity}

For every \(x\in\mathcal D\),

\[
\boxed{ \operatorname{urx}(x) + \operatorname{uxp}(x) = \frac4{\sin 2x} }.
\]

\begin{proof}
From Lemma 4.2,

\[
\operatorname{urx} + \operatorname{uxp} = \frac{1+p}{sc} + \frac{1-p}{sc}.
\]

Therefore,

\[
\operatorname{urx} + \operatorname{uxp} = \frac2{sc}.
\]

Since

\[
\sin 2x=2sc,
\]

we obtain

\[
\frac2{sc} = \frac4{\sin 2x}.
\]

Thus,

\[
\boxed{ \operatorname{urx} + \operatorname{uxp} = \frac4{\sin 2x} }.
\]

\end{proof}

\PCsubsection{Corollary 5.2 --- Inverse UNA sum identity}

Because the UNA sum is nonzero on \(\mathcal D\),

\[
\boxed{ \frac1{\operatorname{urx}(x)+\operatorname{uxp}(x)} = \frac{\sin 2x}{4} }.
\]

This is the exact double-angle closure identity.


\PCsection{SECTION 6 --- UNA Product}

\PCsubsection{Theorem 6.1 --- UNA product identity}

For every \(x\in\mathcal D\),

\[
\boxed{ \operatorname{urx}(x) \operatorname{uxp}(x) = \frac2{|sc|} }.
\]

Equivalently,

\[
\boxed{ \operatorname{urx}(x) \operatorname{uxp}(x) = \frac4{|\sin 2x|} }.
\]

\begin{proof}
Using Lemma 4.2,

\[
\operatorname{urx}\operatorname{uxp} = \frac{(1+p)(1-p)}{s^2c^2}.
\]

Thus,

\[
\operatorname{urx}\operatorname{uxp} = \frac{1-p^2}{s^2c^2}.
\]

From Lemma 4.3,

\[
p^2=1-2|sc|.
\]

Therefore,

\[
1-p^2=2|sc|.
\]

Hence,

\[
\operatorname{urx}\operatorname{uxp} = \frac{2|sc|}{s^2c^2}.
\]

Since

\[
s^2c^2=|sc|^2,
\]

we obtain

\[
\operatorname{urx}\operatorname{uxp} = \frac2{|sc|}.
\]

Finally,

\[
|\sin 2x|=2|sc|,
\]

so

\[
\frac2{|sc|} = \frac4{|\sin 2x|}.
\]


\end{proof}

\PCsection{SECTION 7 --- Flatwave}

\PCsubsection{Theorem 7.1 --- Exact Flatwave theorem}

For every \(x\in\mathcal D\),

\[
\boxed{ \operatorname{Flatwave}(x) = \operatorname{sgn}(\sin 2x) }.
\]

\begin{proof}
By definition,

\[
\operatorname{Flatwave} = \frac1{\operatorname{urx}} + \frac1{\operatorname{uxp}}.
\]

Combining fractions,

\[
\operatorname{Flatwave} = \frac{\operatorname{urx}+\operatorname{uxp}} {\operatorname{urx}\operatorname{uxp}}.
\]

By Theorems 5.1 and 6.1,

\[
\operatorname{Flatwave} = \frac{2/(sc)}{2/|sc|}.
\]

Therefore,

\[
\operatorname{Flatwave} = \frac{|sc|}{sc}.
\]

But

\[
\frac{|sc|}{sc} = \operatorname{sgn}(sc).
\]

Since

\[
\operatorname{sgn}(sc) = \operatorname{sgn}(\sin 2x),
\]

we obtain

\[
\boxed{ \operatorname{Flatwave}(x) = \operatorname{sgn}(\sin 2x) }.
\]

\end{proof}

\PCsubsection{Corollary 7.2 --- Exact Flatwave magnitude}

\[
\boxed{ |\operatorname{Flatwave}(x)|=1 }.
\]

\PCsubsection{Corollary 7.3 --- Quadrant values}

\[
\operatorname{Flatwave}(x) = \begin{cases} +1, & x\text{ lies in Quadrant I or III}, \\[4pt] -1, & x\text{ lies in Quadrant II or IV}. \end{cases}
\]

\PCsubsection{Corollary 7.4 --- Flatwave derivative}

On every open quadrant,

\[
\boxed{ \frac{d}{dx}\operatorname{Flatwave}(x)=0 }.
\]

Flatwave is undefined at the excluded quadrant boundaries.


\PCsection{SECTION 8 --- Saw Partition}

\PCsubsection{Theorem 8.1 --- Signed saw partition}

For every \(x\in\mathcal D\),

\[
\boxed{ \operatorname{saw}_{r}(x) + \operatorname{saw}_{x}(x) = \operatorname{sgn}(\sin 2x) }.
\]


This follows directly from the definition of Flatwave and Theorem 7.1.

\PCsubsection{Corollary 8.2 --- Unsigned saw partition}

For every \(x\in\mathcal D\),

\[
\boxed{ \left|\operatorname{saw}_{r}(x)\right| + \left|\operatorname{saw}_{x}(x)\right| = 1 }.
\]

\begin{proof}
By Theorem 4.5, \(\operatorname{urx}\) and \(\operatorname{uxp}\) have the same sign. Their reciprocals therefore also have the same sign.

Consequently,

\[
\left| \operatorname{saw}_{r} + \operatorname{saw}_{x} \right| = |\operatorname{saw}_{r}| + |\operatorname{saw}_{x}|.
\]


But by Theorem 8.1,

\[
\left| \operatorname{saw}_{r} + \operatorname{saw}_{x} \right| = |\operatorname{sgn}(\sin 2x)| = 1.
\]


Therefore,

\[
\boxed{ |\operatorname{saw}_{r}| + |\operatorname{saw}_{x}| = 1 }.
\]


This establishes an exact signed and unsigned partition of unity within the operator system.

It does not yet establish a physical conservation law.


\end{proof}

\PCsection{SECTION 9 --- Half-Angle Representation}

\PCsubsection{Theorem 9.1 --- Quadrant-I half-angle forms}

For

\[
0<x<\frac{\pi}{2},
\]

we have

\[
\boxed{ \operatorname{srx}(x)=\cot\frac{x}{2} },
\]

\[
\boxed{ \operatorname{sxp}(x)=\tan\frac{x}{2} },
\]

\[
\boxed{ \operatorname{cxp}(x) = \tan\left(\frac{\pi}{4}+\frac{x}{2}\right) },
\]

and

\[
\boxed{ \operatorname{crx}(x) = \tan\left(\frac{\pi}{4}-\frac{x}{2}\right) }.
\]

\begin{proof}
In Quadrant I,

\[
|\csc x|=\csc x, \qquad |\sec x|=\sec x.
\]

The standard half-angle identities give

\[
\csc x+\cot x = \frac{1+\cos x}{\sin x} = \cot\frac{x}{2},
\]

and

\[
\csc x-\cot x = \frac{1-\cos x}{\sin x} = \tan\frac{x}{2}.
\]

Similarly,

\[
\sec x+\tan x = \frac{1+\sin x}{\cos x} = \tan\left(\frac{\pi}{4}+\frac{x}{2}\right),
\]

and

\[
\sec x-\tan x = \frac{1-\sin x}{\cos x} = \tan\left(\frac{\pi}{4}-\frac{x}{2}\right).
\]


\end{proof}

\PCsection{SECTION 10 --- Octant Dominance}

\PCsubsection{Theorem 10.1 --- Dominance in the first octant}

For

\[
0<x<\frac{\pi}{4},
\]

\[
\boxed{ \operatorname{srx}(x) > \operatorname{cxp}(x), \operatorname{sxp}(x), \operatorname{crx}(x) }.
\]

\begin{proof}
In this interval,

\[
c>s>0.
\]

Compare \(\operatorname{srx}\) and \(\operatorname{cxp}\):

\[
\operatorname{srx} = \frac{1+c}{s}, \qquad \operatorname{cxp} = \frac{1+s}{c}.
\]

Because \(sc>0\),

\[
\operatorname{srx}>\operatorname{cxp}
\]

if and only if

\[
c(1+c)>s(1+s).
\]

Subtracting the right side from the left gives

\[
c+c^2-s-s^2 = (c-s)+(c^2-s^2).
\]

Factoring,

\[
(c-s)+(c-s)(c+s) = (c-s)(1+c+s).
\]

Both factors are positive. Therefore,

\[
\operatorname{srx}>\operatorname{cxp}.
\]

Also,

\[
\operatorname{srx}>1,
\]

while

\[
\operatorname{sxp}<1
\]

because they are reciprocal.

Likewise,

\[
\operatorname{crx}<1
\]

in Quadrant I.

Therefore, \(\operatorname{srx}\) is the unique dominant operator in the first octant.

\end{proof}

\PCsubsection{Theorem 10.2 --- Dominance in the second octant}

For

\[
\frac{\pi}{4}<x<\frac{\pi}{2},
\]

\[
\boxed{ \operatorname{cxp}(x) > \operatorname{srx}(x), \operatorname{sxp}(x), \operatorname{crx}(x) }.
\]

The proof is identical, with \(s>c>0\).

\PCsubsection{Lemma 10.3 --- Quarter-turn operator permutation}

For every admissible \(x\),

\[
\operatorname{srx}\left(x+\frac{\pi}{2}\right) = \operatorname{crx}(x),
\]

\[
\operatorname{cxp}\left(x+\frac{\pi}{2}\right) = \operatorname{sxp}(x),
\]

\[
\operatorname{crx}\left(x+\frac{\pi}{2}\right) = \operatorname{srx}(x),
\]

and

\[
\operatorname{sxp}\left(x+\frac{\pi}{2}\right) = \operatorname{cxp}(x).
\]

\PCsubsection{Theorem 10.4 --- Complete octant dominance cycle}

Using Theorems 10.1–10.2 and the quarter-turn permutation,

\[
\boxed{ \operatorname{srx} \rightarrow \operatorname{cxp} \rightarrow \operatorname{crx} \rightarrow \operatorname{sxp} \rightarrow \operatorname{srx} }
\]

is the unique magnitude-dominance cycle over successive open octants.

Explicitly:

\[
\begin{array}{c|c} \text{Octant}&\text{Unique dominant operator}\\ \hline (0,\pi/4)&\operatorname{srx}\\ (\pi/4,\pi/2)&\operatorname{cxp}\\ (\pi/2,3\pi/4)&\operatorname{crx}\\ (3\pi/4,\pi)&\operatorname{sxp}\\ (\pi,5\pi/4)&\operatorname{srx}\\ (5\pi/4,3\pi/2)&\operatorname{cxp}\\ (3\pi/2,7\pi/4)&\operatorname{crx}\\ (7\pi/4,2\pi)&\operatorname{sxp} \end{array}
\]

At odd multiples of \(\pi/4\), adjacent dominant operators tie. Those points are admissible but are not unique-winner points.


\PCsection{CLOSURE DIVISION}

Starting from the Pythagorean identity alone, together with ordinary trigonometric definitions, we have rigorously established:

\begin{enumerate}
\item Exact reciprocal conjugates:
\end{enumerate}
\[
\operatorname{crx}=M-\tan x, \qquad \operatorname{sxp}=N-\cot x.
\]

\begin{enumerate}
\item Positivity of all four primitive displayed operators.
\item Exact reciprocal pairing.
\item Exact signed forms of the UNA operators.
\item Common UNA sign:
\end{enumerate}
\[
\operatorname{sgn}(\operatorname{urx}) = \operatorname{sgn}(\operatorname{uxp}) = \operatorname{sgn}(\sin 2x).
\]

\begin{enumerate}
\item The double-angle closure identity:
\end{enumerate}
\[
\operatorname{urx}+\operatorname{uxp} = \frac4{\sin 2x}.
\]

\begin{enumerate}
\item The UNA product identity:
\end{enumerate}
\[
\operatorname{urx}\operatorname{uxp} = \frac4{|\sin 2x|}.
\]

\begin{enumerate}
\item The exact Flatwave theorem:
\end{enumerate}
\[
\operatorname{Flatwave} = \operatorname{sgn}(\sin 2x).
\]

\begin{enumerate}
\item The signed and unsigned saw partitions:
\end{enumerate}
\[
\operatorname{saw}_{r}+\operatorname{saw}_{x} = \operatorname{sgn}(\sin 2x),
\]

\[
|\operatorname{saw}_{r}|+|\operatorname{saw}_{x}|=1.
\]

\begin{enumerate}
\item The complete octant dominance cycle:
\end{enumerate}
\[
\operatorname{srx} \rightarrow \operatorname{cxp} \rightarrow \operatorname{crx} \rightarrow \operatorname{sxp}.
\]

These results constitute a rigorous Pythagorean core for the bounded phase operator calculus.

They establish the operator geometry.

They do not yet establish its physical interpretation.

% ===== END S0_I =====
% ===== BEGIN S0_II =====
% Source: S0.ii.docx
\PCchapter{S0.II}{The Canonical Generator and Quadrant-Local Closure}{Book 0 - Formal Core}

\PCsubsection{Abstract}

This essay determines the algebraic size of the bounded phase operator calculus established in S0.I.

On each open quadrant, let

\[
z(x)=\operatorname{srx}(x).
\]

Using the exact Flatwave theorem, the second primary envelope is shown to be a Möbius transformation of \(z\):

\[
\operatorname{cxp}(x) = \frac{z(x)+\varepsilon} {\varepsilon z(x)-1},
\]

where

\[
\varepsilon = \operatorname{sgn}(\sin 2x)
\]

is constant on the quadrant. The reciprocal operators, UNA operators, saw responses, and Flatwave then reduce to rational functions of the single generator \(z\).

The phase itself also admits a rational readout:

\[
\sin 2x = \frac{4z(z^2-1)}{(z^2+1)^2}.
\]

The generator obeys the autonomous differential equation

\[
\frac{dz}{dx} = -\frac{1+z^2}{2}.
\]

Consequently, every rational expression in \(z\) remains rational in \(z\) after quadrant-local differentiation. This establishes a precise closure result: the calculus is not a finite algebra of named functions, but a one-generator differential field on each admissible quadrant.

No claim of global analytic continuation, physical interpretation, or uniqueness of the chosen generator is made.


\PCsection{0. Scope and Proof Standard}

\PCsubsection{0.1 Scope}

This essay proves:

\begin{enumerate}
\item the rational dependence of all named operators on one generator;
\item the exact Möbius relation between the two primary envelopes;
\item canonical rational forms for the UNA and saw operators;
\item a rational representation of \(\sin 2x\);
\item the differential equation satisfied by each primitive operator;
\item quadrant-local closure under differentiation.
\end{enumerate}
\PCsubsection{0.2 Non-scope}

This essay does not prove:

\begin{itemize}
\item that \(\operatorname{srx}\) is the only possible generator;
\item that the named operator list is closed under arbitrary operations;
\item that every rational function of the generator is physically meaningful;
\item that the calculus supplies a law of temporal evolution;
\item that differentiation with respect to phase is differentiation with respect to time.
\end{itemize}
\PCsubsection{0.3 Inheritance}

All definitions and results of S0.I are inherited.

In particular, for every admissible \(x\),

\[
\operatorname{srx}>0, \qquad \operatorname{cxp}>0,
\]

\[
\operatorname{sxp} = \frac1{\operatorname{srx}}, \qquad \operatorname{crx} = \frac1{\operatorname{cxp}},
\]

and

\[
\operatorname{Flatwave}(x) = \operatorname{sgn}(\sin 2x).
\]

The admissible domain remains

\[
\mathcal D = \mathbb R \setminus \left\{ \frac{k\pi}{2}:k\in\mathbb Z \right\}.
\]


\PCsection{1. Quadrant-Local Setting}

\PCsubsection{Definition 1.1 --- Open quadrants}

For each integer \(k\), define

\[
Q_k = \left( \frac{k\pi}{2}, \frac{(k+1)\pi}{2} \right).
\]

Every \(Q_k\) is a connected component of \(\mathcal D\).

\PCsubsection{Definition 1.2 --- Quadrant sign}

On a fixed quadrant \(Q_k\), define

\[
\varepsilon_k = \operatorname{sgn}(\sin 2x).
\]

Because \(\sin 2x\) has no zero inside \(Q_k\), \(\varepsilon_k\) is constant throughout the quadrant.

For readability, write simply

\[
\varepsilon=\varepsilon_k
\]

when a quadrant has been fixed.

Thus,

\[
\varepsilon\in\{-1,+1\}.
\]

\PCsubsection{Definition 1.3 --- Canonical generator}

Define

\[
\boxed{ z(x)=\operatorname{srx}(x) }.
\]

Since \(\operatorname{srx}>0\),

\[
z(x)>0
\]

throughout every admissible quadrant.

\PCsubsection{Definition 1.4 --- Second primary envelope}

For temporary convenience, write

\[
w(x)=\operatorname{cxp}(x).
\]

The objective of the following sections is to eliminate \(w\) as an independent algebraic quantity.


\PCsection{2. Rational Form of Flatwave}

\PCsubsection{Lemma 2.1 --- Rational saw forms}

For every \(x\in\mathcal D\),

\[
\operatorname{saw}_{r} = \frac{w}{zw-1},
\]

and

\[
\operatorname{saw}_{x} = \frac{z}{zw-1}.
\]

\begin{proof}
By definition,

\[
\operatorname{urx} = \operatorname{srx}-\operatorname{crx}.
\]

Using

\[
\operatorname{srx}=z \qquad\text{and}\qquad \operatorname{crx}=\frac1w,
\]

we have

\[
\operatorname{urx} = z-\frac1w.
\]

Combining the terms,

\[
\operatorname{urx} = \frac{zw-1}{w}.
\]

Therefore,

\[
\operatorname{saw}_{r} = \frac1{\operatorname{urx}} = \frac{w}{zw-1}.
\]

Similarly,

\[
\operatorname{uxp} = w-\frac1z = \frac{zw-1}{z}.
\]

Hence,

\[
\operatorname{saw}_{x} = \frac1{\operatorname{uxp}} = \frac{z}{zw-1}.
\]

\end{proof}

\PCsubsection{Corollary 2.2 --- Rational Flatwave form}

\[
\boxed{ \operatorname{Flatwave} = \frac{z+w}{zw-1} }.
\]

\begin{proof}
Adding the two saw forms gives

\[
\operatorname{Flatwave} = \frac{w}{zw-1} + \frac{z}{zw-1}.
\]

Therefore,

\[
\operatorname{Flatwave} = \frac{z+w}{zw-1}.
\]


\end{proof}

\PCsection{3. The Möbius Reduction}

\PCsubsection{Theorem 3.1 --- Primary-envelope dependence}

On every open quadrant,

\[
\boxed{ w = \frac{z+\varepsilon} {\varepsilon z-1} }.
\]

Equivalently,

\[
\boxed{ \operatorname{cxp}(x) = \frac{\operatorname{srx}(x)+\varepsilon} {\varepsilon\operatorname{srx}(x)-1} }.
\]

\begin{proof}
By the Flatwave theorem of S0.I,

\[
\operatorname{Flatwave}=\varepsilon.
\]

By Corollary 2.2,

\[
\frac{z+w}{zw-1} = \varepsilon.
\]

Multiplying by \(zw-1\),

\[
z+w = \varepsilon(zw-1).
\]

Expanding the right-hand side,

\[
z+w = \varepsilon zw-\varepsilon.
\]

Collect the terms containing \(w\):

\[
w-\varepsilon zw = -\varepsilon-z.
\]

Factor \(w\):

\[
w(1-\varepsilon z) = -(z+\varepsilon).
\]

Multiplying numerator and denominator by \(-1\),

\[
w = \frac{z+\varepsilon} {\varepsilon z-1}.
\]

It remains to verify that the denominator is nonzero.

If \(\varepsilon=-1\), then

\[
\varepsilon z-1=-z-1<0,
\]

because \(z>0\).

If \(\varepsilon=+1\), the denominator could vanish only if \(z=1\). But \(z=1\) occurs only at an excluded quadrant boundary. Therefore,

\[
\varepsilon z-1\neq0
\]

on every open quadrant.

Hence,

\[
\boxed{ w = \frac{z+\varepsilon} {\varepsilon z-1} }.
\]

\end{proof}

\PCsubsection{Corollary 3.2 --- Reciprocal forms}

The two reciprocal operators are

\[
\boxed{ \operatorname{sxp} = \frac1z },
\]

and

\[
\boxed{ \operatorname{crx} = \frac{\varepsilon z-1} {z+\varepsilon} }.
\]

\begin{proof}
The first identity follows from the definition of \(\operatorname{sxp}\).

For the second,

\[
\operatorname{crx} = \frac1{\operatorname{cxp}} = \frac1w.
\]

Taking the reciprocal of Theorem 3.1 gives

\[
\operatorname{crx} = \frac{\varepsilon z-1} {z+\varepsilon}.
\]

\end{proof}

\PCsubsection{Corollary 3.3 --- One-generator sufficiency}

On a fixed quadrant, every primitive displayed operator is a rational function of \(z\).

Explicitly,

\[
\operatorname{srx}=z,
\]

\[
\operatorname{sxp}=\frac1z,
\]

\[
\operatorname{cxp} = \frac{z+\varepsilon}{\varepsilon z-1},
\]

and

\[
\operatorname{crx} = \frac{\varepsilon z-1}{z+\varepsilon}.
\]

Thus the four displayed operators do not constitute four algebraically independent generators.


\PCsection{4. Generator Range and Direction Sign}

\PCsubsection{Theorem 4.1 --- Generator range}

On every quadrant with \(\varepsilon=+1\),

\[
\boxed{z>1}.
\]

On every quadrant with \(\varepsilon=-1\),

\[
\boxed{0<z<1}.
\]

\begin{proof}
From S0.I, on an upper sine half-cycle,

\[
z=\cot\frac{x}{2},
\]

while on the corresponding lower half-cycle,

\[
z=-\tan\frac{x}{2}.
\]

In Quadrants I and III, \(\sin 2x>0\), so \(\varepsilon=+1\). On both such quadrants, the corresponding positive half-angle expression takes values greater than \(1\).

In Quadrants II and IV, \(\sin 2x<0\), so \(\varepsilon=-1\). On both such quadrants, the corresponding positive half-angle expression lies strictly between \(0\) and \(1\).

The value \(z=1\) occurs only at odd multiples of \(\pi/2\), which are excluded from \(\mathcal D\).

Therefore,

\[
\varepsilon=+1\implies z>1,
\]

and

\[
\varepsilon=-1\implies 0<z<1.
\]

\end{proof}

\PCsubsection{Corollary 4.2 --- Direction sign from the generator}

\[
\boxed{ \varepsilon = \operatorname{sgn}(z-1) }.
\]

Since \(z>0\),

\[
\boxed{ \varepsilon = \operatorname{sgn}(z^2-1) }.
\]

The Flatwave sign therefore need not be supplied independently once the generator value is known.


\PCsection{5. Canonical UNA Forms}

\PCsubsection{Theorem 5.1 --- Canonical Unainverex form}

\[
\boxed{ \operatorname{urx} = \frac{z^2+1}{z+\varepsilon} }.
\]

\begin{proof}
By definition,

\[
\operatorname{urx} = \operatorname{srx}-\operatorname{crx}.
\]

Using the canonical expressions,

\[
\operatorname{urx} = z- \frac{\varepsilon z-1}{z+\varepsilon}.
\]

Place the terms over a common denominator:

\[
\operatorname{urx} = \frac{z(z+\varepsilon)-(\varepsilon z-1)} {z+\varepsilon}.
\]

Expand the numerator:

\[
z(z+\varepsilon)-(\varepsilon z-1) = z^2+\varepsilon z-\varepsilon z+1.
\]

Therefore,

\[
\operatorname{urx} = \frac{z^2+1}{z+\varepsilon}.
\]

\end{proof}

\PCsubsection{Theorem 5.2 --- Canonical Unaexpon form}

\[
\boxed{ \operatorname{uxp} = \frac{z^2+1} {z(\varepsilon z-1)} }.
\]

\begin{proof}
By definition,

\[
\operatorname{uxp} = \operatorname{cxp}-\operatorname{sxp}.
\]

Substitute the canonical expressions:

\[
\operatorname{uxp} = \frac{z+\varepsilon}{\varepsilon z-1} - \frac1z.
\]

Use the common denominator

\[
z(\varepsilon z-1):
\]

\[
\operatorname{uxp} = \frac{z(z+\varepsilon)-(\varepsilon z-1)} {z(\varepsilon z-1)}.
\]

The numerator is again

\[
z^2+\varepsilon z-\varepsilon z+1 = z^2+1.
\]

Thus,

\[
\operatorname{uxp} = \frac{z^2+1} {z(\varepsilon z-1)}.
\]

\end{proof}

\PCsubsection{Corollary 5.3 --- UNA signs}

Because

\[
z^2+1>0,
\]

the signs are determined entirely by the denominators.

Theorem 4.1 therefore gives

\[
\operatorname{sgn}(\operatorname{urx}) = \operatorname{sgn}(\operatorname{uxp}) = \varepsilon.
\]

This recovers the common-sign theorem of S0.I in canonical-generator form.


\PCsection{6. Canonical Saw Forms}

\PCsubsection{Theorem 6.1 --- Canonical inverex saw}

\[
\boxed{ \operatorname{saw}_{r} = \frac{z+\varepsilon}{z^2+1} }.
\]

\begin{proof}
By definition,

\[
\operatorname{saw}_{r} = \frac1{\operatorname{urx}}.
\]

Using Theorem 5.1,

\[
\operatorname{saw}_{r} = \frac{z+\varepsilon}{z^2+1}.
\]

\end{proof}

\PCsubsection{Theorem 6.2 --- Canonical expon saw}

\[
\boxed{ \operatorname{saw}_{x} = \frac{z(\varepsilon z-1)}{z^2+1} }.
\]

\begin{proof}
By definition,

\[
\operatorname{saw}_{x} = \frac1{\operatorname{uxp}}.
\]

Taking the reciprocal of Theorem 5.2 gives

\[
\operatorname{saw}_{x} = \frac{z(\varepsilon z-1)} {z^2+1}.
\]

\end{proof}

\PCsubsection{Corollary 6.3 --- Flatwave reduction}

Adding the two canonical saw forms,

\[
\operatorname{Flatwave} = \frac{z+\varepsilon} {z^2+1} + \frac{z(\varepsilon z-1)} {z^2+1}.
\]

The numerator becomes

\[
z+\varepsilon+\varepsilon z^2-z.
\]

Therefore,

\[
\operatorname{Flatwave} = \frac{\varepsilon(z^2+1)} {z^2+1}.
\]

Hence,

\[
\boxed{ \operatorname{Flatwave}=\varepsilon }.
\]

The exact Flatwave theorem is therefore built directly into the canonical forms.


\PCsection{7. Rational Phase Readout}

\PCsubsection{Theorem 7.1 --- Double-angle phase in generator form}

\[
\boxed{ \sin 2x = \frac{4z(z^2-1)} {(z^2+1)^2} }.
\]

\begin{proof}
From S0.I,

\[
\operatorname{urx} + \operatorname{uxp} = \frac4{\sin 2x}.
\]

Using Theorems 5.1 and 5.2,

\[
\operatorname{urx} + \operatorname{uxp} = \frac{z^2+1}{z+\varepsilon} + \frac{z^2+1}{z(\varepsilon z-1)}.
\]

Factor \(z^2+1\):

\[
\operatorname{urx} + \operatorname{uxp} = (z^2+1) \left[ \frac1{z+\varepsilon} + \frac1{z(\varepsilon z-1)} \right].
\]

Combine the fractions:

\[
\operatorname{urx} + \operatorname{uxp} = (z^2+1) \frac{z(\varepsilon z-1)+(z+\varepsilon)} {z(z+\varepsilon)(\varepsilon z-1)}.
\]

Simplify the numerator:

\[
z(\varepsilon z-1)+(z+\varepsilon) = \varepsilon z^2-z+z+\varepsilon.
\]

Thus,

\[
z(\varepsilon z-1)+(z+\varepsilon) = \varepsilon(z^2+1).
\]

Next simplify

\[
(z+\varepsilon)(\varepsilon z-1).
\]

Expanding,

\[
(z+\varepsilon)(\varepsilon z-1) = \varepsilon z^2-z+\varepsilon^2z-\varepsilon.
\]

Since \(\varepsilon^2=1\),

\[
(z+\varepsilon)(\varepsilon z-1) = \varepsilon z^2-z+z-\varepsilon = \varepsilon(z^2-1).
\]

Therefore,

\[
\operatorname{urx} + \operatorname{uxp} = (z^2+1) \frac{\varepsilon(z^2+1)} {\varepsilon z(z^2-1)}.
\]

Cancel \(\varepsilon\):

\[
\operatorname{urx} + \operatorname{uxp} = \frac{(z^2+1)^2} {z(z^2-1)}.
\]

Since this also equals \(4/\sin 2x\),

\[
\frac4{\sin 2x} = \frac{(z^2+1)^2} {z(z^2-1)}.
\]

Taking reciprocals,

\[
\frac{\sin 2x}{4} = \frac{z(z^2-1)} {(z^2+1)^2}.
\]

Therefore,

\[
\boxed{ \sin 2x = \frac{4z(z^2-1)} {(z^2+1)^2} }.
\]

\end{proof}

\PCsubsection{Corollary 7.2 --- Flatwave sign}

Because \(z>0\) and \((z^2+1)^2>0\),

\[
\operatorname{sgn}(\sin 2x) = \operatorname{sgn}(z^2-1).
\]

Therefore,

\[
\operatorname{Flatwave} = \operatorname{sgn}(z^2-1).
\]

\PCsubsection{Corollary 7.3 --- Boundary identification}

The excluded quadrant boundaries correspond to:

\[
z\rightarrow\infty, \qquad z\rightarrow0, \qquad\text{or}\qquad z\rightarrow1,
\]

depending on the side of approach.

The value \(z=1\) is never attained on an admissible open quadrant.


\PCsection{8. Differential Law of the Generator}

\PCsubsection{Analytic Convention 8.1}

Ordinary real differentiation rules are inherited on each open quadrant, including the chain rule and the derivatives

\[
\frac{d}{du}\tan u=\sec ^2u,
\]

and

\[
\frac{d}{du}\cot u=-\csc ^2u.
\]

No differentiation across an excluded boundary is asserted.

\PCsubsection{Theorem 8.2 --- Generator differential equation}

On every open quadrant,

\[
\boxed{ \frac{dz}{dx} = -\frac{1+z^2}{2} }.
\]

\begin{proof}
From the half-angle forms proved in S0.I, there are two branch types.

On the first branch,

\[
z=\cot\frac{x}{2}.
\]

Differentiate:

\[
\frac{dz}{dx} = -\csc ^2\frac{x}{2}\cdot\frac12.
\]

Using

\[
\csc ^2u=1+\cot ^2u,
\]

we obtain

\[
\frac{dz}{dx} = -\frac12 \left( 1+\cot ^2\frac{x}{2} \right).
\]

Since

\[
z=\cot\frac{x}{2},
\]

this becomes

\[
\frac{dz}{dx} = -\frac{1+z^2}{2}.
\]

On the second branch,

\[
z=-\tan\frac{x}{2}.
\]

Differentiating,

\[
\frac{dz}{dx} = -\sec ^2\frac{x}{2}\cdot\frac12.
\]

Since

\[
\sec ^2u=1+\tan ^2u,
\]

and

\[
z^2=\tan ^2\frac{x}{2},
\]

we again obtain

\[
\frac{dz}{dx} = -\frac{1+z^2}{2}.
\]

Thus the same differential equation holds on every open quadrant:

\[
\boxed{ z' = -\frac{1+z^2}{2} }.
\]

\end{proof}

\PCsubsection{Corollary 8.3 --- Strict monotonicity}

Since

\[
1+z^2>0,
\]

we have

\[
z'<0.
\]

Therefore,

\[
\boxed{ \operatorname{srx}(x) \text{ is strictly decreasing on every open quadrant.} }
\]


\PCsection{9. Differential Laws of the Primitive Family}

\PCsubsection{Theorem 9.1 --- Cosinexpon differential equation}

Let

\[
w=\operatorname{cxp}.
\]

Then

\[
\boxed{ \frac{dw}{dx} = \frac{1+w^2}{2} }.
\]

\begin{proof}
From Theorem 3.1,

\[
w = \frac{z+\varepsilon} {\varepsilon z-1}.
\]

The quadrant sign \(\varepsilon\) is constant on the chosen quadrant. Differentiate with respect to \(z\):

\[
\frac{dw}{dz} = \frac{ (\varepsilon z-1) - \varepsilon(z+\varepsilon) } {(\varepsilon z-1)^2}.
\]

Simplify the numerator:

\[
\varepsilon z-1-\varepsilon z-\varepsilon^2 = -1-1 = -2.
\]

Therefore,

\[
\frac{dw}{dz} = -\frac2{(\varepsilon z-1)^2}.
\]

By Theorem 8.2,

\[
\frac{dz}{dx} = -\frac{1+z^2}{2}.
\]

Using the chain rule,

\[
\frac{dw}{dx} = \frac{dw}{dz}\frac{dz}{dx}.
\]

Hence,

\[
\frac{dw}{dx} = \left[ -\frac2{(\varepsilon z-1)^2} \right] \left[ -\frac{1+z^2}{2} \right].
\]

Thus,

\[
\frac{dw}{dx} = \frac{1+z^2} {(\varepsilon z-1)^2}.
\]

Now compute \(1+w^2\):

\[
1+w^2 = 1+ \left( \frac{z+\varepsilon} {\varepsilon z-1} \right)^2.
\]

Using a common denominator,

\[
1+w^2 = \frac{ (\varepsilon z-1)^2+(z+\varepsilon)^2 } {(\varepsilon z-1)^2}.
\]

Expand the numerator:

\[
(\varepsilon z-1)^2 = z^2-2\varepsilon z+1,
\]

and

\[
(z+\varepsilon)^2 = z^2+2\varepsilon z+1.
\]

Adding them gives

\[
2z^2+2 = 2(1+z^2).
\]

Therefore,

\[
1+w^2 = \frac{2(1+z^2)} {(\varepsilon z-1)^2}.
\]

Hence,

\[
\frac{1+w^2}{2} = \frac{1+z^2} {(\varepsilon z-1)^2}.
\]

This is exactly \(dw/dx\). Therefore,

\[
\boxed{ w' = \frac{1+w^2}{2} }.
\]

\end{proof}

\PCsubsection{Corollary 9.2 --- Primitive derivative table}

The four primitive operators satisfy:

\[
\boxed{ \operatorname{srx}' = -\frac{1+\operatorname{srx}^2}{2} },
\]

\[
\boxed{ \operatorname{sxp}' = \frac{1+\operatorname{sxp}^2}{2} },
\]

\[
\boxed{ \operatorname{cxp}' = \frac{1+\operatorname{cxp}^2}{2} },
\]

and

\[
\boxed{ \operatorname{crx}' = -\frac{1+\operatorname{crx}^2}{2} }.
\]

\begin{proof}
The first and third identities are Theorems 8.2 and 9.1.

For

\[
\operatorname{sxp}=\frac1z,
\]

differentiate:

\[
\operatorname{sxp}' = -\frac{z'}{z^2}.
\]

Using

\[
z'=-\frac{1+z^2}{2},
\]

we obtain

\[
\operatorname{sxp}' = \frac{1+z^2}{2z^2}.
\]

Therefore,

\[
\operatorname{sxp}' = \frac12 \left( 1+\frac1{z^2} \right).
\]

Since

\[
\operatorname{sxp}=\frac1z,
\]

it follows that

\[
\operatorname{sxp}' = \frac{1+\operatorname{sxp}^2}{2}.
\]

The proof for

\[
\operatorname{crx}=\frac1w
\]

is identical:

\[
\operatorname{crx}' = -\frac{w'}{w^2} = -\frac{1+w^2}{2w^2} = -\frac{1+\operatorname{crx}^2}{2}.
\]

\end{proof}

\PCsubsection{Corollary 9.3 --- Primitive monotonicity}

On every open quadrant:

\[
\operatorname{srx}'<0, \qquad \operatorname{crx}'<0,
\]

while

\[
\operatorname{sxp}'>0, \qquad \operatorname{cxp}'>0.
\]

Thus the two inverex operators decrease with phase, while the two expon operators increase with phase.

This is an exact local phase statement, not yet a temporal growth or decay law.


\PCsection{10. Quadrant-Local Differential Closure}

\PCsubsection{Definition 10.1 --- Rational expression field}

Fix an open quadrant and its constant sign \(\varepsilon\).

Let

\[
\mathbb R(z)
\]

denote the set of rational functions

\[
F(z)=\frac{P(z)}{Q(z)},
\]

where

\[
P,Q\in\mathbb R[z]
\]

and \(Q(z)\neq0\) on the declared evaluation interval.

\PCsubsection{Theorem 10.2 --- Algebraic containment}

Every named operator and named composite introduced in S0.I belongs to

\[
\mathbb R(z)
\]

on each open quadrant.

\begin{proof}
The canonical forms derived above are:

\[
\operatorname{srx}=z,
\]

\[
\operatorname{sxp}=\frac1z,
\]

\[
\operatorname{cxp} = \frac{z+\varepsilon}{\varepsilon z-1},
\]

\[
\operatorname{crx} = \frac{\varepsilon z-1}{z+\varepsilon},
\]

\[
\operatorname{urx} = \frac{z^2+1}{z+\varepsilon},
\]

\[
\operatorname{uxp} = \frac{z^2+1}{z(\varepsilon z-1)},
\]

\[
\operatorname{saw}_{r} = \frac{z+\varepsilon}{z^2+1},
\]

\[
\operatorname{saw}_{x} = \frac{z(\varepsilon z-1)}{z^2+1},
\]

and

\[
\operatorname{Flatwave}=\varepsilon.
\]

Each is a rational function of \(z\) with real coefficients once \(\varepsilon\) is fixed.

Therefore, every named operator belongs to \(\mathbb R(z)\).

\end{proof}

\PCsubsection{Theorem 10.3 --- Closure under differentiation}

If

\[
F(z)\in\mathbb R(z),
\]

then

\[
\frac{d}{dx}F(z(x)) \in\mathbb R(z)
\]

on every open quadrant where \(F\) is defined.

\begin{proof}
Let

\[
F(z)=\frac{P(z)}{Q(z)},
\]

with \(P,Q\in\mathbb R[z]\).

Differentiating with respect to \(z\),

\[
\frac{dF}{dz} = \frac{P'(z)Q(z)-P(z)Q'(z)} {Q(z)^2}.
\]

The numerator and denominator are polynomials in \(z\). Therefore,

\[
\frac{dF}{dz}\in\mathbb R(z).
\]

By the chain rule,

\[
\frac{dF}{dx} = \frac{dF}{dz}\frac{dz}{dx}.
\]

From Theorem 8.2,

\[
\frac{dz}{dx} = -\frac{1+z^2}{2},
\]

which is a polynomial in \(z\).

The product of two rational functions of \(z\) is again a rational function of \(z\). Therefore,

\[
\frac{dF}{dx}\in\mathbb R(z).
\]

\end{proof}

\PCsubsection{Corollary 10.4 --- Closure under repeated differentiation}

For every nonnegative integer \(n\),

\[
\boxed{ \frac{d^n}{dx^n}F(z(x)) \in\mathbb R(z) }.
\]

\begin{proof}
The result follows by induction.

The case \(n=0\) is the assumption that \(F\in\mathbb R(z)\).

If

\[
\frac{d^nF}{dx^n}\in\mathbb R(z),
\]

then Theorem 10.3 implies

\[
\frac{d}{dx} \left( \frac{d^nF}{dx^n} \right) = \frac{d^{n+1}F}{dx^{n+1}} \in\mathbb R(z).
\]

Therefore, the result holds for every \(n\geq0\).


\end{proof}

\PCsection{11. Meaning of Closure}

\PCsubsection{Theorem 11.1 --- One-generator differential-field closure}

On every open quadrant, the bounded phase operator calculus generated by the named R-operators is contained in the differential field

\[
\boxed{\mathbb R(z)}
\]

with derivation determined by

\[
\boxed{ z' = -\frac{1+z^2}{2} }.
\]

\PCsubsection{Clarification 11.2 --- What has been proved}

The following statement is proved:

Every named operator, every rational combination of named operators, and every quadrant-local derivative of such a combination reduces to a rational function of one generator.

\PCsubsection{Clarification 11.3 --- What has not been proved}

The following stronger statements are not proved:

\begin{itemize}
\item that every rational function of \(z\) should be designated an R-operator;
\item that the finite named list is closed under all rational operations;
\item that no alternative generator exists;
\item that every admissible phase function belongs to \(\mathbb R(z)\);
\item that the expression field has a physical interpretation.
\end{itemize}
The closure is a mathematical containment and reduction result.


\PCsection{12. Minimality of Generator Count}

\PCsubsection{Theorem 12.1 --- One nonconstant generator is sufficient}

The named quadrant-local calculus requires no more than one nonconstant algebraic generator.

\begin{proof}
Theorem 10.2 explicitly expresses every named operator as a rational function of \(z\).

Therefore, one nonconstant generator is sufficient.

\end{proof}

\PCsubsection{Theorem 12.2 --- Zero nonconstant generators are insufficient}

No field containing only real constants can represent the named calculus.

\begin{proof}
A field generated by no nonconstant element over \(\mathbb R\) contains only constant functions.

But

\[
z(x)=\operatorname{srx}(x)
\]

is strictly decreasing on every open quadrant and is therefore nonconstant.

Hence, a constant-only field cannot contain \(z\), and therefore cannot contain the full named calculus.

At least one nonconstant generator is required.

\end{proof}

\PCsubsection{Corollary 12.3 --- Minimal generator count}

Exactly one nonconstant generator is necessary and sufficient for the named calculus on each open quadrant.

This proves minimality of the generator count.

It does not prove that \(\operatorname{srx}\) is the unique possible generator.


\PCsection{13. Terminal Statement}

S0.I established the Pythagorean identities of the bounded phase operators.

S0.II establishes their algebraic organization.

On each open quadrant:

\[
z=\operatorname{srx}
\]

is sufficient to generate the entire named operator family.

The second primary envelope is not independent:

\[
\operatorname{cxp} = \frac{z+\varepsilon} {\varepsilon z-1}.
\]

The reciprocal, UNA, saw, and Flatwave operators are rational functions of \(z\).

The double-angle phase readout is

\[
\sin 2x = \frac{4z(z^2-1)} {(z^2+1)^2}.
\]

The generator obeys

\[
z' = -\frac{1+z^2}{2}.
\]

Therefore, the quadrant-local expression field is closed under differentiation:

\[
F(z)\in\mathbb R(z) \quad\Longrightarrow\quad \frac{dF}{dx}\in\mathbb R(z).
\]

The result is a precise form of closure:

\[
\boxed{ \text{one generator, rational reduction, differential closure.} }
\]

No additional primitive function is required to express or differentiate the named bounded phase operators within an open quadrant.

% ===== END S0_II =====
% ===== BEGIN S0_III =====
% Source: S0.iii.docx
\PCchapter{S0.III}{Tangential Geometry and Bounded Transfer}{Book 0 - Formal Core}

\PCsubsection{Abstract}

This essay develops the quadrant-local differential geometry of the bounded phase operators established in S0.I and reduced to one generator in S0.II.

The phase derivative is introduced as a derived operation rather than a new primitive. Each displayed operator satisfies an autonomous Riccati equation. Its normalized derivative reduces exactly to an absolute reciprocal trigonometric function:

\[
\widehat T[\operatorname{srx}] = -|\csc x|, \qquad \widehat T[\operatorname{sxp}] = |\csc x|,
\]

\[
\widehat T[\operatorname{cxp}] = |\sec x|, \qquad \widehat T[\operatorname{crx}] = -|\sec x|.
\]

Reciprocal pairing therefore reverses normalized tangential direction without introducing a new rate magnitude.

A pair of bounded auxiliary coordinates \(p\) and \(q\) is then defined. They satisfy

\[
p'=-q, \qquad q'=p, \qquad p^2+q^2=2.
\]

These coordinates yield a canonical transfer variable

\[
\lambda=\frac{1-p}{2},
\]

which traverses strictly from \(0\) to \(1\) across every open quadrant. The two saw responses become

\[
\operatorname{saw}_{r} = \varepsilon\lambda, \qquad \operatorname{saw}_{x} = \varepsilon(1-\lambda),
\]

where

\[
\varepsilon=\operatorname{sgn}(\sin 2x).
\]

Their derivatives are equal and opposite:

\[
\operatorname{saw}_{r}' = \varepsilon\frac q2, \qquad \operatorname{saw}_{x}' = -\varepsilon\frac q2.
\]

The total Flatwave derivative is therefore zero on every open quadrant. Transfer occurs entirely between the two saw components while their signed total remains fixed.

The transfer rate is bounded:

\[
\frac12 < \left| \operatorname{saw}_{r}' \right| = \left| \operatorname{saw}_{x}' \right| \leq \frac1{\sqrt2}.
\]

Each saw changes curvature only at the admissible midpoint of a quadrant, so curvature sign is fixed within every open octant.

No temporal, physical, or causal interpretation is assumed.


\PCsection{0. Scope and Method}

\PCsubsection{0.1 Scope}

This essay proves:

\begin{enumerate}
\item the autonomous phase-derivative equations of the four primitive operators;
\item their normalized tangential responses;
\item exact reversal of normalized response under reciprocal pairing;
\item the location of primary rate balance;
\item a bounded transfer coordinate for the two saw responses;
\item equal-and-opposite saw transfer;
\item the boundedness of the saw-transfer rate;
\item the curvature structure of the saw responses;
\item corresponding differential forms for the UNA operators;
\item closure of the tangential calculus.
\end{enumerate}
\PCsubsection{0.2 Non-scope}

This essay does not:

\begin{itemize}
\item identify phase \(x\) with time;
\item specify a traversal law \(x(\tau)\);
\item introduce force, energy, mass, or information;
\item interpret derivative sign as physical causation;
\item assert a physical conservation law;
\item integrate across excluded quadrant boundaries.
\end{itemize}
\PCsubsection{0.3 Inheritance}

All definitions and proved results of S0.I and S0.II are inherited.

In particular, the admissible domain is

\[
\mathcal D = \mathbb R \setminus \left\{ \frac{k\pi}{2}:k\in\mathbb Z \right\}.
\]

The primitive operators are positive on \(\mathcal D\), and the Flatwave satisfies

\[
\operatorname{Flatwave}(x) = \operatorname{sgn}(\sin 2x).
\]

On each open quadrant,

\[
z(x)=\operatorname{srx}(x)
\]

satisfies

\[
z' = -\frac{1+z^2}{2}.
\]

Every named operator is a rational function of \(z\) on that quadrant.

\PCsubsection{0.4 Analytic convention}

Ordinary real differentiation is inherited on each connected component of \(\mathcal D\).

No derivative is defined through an excluded boundary.


\PCsection{1. Tangential Notation}

\PCsubsection{Definition 1.1 --- Phase derivative}

For a differentiable function \(f\) on an open quadrant, define its phase-tangential derivative by

\[
\boxed{ Tf=\frac{df}{dx} }.
\]

The operator \(T\) is derived from ordinary differentiation. It is not a new primitive member of the R-operator family.

\PCsubsection{Definition 1.2 --- Normalized tangential response}

For a differentiable positive function \(f\), define

\[
\boxed{ \widehat T[f] = \frac{f'(x)}{f(x)} }.
\]

Since the primitive operators are positive,

\[
\widehat T[f] = \frac{d}{dx}\ln f
\]

for every primitive operator \(f\).

\PCsubsection{Definition 1.3 --- Tangential direction}

Define

\[
\operatorname{dir}(f) = \operatorname{sgn}\bigl(\widehat T[f]\bigr).
\]

\PCsubsection{Definition 1.4 --- Tangential speed}

Define

\[
V_f = \left| \widehat T[f] \right|.
\]

These quantities describe change with respect to phase position \(x\), not change with respect to physical time.


\PCsection{2. Primitive Differential Equations}

\PCsubsection{Theorem 2.1 --- Sininverex differential equation}

For every admissible \(x\),

\[
\boxed{ \operatorname{srx}'(x) = -\frac{1+\operatorname{srx}(x)^2}{2} }.
\]

\begin{proof}
By S0.II, on each open quadrant the generator

\[
z=\operatorname{srx}
\]

satisfies

\[
z' = -\frac{1+z^2}{2}.
\]

Substitution gives the stated identity.

\end{proof}

\PCsubsection{Theorem 2.2 --- Sinexpon differential equation}

For every admissible \(x\),

\[
\boxed{ \operatorname{sxp}'(x) = \frac{1+\operatorname{sxp}(x)^2}{2} }.
\]

\begin{proof}
Since

\[
\operatorname{sxp} = \frac1{\operatorname{srx}},
\]

differentiate:

\[
\operatorname{sxp}' = -\frac{\operatorname{srx}'} {\operatorname{srx}^2}.
\]

Using Theorem 2.1,

\[
\operatorname{sxp}' = -\frac{ -\frac12(1+\operatorname{srx}^2) }{ \operatorname{srx}^2 }.
\]

Thus,

\[
\operatorname{sxp}' = \frac12 \left( 1+\frac1{\operatorname{srx}^2} \right).
\]

Because

\[
\operatorname{sxp} = \frac1{\operatorname{srx}},
\]

we obtain

\[
\boxed{ \operatorname{sxp}' = \frac{1+\operatorname{sxp}^2}{2} }.
\]

\end{proof}

\PCsubsection{Theorem 2.3 --- Cosinexpon differential equation}

For every admissible \(x\),

\[
\boxed{ \operatorname{cxp}'(x) = \frac{1+\operatorname{cxp}(x)^2}{2} }.
\]

\begin{proof}
This was established from the canonical Möbius representation in S0.II. It may also be verified directly.

Let

\[
M=|\sec x|.
\]

On a fixed quadrant, let

\[
\beta=\operatorname{sgn}(\cos x).
\]

Then

\[
M=\beta\sec x,
\]

because \(\beta\) is constant on that quadrant.

Differentiate:

\[
M' = \beta\sec x\tan x = M\tan x.
\]

Since

\[
\operatorname{cxp} = M+\tan x,
\]

we have

\[
\operatorname{cxp}' = M\tan x+\sec^2x.
\]

Using

\[
\sec^2x=M^2,
\]

this becomes

\[
\operatorname{cxp}' = M\tan x+M^2 = M(M+\tan x).
\]

Therefore,

\[
\operatorname{cxp}' = M\operatorname{cxp}.
\]

From the reciprocal-conjugate identity,

\[
\operatorname{cxp} + \operatorname{crx} = 2M.
\]

Since

\[
\operatorname{crx} = \frac1{\operatorname{cxp}},
\]

we obtain

\[
2M = \operatorname{cxp} + \frac1{\operatorname{cxp}}.
\]

Multiplying by \(\operatorname{cxp}/2\),

\[
M\operatorname{cxp} = \frac{1+\operatorname{cxp}^2}{2}.
\]

Hence,

\[
\boxed{ \operatorname{cxp}' = \frac{1+\operatorname{cxp}^2}{2} }.
\]

\end{proof}

\PCsubsection{Theorem 2.4 --- Cosininverex differential equation}

For every admissible \(x\),

\[
\boxed{ \operatorname{crx}'(x) = -\frac{1+\operatorname{crx}(x)^2}{2} }.
\]

\begin{proof}
Since

\[
\operatorname{crx} = \frac1{\operatorname{cxp}},
\]

differentiate:

\[
\operatorname{crx}' = -\frac{\operatorname{cxp}'} {\operatorname{cxp}^2}.
\]

Using Theorem 2.3,

\[
\operatorname{crx}' = -\frac{ \frac12(1+\operatorname{cxp}^2) }{ \operatorname{cxp}^2 }.
\]

Thus,

\[
\operatorname{crx}' = -\frac12 \left( 1+\frac1{\operatorname{cxp}^2} \right).
\]

Since

\[
\operatorname{crx} = \frac1{\operatorname{cxp}},
\]

we conclude that

\[
\boxed{ \operatorname{crx}' = -\frac{1+\operatorname{crx}^2}{2} }.
\]

\end{proof}

\PCsubsection{Corollary 2.5 --- Primitive direction table}

On every open quadrant,

\[
\operatorname{srx}'<0, \qquad \operatorname{crx}'<0,
\]

and

\[
\operatorname{sxp}'>0, \qquad \operatorname{cxp}'>0.
\]

Thus the inverex operators decrease with increasing phase, while the expon operators increase with increasing phase.

This is an ordering statement in \(x\), not yet a physical growth or decay law.


\PCsection{3. Normalized Primitive Responses}

\PCsubsection{Lemma 3.1 --- Sine-family sum}

\[
\boxed{ \operatorname{srx} + \operatorname{sxp} = 2N }.
\]

\begin{proof}
By the reciprocal-conjugate forms,

\[
\operatorname{srx} = N+\cot x,
\]

and

\[
\operatorname{sxp} = N-\cot x.
\]

Adding,

\[
\operatorname{srx} + \operatorname{sxp} = 2N.
\]

\end{proof}

\PCsubsection{Lemma 3.2 --- Cosine-family sum}

\[
\boxed{ \operatorname{cxp} + \operatorname{crx} = 2M }.
\]

\begin{proof}
Using

\[
\operatorname{cxp} = M+\tan x
\]

and

\[
\operatorname{crx} = M-\tan x,
\]

addition gives

\[
\operatorname{cxp} + \operatorname{crx} = 2M.
\]

\end{proof}

\PCsubsection{Theorem 3.3 --- Sininverex normalized response}

\[
\boxed{ \widehat T[\operatorname{srx}] = -N = -|\csc x| }.
\]

\begin{proof}
From Theorem 2.1,

\[
\widehat T[\operatorname{srx}] = \frac{\operatorname{srx}'} {\operatorname{srx}} = -\frac{1+\operatorname{srx}^2} {2\operatorname{srx}}.
\]

Therefore,

\[
\widehat T[\operatorname{srx}] = -\frac12 \left( \operatorname{srx} + \frac1{\operatorname{srx}} \right).
\]

Since

\[
\frac1{\operatorname{srx}} = \operatorname{sxp},
\]

we have

\[
\widehat T[\operatorname{srx}] = -\frac12 \left( \operatorname{srx} + \operatorname{sxp} \right).
\]

By Lemma 3.1,

\[
\widehat T[\operatorname{srx}] = -\frac12(2N) = -N.
\]

Thus,

\[
\boxed{ \widehat T[\operatorname{srx}] = -|\csc x| }.
\]

\end{proof}

\PCsubsection{Theorem 3.4 --- Sinexpon normalized response}

\[
\boxed{ \widehat T[\operatorname{sxp}] = N = |\csc x| }.
\]

\begin{proof}
Since

\[
\operatorname{sxp} = \frac1{\operatorname{srx}},
\]

we have

\[
\ln\operatorname{sxp} = -\ln\operatorname{srx}.
\]

Differentiating,

\[
\widehat T[\operatorname{sxp}] = -\widehat T[\operatorname{srx}].
\]

Using Theorem 3.3,

\[
\widehat T[\operatorname{sxp}] = N.
\]

\end{proof}

\PCsubsection{Theorem 3.5 --- Cosinexpon normalized response}

\[
\boxed{ \widehat T[\operatorname{cxp}] = M = |\sec x| }.
\]

\begin{proof}
From Theorem 2.3,

\[
\widehat T[\operatorname{cxp}] = \frac{1+\operatorname{cxp}^2} {2\operatorname{cxp}}.
\]

Hence,

\[
\widehat T[\operatorname{cxp}] = \frac12 \left( \operatorname{cxp} + \frac1{\operatorname{cxp}} \right).
\]

Since

\[
\frac1{\operatorname{cxp}} = \operatorname{crx},
\]

Lemma 3.2 gives

\[
\widehat T[\operatorname{cxp}] = \frac12(2M) = M.
\]

Therefore,

\[
\boxed{ \widehat T[\operatorname{cxp}] = |\sec x| }.
\]

\end{proof}

\PCsubsection{Theorem 3.6 --- Cosininverex normalized response}

\[
\boxed{ \widehat T[\operatorname{crx}] = -M = -|\sec x| }.
\]

\begin{proof}
Because

\[
\operatorname{crx} = \frac1{\operatorname{cxp}},
\]

normalized differentiation reverses sign:

\[
\widehat T[\operatorname{crx}] = -\widehat T[\operatorname{cxp}].
\]

Theorem 3.5 then gives

\[
\widehat T[\operatorname{crx}] = -M.
\]

\end{proof}

\PCsubsection{Corollary 3.7 --- Complete normalized-response table}

\[
\boxed{ \begin{array}{c|c} \text{Operator} & \text{Normalized tangential response} \\ \hline \operatorname{srx} & -N \\ \operatorname{sxp} & +N \\ \operatorname{cxp} & +M \\ \operatorname{crx} & -M \end{array} }
\]


\PCsection{4. Reciprocal Tangential Symmetry}

\PCsubsection{Theorem 4.1 --- Reciprocal reversal}

Let \(f>0\) be differentiable and define

\[
f^\ast=\frac1f.
\]

Then

\[
\boxed{ \widehat T[f^\ast] = -\widehat T[f] }.
\]

\begin{proof}
Differentiate:

\[
(f^\ast)' = -\frac{f'}{f^2}.
\]

Then

\[
\widehat T[f^\ast] = \frac{(f^\ast)'}{f^\ast} = \frac{-f'/f^2}{1/f}.
\]

Therefore,

\[
\widehat T[f^\ast] = -\frac{f'}f = -\widehat T[f].
\]

\end{proof}

\PCsubsection{Corollary 4.2 --- Mirror rate magnitude}

Reciprocal pairing preserves normalized speed:

\[
\boxed{ V_{f^\ast}=V_f }.
\]

It reverses only normalized direction.

\PCsubsection{Corollary 4.3 --- Family rates}

The sine reciprocal pair shares rate magnitude

\[
V_{\sin}=N,
\]

while the cosine reciprocal pair shares rate magnitude

\[
V_{\cos}=M.
\]

Thus reciprocal access creates no independent rate scale.


\PCsection{5. Primary Rate Comparison}

\PCsubsection{Definition 5.1 --- Rate sum}

Define the positive rate sum

\[
\Sigma_T = \widehat T[\operatorname{cxp}] - \widehat T[\operatorname{srx}].
\]

Using Section 3,

\[
\boxed{ \Sigma_T=M+N }.
\]

\PCsubsection{Definition 5.2 --- Rate contrast}

Define the signed rate contrast

\[
\Delta_T = \widehat T[\operatorname{cxp}] + \widehat T[\operatorname{srx}].
\]

Therefore,

\[
\boxed{ \Delta_T=M-N }.
\]

The sum appears here because the two primary normalized directions are opposite.

\PCsubsection{Theorem 5.3 --- Positivity of total rate separation}

For every \(x\in\mathcal D\),

\[
\boxed{ \Sigma_T=M+N>0 }.
\]

Thus the two primary normalized directions never coincide.

\PCsubsection{Theorem 5.4 --- Unique rate balance within each quadrant}

The primary rate magnitudes are equal exactly when

\[
M=N.
\]

This occurs exactly at

\[
\boxed{ x=\frac{(2k+1)\pi}{4}, \qquad k\in\mathbb Z }.
\]

\begin{proof}
The condition

\[
M=N
\]

is

\[
\frac1{|\cos x|} = \frac1{|\sin x|}.
\]

Since both denominators are positive,

\[
|\sin x|=|\cos x|.
\]

Squaring,

\[
\sin^2x=\cos^2x.
\]

Using

\[
\sin^2x+\cos^2x=1,
\]

we obtain

\[
\sin^2x=\cos^2x=\frac12.
\]

This occurs exactly at the odd multiples of \(\pi/4\):

\[
x=\frac{(2k+1)\pi}{4}.
\]

Each open quadrant contains exactly one such point.

\end{proof}

\PCsubsection{Corollary 5.5 --- Rate balance and octant division}

The unique primary rate-balance point in each quadrant is its midpoint and is also the boundary between its two octants.

At this point,

\[
M=N=\sqrt2.
\]

The neighboring dominant primitive operators tie there.


\PCsection{6. Bounded Transfer Coordinates}

\PCsubsection{Definition 6.1 --- Quadrant signs}

On a fixed open quadrant, define

\[
\alpha=\operatorname{sgn}(\sin x), \qquad \beta=\operatorname{sgn}(\cos x).
\]

Both are constant on that quadrant.

Define also

\[
\varepsilon=\alpha\beta = \operatorname{sgn}(\sin x\cos x) = \operatorname{sgn}(\sin 2x).
\]

\PCsubsection{Definition 6.2 --- Transfer coordinates}

Define

\[
\boxed{ p=\alpha\cos x-\beta\sin x }
\]

and

\[
\boxed{ q=\alpha\sin x+\beta\cos x }.
\]

Since

\[
\alpha\sin x=|\sin x|
\]

and

\[
\beta\cos x=|\cos x|,
\]

we have

\[
\boxed{ q=|\sin x|+|\cos x| }.
\]

Therefore,

\[
q>0
\]

throughout every open quadrant.

\PCsubsection{Theorem 6.3 --- Transfer-circle identity}

\[
\boxed{ p^2+q^2=2 }.
\]

\begin{proof}
Expand:

\[
p^2 = (\alpha c-\beta s)^2 = c^2+s^2-2\alpha\beta sc.
\]

Similarly,

\[
q^2 = (\alpha s+\beta c)^2 = s^2+c^2+2\alpha\beta sc.
\]

Adding,

\[
p^2+q^2 = 2(s^2+c^2).
\]

By the Pythagorean identity,

\[
s^2+c^2=1.
\]

Therefore,

\[
\boxed{ p^2+q^2=2 }.
\]

\end{proof}

\PCsubsection{Theorem 6.4 --- Transfer-coordinate differential system}

\[
\boxed{ p'=-q }
\]

and

\[
\boxed{ q'=p }.
\]

\begin{proof}
Because \(\alpha\) and \(\beta\) are constant on the selected open quadrant,

\[
p' = -\alpha\sin x-\beta\cos x.
\]

Thus,

\[
p'=-q.
\]

Likewise,

\[
q' = \alpha\cos x-\beta\sin x = p.
\]

\end{proof}

\PCsubsection{Corollary 6.5 --- Harmonic closure}

\[
\boxed{ p''=-p }
\]

and

\[
\boxed{ q''=-q }.
\]

The transfer coordinates therefore form a closed two-dimensional differential system.

\PCsubsection{Theorem 6.6 --- Quadrant traversal of \(p\)}

On every open quadrant,

\[
-1<p<1,
\]

and

\[
p'<0.
\]

Moreover, under one-sided boundary limits,

\[
p\longrightarrow1
\]

at the beginning of the quadrant, while

\[
p\longrightarrow-1
\]

at its end.

\begin{proof}
Since \(q>0\), Theorem 6.4 gives

\[
p'=-q<0.
\]

Thus \(p\) is strictly decreasing.

At the left boundary of any quadrant, one of \(|\sin x|\) or \(|\cos x|\) approaches \(0\), while the other approaches \(1\). Direct substitution into the signed definition of \(p\) gives the left limit \(1\).

At the right boundary, the roles reverse and the limit is \(-1\).

Because the boundary points are excluded,

\[
-1<p<1
\]

inside the quadrant.


\end{proof}

\PCsection{7. Canonical Transfer Variable}

\PCsubsection{Definition 7.1 --- Bounded transfer fraction}

Define

\[
\boxed{ \lambda = \frac{1-p}{2} }.
\]

\PCsubsection{Theorem 7.2 --- Transfer range and monotonicity}

On every open quadrant,

\[
\boxed{ 0<\lambda<1 }
\]

and

\[
\boxed{ \lambda'>0 }.
\]

Under one-sided boundary limits,

\[
\lambda\longrightarrow0
\]

at the beginning of the quadrant and

\[
\lambda\longrightarrow1
\]

at its end.

\begin{proof}
Since

\[
-1<p<1,
\]

we have

\[
0<1-p<2.
\]

Dividing by \(2\),

\[
0<\lambda<1.
\]

Differentiate:

\[
\lambda' = -\frac{p'}2.
\]

Using \(p'=-q\),

\[
\boxed{ \lambda'=\frac q2>0 }.
\]

The boundary limits follow directly from the limits \(p\to1\) and \(p\to-1\).

\end{proof}

\PCsubsection{Corollary 7.3 --- Transfer acceleration}

\[
\boxed{ \lambda'' = \frac p2 = \frac12-\lambda }.
\]

\begin{proof}
Differentiate

\[
\lambda'=\frac q2.
\]

Then

\[
\lambda'' = \frac{q'}2 = \frac p2.
\]

Since

\[
p=1-2\lambda,
\]

we have

\[
\lambda'' = \frac{1-2\lambda}{2} = \frac12-\lambda.
\]

\end{proof}

\PCsubsection{Corollary 7.4 --- Transfer differential equation}

\[
\boxed{ \lambda''+\lambda=\frac12 }.
\]

This is a phase-domain identity. It is not yet an equation of physical motion.


\PCsection{8. Saw Responses as Complementary Transfer Channels}

\PCsubsection{Theorem 8.1 --- Saw representation by \(p\)}

The two saw responses satisfy

\[
\boxed{ \operatorname{saw}_{r} = \varepsilon\frac{1-p}{2} }
\]

and

\[
\boxed{ \operatorname{saw}_{x} = \varepsilon\frac{1+p}{2} }.
\]

\begin{proof}
From S0.I, the UNA branch forms are

\[
\operatorname{urx} = \frac{1+p}{sc},
\]

and

\[
\operatorname{uxp} = \frac{1-p}{sc}.
\]

Also,

\[
sc=\varepsilon|sc|.
\]

From the proof of the UNA product identity,

\[
1-p^2=2|sc|.
\]

Therefore,

\[
sc = \varepsilon\frac{1-p^2}{2}.
\]

Now,

\[
\operatorname{saw}_{r} = \frac1{\operatorname{urx}} = \frac{sc}{1+p}.
\]

Substitute the expression for \(sc\):

\[
\operatorname{saw}_{r} = \varepsilon \frac{1-p^2}{2(1+p)}.
\]

Since

\[
1-p^2=(1-p)(1+p),
\]

and \(1+p>0\),

\[
\operatorname{saw}_{r} = \varepsilon\frac{1-p}{2}.
\]

Similarly,

\[
\operatorname{saw}_{x} = \frac{sc}{1-p} = \varepsilon\frac{1-p^2}{2(1-p)} = \varepsilon\frac{1+p}{2}.
\]

\end{proof}

\PCsubsection{Corollary 8.2 --- Saw representation by \(\lambda\)}

Since

\[
\lambda=\frac{1-p}{2},
\]

we have

\[
\boxed{ \operatorname{saw}_{r} = \varepsilon\lambda }
\]

and

\[
\boxed{ \operatorname{saw}_{x} = \varepsilon(1-\lambda) }.
\]

\PCsubsection{Corollary 8.3 --- Signed and unsigned partition}

\[
\operatorname{saw}_{r} + \operatorname{saw}_{x} = \varepsilon,
\]


and

\[
|\operatorname{saw}_{r}| + |\operatorname{saw}_{x}| = 1.
\]

The transfer fraction \(\lambda\) gives the exact distribution of the signed Flatwave total between the two saw channels.


\PCsection{9. Tangential Saw Transfer}

\PCsubsection{Theorem 9.1 --- Equal-and-opposite transfer rates}

\[
\boxed{ \operatorname{saw}_{r}' = \varepsilon\frac q2 }
\]

and

\[
\boxed{ \operatorname{saw}_{x}' = -\varepsilon\frac q2 }.
\]

\begin{proof}
From Corollary 8.2,

\[
\operatorname{saw}_{r} = \varepsilon\lambda.
\]

Since \(\varepsilon\) is constant on the quadrant,

\[
\operatorname{saw}_{r}' = \varepsilon\lambda'.
\]

Using

\[
\lambda'=\frac q2,
\]

we get

\[
\operatorname{saw}_{r}' = \varepsilon\frac q2.
\]

Likewise,

\[
\operatorname{saw}_{x} = \varepsilon(1-\lambda),
\]

so

\[
\operatorname{saw}_{x}' = -\varepsilon\lambda' = -\varepsilon\frac q2.
\]

\end{proof}

\PCsubsection{Corollary 9.2 --- Flatwave tangential invariance}

\[
\boxed{ \operatorname{Flatwave}' = \operatorname{saw}_{r}' + \operatorname{saw}_{x}' = 0 }.
\]



The Flatwave is not itself a varying transfer channel within a quadrant. It is the fixed signed total under which the two saw components exchange magnitude.

\PCsubsection{Corollary 9.3 --- Transfer-rate equality}

\[
\boxed{ \left| \operatorname{saw}_{r}' \right| = \left| \operatorname{saw}_{x}' \right| = \frac q2 }.
\]

The two channels transfer at exactly equal magnitude and opposite direction.


\PCsection{10. Boundedness of Saw Transfer}

\PCsubsection{Lemma 10.1 --- Bounds on \(q\)}

On every open quadrant,

\[
\boxed{ 1<q\leq\sqrt2 }.
\]

The upper bound is attained at the quadrant midpoint.

The lower value \(1\) is approached only at excluded quadrant boundaries.

\begin{proof}
Since

\[
q=|\sin x|+|\cos x|,
\]

we have

\[
q^2 = \sin^2x+\cos^2x+2|\sin x\cos x|.
\]

Therefore,

\[
q^2 = 1+2|sc|.
\]

Inside an open quadrant,

\[
|sc|>0,
\]

so

\[
q^2>1.
\]

Also,

\[
2|sc|\leq s^2+c^2=1.
\]

Hence,

\[
q^2\leq2.
\]

Since \(q>0\),

\[
1<q\leq\sqrt2.
\]

Equality \(q=\sqrt2\) occurs exactly when

\[
|\sin x|=|\cos x|=\frac1{\sqrt2},
\]

which is the quadrant midpoint.

\end{proof}

\PCsubsection{Theorem 10.2 --- Bounded saw-transfer speed}

\[
\boxed{ \frac12 < \left| \operatorname{saw}_{r}' \right| = \left| \operatorname{saw}_{x}' \right| \leq \frac1{\sqrt2} }.
\]

\begin{proof}
By Corollary 9.3,

\[
\left| \operatorname{saw}_{r}' \right| = \left| \operatorname{saw}_{x}' \right| = \frac q2.
\]

Applying Lemma 10.1 gives

\[
\frac12 < \frac q2 \leq \frac{\sqrt2}{2}.
\]

Since

\[
\frac{\sqrt2}{2} = \frac1{\sqrt2},
\]

the result follows.

\end{proof}

\PCsubsection{Corollary 10.3 --- Maximum transfer point}

The transfer-rate magnitude reaches its unique maximum at each quadrant midpoint:

\[
x=\frac{(2k+1)\pi}{4}.
\]

At that point,

\[
\lambda=\frac12,
\]

\[
\operatorname{saw}_{r} = \operatorname{saw}_{x} = \frac{\varepsilon}{2},
\]

and

\[
\left| \operatorname{saw}_{r}' \right| = \left| \operatorname{saw}_{x}' \right| = \frac1{\sqrt2}.
\]

Thus equality of the saw components and maximum transfer occur at the same admissible phase point.


\PCsection{11. Saw Curvature}

\PCsubsection{Theorem 11.1 --- Saw second derivatives}

\[
\boxed{ \operatorname{saw}_{r}'' = \varepsilon\frac p2 }
\]

and

\[
\boxed{ \operatorname{saw}_{x}'' = -\varepsilon\frac p2 }.
\]

\begin{proof}
Differentiate the first-derivative forms:

\[
\operatorname{saw}_{r}' = \varepsilon\frac q2.
\]

Therefore,

\[
\operatorname{saw}_{r}'' = \varepsilon\frac{q'}2.
\]

Since \(q'=p\),

\[
\operatorname{saw}_{r}'' = \varepsilon\frac p2.
\]

The second identity follows similarly.

\end{proof}

\PCsubsection{Corollary 11.2 --- Saw differential equations}

Both saw responses satisfy

\[
\boxed{ y''+y=\frac{\varepsilon}{2} }.
\]

\begin{proof}[Proof for \(\operatorname{saw}_{r}\)]
From

\[
\operatorname{saw}_{r} = \varepsilon\frac{1-p}{2}
\]

and

\[
\operatorname{saw}_{r}'' = \varepsilon\frac p2,
\]

we obtain

\[
\operatorname{saw}_{r}'' + \operatorname{saw}_{r} = \varepsilon\frac p2 + \varepsilon\frac{1-p}{2} = \frac{\varepsilon}{2}.
\]


The proof for \(\operatorname{saw}_{x}\) is identical.

\end{proof}

\PCsubsection{Theorem 11.3 --- Unique inflection point}

Each saw response has exactly one inflection point in every open quadrant, located at

\[
\boxed{ x=\frac{(2k+1)\pi}{4} }.
\]

\begin{proof}
An inflection requires

\[
\operatorname{saw}_{r}''=0
\]

or equivalently

\[
p=0.
\]

Since \(p\) decreases strictly from \(1\) to \(-1\), it has exactly one zero in the quadrant.

The condition \(p=0\) is equivalent to

\[
\alpha\cos x=\beta\sin x.
\]

Taking magnitudes,

\[
|\cos x|=|\sin x|.
\]

This occurs at the quadrant midpoint

\[
x=\frac{(2k+1)\pi}{4}.
\]

\end{proof}

\PCsubsection{Corollary 11.4 --- Octant-local curvature}

Within each open octant, \(p\) has fixed sign. Therefore, each saw response has fixed curvature sign on every open octant.

The phrase “quasi-linear” may be used only if it is explicitly defined as this octant-local monotonicity with fixed curvature sign. It does not mean that the saw functions are linear.


\PCsection{12. UNA Operators in Transfer Coordinates}

\PCsubsection{Theorem 12.1 --- Canonical UNA transfer forms}

\[
\boxed{ \operatorname{urx} = \frac{\varepsilon}{\lambda} }
\]

and

\[
\boxed{ \operatorname{uxp} = \frac{\varepsilon}{1-\lambda} }.
\]

\begin{proof}
Since

\[
\operatorname{saw}_{r} = \frac1{\operatorname{urx}} = \varepsilon\lambda,
\]

taking reciprocals gives

\[
\operatorname{urx} = \frac1{\varepsilon\lambda}.
\]

Because

\[
\varepsilon^2=1,
\]

we have

\[
\frac1\varepsilon=\varepsilon.
\]

Thus,

\[
\operatorname{urx} = \frac{\varepsilon}{\lambda}.
\]

The proof for \(\operatorname{uxp}\) is identical.

\end{proof}

\PCsubsection{Corollary 12.2 --- UNA sum}

\[
\operatorname{urx} + \operatorname{uxp} = \varepsilon \left( \frac1\lambda + \frac1{1-\lambda} \right).
\]

Therefore,

\[
\boxed{ \operatorname{urx} + \operatorname{uxp} = \frac{\varepsilon} {\lambda(1-\lambda)} }.
\]

\PCsubsection{Lemma 12.3 --- Transfer product relation}

\[
\boxed{ \lambda(1-\lambda) = \frac{|\sin 2x|}{4} }.
\]

\begin{proof}
Using

\[
\lambda=\frac{1-p}{2},
\]

we obtain

\[
\lambda(1-\lambda) = \frac{1-p}{2} \frac{1+p}{2}.
\]

Therefore,

\[
\lambda(1-\lambda) = \frac{1-p^2}{4}.
\]

From S0.I,

\[
1-p^2=2|sc|.
\]

Thus,

\[
\lambda(1-\lambda) = \frac{|sc|}{2}.
\]

Since

\[
|\sin 2x|=2|sc|,
\]

we conclude that

\[
\lambda(1-\lambda) = \frac{|\sin 2x|}{4}.
\]

\end{proof}

\PCsubsection{Corollary 12.4 --- Recovery of double-angle closure}

Since

\[
\varepsilon = \operatorname{sgn}(\sin 2x),
\]

we have

\[
\frac{\varepsilon}{|\sin 2x|} = \frac1{\sin 2x}.
\]

Therefore,

\[
\operatorname{urx} + \operatorname{uxp} = \frac{4}{\sin 2x}.
\]

The double-angle theorem is thus equivalent to the bounded transfer-product relation.


\PCsection{13. UNA Tangential Responses}

\PCsubsection{Theorem 13.1 --- Unainverex phase derivative}

\[
\boxed{ \operatorname{urx}' = -\varepsilon \frac{q}{2\lambda^2} }.
\]

\begin{proof}
From Theorem 12.1,

\[
\operatorname{urx} = \varepsilon\lambda^{-1}.
\]

Differentiate:

\[
\operatorname{urx}' = -\varepsilon \lambda^{-2}\lambda'.
\]

Since

\[
\lambda'=\frac q2,
\]

we obtain

\[
\operatorname{urx}' = -\varepsilon \frac{q}{2\lambda^2}.
\]

\end{proof}

\PCsubsection{Theorem 13.2 --- Unaexpon phase derivative}

\[
\boxed{ \operatorname{uxp}' = \varepsilon \frac{q}{2(1-\lambda)^2} }.
\]

\begin{proof}
From

\[
\operatorname{uxp} = \varepsilon(1-\lambda)^{-1},
\]

differentiate:

\[
\operatorname{uxp}' = \varepsilon (1-\lambda)^{-2}\lambda'.
\]

Using \(\lambda'=q/2\) gives the result.

\end{proof}

\PCsubsection{Corollary 13.3 --- UNA derivative signs}

\[
\boxed{ \operatorname{sgn}(\operatorname{urx}') = -\varepsilon }
\]

and

\[
\boxed{ \operatorname{sgn}(\operatorname{uxp}') = \varepsilon }.
\]

The two UNA operators change in opposite phase directions within each quadrant.

\PCsubsection{Theorem 13.4 --- Normalized UNA responses}

\[
\boxed{ \widehat T[\operatorname{urx}] = -\frac{q}{2\lambda} }
\]

and

\[
\boxed{ \widehat T[\operatorname{uxp}] = \frac{q}{2(1-\lambda)} }.
\]

\begin{proof}
Divide Theorem 13.1 by

\[
\operatorname{urx} = \frac{\varepsilon}{\lambda}.
\]

Then

\[
\widehat T[\operatorname{urx}] = \frac{ -\varepsilon q/(2\lambda^2) }{ \varepsilon/\lambda } = -\frac{q}{2\lambda}.
\]

The proof for \(\operatorname{uxp}\) is identical.

\end{proof}

\PCsubsection{Corollary 13.5 --- UNA rate balance}

The normalized UNA rate magnitudes are equal exactly when

\[
\lambda=\frac12.
\]

This occurs at the quadrant midpoint, where both UNA magnitudes equal \(2\) and both saw components equal \(\varepsilon/2\).


\PCsection{14. Differential Closure}

\PCsubsection{Definition 14.1 --- Generator derivation}

On a fixed open quadrant, define the derivation

\[
\boxed{ \mathcal D_z = -\frac{1+z^2}{2}\frac{d}{dz} }.
\]

For every rational function \(F(z)\),

\[
\frac{d}{dx}F(z(x)) = \mathcal D_zF.
\]

\PCsubsection{Theorem 14.2 --- Rational tangential closure}

If

\[
F(z)\in\mathbb R(z),
\]

then

\[
T[F]\in\mathbb R(z).
\]

This result was established in S0.II and applies to all named primitive, UNA, saw, and Flatwave expressions.

\PCsubsection{Definition 14.3 --- Transfer vector space}

Define

\[
\mathcal V = \operatorname{span}_{\mathbb R}\{1,p,q\}.
\]

\PCsubsection{Theorem 14.4 --- Finite transfer closure}

The vector space \(\mathcal V\) is closed under phase differentiation.

\begin{proof}
The derivatives of its basis elements are

\[
1'=0,
\]

\[
p'=-q,
\]

and

\[
q'=p.
\]

Each derivative belongs to

\[
\operatorname{span}_{\mathbb R}\{1,p,q\}.
\]

Therefore,

\[
T[\mathcal V]\subseteq\mathcal V.
\]

\end{proof}

\PCsubsection{Corollary 14.5 --- Saw derivative closure}

Because both saw responses are affine functions of \(p\), every derivative of either saw response belongs to \(\mathcal V\).

Their derivative sequence cycles through \(p\) and \(q\), while their constant midpoint is fixed by \(\varepsilon/2\).


\PCsection{15. Structural Interpretation}

The exact mathematical results permit the following restricted statements.

\PCsubsection{15.1 Primitive envelopes}

The four primitive operators have unbounded phase slopes near their inherited singular boundaries. Their normalized rates are exactly

\[
\pm|\csc x| \quad\text{or}\quad \pm|\sec x|.
\]

\PCsubsection{15.2 Reciprocal access}

Reciprocal pairing reverses normalized tangential direction while preserving normalized rate magnitude.

\PCsubsection{15.3 Flatwave}

Flatwave has no interior phase slope:

\[
\operatorname{Flatwave}'=0.
\]

It is a signed quadrant discriminator, not a continuously varying rate channel.

\PCsubsection{15.4 Saw transfer}

The two saw components provide the bounded transfer geometry:

\[
\operatorname{saw}_{r} = \varepsilon\lambda, \qquad \operatorname{saw}_{x} = \varepsilon(1-\lambda).
\]

Their transfer is equal, opposite, monotone, and bounded.

\PCsubsection{15.5 Octant structure}

The saw curvature changes sign only at the quadrant midpoint. Therefore, curvature sign is fixed within each octant.

These are structural statements about phase differentiation. Physical interpretation remains deferred.


\PCsection{16. Terminal Statement}

S0.I derived the operator identities from the Pythagorean relation.

S0.II reduced the named calculus to one rational generator.

S0.III establishes the tangential geometry.

The primitive normalized responses are

\[
\boxed{ \widehat T[\operatorname{srx}] = -|\csc x| }
\]

\[
\boxed{ \widehat T[\operatorname{sxp}] = |\csc x| }
\]

\[
\boxed{ \widehat T[\operatorname{cxp}] = |\sec x| }
\]

\[
\boxed{ \widehat T[\operatorname{crx}] = -|\sec x| }.
\]

The bounded transfer coordinate satisfies

\[
\boxed{ 0<\lambda<1, \qquad \lambda'>0, \qquad \lambda''+\lambda=\frac12 }.
\]

The saw channels satisfy

\[
\boxed{ \operatorname{saw}_{r} = \varepsilon\lambda }
\]

and

\[
\boxed{ \operatorname{saw}_{x} = \varepsilon(1-\lambda) }.
\]

Their rates obey

\[
\boxed{ \operatorname{saw}_{r}' = -\operatorname{saw}_{x}' }
\]

and

\[
\boxed{ \frac12 < \left| \operatorname{saw}_{r}' \right| = \left| \operatorname{saw}_{x}' \right| \leq \frac1{\sqrt2} }.
\]

Their sum remains fixed:

\[
\boxed{ \operatorname{Flatwave}'=0 }.
\]

Thus the tangential core is:

\[
\boxed{ \text{unbounded primitive envelopes, reciprocal rate reversal, and bounded complementary transfer.} }
\]

% ===== END S0_III =====
% ===== BEGIN S0_IV =====
% Source: S0.iv.docx
\PCchapter{S0.IV}{Phase Symmetry, Operator Orbits, and Octant Rotation}{Book 0 - Formal Core}

\PCsubsection{Abstract}

This essay determines the discrete symmetry structure of the bounded phase operators.

Two transformations are introduced on the quotient phase circle

\[
\mathbb T=\mathbb R/(2\pi\mathbb Z):
\]

\[
R(x)=x+\frac{\pi}{2},
\]

representing a quarter-turn, and

\[
J(x)=-x,
\]

representing reflection through the phase origin.

These transformations generate the dihedral group of the four-quadrant phase wheel. Throughout this corpus, \(D_4\) denotes the order-eight group (called \(D_8\) in some alternate conventions):

\[
D_4 = \langle R,J\mid R^4=J^2=e,\ JRJ=R^{-1}\rangle.
\]

Their action on the four primitive operators is derived exactly. A quarter-turn exchanges the sine and cosine families while preserving reciprocal position:

\[
\operatorname{srx}\leftrightarrow\operatorname{crx}, \qquad \operatorname{sxp}\leftrightarrow\operatorname{cxp}.
\]

Phase reflection realizes reciprocal inversion:

\[
\operatorname{srx}\leftrightarrow\operatorname{sxp}, \qquad \operatorname{cxp}\leftrightarrow\operatorname{crx}.
\]

The resulting permutation representation is not faithful. A half-turn acts trivially on every primitive operator because the complete operator family is \(\pi\)-periodic. The eight phase symmetries therefore reduce to a four-element operator symmetry group isomorphic to the Klein group:

\[
D_4/\{e,R^2\}\cong V_4.
\]

The actions on the UNA operators, saw responses, Flatwave, and bounded transfer coordinate are derived. Flatwave transforms as a sign representation: quarter-turn and phase reflection reverse its sign, while their composition preserves it.

Finally, the eighth-turn

\[
H(x)=x+\frac{\pi}{4}
\]

is shown to generate a cyclic action on the eight octants. It advances the dominant-operator label through

\[
\operatorname{srx} \rightarrow \operatorname{cxp} \rightarrow \operatorname{crx} \rightarrow \operatorname{sxp} \rightarrow \operatorname{srx}.
\]

This octant action is combinatorial rather than an exact equality among shifted operator functions. The distinction between exact operator symmetry and dominance-label rotation is made explicit.

No physical, temporal, or observer interpretation is introduced.


\PCsection{0. Scope and Inheritance}

\PCsubsection{0.1 Scope}

This essay proves:

\begin{enumerate}
\item the exact quarter-turn transformation of the primitive operators;
\item the exact reflection transformation of the primitive operators;
\item the dihedral relations satisfied by the phase transformations;
\item the induced operator-permutation group;
\item the \(\pi\)-periodicity and half-turn degeneracy of the operator family;
\item transformation laws for UNA, saw, and Flatwave expressions;
\item transformation laws for the bounded transfer variables;
\item the cyclic octant action on dominant-operator labels;
\item the distinction between exact function symmetry and dominance rotation.
\end{enumerate}
\PCsubsection{0.2 Non-scope}

This essay does not:

\begin{itemize}
\item interpret reflection as a physical mirror universe;
\item identify phase rotation with time evolution;
\item introduce observers or reference frames;
\item claim that an eighth-turn permutes the operator functions exactly;
\item assign physical meaning to Flatwave sign;
\item extend any function through an excluded boundary.
\end{itemize}
\PCsubsection{0.3 Inheritance}

All definitions and theorems of S0.I–S0.III are inherited.

The admissible domain is

\[
\mathcal D = \mathbb R \setminus \left\{ \frac{k\pi}{2}:k\in\mathbb Z \right\}.
\]

The primitive operators are

\[
\operatorname{srx}(x)=|\csc x|+\cot x,
\]

\[
\operatorname{sxp}(x)=|\csc x|-\cot x,
\]

\[
\operatorname{cxp}(x)=|\sec x|+\tan x,
\]

\[
\operatorname{crx}(x)=|\sec x|-\tan x.
\]

The UNA operators are

\[
\operatorname{urx} = \operatorname{srx}-\operatorname{crx},
\]

and

\[
\operatorname{uxp} = \operatorname{cxp}-\operatorname{sxp}.
\]

The saw responses are

\[
\operatorname{saw}_{r} = \frac1{\operatorname{urx}}, \qquad \operatorname{saw}_{x} = \frac1{\operatorname{uxp}}.
\]

Flatwave is

\[
\operatorname{Flatwave} = \operatorname{saw}_{r}+\operatorname{saw}_{x} = \operatorname{sgn}(\sin 2x).
\]



\PCsection{1. Elementary Rotation Formulas}

\PCsubsection{Inherited theorem 1.1 --- Angle addition}

For real \(x\) and \(y\),

\[
\sin(x+y) = \sin x\cos y+\cos x\sin y,
\]

and

\[
\cos(x+y) = \cos x\cos y-\sin x\sin y.
\]

\PCsubsection{Lemma 1.2 --- Quarter-turn formulas}

For every admissible \(x\),

\[
\boxed{ \sin\left(x+\frac{\pi}{2}\right)=\cos x }
\]

and

\[
\boxed{ \cos\left(x+\frac{\pi}{2}\right)=-\sin x }.
\]

\begin{proof}
Using

\[
\cos\frac{\pi}{2}=0, \qquad \sin\frac{\pi}{2}=1,
\]

the addition formulas give

\[
\sin\left(x+\frac{\pi}{2}\right) = \sin x\cdot0+\cos x\cdot1 = \cos x,
\]

and

\[
\cos\left(x+\frac{\pi}{2}\right) = \cos x\cdot0-\sin x\cdot1 = -\sin x.
\]

\end{proof}

\PCsubsection{Corollary 1.3 --- Quarter-turn quotient formulas}

\[
\boxed{ \tan\left(x+\frac{\pi}{2}\right) = -\cot x }
\]

and

\[
\boxed{ \cot\left(x+\frac{\pi}{2}\right) = -\tan x }.
\]

\begin{proof}
Since \(x\in\mathcal D\),

\[
\tan\left(x+\frac{\pi}{2}\right) = \frac{\cos x}{-\sin x} = -\cot x.
\]

The second identity follows by reciprocal exchange.

\end{proof}

\PCsubsection{Corollary 1.4 --- Quarter-turn absolute reciprocal formulas}

\[
\boxed{ \left| \csc\left(x+\frac{\pi}{2}\right) \right| = |\sec x| }
\]

and

\[
\boxed{ \left| \sec\left(x+\frac{\pi}{2}\right) \right| = |\csc x| }.
\]


\PCsection{2. Quarter-Turn Operator Symmetry}

\PCsubsection{Definition 2.1 --- Quarter-turn transformation}

Define the quotient action

\[
\boxed{ R([x])=\left[x+\frac{\pi}{2}\right] }
\]

on \(\mathbb T\). In formulas involving the periodic operators, the shorter notation \(R x=x+\pi/2\pmod{2\pi}\) is retained.

\PCsubsection{Theorem 2.2 --- Quarter-turn action on Sininverex}

\[
\boxed{ \operatorname{srx}(Rx) = \operatorname{crx}(x) }.
\]

\begin{proof}
By definition,

\[
\operatorname{srx}(Rx) = \left| \csc\left(x+\frac{\pi}{2}\right) \right| + \cot\left(x+\frac{\pi}{2}\right).
\]

Using Corollaries 1.3–1.4,

\[
\operatorname{srx}(Rx) = |\sec x|-\tan x.
\]

But

\[
|\sec x|-\tan x = \operatorname{crx}(x).
\]

Therefore,

\[
\operatorname{srx}(Rx) = \operatorname{crx}(x).
\]

\end{proof}

\PCsubsection{Theorem 2.3 --- Complete quarter-turn action}

\[
\boxed{ \begin{aligned} \operatorname{srx}(Rx)&=\operatorname{crx}(x),\\ \operatorname{sxp}(Rx)&=\operatorname{cxp}(x),\\ \operatorname{cxp}(Rx)&=\operatorname{sxp}(x),\\ \operatorname{crx}(Rx)&=\operatorname{srx}(x). \end{aligned} }
\]

\begin{proof}
The first identity is Theorem 2.2.

For Sinexpon,

\[
\operatorname{sxp}(Rx) = \left|\csc\left(x+\frac{\pi}{2}\right)\right| - \cot\left(x+\frac{\pi}{2}\right).
\]


Thus,

\[
\operatorname{sxp}(Rx) = |\sec x|+\tan x = \operatorname{cxp}(x).
\]

Similarly,

\[
\operatorname{cxp}(Rx) = \left| \sec\left(x+\frac{\pi}{2}\right) \right| + \tan\left(x+\frac{\pi}{2}\right)
\]

becomes

\[
\operatorname{cxp}(Rx) = |\csc x|-\cot x = \operatorname{sxp}(x).
\]

Finally,

\[
\operatorname{crx}(Rx) = |\csc x|+\cot x = \operatorname{srx}(x).
\]

\end{proof}

\PCsubsection{Corollary 2.4 --- Quarter-turn permutation}

On the ordered primitive set

\[
\mathcal O = \{\operatorname{srx},\operatorname{sxp},\operatorname{cxp},\operatorname{crx}\}.
\]


the quarter-turn induces the permutation

\[
\boxed{ \rho(R) = (\operatorname{srx}\ \operatorname{crx}) (\operatorname{sxp}\ \operatorname{cxp}) }.
\]

This induced permutation has order \(2\), even though the phase transformation \(R\) has order \(4\).


\PCsection{3. Half-Turn Invariance}

\PCsubsection{Lemma 3.1 --- Half-turn trigonometric formulas}

\[
\sin(x+\pi)=-\sin x,
\]

\[
\cos(x+\pi)=-\cos x,
\]

\[
\tan(x+\pi)=\tan x,
\]

and

\[
\cot(x+\pi)=\cot x.
\]

Furthermore,

\[
|\csc(x+\pi)|=|\csc x|,
\]

and

\[
|\sec(x+\pi)|=|\sec x|.
\]

\PCsubsection{Theorem 3.2 --- Primitive half-turn invariance}

For every primitive operator \(f\in\mathcal O\),

\[
\boxed{ f(x+\pi)=f(x). }
\]

\begin{proof}
Each primitive consists of:

\begin{itemize}
\item an absolute reciprocal function, which is \(\pi\)-periodic; and
\item either tangent or cotangent, which is also \(\pi\)-periodic.
\end{itemize}
Therefore each complete primitive is \(\pi\)-periodic.

\end{proof}

\PCsubsection{Corollary 3.3 --- Operator period}

The displayed operator family has fundamental phase repetition no greater than \(\pi\):

\[
\boxed{ \mathcal O(x+\pi)=\mathcal O(x). }
\]

\PCsubsection{Corollary 3.4 --- Half-turn kernel}

Although

\[
R^2(x)=x+\pi
\]

is not the identity transformation on the \(2\pi\)-phase circle, it acts trivially on every primitive operator:

\[
\boxed{ \rho(R^2)=e. }
\]

The operator representation therefore does not distinguish phase positions separated by \(\pi\).


\PCsection{4. Reflection and Reciprocal Inversion}

\PCsubsection{Definition 4.1 --- Phase reflection}

Define the quotient action

\[
\boxed{ J([x])=[-x] }
\]

on \(\mathbb T\). The shorter notation \(Jx=-x\pmod{2\pi}\) is retained in operator formulas.

\PCsubsection{Lemma 4.2 --- Reflection formulas}

\[
\sin(-x)=-\sin x, \qquad \cos(-x)=\cos x,
\]

\[
\tan(-x)=-\tan x, \qquad \cot(-x)=-\cot x.
\]

Absolute reciprocal magnitudes satisfy

\[
|\csc(-x)|=|\csc x|,
\]

and

\[
|\sec(-x)|=|\sec x|.
\]

\PCsubsection{Theorem 4.3 --- Reflection action on the primitive operators}

\[
\boxed{ \begin{aligned} \operatorname{srx}(Jx)&=\operatorname{sxp}(x),\\ \operatorname{sxp}(Jx)&=\operatorname{srx}(x),\\ \operatorname{cxp}(Jx)&=\operatorname{crx}(x),\\ \operatorname{crx}(Jx)&=\operatorname{cxp}(x). \end{aligned} }
\]

\begin{proof}
For Sininverex,

\[
\operatorname{srx}(-x) = |\csc(-x)|+\cot(-x).
\]

Therefore,

\[
\operatorname{srx}(-x) = |\csc x|-\cot x = \operatorname{sxp}(x).
\]

The remaining identities follow by the same substitution.

\end{proof}

\PCsubsection{Corollary 4.4 --- Reflection permutation}

Reflection induces

\[
\boxed{ \rho(J) = (\operatorname{srx}\ \operatorname{sxp}) (\operatorname{cxp}\ \operatorname{crx}) }.
\]

\PCsubsection{Corollary 4.5 --- Reciprocal inversion is geometrically realized}

Since

\[
\operatorname{sxp} = \frac1{\operatorname{srx}}
\]

and

\[
\operatorname{crx} = \frac1{\operatorname{cxp}},
\]

phase reflection realizes reciprocal exchange:

\[
\boxed{ f(-x)=\frac1{f(x)} }
\]

for each primary envelope \(f\in\{\operatorname{srx},\operatorname{cxp}\}\).

This statement applies to the paired primary envelopes. It is not a general rule for every composite expression.


\PCsection{5. The Phase Symmetry Group}

\PCsubsection{Theorem 5.1 --- Dihedral relations}

The phase transformations \(R\) and \(J\) satisfy

\[
\boxed{ R^4=e, \qquad J^2=e, \qquad JRJ=R^{-1}. }
\]

\begin{proof}
Four quarter-turns give

\[
R^4(x) = x+4\frac{\pi}{2} = x+2\pi \equiv x.
\]


Thus,

\[
R^4=e.
\]

Two reflections give

\[
J^2(x) = -(-x) = x,
\]


so

\[
J^2=e.
\]

Finally,

\[
JRJ(x) = J\left(R(-x)\right).
\]


Now,

\[
R(-x) = -x+\frac{\pi}{2}.
\]


Applying \(J\),

\[
J\left(-x+\frac{\pi}{2}\right) = x-\frac{\pi}{2}.
\]


But

\[
x-\frac{\pi}{2}=R^{-1}(x).
\]


Therefore,

\[
JRJ=R^{-1}.
\]

\end{proof}

\PCsubsection{Corollary 5.2 --- Quadrant phase group}

The transformations generated by \(R\) and \(J\) form the dihedral group

\[
\boxed{ D_4 = \langle R,J\mid R^4=J^2=e,\ JRJ=R^{-1}\rangle }.
\]


Under the convention fixed above, it contains eight phase transformations:

\[
e,\ R,\ R^2,\ R^3,\ J,\ RJ,\ R^2J,\ R^3J.
\]


\PCsection{6. The Induced Operator Group}

\PCsubsection{Definition 6.1 --- Operator representation}

Let

\[
\rho:D_4\longrightarrow\operatorname{Sym}(\mathcal O)
\]

assign to each phase symmetry its induced permutation of the primitive operator set.

\PCsubsection{Lemma 6.2 --- Generator images}

\[
\rho(R) = (\operatorname{srx}\ \operatorname{crx}) (\operatorname{sxp}\ \operatorname{cxp}),
\]


and

\[
\rho(J) = (\operatorname{srx}\ \operatorname{sxp}) (\operatorname{cxp}\ \operatorname{crx}).
\]


Both induced permutations have order \(2\).

\PCsubsection{Theorem 6.3 --- Induced operator group}

The image of \(\rho\) is

\[
\boxed{ \operatorname{im}\rho = \{e,Q,I,QI\} }.
\]

where

\[
Q = (\operatorname{srx}\ \operatorname{crx}) (\operatorname{sxp}\ \operatorname{cxp}),
\]

\[
I = (\operatorname{srx}\ \operatorname{sxp}) (\operatorname{cxp}\ \operatorname{crx}),
\]

and

\[
QI = (\operatorname{srx}\ \operatorname{cxp}) (\operatorname{sxp}\ \operatorname{crx}).
\]


These three nonidentity permutations satisfy

\[
Q^2=I^2=(QI)^2=e
\]

and commute.

Therefore,

\[
\boxed{ \operatorname{im}\rho\cong V_4, }
\]

the Klein four-group.

\begin{proof}
The two generator images \(Q\) and \(I\) are distinct commuting involutions.

Their product is

\[
QI = (\operatorname{srx}\ \operatorname{cxp}) (\operatorname{sxp}\ \operatorname{crx}).
\]

Thus the generated group contains exactly

\[
e,\ Q,\ I,\ QI.
\]

Every nonidentity element has order \(2\). This is the defining structure of \(V_4\).

\end{proof}

\PCsubsection{Theorem 6.4 --- Representation kernel}

\[
\boxed{ \ker\rho=\{e,R^2\}. }
\]

\begin{proof}
By half-turn invariance,

\[
\rho(R^2)=e.
\]

Thus,

\[
\{e,R^2\} \subseteq \ker\rho.
\]

The phase group has eight elements, while its image has four. By the homomorphism counting relation,

\[
|\ker\rho| = \frac{|D_4|}{|\operatorname{im}\rho|} = \frac84 = 2.
\]


Therefore the kernel contains exactly the two already identified elements:

\[
\ker\rho = \{e,R^2\}.
\]

\end{proof}

\PCsubsection{Corollary 6.5 --- Quotient symmetry}

\[
\boxed{ D_4/\{e,R^2\} \cong V_4. }
\]

The primitive operators retain only half of the full phase-wheel symmetry information.


\PCsection{7. The Third Primitive Symmetry}

\PCsubsection{Definition 7.1 --- Diagonal reflection}

Define the composition

\[
K=RJ,
\]

with \(J\) applied first.

Then

\[
\boxed{ K(x)=\frac{\pi}{2}-x }.
\]


\PCsubsection{Theorem 7.2 --- Diagonal-reflection action}

\[
\boxed{ \begin{aligned} \operatorname{srx}(Kx)&=\operatorname{cxp}(x),\\ \operatorname{sxp}(Kx)&=\operatorname{crx}(x),\\ \operatorname{cxp}(Kx)&=\operatorname{srx}(x),\\ \operatorname{crx}(Kx)&=\operatorname{sxp}(x). \end{aligned} }
\]

\begin{proof}
The induced permutation is

\[
\rho(K) = \rho(R)\rho(J) = QI.
\]

The explicit permutation \(QI\) exchanges

\[
\operatorname{srx}\leftrightarrow\operatorname{cxp}
\]

and

\[
\operatorname{sxp}\leftrightarrow\operatorname{crx}.
\]

The same result follows by direct substitution using

\[
\sin\left(\frac{\pi}{2}-x\right)=\cos x
\]

and

\[
\cos\left(\frac{\pi}{2}-x\right)=\sin x.
\]

\end{proof}

\PCsubsection{Corollary 7.3 --- Operator orbit}

The orbit of any primitive operator under the phase symmetry group contains all four primitive operators.

For example,

\[
\operatorname{srx} \overset{J}{\longmapsto} \operatorname{sxp},
\]

\[
\operatorname{srx} \overset{R}{\longmapsto} \operatorname{crx},
\]

and

\[
\operatorname{srx} \overset{K}{\longmapsto} \operatorname{cxp}.
\]

Thus the primitive family is one symmetry orbit.


\PCsection{8. Transformation of the UNA Operators}

\PCsubsection{Theorem 8.1 --- Quarter-turn UNA reversal}

\[
\boxed{ \operatorname{urx}(Rx) = -\operatorname{urx}(x) }
\]

and

\[
\boxed{ \operatorname{uxp}(Rx) = -\operatorname{uxp}(x). }
\]

\begin{proof}
Using the primitive quarter-turn transformations,

\[
\operatorname{urx}(Rx) = \operatorname{srx}(Rx) - \operatorname{crx}(Rx).
\]

Therefore,

\[
\operatorname{urx}(Rx) = \operatorname{crx}(x) - \operatorname{srx}(x) = -\operatorname{urx}(x).
\]

Likewise,

\[
\operatorname{uxp}(Rx) = \operatorname{sxp}(x) - \operatorname{cxp}(x) = -\operatorname{uxp}(x).
\]

\end{proof}

\PCsubsection{Theorem 8.2 --- Reflection exchanges the UNA operators with sign reversal}

\[
\boxed{ \operatorname{urx}(Jx) = -\operatorname{uxp}(x) }
\]

and

\[
\boxed{ \operatorname{uxp}(Jx) = -\operatorname{urx}(x). }
\]

\begin{proof}
Under reflection,

\[
\operatorname{srx}\mapsto\operatorname{sxp}, \qquad \operatorname{crx}\mapsto\operatorname{cxp}.
\]

Therefore,

\[
\operatorname{urx}(Jx) = \operatorname{sxp}-\operatorname{cxp} = -\operatorname{uxp}.
\]

The second identity follows similarly.

\end{proof}

\PCsubsection{Theorem 8.3 --- Diagonal reflection exchanges the UNA operators}

\[
\boxed{ \operatorname{urx}(Kx) = \operatorname{uxp}(x) }
\]

and

\[
\boxed{ \operatorname{uxp}(Kx) = \operatorname{urx}(x). }
\]

There is no sign reversal under \(K\).


\PCsection{9. Transformation of the Saw Responses}

\PCsubsection{Theorem 9.1 --- Quarter-turn saw reversal}

\[
\boxed{ \operatorname{saw}_{r}(Rx) = -\operatorname{saw}_{r}(x) }
\]

and

\[
\boxed{ \operatorname{saw}_{x}(Rx) = -\operatorname{saw}_{x}(x). }
\]

\begin{proof}
Quarter-turn negates both UNA denominators. Taking reciprocals preserves that negative factor.

\end{proof}

\PCsubsection{Theorem 9.2 --- Reflection exchanges the saw channels with sign reversal}

\[
\boxed{ \operatorname{saw}_{r}(Jx) = -\operatorname{saw}_{x}(x) }
\]

and

\[
\boxed{ \operatorname{saw}_{x}(Jx) = -\operatorname{saw}_{r}(x). }
\]

\PCsubsection{Theorem 9.3 --- Diagonal reflection exchanges the saw channels}

\[
\boxed{ \operatorname{saw}_{r}(Kx) = \operatorname{saw}_{x}(x) }
\]

and

\[
\boxed{ \operatorname{saw}_{x}(Kx) = \operatorname{saw}_{r}(x). }
\]


\PCsection{10. Flatwave as a Sign Representation}

\PCsubsection{Theorem 10.1 --- Quarter-turn Flatwave reversal}

\[
\boxed{ \operatorname{Flatwave}(Rx) = -\operatorname{Flatwave}(x). }
\]

\begin{proof}
Since

\[
\sin\left(2x+\pi\right) = -\sin 2x,
\]

we have

\[
\operatorname{sgn} \left( \sin 2Rx \right) = -\operatorname{sgn}(\sin 2x).
\]

The result follows from the exact Flatwave theorem.

\end{proof}

\PCsubsection{Theorem 10.2 --- Reflection Flatwave reversal}

\[
\boxed{ \operatorname{Flatwave}(Jx) = -\operatorname{Flatwave}(x). }
\]

\begin{proof}
\[
\sin(-2x)=-\sin 2x.
\]

Therefore its sign reverses.

\end{proof}

\PCsubsection{Theorem 10.3 --- Diagonal-reflection invariance}

\[
\boxed{ \operatorname{Flatwave}(Kx) = \operatorname{Flatwave}(x). }
\]

\begin{proof}
Since

\[
Kx=\frac{\pi}{2}-x,
\]

we have

\[
\sin(2Kx) = \sin(\pi-2x) = \sin 2x.
\]

Therefore Flatwave is unchanged.

\end{proof}

\PCsubsection{Definition 10.4 --- Flatwave character}

Define

\[
\chi:D_4\rightarrow\{-1,+1\}
\]

by

\[
\operatorname{Flatwave}(gx) = \chi(g)\operatorname{Flatwave}(x).
\]

The generators satisfy

\[
\chi(R)=-1, \qquad \chi(J)=-1.
\]

Therefore,

\[
\boxed{ \chi(R^aJ^b)=(-1)^{a+b}. }
\]

\PCsubsection{Corollary 10.5 --- Flatwave-preserving subgroup}

The phase symmetries that preserve Flatwave form the four-element subgroup

\[
\boxed{ \ker\chi = \{e,\ R^2,\ RJ,\ R^3J\}. }
\]

Quarter-turn and origin reflection reverse quadrant orientation; their composition preserves it.


\PCsection{11. Transformation of the Transfer Coordinates}

Recall from S0.III:

\[
\alpha=\operatorname{sgn}(\sin x), \qquad \beta=\operatorname{sgn}(\cos x),
\]

\[
\varepsilon=\alpha\beta,
\]

\[
p=\alpha\cos x-\beta\sin x,
\]

\[
q=\alpha\sin x+\beta\cos x,
\]

and

\[
\lambda=\frac{1-p}{2}.
\]

\PCsubsection{Theorem 11.1 --- Quarter-turn transfer invariance}

Under \(R\),

\[
\boxed{ p(Rx)=p(x), }
\]

\[
\boxed{ q(Rx)=q(x), }
\]

\[
\boxed{ \lambda(Rx)=\lambda(x), }
\]

while

\[
\boxed{ \varepsilon(Rx)=-\varepsilon(x). }
\]

\begin{proof}
Under the quarter-turn,

\[
\sin Rx=\cos x, \qquad \cos Rx=-\sin x.
\]

Therefore,

\[
\alpha_R=\beta, \qquad \beta_R=-\alpha.
\]

Then

\[
p(Rx) = \alpha_R\cos Rx-\beta_R\sin Rx.
\]

Substituting,

\[
p(Rx) = \beta(-\sin x)-(-\alpha)(\cos x).
\]

Thus,

\[
p(Rx) = \alpha\cos x-\beta\sin x = p(x).
\]

Similarly,

\[
q(Rx) = \beta\cos x+(-\alpha)(-\sin x) = q(x).
\]

Because \(\lambda\) depends only on \(p\), it is unchanged.

Finally,

\[
\varepsilon_R = \alpha_R\beta_R = \beta(-\alpha) = -\varepsilon.
\]

\end{proof}

\PCsubsection{Theorem 11.2 --- Reflection transfer complement}

Under \(J\),

\[
\boxed{ p(Jx)=-p(x), }
\]

\[
\boxed{ q(Jx)=q(x), }
\]

\[
\boxed{ \lambda(Jx)=1-\lambda(x), }
\]

and

\[
\boxed{ \varepsilon(Jx)=-\varepsilon(x). }
\]

\begin{proof}
Under reflection,

\[
\sin(-x)=-\sin x, \qquad \cos(-x)=\cos x,
\]

so

\[
\alpha_J=-\alpha, \qquad \beta_J=\beta.
\]

Therefore,

\[
p(Jx) = (-\alpha)\cos x-\beta(-\sin x) = -p(x).
\]

Likewise,

\[
q(Jx) = (-\alpha)(-\sin x)+\beta\cos x = q(x).
\]

Hence,

\[
\lambda(Jx) = \frac{1+p}{2} = 1-\lambda(x).
\]

Also,

\[
\varepsilon_J = (-\alpha)\beta = -\varepsilon.
\]

\end{proof}

\PCsubsection{Corollary 11.3 --- Diagonal-reflection transfer complement}

Under \(K=RJ\),

\[
\boxed{ p(Kx)=-p(x), }
\]

\[
\boxed{ q(Kx)=q(x), }
\]

\[
\boxed{ \lambda(Kx)=1-\lambda(x), }
\]

while

\[
\boxed{ \varepsilon(Kx)=\varepsilon(x). }
\]

Thus \(K\) exchanges the two transfer channels without changing their signed total.


\PCsection{12. Exact Symmetry and Symmetric Invariants}

\PCsubsection{Theorem 12.1 --- Primitive-set invariance}

Every phase symmetry in \(D_4\) permutes the primitive operator set without introducing a new primitive function.

Therefore, the unordered set

\[
\boxed{ \{ \operatorname{srx}, \operatorname{sxp}, \operatorname{cxp}, \operatorname{crx} \} }
\]

is invariant under \(D_4\).

\PCsubsection{Corollary 12.2 --- Symmetric-expression invariance}

Every symmetric rational expression in the four primitive operators is invariant under the induced operator group.

Examples include

\[
\operatorname{srx} + \operatorname{sxp} + \operatorname{cxp} + \operatorname{crx},
\]

and

\[
\operatorname{srx} \operatorname{sxp} \operatorname{cxp} \operatorname{crx}.
\]

\PCsubsection{Corollary 12.3 --- Total primitive product}

\[
\boxed{ \operatorname{srx}\, \operatorname{sxp}\, \operatorname{cxp}\, \operatorname{crx} = 1 }.
\]


This follows from the two reciprocal pairs and is invariant under every phase symmetry.


\PCsection{13. Octant Rotation}

\PCsubsection{Definition 13.1 --- Open octants}

For \(j\in\mathbb Z/8\mathbb Z\), define

\[
O_j = \left( \frac{j\pi}{4}, \frac{(j+1)\pi}{4} \right).
\]

\PCsubsection{Definition 13.2 --- Eighth-turn transformation}

Define

\[
\boxed{ H(x)=x+\frac{\pi}{4} \pmod{2\pi}. }
\]

Then

\[
H(O_j)=O_{j+1}.
\]

\PCsubsection{Theorem 13.3 --- Octant rotation group}

The transformation \(H\) satisfies

\[
H^8=e.
\]

It generates a cyclic action on the eight octants:

\[
\boxed{ \langle H\rangle\cong C_8. }
\]

\begin{proof}
Eight eighth-turns give

\[
H^8(x) = x+8\frac{\pi}{4} = x+2\pi \equiv x.
\]

No smaller positive power fixes every phase position modulo \(2\pi\). Therefore \(H\) has order \(8\).


\end{proof}

\PCsection{14. Dominance-Label Rotation}

\PCsubsection{Definition 14.1 --- Dominant-operator map}

For each open octant, let \(W(O_j)\) denote its unique magnitude-dominant primitive operator.

From S0.I,

\[
\begin{aligned} W(O_0)&=\operatorname{srx},\\ W(O_1)&=\operatorname{cxp},\\ W(O_2)&=\operatorname{crx},\\ W(O_3)&=\operatorname{sxp}, \end{aligned}
\]

with repetition after four octants.

\PCsubsection{Definition 14.2 --- Winner-cycle permutation}

Define

\[
\boxed{ C = (\operatorname{srx}\ \operatorname{cxp}\ \operatorname{crx}\ \operatorname{sxp}) }.
\]


Thus,

\[
C(\operatorname{srx})=\operatorname{cxp},
\]

\[
C(\operatorname{cxp})=\operatorname{crx},
\]

\[
C(\operatorname{crx})=\operatorname{sxp},
\]

and

\[
C(\operatorname{sxp})=\operatorname{srx}.
\]

\PCsubsection{Theorem 14.3 --- Dominance advancement}

For every open octant,

\[
\boxed{ W(HO_j) = C\bigl(W(O_j)\bigr). }
\]

\begin{proof}
Since

\[
H(O_j)=O_{j+1},
\]

the result follows directly from the proved octant sequence:

\[
\operatorname{srx} \rightarrow \operatorname{cxp} \rightarrow \operatorname{crx} \rightarrow \operatorname{sxp} \rightarrow \operatorname{srx}.
\]

\end{proof}

\PCsubsection{Corollary 14.4 --- Dominance quotient}

The eighth-turn phase action has order \(8\), while the winner-cycle permutation has order \(4\):

\[
C^4=e.
\]

The half-turn

\[
H^4(x)=x+\pi
\]

acts trivially on winner labels.

Therefore, the dominance-label action factors through

\[
\boxed{ C_8/\{e,H^4\}\cong C_4. }
\]


\PCsection{15. Exact Operator Symmetry Versus Dominance Rotation}

\PCsubsection{Theorem 15.1 --- Quarter-turn consistency}

A quarter-turn equals two eighth-turns:

\[
R=H^2.
\]

The dominance permutation generated by two octant steps is

\[
C^2 = (\operatorname{srx}\ \operatorname{crx}) (\operatorname{cxp}\ \operatorname{sxp}).
\]

This is exactly the primitive quarter-turn permutation \(\rho(R)\).

Thus exact operator symmetry and dominance-label rotation agree at the quarter-turn scale.

\PCsubsection{Theorem 15.2 --- Eighth-turn is not generally an exact primitive permutation}

In general,

\[
\boxed{ \operatorname{srx}\left(x+\frac{\pi}{4}\right) \neq \operatorname{cxp}(x). }
\]

Therefore, the winner-cycle permutation \(C\) is not an exact function-substitution symmetry.

\begin{proof}
In the first quadrant,

\[
\operatorname{srx}\left(x+\frac{\pi}{4}\right) = \cot\left( \frac x2+\frac{\pi}{8} \right),
\]

while

\[
\operatorname{cxp}(x) = \tan\left( \frac{\pi}{4}+\frac x2 \right).
\]

Equality would require

\[
\cot\left( \frac x2+\frac{\pi}{8} \right) = \tan\left( \frac{\pi}{4}+\frac x2 \right).
\]

Using

\[
\cot u=\tan\left(\frac{\pi}{2}-u\right),
\]

this becomes

\[
\tan\left( \frac{3\pi}{8}-\frac x2 \right) = \tan\left( \frac{\pi}{4}+\frac x2 \right).
\]

Within the relevant interval, equality requires

\[
\frac{3\pi}{8}-\frac x2 = \frac{\pi}{4}+\frac x2.
\]

Therefore,

\[
x=\frac{\pi}{8}.
\]

Thus the equality holds only at the midpoint of the first octant, not throughout the octant.

Hence an eighth-turn advances the dominant label but does not generally permute the primitive functions exactly.

\end{proof}

\PCsubsection{Corollary 15.3 --- Two distinct symmetry layers}

The bounded phase calculus contains two related but distinct discrete structures:

\begin{enumerate}
\item Exact operator symmetry: generated by quarter-turn and reflection, with induced group \(V_4\).
\item Dominance-label rotation: generated by the eighth-turn, with induced cycle \(C_4\).
\end{enumerate}
They agree when an eighth-turn is applied twice, but not after a single eighth-turn.


\PCsection{16. Terminal Statement}

S0.IV establishes the discrete symmetry architecture of the bounded phase operators.

The phase wheel carries the dihedral symmetry

\[
\boxed{ D_4 = \langle R,J \mid R^4=J^2=e,\ JRJ=R^{-1} \rangle. }
\]

The primitive operators provide a nonfaithful representation of this group. The half-turn acts trivially, so the induced operator group is

\[
\boxed{ D_4/\{e,R^2\} \cong V_4. }
\]

Quarter-turn acts by

\[
\boxed{ \operatorname{srx}\leftrightarrow\operatorname{crx}, \qquad \operatorname{sxp}\leftrightarrow\operatorname{cxp}. }
\]

Reflection acts by reciprocal exchange:

\[
\boxed{ \operatorname{srx}\leftrightarrow\operatorname{sxp}, \qquad \operatorname{cxp}\leftrightarrow\operatorname{crx}. }
\]

Diagonal reflection exchanges the sine and cosine families:

\[
\boxed{ \operatorname{srx}\leftrightarrow\operatorname{cxp}, \qquad \operatorname{sxp}\leftrightarrow\operatorname{crx}. }
\]

Flatwave transforms according to a sign character:

\[
\boxed{ \operatorname{Flatwave}(Rx) = -\operatorname{Flatwave}(x), }
\]

\[
\boxed{ \operatorname{Flatwave}(Jx) = -\operatorname{Flatwave}(x), }
\]

and

\[
\boxed{ \operatorname{Flatwave}(Kx) = \operatorname{Flatwave}(x). }
\]

The transfer fraction obeys

\[
\boxed{ \lambda(Rx)=\lambda(x), }
\]

\[
\boxed{ \lambda(Jx)=1-\lambda(x), }
\]

and

\[
\boxed{ \lambda(Kx)=1-\lambda(x). }
\]

Finally, the eighth-turn advances the dominant operator through

\[
\boxed{ \operatorname{srx} \rightarrow \operatorname{cxp} \rightarrow \operatorname{crx} \rightarrow \operatorname{sxp} \rightarrow \operatorname{srx}. }
\]

This is an exact rotation of dominance labels, but not generally an exact permutation of the shifted primitive functions.

The resulting structure is:

\[
\boxed{ \text{dihedral phase symmetry, Klein operator symmetry, and cyclic octant dominance.} }
\]

% ===== END S0_IV =====
% ===== BEGIN S0_V =====
% Source: S0.v.docx
\PCchapter{S0.V}{Boundary Structure, Singular Limits, and Admissible Extension}{Book 0 - Formal Core}

\PCsubsection{Abstract}

The bounded phase operators are defined on the punctured phase domain

\[
\mathcal D = \mathbb R \setminus \left\{ \frac{k\pi}{2}:k\in\mathbb Z \right\}.
\]

The excluded points divide the phase line into open quadrants. The preceding essays established exact algebraic, differential, transfer, and symmetry structures inside those quadrants. This essay determines what the core calculus establishes at their boundaries.

The primitive operators exhibit alternating zero–pole exchanges. At a sine-zero boundary, the Sininverex–Sinexpon pair exchanges \(0\) and \(+\infty\), while the cosine-family pair tends to \(1\). At a cosine-zero boundary, the Cosinexpon–Cosininverex pair exchanges \(+\infty\) and \(0\), while the sine-family pair tends to \(1\).

Within every open quadrant, the transfer coordinate

\[
\lambda=\frac{1-p}{2}
\]

increases continuously from \(0\) to \(1\). It therefore extends continuously to the corresponding closed quadrant. The saw responses extend with it:

\[
\operatorname{saw}_{r} = \varepsilon\lambda, \qquad \operatorname{saw}_{x} = \varepsilon(1-\lambda),
\]

so the signed unit total is transferred from the expon saw channel to the inverex saw channel.

At a common boundary between adjacent quadrants, however, the Flatwave sign reverses. The one-sided limits satisfy

\[
\lim_{x\to b^-}\operatorname{Flatwave}(x) = - \lim_{x\to b^+}\operatorname{Flatwave}(x).
\]


Consequently, Flatwave has no continuous extension to the ordinary phase circle. The individual saw responses also possess finite but unequal one-sided limits at each boundary.

A split-boundary extension domain is introduced in which every boundary has a distinct entering and exiting copy. All bounded transfer quantities extend continuously to this enlarged domain. Identifying those copies requires an additional seam rule. The resulting countable disjoint union is not asserted to be compact as a whole.

The core calculus therefore determines one-sided limits but does not determine boundary-crossing dynamics.


\PCsection{0. Scope and Inheritance}

\PCsubsection{0.1 Scope}

This essay proves:

\begin{enumerate}
\item the complete one-sided boundary limits of the primitive operators;
\item the first-order asymptotic behavior of the zero–pole exchanges;
\item the one-sided limits of the UNA operators;
\item the finite boundary limits of the saw responses;
\item continuous compactification of the transfer coordinate on each quadrant;
\item the Flatwave sign-jump theorem;
\item the impossibility of a global continuous Flatwave extension on the ordinary phase circle;
\item continuous extension on a split-boundary phase domain, without claiming global compactness;
\item the underdetermination of any boundary-crossing law.
\end{enumerate}
\PCsubsection{0.2 Non-scope}

This essay does not:

\begin{itemize}
\item define a physical phase transition;
\item prescribe what happens when a trajectory reaches an excluded point;
\item identify divergence with infinite physical energy or mass;
\item identify a sign reversal with annihilation, creation, or temporal reversal;
\item introduce a new evolution law;
\item remove the original exclusions from the analytic definitions.
\end{itemize}
\PCsubsection{0.3 Inheritance}

All definitions and theorems of S0.I–S0.IV are inherited.

In particular, the primitive operators are

\[
\operatorname{srx}(x) = |\csc x|+\cot x,
\]

\[
\operatorname{sxp}(x) = |\csc x|-\cot x = \frac1{\operatorname{srx}(x)},
\]

\[
\operatorname{cxp}(x) = |\sec x|+\tan x,
\]

and

\[
\operatorname{crx}(x) = |\sec x|-\tan x = \frac1{\operatorname{cxp}(x)}.
\]

The UNA operators are

\[
\operatorname{urx} = \operatorname{srx}-\operatorname{crx},
\]

and

\[
\operatorname{uxp} = \operatorname{cxp}-\operatorname{sxp}.
\]

The saw responses are

\[
\operatorname{saw}_{r} = \frac1{\operatorname{urx}}, \qquad \operatorname{saw}_{x} = \frac1{\operatorname{uxp}}.
\]

Flatwave is

\[
\operatorname{Flatwave} = \operatorname{saw}_{r}+\operatorname{saw}_{x} = \operatorname{sgn}(\sin 2x).
\]


On every open quadrant,

\[
\operatorname{saw}_{r} = \varepsilon\lambda,
\]

and

\[
\operatorname{saw}_{x} = \varepsilon(1-\lambda),
\]

where

\[
\varepsilon=\operatorname{sgn}(\sin 2x)
\]

is constant on that quadrant and

\[
0<\lambda<1.
\]


\PCsection{1. Boundary Indexing}

\PCsubsection{Definition 1.1 --- Boundary sequence}

For each integer \(k\), define

\[
\boxed{ b_k=\frac{k\pi}{2}. }
\]

The forbidden set is

\[
\mathbb B = \{b_k:k\in\mathbb Z\}.
\]

\PCsubsection{Definition 1.2 --- Open quadrants}

Define

\[
\boxed{ Q_k=(b_k,b_{k+1}). }
\]

Then

\[
\mathcal D = \bigcup_{k\in\mathbb Z}Q_k.
\]

The quadrants are pairwise disjoint connected components of \(\mathcal D\).

\PCsubsection{Definition 1.3 --- Quadrant sign}

For \(x\in Q_k\), define

\[
\varepsilon_k = \operatorname{sgn}(\sin 2x).
\]

Since

\[
2x\in(k\pi,(k+1)\pi),
\]

the sine function has constant sign on that interval. Therefore,

\[
\boxed{ \varepsilon_k=(-1)^k. }
\]

Adjacent quadrants have opposite signs:

\[
\boxed{ \varepsilon_{k+1}=-\varepsilon_k. }
\]

\PCsubsection{Definition 1.4 --- Boundary types}

A boundary \(b_k\) is called a sine boundary when \(k\) is even. Then

\[
b_k=m\pi
\]

for some integer \(m\), and

\[
\sin b_k=0.
\]

A boundary \(b_k\) is called a cosine boundary when \(k\) is odd. Then

\[
b_k=\frac{\pi}{2}+m\pi,
\]

and

\[
\cos b_k=0.
\]


\PCsection{2. Standard Small-Angle Limits}

\PCsubsection{Inherited limit 2.1}

As \(h\to0\),

\[
\frac{\sin h}{h}\to1.
\]

Equivalently,

\[
\sin h\sim h.
\]

\PCsubsection{Inherited limit 2.2}

As \(h\to0\),

\[
1-\cos h\sim\frac{h^2}{2},
\]

while

\[
1+\cos h\to2.
\]

\PCsubsection{Definition 2.3 --- Asymptotic equivalence}

For functions \(f\) and \(g\), the notation

\[
f(h)\sim g(h) \quad\text{as }h\to0
\]

means

\[
\lim_{h\to0}\frac{f(h)}{g(h)}=1.
\]

These standard limits will be used to determine the order of each boundary zero or pole.


\PCsection{3. Primitive Limits at Sine Boundaries}

\PCsubsection{Theorem 3.1 --- Sine-family one-sided limits}

Let

\[
b=m\pi
\]

be a sine boundary. Then

\[
\boxed{ \lim_{x\to b^-}\operatorname{srx}(x)=0, }
\]

\[
\boxed{ \lim_{x\to b^+}\operatorname{srx}(x)=+\infty, }
\]

\[
\boxed{ \lim_{x\to b^-}\operatorname{sxp}(x)=+\infty, }
\]

and

\[
\boxed{ \lim_{x\to b^+}\operatorname{sxp}(x)=0. }
\]

\begin{proof}
Because all primitive operators are \(\pi\)-periodic, it is sufficient to examine \(b=0\).

First approach from the right. Let

\[
x=h, \qquad h\to0^+.
\]

Then

\[
\operatorname{srx}(h) = \csc h+\cot h = \frac{1+\cos h}{\sin h}.
\]

Therefore,

\[
\operatorname{srx}(h) \sim \frac2h,
\]

so

\[
\operatorname{srx}(h)\to+\infty.
\]

Since

\[
\operatorname{sxp} = \frac1{\operatorname{srx}},
\]

it follows that

\[
\operatorname{sxp}(h)\to0.
\]

Now approach from the left. Let

\[
x=-h, \qquad h\to0^+.
\]

Then

\[
|\csc(-h)|=\csc h
\]

and

\[
\cot(-h)=-\cot h.
\]

Thus,

\[
\operatorname{srx}(-h) = \csc h-\cot h = \frac{1-\cos h}{\sin h}.
\]

Using the standard limits,

\[
\operatorname{srx}(-h) \sim \frac{h^2/2}{h} = \frac h2.
\]

Therefore,

\[
\operatorname{srx}(-h)\to0.
\]

Taking reciprocals gives

\[
\operatorname{sxp}(-h)\to+\infty.
\]

The same limits hold at every \(m\pi\) by \(\pi\)-periodicity.

\end{proof}

\PCsubsection{Theorem 3.2 --- Cosine-family limits at sine boundaries}

At every sine boundary \(b=m\pi\),

\[
\boxed{ \lim_{x\to b^\pm}\operatorname{cxp}(x)=1 }
\]

and

\[
\boxed{ \lim_{x\to b^\pm}\operatorname{crx}(x)=1. }
\]

\begin{proof}
As \(x\to m\pi\),

\[
|\sec x|\to1, \qquad \tan x\to0.
\]

Therefore,

\[
\operatorname{cxp}(x) = |\sec x|+\tan x \to1.
\]

Likewise,

\[
\operatorname{crx}(x) = |\sec x|-\tan x \to1.
\]

\end{proof}

\PCsubsection{Corollary 3.3 --- Sine-boundary primitive table}

At every sine boundary,

\[
\boxed{ \begin{array}{c|cc} & x\to b^- & x\to b^+ \\ \hline \operatorname{srx} & 0 & +\infty \\ \operatorname{sxp} & +\infty & 0 \\ \operatorname{cxp} & 1 & 1 \\ \operatorname{crx} & 1 & 1 \end{array} }
\]


\PCsection{4. Primitive Limits at Cosine Boundaries}

\PCsubsection{Theorem 4.1 --- Cosine-family one-sided limits}

Let

\[
b=\frac{\pi}{2}+m\pi
\]

be a cosine boundary. Then

\[
\boxed{ \lim_{x\to b^-}\operatorname{cxp}(x)=+\infty, }
\]

\[
\boxed{ \lim_{x\to b^+}\operatorname{cxp}(x)=0, }
\]

\[
\boxed{ \lim_{x\to b^-}\operatorname{crx}(x)=0, }
\]

and

\[
\boxed{ \lim_{x\to b^+}\operatorname{crx}(x)=+\infty. }
\]

\begin{proof}
By \(\pi\)-periodicity, it is sufficient to examine \(b=\pi/2\).

Approach from the left. Let

\[
x=\frac{\pi}{2}-h, \qquad h\to0^+.
\]

Then

\[
\cos x=\sin h, \qquad \sin x=\cos h.
\]

Therefore,

\[
\operatorname{cxp}(x) = \frac{1+\sin x}{\cos x} = \frac{1+\cos h}{\sin h}.
\]

Hence,

\[
\operatorname{cxp}(x)\sim\frac2h,
\]

so

\[
\operatorname{cxp}(x)\to+\infty.
\]

Its reciprocal satisfies

\[
\operatorname{crx}(x)\to0.
\]

Approach from the right. Let

\[
x=\frac{\pi}{2}+h.
\]

Then

\[
\cos x=-\sin h, \qquad \sin x=\cos h.
\]

Since

\[
|\sec x|=\frac1{\sin h}
\]

and

\[
\tan x=-\frac{\cos h}{\sin h},
\]

we have

\[
\operatorname{cxp}(x) = \frac{1-\cos h}{\sin h}.
\]

Thus,

\[
\operatorname{cxp}(x) \sim \frac h2,
\]

so

\[
\operatorname{cxp}(x)\to0.
\]

Its reciprocal therefore diverges:

\[
\operatorname{crx}(x)\to+\infty.
\]

\end{proof}

\PCsubsection{Theorem 4.2 --- Sine-family limits at cosine boundaries}

At every cosine boundary,

\[
\boxed{ \lim_{x\to b^\pm}\operatorname{srx}(x)=1 }
\]

and

\[
\boxed{ \lim_{x\to b^\pm}\operatorname{sxp}(x)=1. }
\]

\begin{proof}
As \(x\to\pi/2+m\pi\),

\[
|\csc x|\to1, \qquad \cot x\to0.
\]

Therefore,

\[
\operatorname{srx}\to1
\]

and

\[
\operatorname{sxp}\to1.
\]

\end{proof}

\PCsubsection{Corollary 4.3 --- Cosine-boundary primitive table}

At every cosine boundary,

\[
\boxed{ \begin{array}{c|cc} & x\to b^- & x\to b^+ \\ \hline \operatorname{srx} & 1 & 1 \\ \operatorname{sxp} & 1 & 1 \\ \operatorname{cxp} & +\infty & 0 \\ \operatorname{crx} & 0 & +\infty \end{array} }
\]


\PCsection{5. First-Order Primitive Asymptotics}

\PCsubsection{Theorem 5.1 --- Sine-boundary zero–pole scale}

Let \(b=m\pi\) and \(h\to0^+\). Then

\[
\boxed{ \operatorname{srx}(b-h)\sim\frac h2, }
\]

\[
\boxed{ \operatorname{sxp}(b-h)\sim\frac2h, }
\]

\[
\boxed{ \operatorname{srx}(b+h)\sim\frac2h, }
\]

and

\[
\boxed{ \operatorname{sxp}(b+h)\sim\frac h2. }
\]

\begin{proof}
These asymptotic relations were obtained in the proof of Theorem 3.1. The reciprocal relations follow because

\[
\operatorname{srx}\operatorname{sxp}=1.
\]

\end{proof}

\PCsubsection{Theorem 5.2 --- Cosine-boundary zero–pole scale}

Let

\[
b=\frac{\pi}{2}+m\pi
\]

and \(h\to0^+\). Then

\[
\boxed{ \operatorname{cxp}(b-h)\sim\frac2h, }
\]

\[
\boxed{ \operatorname{crx}(b-h)\sim\frac h2, }
\]

\[
\boxed{ \operatorname{cxp}(b+h)\sim\frac h2, }
\]

and

\[
\boxed{ \operatorname{crx}(b+h)\sim\frac2h. }
\]

\PCsubsection{Corollary 5.3 --- Reciprocal pole order}

Every primitive divergence at an excluded boundary is a first-order pole in the distance to the boundary.

Every reciprocal zero vanishes linearly in that same distance.

The products

\[
\operatorname{srx}\operatorname{sxp}=1
\]

and

\[
\operatorname{cxp}\operatorname{crx}=1
\]

remain exact throughout the approach, even though neither factor is defined at the boundary itself.


\PCsection{6. Transfer Coordinate at Quadrant Boundaries}

Fix an open quadrant

\[
Q_k=(b_k,b_{k+1})
\]

and write

\[
\varepsilon=\varepsilon_k.
\]

\PCsubsection{Theorem 6.1 --- Transfer-coordinate endpoint limits}

Within \(Q_k\),

\[
\boxed{ \lim_{x\to b_k^+}\lambda(x)=0 }
\]

and

\[
\boxed{ \lim_{x\to b_{k+1}^-}\lambda(x)=1. }
\]

\begin{proof}
From S0.III,

\[
\lambda=\frac{1-p}{2},
\]

where \(p\) decreases continuously from \(1\) to \(-1\) across every open quadrant.

Therefore,

\[
p\to1 \quad\text{as }x\to b_k^+,
\]

which gives

\[
\lambda\to0.
\]

Likewise,

\[
p\to-1 \quad\text{as }x\to b_{k+1}^-,
\]

which gives

\[
\lambda\to1.
\]

\end{proof}

\PCsubsection{Theorem 6.2 --- First-order transfer asymptotics}

Let \(h\to0^+\). Then

\[
\boxed{ \lambda(b_k+h)\sim\frac h2 }
\]

and

\[
\boxed{ 1-\lambda(b_{k+1}-h)\sim\frac h2. }
\]

\begin{proof}
From S0.III,

\[
\lambda'=\frac q2,
\]

where

\[
q=|\sin x|+|\cos x|.
\]

At either boundary of a quadrant,

\[
q\to1.
\]

Therefore,

\[
\lambda'\to\frac12
\]

at the entering boundary.

Since \(\lambda\to0\) there,

\[
\lambda(b_k+h)\sim\frac h2.
\]

At the exiting boundary,

\[
(1-\lambda)'=-\lambda'\to-\frac12.
\]

Since \(1-\lambda\to0\),

\[
1-\lambda(b_{k+1}-h)\sim\frac h2.
\]


\end{proof}

\PCsection{7. UNA Boundary Structure}

\PCsubsection{Theorem 7.1 --- Quadrant-relative UNA limits}

Within the quadrant \(Q_k\), with sign \(\varepsilon=\varepsilon_k\),

\[
\boxed{ \lim_{x\to b_k^+}\operatorname{urx}(x) = \varepsilon\infty, }
\]

\[
\boxed{ \lim_{x\to b_k^+}\operatorname{uxp}(x) = \varepsilon, }
\]

\[
\boxed{ \lim_{x\to b_{k+1}^-}\operatorname{urx}(x) = \varepsilon, }
\]

and

\[
\boxed{ \lim_{x\to b_{k+1}^-}\operatorname{uxp}(x) = \varepsilon\infty. }
\]

Here \(\varepsilon\infty\) denotes divergence with sign \(\varepsilon\).

\begin{proof}
From S0.III,

\[
\operatorname{urx} = \frac{\varepsilon}{\lambda},
\]

and

\[
\operatorname{uxp} = \frac{\varepsilon}{1-\lambda}.
\]

At the entering boundary,

\[
\lambda\to0^+, \qquad 1-\lambda\to1.
\]

Thus,

\[
\operatorname{urx}\to\varepsilon\infty
\]

and

\[
\operatorname{uxp}\to\varepsilon.
\]

At the exiting boundary,

\[
\lambda\to1, \qquad 1-\lambda\to0^+.
\]

Therefore,

\[
\operatorname{urx}\to\varepsilon
\]

and

\[
\operatorname{uxp}\to\varepsilon\infty.
\]

\end{proof}

\PCsubsection{Corollary 7.2 --- UNA divergence transfer}

Across every open quadrant, the first-order pole transfers from Unainverex to Unaexpon:

\[
\boxed{ (\varepsilon\infty,\varepsilon) \longrightarrow (\varepsilon,\varepsilon\infty). }
\]

The two operators retain the same sign throughout the quadrant.

\PCsubsection{Theorem 7.3 --- UNA first-order asymptotics}

Let \(h\to0^+\). Within \(Q_k\),

\[
\boxed{ \operatorname{urx}(b_k+h)\sim\frac{2\varepsilon}{h}, }
\]

and

\[
\boxed{ \operatorname{uxp}(b_{k+1}-h)\sim\frac{2\varepsilon}{h}. }
\]

\begin{proof}
Use

\[
\operatorname{urx} = \frac{\varepsilon}{\lambda}
\]

with

\[
\lambda(b_k+h)\sim\frac h2.
\]

Then

\[
\operatorname{urx}(b_k+h) \sim \frac{\varepsilon}{h/2} = \frac{2\varepsilon}{h}.
\]

The proof for \(\operatorname{uxp}\) is identical.


\end{proof}

\PCsection{8. Saw Boundary Structure}

\PCsubsection{Theorem 8.1 --- Quadrant-relative saw limits}

Within \(Q_k\),

\[
\boxed{ \lim_{x\to b_k^+} \operatorname{saw}_{r}(x)=0, }
\]

\[
\boxed{ \lim_{x\to b_k^+} \operatorname{saw}_{x}(x)=\varepsilon_k, }
\]

\[
\boxed{ \lim_{x\to b_{k+1}^-} \operatorname{saw}_{r}(x)=\varepsilon_k, }
\]

and

\[
\boxed{ \lim_{x\to b_{k+1}^-} \operatorname{saw}_{x}(x)=0. }
\]

\begin{proof}
Use

\[
\operatorname{saw}_{r} = \varepsilon_k\lambda
\]

and

\[
\operatorname{saw}_{x} = \varepsilon_k(1-\lambda).
\]

The result follows from the endpoint limits of \(\lambda\).

\end{proof}

\PCsubsection{Corollary 8.2 --- Signed unit transfer}

Across every open quadrant,

\[
\boxed{ (\operatorname{saw}_{r},\operatorname{saw}_{x}) : (0,\varepsilon_k) \longrightarrow (\varepsilon_k,0). }
\]

Thus the complete signed unit weight passes continuously from the expon saw channel to the inverex saw channel.

\PCsubsection{Theorem 8.3 --- Saw first-order asymptotics}

Let \(h\to0^+\). Then

\[
\operatorname{saw}_{r}(b_k+h) \sim \varepsilon_k\frac h2,
\]

\[
\operatorname{saw}_{x}(b_k+h) = \varepsilon_k \left( 1-\frac h2+o(h) \right),
\]

\[
\operatorname{saw}_{r}(b_{k+1}-h) = \varepsilon_k\left(1-\frac h2+o(h)\right),
\]


and

\[
\operatorname{saw}_{x}(b_{k+1}-h) \sim \varepsilon_k\frac h2.
\]

\PCsubsection{Corollary 8.4 --- No saw divergence}

Neither saw response diverges as an excluded boundary is approached from within a quadrant.

The saw expressions remain undefined at the boundary under their original reciprocal definitions, but their quadrant-relative limits are finite.


\PCsection{9. Two-Sided Saw Limits at a Common Boundary}

Let \(b_k\) be a common boundary of \(Q_{k-1}\) and \(Q_k\).

Since

\[
\varepsilon_{k-1}=-\varepsilon_k,
\]

the two neighboring quadrants carry opposite Flatwave signs.

\PCsubsection{Theorem 9.1 --- Two-sided saw limit table}

At \(b_k\),

\[
\boxed{ \lim_{x\to b_k^-}\operatorname{saw}_{r}(x) = \varepsilon_{k-1} = -\varepsilon_k }.
\]


while

\[
\boxed{ \lim_{x\to b_k^+}\operatorname{saw}_{r}(x)=0. }
\]

Likewise,

\[
\boxed{ \lim_{x\to b_k^-}\operatorname{saw}_{x}(x)=0, }
\]

while

\[
\boxed{ \lim_{x\to b_k^+}\operatorname{saw}_{x}(x) = \varepsilon_k. }
\]

\begin{proof}
The left side of \(b_k\) is the exiting boundary of \(Q_{k-1}\). Apply Theorem 8.1 with sign \(\varepsilon_{k-1}\).

The right side is the entering boundary of \(Q_k\). Apply Theorem 8.1 with sign \(\varepsilon_k\).

\end{proof}

\PCsubsection{Corollary 9.2 --- No global continuous saw extension}

Neither individual saw response has equal left and right limits at a boundary.

Therefore, neither

\[
\operatorname{saw}_{r}
\]

nor

\[
\operatorname{saw}_{x}
\]

admits a continuous extension to the ordinary phase circle at any excluded boundary.

Their exclusions are removable only relative to a chosen closed quadrant, not globally after neighboring quadrants are identified.


\PCsection{10. Flatwave Boundary Discontinuity}

\PCsubsection{Theorem 10.1 --- Flatwave one-sided limits}

At every boundary \(b_k\),

\[
\boxed{ \lim_{x\to b_k^-}\operatorname{Flatwave}(x) = \varepsilon_{k-1}, }
\]

and

\[
\boxed{ \lim_{x\to b_k^+}\operatorname{Flatwave}(x) = \varepsilon_k. }
\]

Since

\[
\varepsilon_k=-\varepsilon_{k-1},
\]

we have

\[
\boxed{ \lim_{x\to b_k^-}\operatorname{Flatwave}(x) = - \lim_{x\to b_k^+}\operatorname{Flatwave}(x) }.
\]


\begin{proof}
Flatwave is constant on every open quadrant:

\[
\operatorname{Flatwave}(x)=\varepsilon_j \qquad \text{for }x\in Q_j.
\]

The left-hand limit at \(b_k\) is therefore \(\varepsilon_{k-1}\), while the right-hand limit is \(\varepsilon_k\).

Adjacent signs are opposite.

\end{proof}

\PCsubsection{Corollary 10.2 --- Nonexistence of the two-sided Flatwave limit}

For every \(b_k\in\mathbb B\),

\[
\boxed{ \lim_{x\to b_k}\operatorname{Flatwave}(x) \text{ does not exist.} }
\]

\PCsubsection{Corollary 10.3 --- No continuous real-valued extension}

There is no function

\[
\overline{\operatorname{Flatwave}} : \mathbb R\to\mathbb R
\]

that is continuous at every \(b_k\) and agrees with Flatwave on \(\mathcal D\).

\begin{proof}
A continuous extension at \(b_k\) would require equal left and right limits. Theorem 10.1 proves that these limits are opposite.

\end{proof}

\PCsubsection{Definition 10.4 --- Flatwave phase reversal}

The exact mathematical event at a boundary is the sign reversal

\[
\varepsilon_{k-1}\longrightarrow\varepsilon_k=-\varepsilon_{k-1}.
\]

This may be called a Flatwave phase reversal.

The term denotes a discontinuity of the signed quadrant invariant. It does not, by itself, identify a physical transition.


\PCsection{11. Quadrant-Wise Compactification}

\PCsubsection{Definition 11.1 --- Closed quadrant}

For each integer \(k\), define

\[
\overline Q_k=[b_k,b_{k+1}].
\]

The ordinary operator formulas remain defined only on the interior \(Q_k\).

\PCsubsection{Definition 11.2 --- Extended transfer coordinate}

Define

\[
\overline\lambda_k: \overline Q_k\to[0,1]
\]

by

\[
\overline\lambda_k(x) = \lambda(x) \qquad \text{for }x\in Q_k,
\]

together with

\[
\overline\lambda_k(b_k)=0
\]

and

\[
\overline\lambda_k(b_{k+1})=1.
\]

\PCsubsection{Theorem 11.3 --- Continuous transfer compactification}

The function

\[
\overline\lambda_k
\]

is continuous on the closed quadrant \(\overline Q_k\).

\begin{proof}
The original \(\lambda\) is continuous on \(Q_k\).

By Theorem 6.1,

\[
\lim_{x\to b_k^+}\lambda(x)=0 = \overline\lambda_k(b_k),
\]

and

\[
\lim_{x\to b_{k+1}^-}\lambda(x)=1 = \overline\lambda_k(b_{k+1}).
\]

Therefore, the endpoint assignments produce a continuous extension.

\end{proof}

\PCsubsection{Definition 11.4 --- Extended saw responses}

Define

\[
\overline{\operatorname{saw}}_{r,k} = \varepsilon_k\overline\lambda_k,
\]

and

\[
\overline{\operatorname{saw}}_{x,k} = \varepsilon_k(1-\overline\lambda_k).
\]

\PCsubsection{Corollary 11.5 --- Continuous saw compactification}

Both extended saw responses are continuous on \(\overline Q_k\).

Their endpoint values are

\[
\overline{\operatorname{saw}}_{r,k}(b_k)=0,
\]

\[
\overline{\operatorname{saw}}_{x,k}(b_k)=\varepsilon_k,
\]

\[
\overline{\operatorname{saw}}_{r,k}(b_{k+1})=\varepsilon_k,
\]

and

\[
\overline{\operatorname{saw}}_{x,k}(b_{k+1})=0.
\]

\PCsubsection{Definition 11.6 --- Quadrant-relative Flatwave extension}

Define

\[
\overline{\operatorname{Flatwave}}_k(x)=\varepsilon_k
\]

for every

\[
x\in\overline Q_k.
\]

Then

\[
\overline{\operatorname{Flatwave}}_k = \overline{\operatorname{saw}}_{r,k} + \overline{\operatorname{saw}}_{x,k}
\]


on the entire closed quadrant.

\PCsubsection{Corollary 11.7 --- Local but incompatible extensions}

Every bounded transfer expression extends continuously to each closed quadrant separately.

However, adjacent closed quadrants assign incompatible values to their shared boundary.

For example,

\[
\overline{\operatorname{Flatwave}}_{k-1}(b_k) = \varepsilon_{k-1},
\]

while

\[
\overline{\operatorname{Flatwave}}_k(b_k) = \varepsilon_k = -\varepsilon_{k-1}.
\]

Thus the obstruction is not failure of one-sided compactification. It is incompatibility under boundary identification.


\PCsection{12. Split-Boundary Phase Domain}

\PCsubsection{Definition 12.1 --- Split-boundary extension domain}

Define

\[
\boxed{ \widehat{\mathcal D} = \bigsqcup_{k\in\mathbb Z}\overline Q_k, }
\]

the disjoint union of all closed quadrants.

Each component \(\overline Q_k\) is compact, but the countable disjoint union \(\widehat{\mathcal D}\) is not thereby a compact space. Accordingly, \emph{compactification} in this essay applies only componentwise to \(Q_k\subset\overline Q_k\); the global object is the split-boundary extension domain.

In this space, the exiting copy of a boundary and the entering copy of the same numerical phase are distinct points.

For example, the numerical boundary \(b_k\) appears as:

\[
b_k^{\mathrm{exit}} \in\overline Q_{k-1},
\]

and

\[
b_k^{\mathrm{entry}} \in\overline Q_k.
\]

\PCsubsection{Theorem 12.2 --- Continuous bounded extension on the split domain}

The functions

\[
\lambda, \qquad \operatorname{saw}_{r}, \qquad \operatorname{saw}_{x}, \qquad \operatorname{Flatwave}
\]


extend continuously to \(\widehat{\mathcal D}\).

\begin{proof}
Each component of \(\widehat{\mathcal D}\) is one closed quadrant \(\overline Q_k\).

Section 11 constructs a continuous extension of every listed function to each \(\overline Q_k\).

Because the components are disjoint, no equality between entering and exiting boundary values is required.

Therefore, these component-wise extensions define continuous functions on the disjoint union.

\end{proof}

\PCsubsection{Corollary 12.3 --- Boundary state distinction}

The core bounded variables distinguish an exiting boundary state from an entering boundary state.

At \(b_k\),

\[
\left( \lambda, \operatorname{saw}_{r}, \operatorname{saw}_{x}, \operatorname{Flatwave} \right)_{\mathrm{exit}} = \left(1,\varepsilon_{k-1},0,\varepsilon_{k-1}\right),
\]


while

\[
\left( \lambda, \operatorname{saw}_{r}, \operatorname{saw}_{x}, \operatorname{Flatwave} \right)_{\mathrm{entry}} = \left(0,0,\varepsilon_k,\varepsilon_k\right).
\]


Since

\[
\varepsilon_k=-\varepsilon_{k-1},
\]

these are distinct boundary states.


\PCsection{13. Classification of Boundary Behavior}

\PCsubsection{Definition 13.1 --- Divergent boundary}

A function has a divergent boundary at \(b\) from one side if its magnitude tends to \(+\infty\) from that side.

\PCsubsection{Definition 13.2 --- Finite one-sided boundary}

A function has a finite one-sided boundary value if the corresponding one-sided limit exists in \(\mathbb R\).

\PCsubsection{Definition 13.3 --- Globally removable exclusion}

An excluded point is globally removable for a real-valued function when the left and right limits exist, are finite, and are equal.

\PCsubsection{Definition 13.4 --- Quadrant-removable exclusion}

An excluded endpoint is quadrant-removable when the interior function has a finite limit from within that quadrant, allowing a continuous extension to that closed quadrant.

\PCsubsection{Theorem 13.5 --- Boundary classification}

The operator families satisfy the following classification.

\PCsubsubsection{Primitive operators}

At every excluded boundary:

\begin{itemize}
\item one reciprocal pair undergoes a zero–pole exchange;
\item the other reciprocal pair approaches \(1\);
\item at least one primitive diverges from one side.
\end{itemize}
Therefore, the primitive family does not admit a finite global extension.

\PCsubsubsection{UNA operators}

Within each quadrant:

\begin{itemize}
\item Unainverex diverges at entry and is finite at exit;
\item Unaexpon is finite at entry and diverges at exit.
\end{itemize}
Therefore, neither UNA operator extends finitely to both endpoints of a closed quadrant.

\PCsubsubsection{Saw operators}

Each saw has finite one-sided limits at every boundary and extends continuously to every closed quadrant.

However, its two-sided limits are unequal.

Therefore, each saw exclusion is quadrant-removable but not globally removable.

\PCsubsubsection{Flatwave}

Flatwave extends continuously to every closed quadrant as a constant.

Its neighboring one-sided values are opposite.

Therefore, its boundary exclusion is quadrant-removable but globally nonremovable.

\PCsubsubsection{Transfer coordinate}

The transfer coordinate extends continuously to every closed quadrant from \(0\) to \(1\).

Its exit value and the next quadrant's entry value differ:

\[
1\neq0.
\]

Therefore, global identification requires an explicit reset or gluing rule.


\PCsection{14. Seam Maps}

\PCsubsection{Definition 14.1 --- Exit and entry sets}

Let

\[
E_k = b_{k+1}^{\mathrm{exit}} \in\overline Q_k
\]

denote the exit endpoint of \(Q_k\).

Let

\[
I_{k+1} = b_{k+1}^{\mathrm{entry}} \in\overline Q_{k+1}
\]

denote the entering endpoint of the next quadrant.

\PCsubsection{Definition 14.2 --- Seam map}

A seam map is an additional rule assigning an entering state to an exiting state:

\[
\Sigma:E\to I,
\]

where \(E\) is the set of exit endpoints and \(I\) is the set of entry endpoints.

A seam map may specify, for example:

\begin{itemize}
\item which quadrant follows an exit;
\item whether phase direction is preserved or reversed;
\item whether \(\lambda\) resets from \(1\) to \(0\);
\item whether saw labels are exchanged;
\item whether Flatwave sign reverses;
\item whether additional state variables are retained.
\end{itemize}
\PCsubsection{Proposition 14.3 --- Interior independence of seam choice}

The identities proved in S0.I–S0.IV and the one-sided limits proved in this essay hold independently of the choice of seam map.

\begin{proof}
Every earlier theorem is stated either:

\begin{itemize}
\item at an interior point \(x\in Q_k\); or
\item as a one-sided limit taken from within a specified quadrant.
\end{itemize}
A seam map acts only after an exit endpoint has been reached in the split-boundary description.

Changing the seam map does not alter any value or derivative at an interior point, nor does it alter a one-sided limit already determined from within the quadrant.

Therefore, all core theorems remain valid under more than one possible seam map.

\end{proof}

\PCsubsection{Theorem 14.4 --- No unique boundary-crossing law}

The core bounded phase calculus does not uniquely determine a seam map.

\begin{proof}
The core calculus determines:

\[
\lambda\to1, \qquad \operatorname{saw}_{r}\to\varepsilon_k, \qquad \operatorname{saw}_{x}\to0
\]


at the exit of \(Q_k\).

It also determines the possible entry data of every quadrant:

\[
\lambda=0, \qquad \operatorname{saw}_{r}=0, \qquad \operatorname{saw}_{x}=\varepsilon_j
\]


at the entry of \(Q_j\).

Nothing in the interior identities selects which entry copy must be paired with a given exit copy.

At minimum, one may define:

\begin{enumerate}
\item a forward phase seam pairing the exit of \(Q_k\) with the entry of \(Q_{k+1}\);
\item a reflected seam pairing that exit with an entry selected by a phase reflection;
\item a periodic or lifted seam retaining additional winding information;
\item a terminating rule under which no continuation is defined.
\end{enumerate}
Each choice leaves all interior identities and one-sided limits unchanged.

Therefore, no unique seam map follows from the core calculus.

\end{proof}

\PCsubsection{Corollary 14.5 --- Boundary dynamics require a new assumption}

Any statement specifying what happens after a boundary is reached is additional to Book 0's interior operator calculus.

Such a statement must be labeled as one of:

\begin{itemize}
\item a boundary axiom;
\item an evolution rule;
\item a projection convention;
\item or a model-specific hypothesis.
\end{itemize}
It may not be presented as a consequence of the Flatwave identity alone.


\PCsection{15. Boundary Ledger}

\PCsubsection{15.1 Sine boundary}

For

\[
b=m\pi,
\]

the primitive limits are

\[
\begin{array}{c|cc} & b^- & b^+ \\ \hline \operatorname{srx} & 0 & +\infty \\ \operatorname{sxp} & +\infty & 0 \\ \operatorname{cxp} & 1 & 1 \\ \operatorname{crx} & 1 & 1 \end{array}
\]

\PCsubsection{15.2 Cosine boundary}

For

\[
b=\frac{\pi}{2}+m\pi,
\]

the primitive limits are

\[
\begin{array}{c|cc} & b^- & b^+ \\ \hline \operatorname{srx} & 1 & 1 \\ \operatorname{sxp} & 1 & 1 \\ \operatorname{cxp} & +\infty & 0 \\ \operatorname{crx} & 0 & +\infty \end{array}
\]

\PCsubsection{15.3 General seam \(b_k\)}

Let

\[
\varepsilon_k=(-1)^k.
\]

Then

\[
\begin{array}{c|cc} & b_k^- & b_k^+ \\ \hline \operatorname{urx} & -\varepsilon_k & \varepsilon_k\infty \\ \operatorname{uxp} & -\varepsilon_k\infty & \varepsilon_k \\ \operatorname{saw}_{r} & -\varepsilon_k & 0 \\ \operatorname{saw}_{x} & 0 & \varepsilon_k \\ \operatorname{Flatwave} & -\varepsilon_k & \varepsilon_k \end{array}
\]


This table records one-sided limits only. No row assigns an actual value at \(b_k\) on the original domain.


\PCsection{16. Terminal Statement}

S0.I established the Pythagorean operator identities.

S0.II established one-generator rational reduction.

S0.III established tangential transfer.

S0.IV established discrete phase and operator symmetry.

S0.V establishes the boundary structure.

At a sine boundary, the sine reciprocal pair exchanges

\[
0 \longleftrightarrow +\infty,
\]

while the cosine pair approaches

\[
1.
\]

At a cosine boundary, the cosine reciprocal pair exchanges

\[
+\infty \longleftrightarrow 0,
\]

while the sine pair approaches

\[
1.
\]

Within each quadrant,

\[
\boxed{ \lambda:0\longrightarrow1 }
\]

and

\[
\boxed{ (\operatorname{saw}_{r},\operatorname{saw}_{x}) : (0,\varepsilon) \longrightarrow (\varepsilon,0). }
\]

The saw responses have finite one-sided limits, but their neighboring limits do not agree.

Flatwave obeys

\[
\boxed{ \lim_{x\to b^-}\operatorname{Flatwave}(x) = - \lim_{x\to b^+}\operatorname{Flatwave}(x) }.
\]


Therefore,

\[
\boxed{ \text{Flatwave has no continuous extension across an identified boundary.} }
\]

A continuous bounded extension exists only when entering and exiting boundary copies remain distinct:

\[
\boxed{ \widehat{\mathcal D} = \bigsqcup_k\overline Q_k. }
\]

The core calculus determines the approach to every boundary but does not determine passage through it.

Thus the boundary result is:

\[
\boxed{\substack{\text{exact one-sided limits and componentwise endpoint compactification,}\\[2pt]\text{with no canonical seam dynamics.}}}
\]

% ===== END S0_V =====
% ===== BEGIN S0_VI =====
% Source: S0.vi.docx
\PCchapter{S0.VI}{Axiomatic Status, Minimality, and the Book 0 Theorem Ledger}{Book 0 - Formal Core}

\PCsubsection{Abstract}

This essay closes Book 0 by auditing the logical status of the bounded phase operator calculus developed in S0.I–S0.V.

The audit distinguishes inherited mathematical background, definitions, conventions, proved theorems, corollaries, conditional statements, conjectures, and interpretations. It shows that, relative to ordinary real trigonometry and elementary real analysis, Book 0 requires no independent R-specific mathematical axiom. The R-operators are definitions; their reciprocal, differential, transfer, symmetry, and boundary properties are derived consequences.

The four displayed primitive operators are not algebraically independent. They form two reciprocal pairs, and on every open quadrant the entire named calculus reduces to rational functions of one nonconstant generator,

\[
z=\operatorname{srx}.
\]

One nonconstant generator is therefore necessary and sufficient for quadrant-local generation. This is a theorem about generator count, not a uniqueness theorem for the particular choice of generator.

The finite named list is not closed under unrestricted rational operations. The correct closure object is the quadrant-local rational differential field

\[
\mathbb R(z),
\]

with derivation

\[
\mathcal D_z = -\frac{1+z^2}{2}\frac{d}{dz}.
\]

Globally, the calculus is piecewise rational and requires branch information identifying the open quadrant or the sign

\[
\varepsilon=\operatorname{sgn}(\sin 2x).
\]

The exact Flatwave theorem,

\[
\operatorname{Flatwave}(x) = \operatorname{sgn}(\sin 2x),
\]

proves a signed partition invariant. It does not by itself establish conservation of physical mass, energy, or information. Likewise, the operator calculus determines one-sided boundary limits but does not determine a seam-crossing or evolution law.

The final result of Book 0 is therefore a bounded phase operator calculus with exact algebraic identities, one-generator local reduction, differential closure, complementary transfer, discrete symmetry, and rigorously classified boundary behavior.


\PCsection{0. Purpose of the Closing Essay}

Book 0 has introduced and analyzed the mathematical core of the S Essay series.

Its purpose has not been to produce a physical theory. Its purpose has been to determine exactly what follows from the bounded phase operators before any external interpretation, projection, calibration, or evolution law is introduced.

This closing essay answers six questions.

\begin{enumerate}
\item What mathematical background was inherited?
\item What objects were merely defined?
\item Which statements were proved?
\item In what sense is the calculus minimal?
\item In what sense is the calculus closed?
\item What claims remain unavailable at the end of Book 0?
\end{enumerate}
No new operator is introduced in this essay.

No new dynamical rule is introduced.

No physical interpretation is promoted to theorem status.


\PCsection{1. Status Vocabulary}

\PCsubsection{Definition 1.1 --- Background result}

A background result is a standard theorem or rule inherited from ordinary mathematics rather than proved within Book 0.

Examples include:

\[
\sin^2x+\cos^2x=1,
\]

the angle-addition formulas, the ordinary derivative rules for sine and cosine, and standard one-sided limits.

A background result is not an R-specific axiom.

\PCsubsection{Definition 1.2 --- Convention}

A convention selects notation, coordinates, orientation, or presentation without asserting a new mathematical fact.

Examples include:

\begin{itemize}
\item measuring phase in radians;
\item writing phase modulo \(2\pi\);
\item naming an operator Sininverex;
\item choosing the sign convention
\end{itemize}
\[
\operatorname{uxp} = \operatorname{cxp}-\operatorname{sxp}.
\]

\PCsubsection{Definition 1.3 --- Definition}

A definition declares a mathematical object.

A definition is neither true nor false in the manner of a theorem. It is judged by consistency, clarity, and usefulness.

\PCsubsection{Definition 1.4 --- Axiom}

An axiom is an assumed proposition used without derivation inside the formal system.

A genuine axiom must assert more than a choice of notation or the definition of an object.

\PCsubsection{Definition 1.5 --- Lemma}

A lemma is a proved supporting proposition used primarily to establish a later result.

\PCsubsection{Definition 1.6 --- Theorem}

A theorem is a principal proposition derived from earlier definitions, background results, axioms, lemmas, or theorems.

\PCsubsection{Definition 1.7 --- Corollary}

A corollary is an immediate or nearly immediate consequence of a theorem.

\PCsubsection{Definition 1.8 --- Conditional theorem}

A conditional theorem is proved only after an additional condition, candidate class, branch selection, or model rule has been declared.

Its conclusion may not be generalized beyond those conditions.

\PCsubsection{Definition 1.9 --- Conjecture}

A conjecture is a mathematically precise claim for which no proof has yet been supplied.

\PCsubsection{Definition 1.10 --- Interpretation}

An interpretation assigns conceptual, physical, informational, biological, economic, or metaphysical meaning to a mathematical object.

An interpretation is not a theorem unless a separate projection and validation argument establishes it.


\PCsection{2. Inherited Mathematical Background}

Book 0 depends on ordinary real mathematics.

The background may be stated compactly as follows.

\PCsubsection{Background B.1 --- Real-number structure}

The real numbers are assumed with their ordinary field operations, order relation, absolute value, and limits.

\PCsubsection{Background B.2 --- Real trigonometric functions}

The functions

\[
\sin x \qquad\text{and}\qquad \cos x
\]

are inherited with their ordinary real definitions.

The quotient and reciprocal functions are defined where their denominators are nonzero:

\[
\tan x=\frac{\sin x}{\cos x},
\]

\[
\cot x=\frac{\cos x}{\sin x},
\]

\[
\sec x=\frac1{\cos x},
\]

and

\[
\csc x=\frac1{\sin x}.
\]

\PCsubsection{Background B.3 --- Pythagorean identity}

\[
\boxed{ \sin^2x+\cos^2x=1. }
\]

This is the primary algebraic background theorem of S0.I.

\PCsubsection{Background B.4 --- Angle-addition identities}

\[
\sin(x+y) = \sin x\cos y+\cos x\sin y,
\]

and

\[
\cos(x+y) = \cos x\cos y-\sin x\sin y.
\]

These support the rotation and reflection analysis in S0.IV.

\PCsubsection{Background B.5 --- Elementary differentiation}

On open intervals where the functions are defined,

\[
\frac{d}{dx}\sin x=\cos x,
\]

\[
\frac{d}{dx}\cos x=-\sin x,
\]

together with the product, quotient, reciprocal, and chain rules.

These support S0.II and S0.III.

\PCsubsection{Background B.6 --- Standard local limits}

As \(h\to0\),

\[
\frac{\sin h}{h}\to1,
\]

and

\[
\frac{1-\cos h}{h^2/2}\to1.
\]

These support the boundary asymptotics of S0.V.

\PCsubsection{Theorem 2.1 --- No independent R-specific axiom is required}

Relative to Background B.1–B.6, all Book 0 results follow from definitions and derivations.

\begin{proof}
The operators are introduced by explicit definitions using ordinary trigonometric functions, absolute value, arithmetic operations, and reciprocation.

Their algebraic identities are derived from the Pythagorean identity.

Their symmetry properties are derived from angle-addition identities.

Their differential properties are derived from ordinary differentiation.

Their boundary properties are derived from standard limits.

No additional unproved proposition concerning the R-operators is needed to obtain the Book 0 theorem set.

Therefore, relative to the inherited mathematical background, Book 0 requires no independent R-specific axiom.

\end{proof}

\PCsubsection{Clarification 2.2 --- Earlier uses of the word “axiom”}

Earlier essays occasionally used “axiom” for declarations such as periodic phase, boundary exclusion, or absolute-value placement.

In the final logical classification:

\begin{itemize}
\item periodic phase is a coordinate convention inherited from trigonometric periodicity;
\item boundary exclusion is part of the domain definition;
\item absolute-value placement is part of the operator definitions;
\item the UNA sign convention is a naming convention.
\end{itemize}
None is an independent mathematical proposition comparable to the Pythagorean identity.


\PCsection{3. Canonical Definitions of Book 0}

\PCsubsection{Definition 3.1 --- Phase domain}

The analytic carrier is the lifted real line

\[
x\in\mathbb R.
\]

The operator family is \(2\pi\)-periodic. The identified phase circle is

\[
\mathbb T=\mathbb R/(2\pi\mathbb Z),
\]

with quotient relation

\[
x\equiv x+2\pi.
\]

Analytic expressions, domains, derivatives, and one-sided limits use the lifted coordinate in radians; finite rotation groups act on \(\mathbb T\).

\PCsubsection{Definition 3.2 --- Forbidden set}

\[
\mathbb B = \left\{ \frac{k\pi}{2}:k\in\mathbb Z \right\}.
\]

\PCsubsection{Definition 3.3 --- Admissible domain}

\[
\boxed{ \mathcal D = \mathbb R\setminus\mathbb B. }
\]

Thus, for \(x\in\mathcal D\),

\[
\sin x\neq0 \qquad\text{and}\qquad \cos x\neq0.
\]

\PCsubsection{Definition 3.4 --- Absolute reciprocal magnitudes}

\[
M(x)=|\sec x| = \frac1{|\cos x|},
\]

and

\[
N(x)=|\csc x| = \frac1{|\sin x|}.
\]

\PCsubsection{Definition 3.5 --- Primitive displayed operators}

\[
\boxed{ \operatorname{srx}(x)=N(x)+\cot x }
\]

\[
\boxed{ \operatorname{sxp}(x)=\frac1{\operatorname{srx}(x)} }
\]

\[
\boxed{ \operatorname{cxp}(x)=M(x)+\tan x }
\]

\[
\boxed{ \operatorname{crx}(x)=\frac1{\operatorname{cxp}(x)}. }
\]

\PCsubsection{Definition 3.6 --- UNA operators}

\[
\boxed{ \operatorname{urx} = \operatorname{srx}-\operatorname{crx} }
\]

and

\[
\boxed{ \operatorname{uxp} = \operatorname{cxp}-\operatorname{sxp}. }
\]

\PCsubsection{Definition 3.7 --- Saw responses}

\[
\boxed{ \operatorname{saw}_{r} = \frac1{\operatorname{urx}} }
\]

and

\[
\boxed{ \operatorname{saw}_{x} = \frac1{\operatorname{uxp}}. }
\]

\PCsubsection{Definition 3.8 --- Flatwave}

\[
\boxed{ \operatorname{Flatwave} = \operatorname{saw}_{r}+\operatorname{saw}_{x} }.
\]


\PCsubsection{Definition 3.9 --- Open quadrants}

\[
Q_k = \left( \frac{k\pi}{2}, \frac{(k+1)\pi}{2} \right).
\]

\PCsubsection{Definition 3.10 --- Quadrant sign}

For \(x\in Q_k\),

\[
\boxed{ \varepsilon_k = \operatorname{sgn}(\sin 2x) = (-1)^k. }
\]

\PCsubsection{Definition 3.11 --- Canonical generator}

On an open quadrant,

\[
\boxed{ z=\operatorname{srx}. }
\]

\PCsubsection{Definition 3.12 --- Transfer coordinates}

Let

\[
\alpha=\operatorname{sgn}(\sin x), \qquad \beta=\operatorname{sgn}(\cos x).
\]

Define

\[
p=\alpha\cos x-\beta\sin x,
\]

\[
q=\alpha\sin x+\beta\cos x,
\]

and

\[
\boxed{ \lambda=\frac{1-p}{2}. }
\]

\PCsubsection{Definition 3.13 --- Phase-tangential derivative}

For a differentiable function \(f\) on an open quadrant,

\[
T[f]=\frac{df}{dx}.
\]

For \(f>0\), define

\[
\widehat T[f] = \frac{f'}f.
\]

These definitions exhaust the canonical mathematical vocabulary of Book 0.


\PCsection{4. Primary Dependency Chain}

The central proof structure of Book 0 is:

\[
\boxed{ \text{Pythagorean identity} }
\]

\[
\Downarrow
\]

\[
\boxed{ \text{secant–tangent and cosecant–cotangent identities} }
\]

\[
\Downarrow
\]

\[
\boxed{ \text{reciprocal-conjugate forms} }
\]

\[
\Downarrow
\]

\[
\boxed{ \text{positivity and reciprocal pairing} }
\]

\[
\Downarrow
\]

\[
\boxed{ \text{UNA sum, product, and common sign} }
\]

\[
\Downarrow
\]

\[
\boxed{ \text{Flatwave and saw partition} }
\]

\[
\Downarrow
\]

\[
\boxed{ \text{one-generator reduction} }
\]

\[
\Downarrow
\]

\[
\boxed{ \text{differential and transfer structure} }
\]

\[
\Downarrow
\]

\[
\boxed{ \text{symmetry and boundary classification}. }
\]

No physical premise occurs in this chain.


\PCsection{5. Theorem Ledger for S0.I}

\PCsubsection{Theorem I.1 --- Reciprocal Pythagorean identities}

\[
\boxed{ M^2-\tan^2x=1 }
\]

and

\[
\boxed{ N^2-\cot^2x=1. }
\]

Status: Proved theorem.

Dependencies: Pythagorean identity and ordinary reciprocal definitions.

\PCsubsection{Theorem I.2 --- Reciprocal-conjugate forms}

\[
\boxed{ \operatorname{crx} = M-\tan x }
\]

and

\[
\boxed{ \operatorname{sxp} = N-\cot x. }
\]

Status: Proved theorem.

\PCsubsection{Corollary I.3 --- Exact reciprocal pairs}

\[
\boxed{ \operatorname{cxp}\operatorname{crx}=1 }
\]

and

\[
\boxed{ \operatorname{srx}\operatorname{sxp}=1. }
\]

Status: Proved corollary.

\PCsubsection{Corollary I.4 --- Primitive positivity}

\[
\boxed{ \operatorname{srx}, \operatorname{sxp}, \operatorname{cxp}, \operatorname{crx}>0 }
\]

on \(\mathcal D\).

Status: Proved corollary.

\PCsubsection{Theorem I.5 --- Common UNA sign}

\[
\boxed{ \operatorname{sgn}(\operatorname{urx}) = \operatorname{sgn}(\operatorname{uxp}) = \operatorname{sgn}(\sin 2x) }.
\]


Status: Proved theorem.

\PCsubsection{Theorem I.6 --- UNA sum}

\[
\boxed{ \operatorname{urx} + \operatorname{uxp} = \frac4{\sin 2x}. }
\]

Status: Proved theorem.

\PCsubsection{Corollary I.7 --- Inverse UNA sum}

\[
\boxed{ \frac1{\operatorname{urx}+\operatorname{uxp}} = \frac{\sin 2x}{4}. }
\]

Status: Proved corollary.

\PCsubsection{Theorem I.8 --- UNA product}

\[
\boxed{ \operatorname{urx}\operatorname{uxp} = \frac4{|\sin 2x|}. }
\]

Status: Proved theorem.

\PCsubsection{Theorem I.9 --- Exact Flatwave}

\[
\boxed{ \operatorname{Flatwave} = \operatorname{sgn}(\sin 2x). }
\]

Status: Proved theorem.

\PCsubsection{Corollary I.10 --- Flatwave magnitude}

\[
\boxed{ |\operatorname{Flatwave}|=1. }
\]

Status: Proved corollary.

\PCsubsection{Corollary I.11 --- Signed saw partition}

\[
\boxed{ \operatorname{saw}_{r} + \operatorname{saw}_{x} = \operatorname{sgn}(\sin 2x) }.
\]


Status: Proved corollary.

\PCsubsection{Corollary I.12 --- Unsigned saw partition}


\[
\boxed{ |\operatorname{saw}_{r}|+|\operatorname{saw}_{x}|=1 }.
\]


Status: Proved corollary.

\PCsubsection{Theorem I.13 --- Octant dominance cycle}

The unique dominant primitive on successive open octants follows

\[
\boxed{ \operatorname{srx} \rightarrow \operatorname{cxp} \rightarrow \operatorname{crx} \rightarrow \operatorname{sxp} \rightarrow \operatorname{srx}. }
\]

Status: Proved theorem.

Qualification: At odd multiples of \(\pi/4\), adjacent winners tie. Those points are admissible but are not open-octant interiors.


\PCsection{6. Theorem Ledger for S0.II}

\PCsubsection{Theorem II.1 --- One-generator primary reduction}

On a fixed open quadrant,

\[
\boxed{ \operatorname{cxp} = \frac{z+\varepsilon} {\varepsilon z-1}, }
\]

where

\[
z=\operatorname{srx}
\]

and

\[
\varepsilon=\operatorname{sgn}(\sin 2x).
\]

Status: Proved quadrant-local theorem.

\PCsubsection{Corollary II.2 --- Canonical primitive forms}

\[
\operatorname{srx}=z,
\]

\[
\operatorname{sxp}=\frac1z,
\]

\[
\operatorname{cxp} = \frac{z+\varepsilon}{\varepsilon z-1},
\]

and

\[
\operatorname{crx} = \frac{\varepsilon z-1}{z+\varepsilon}.
\]

Status: Proved quadrant-local corollary.

\PCsubsection{Theorem II.3 --- Generator range}

\[
\varepsilon=+1 \implies z>1,
\]

while

\[
\varepsilon=-1 \implies 0<z<1.
\]

Status: Proved theorem.

\PCsubsection{Corollary II.4 --- Branch sign from the generator}

\[
\boxed{ \varepsilon = \operatorname{sgn}(z-1) = \operatorname{sgn}(z^2-1). }
\]

Status: Proved corollary.

\PCsubsection{Theorem II.5 --- Canonical UNA forms}

\[
\boxed{ \operatorname{urx} = \frac{z^2+1}{z+\varepsilon} }
\]

and

\[
\boxed{ \operatorname{uxp} = \frac{z^2+1} {z(\varepsilon z-1)}. }
\]

Status: Proved theorem.

\PCsubsection{Theorem II.6 --- Canonical saw forms}

\[
\boxed{ \operatorname{saw}_{r} = \frac{z+\varepsilon}{z^2+1} }
\]

and

\[
\boxed{ \operatorname{saw}_{x} = \frac{z(\varepsilon z-1)}{z^2+1}. }
\]

Status: Proved theorem.

\PCsubsection{Theorem II.7 --- Rational double-angle readout}

\[
\boxed{ \sin 2x = \frac{4z(z^2-1)} {(z^2+1)^2}. }
\]

Status: Proved theorem.

\PCsubsection{Theorem II.8 --- Generator differential equation}

\[
\boxed{ z' = -\frac{1+z^2}{2}. }
\]

Status: Proved quadrant-local theorem.

\PCsubsection{Theorem II.9 --- Quadrant-local rational closure}

Every named Book 0 operator belongs to

\[
\boxed{ \mathbb R(z) }
\]

on a fixed open quadrant.

Status: Proved theorem.

\PCsubsection{Theorem II.10 --- Differential closure}

If

\[
F(z)\in\mathbb R(z),
\]

then

\[
\boxed{ \frac{d}{dx}F(z(x)) \in\mathbb R(z). }
\]

Status: Proved theorem.

\PCsubsection{Corollary II.11 --- Repeated differential closure}

For every nonnegative integer \(n\),

\[
\boxed{ \frac{d^n}{dx^n}F(z(x)) \in\mathbb R(z). }
\]

Status: Proved corollary.


\PCsection{7. Theorem Ledger for S0.III}

\PCsubsection{Theorem III.1 --- Primitive Riccati equations}

\[
\boxed{ \operatorname{srx}' = -\frac{1+\operatorname{srx}^2}{2} }
\]

\[
\boxed{ \operatorname{sxp}' = \frac{1+\operatorname{sxp}^2}{2} }
\]

\[
\boxed{ \operatorname{cxp}' = \frac{1+\operatorname{cxp}^2}{2} }
\]

\[
\boxed{ \operatorname{crx}' = -\frac{1+\operatorname{crx}^2}{2}. }
\]

Status: Proved quadrant-local theorems.

\PCsubsection{Theorem III.2 --- Normalized primitive rates}

\[
\boxed{ \widehat T[\operatorname{srx}] = -|\csc x| }
\]

\[
\boxed{ \widehat T[\operatorname{sxp}] = |\csc x| }
\]

\[
\boxed{ \widehat T[\operatorname{cxp}] = |\sec x| }
\]

\[
\boxed{ \widehat T[\operatorname{crx}] = -|\sec x|. }
\]

Status: Proved theorems.

\PCsubsection{Theorem III.3 --- Reciprocal rate reversal}

For \(f>0\),

\[
\boxed{ \widehat T[1/f] = -\widehat T[f]. }
\]

Status: Proved general theorem.

\PCsubsection{Theorem III.4 --- Transfer-circle identity}

\[
\boxed{ p^2+q^2=2. }
\]

Status: Proved theorem.

\PCsubsection{Theorem III.5 --- Transfer differential system}

\[
\boxed{ p'=-q }
\]

and

\[
\boxed{ q'=p. }
\]

Status: Proved theorem.

\PCsubsection{Corollary III.6 --- Harmonic transfer closure}

\[
\boxed{ p''=-p }
\]

and

\[
\boxed{ q''=-q. }
\]

Status: Proved corollary.

\PCsubsection{Theorem III.7 --- Transfer interval}

\[
\boxed{ 0<\lambda<1 }
\]

and

\[
\boxed{ \lambda'>0 }
\]

on every open quadrant.

Status: Proved theorem.

\PCsubsection{Corollary III.8 --- Transfer equation}

\[
\boxed{ \lambda''+\lambda=\frac12. }
\]

Status: Proved corollary.

\PCsubsection{Theorem III.9 --- Saw transfer representation}

\[
\boxed{ \operatorname{saw}_{r} = \varepsilon\lambda }
\]

and

\[
\boxed{ \operatorname{saw}_{x} = \varepsilon(1-\lambda). }
\]

Status: Proved theorem.

\PCsubsection{Theorem III.10 --- Equal-and-opposite saw rates}

\[
\boxed{ \operatorname{saw}_{r}' = -\operatorname{saw}_{x}'. }
\]

More exactly,

\[
\operatorname{saw}_{r}' = \varepsilon\frac q2,
\]

and

\[
\operatorname{saw}_{x}' = -\varepsilon\frac q2.
\]

Status: Proved theorem.

\PCsubsection{Corollary III.11 --- Flatwave tangential invariance}

\[
\boxed{ \operatorname{Flatwave}'=0 }
\]

inside every open quadrant.

Status: Proved corollary.

\PCsubsection{Theorem III.12 --- Bounded saw-transfer rate}

\[
\boxed{ \frac12 < |\operatorname{saw}_{r}'| = |\operatorname{saw}_{x}'| \leq \frac1{\sqrt2}. }
\]

Status: Proved theorem.

\PCsubsection{Theorem III.13 --- Saw curvature and midpoint inflection}

Each saw has exactly one inflection point within each open quadrant, located at its midpoint,

\[
x=\frac{(2k+1)\pi}{4}.
\]

Status: Proved theorem.


\PCsection{8. Theorem Ledger for S0.IV}

\PCsubsection{Theorem IV.1 --- Quarter-turn primitive action}

For

\[
R(x)=x+\frac{\pi}{2},
\]

\[
\boxed{ \operatorname{srx}(Rx)=\operatorname{crx}(x) }
\]

\[
\boxed{ \operatorname{sxp}(Rx)=\operatorname{cxp}(x) }
\]

\[
\boxed{ \operatorname{cxp}(Rx)=\operatorname{sxp}(x) }
\]

\[
\boxed{ \operatorname{crx}(Rx)=\operatorname{srx}(x). }
\]

Status: Proved theorem.

\PCsubsection{Theorem IV.2 --- Half-turn invariance}

For every primitive \(f\),

\[
\boxed{ f(x+\pi)=f(x). }
\]

Status: Proved theorem.

\PCsubsection{Theorem IV.3 --- Reflection as reciprocal exchange}

For

\[
J(x)=-x,
\]

\[
\boxed{ \operatorname{srx}(Jx)=\operatorname{sxp}(x) }
\]

and

\[
\boxed{ \operatorname{cxp}(Jx)=\operatorname{crx}(x), }
\]

with the reciprocal reverse relations also holding.

Status: Proved theorem.

\PCsubsection{Theorem IV.4 --- Phase symmetry group}

On \(\mathbb T\), and using the convention that \(D_4\) has eight elements, the quarter-turn and reflection generate

\[
\boxed{ D_4 = \langle R,J \mid R^4=J^2=e,\ JRJ=R^{-1} \rangle. }
\]

Under the convention used in Book 0, \(D_4\) has eight elements.

Status: Proved theorem.

\PCsubsection{Theorem IV.5 --- Induced operator group}

The primitive operator action has kernel

\[
\{e,R^2\}
\]

and image isomorphic to the Klein four-group:

\[
\boxed{ D_4/\{e,R^2\} \cong V_4. }
\]

Status: Proved theorem.

\PCsubsection{Theorem IV.6 --- Flatwave sign character}

\[
\boxed{ \operatorname{Flatwave}(Rx) = -\operatorname{Flatwave}(x) }
\]

\[
\boxed{ \operatorname{Flatwave}(Jx) = -\operatorname{Flatwave}(x) }.
\]


while diagonal reflection preserves Flatwave.

Status: Proved theorem.

\PCsubsection{Theorem IV.7 --- Transfer transformation}

Under quarter-turn,

\[
\lambda(Rx)=\lambda(x).
\]

Under reflection,

\[
\lambda(Jx)=1-\lambda(x).
\]

Status: Proved theorem.

\PCsubsection{Theorem IV.8 --- Octant-label rotation}

The eighth-turn

\[
H(x)=x+\frac{\pi}{4}
\]

advances the dominant label by

\[
\operatorname{srx} \rightarrow \operatorname{cxp} \rightarrow \operatorname{crx} \rightarrow \operatorname{sxp}.
\]

Status: Proved combinatorial theorem.

\PCsubsection{Negative result IV.9 --- Eighth-turn nonidentity}

In general,

\[
\operatorname{srx}\left(x+\frac{\pi}{4}\right) \neq \operatorname{cxp}(x).
\]

Status: Proved negative theorem.

Meaning: Dominance-label rotation is not the same as exact function permutation.


\PCsection{9. Theorem Ledger for S0.V}

\PCsubsection{Theorem V.1 --- Primitive zero–pole exchange}

At sine boundaries, the sine pair exchanges

\[
0 \longleftrightarrow +\infty,
\]

while the cosine pair approaches \(1\).

At cosine boundaries, the cosine pair exchanges

\[
+\infty \longleftrightarrow 0,
\]

while the sine pair approaches \(1\).

Status: Proved one-sided limit theorem.

\PCsubsection{Theorem V.2 --- First-order asymptotics}

Every primitive pole is first order in boundary distance, with leading behavior

\[
\frac2h.
\]

Every paired zero is first order, with leading behavior

\[
\frac h2.
\]

Status: Proved asymptotic theorem.

\PCsubsection{Theorem V.3 --- Transfer endpoint limits}

Within every open quadrant,

\[
\boxed{ \lambda:0\longrightarrow1. }
\]

Status: Proved one-sided limit theorem.

\PCsubsection{Theorem V.4 --- UNA boundary transfer}

Within a quadrant of sign \(\varepsilon\),

\[
(\operatorname{urx},\operatorname{uxp}) : (\varepsilon\infty,\varepsilon) \longrightarrow (\varepsilon,\varepsilon\infty).
\]

Status: Proved theorem.

\PCsubsection{Theorem V.5 --- Saw boundary transfer}

Within a quadrant,

\[
\boxed{ (\operatorname{saw}_{r},\operatorname{saw}_{x}) : (0,\varepsilon) \longrightarrow (\varepsilon,0). }
\]

Status: Proved theorem.

\PCsubsection{Theorem V.6 --- Flatwave sign jump}

At every excluded boundary \(b\),

\[
\boxed{ \lim_{x\to b^-}\operatorname{Flatwave}(x) = - \lim_{x\to b^+}\operatorname{Flatwave}(x) }.
\]


Status: Proved theorem.

\PCsubsection{Corollary V.7 --- No continuous global Flatwave extension}

Flatwave has no continuous real-valued extension across an identified quadrant boundary.

Status: Proved corollary.

\PCsubsection{Theorem V.8 --- Componentwise quadrant compactification}

The bounded transfer coordinate and saw responses extend continuously to each closed quadrant. This is a compactification of each individual component, not of the full lifted domain.

Status: Proved theorem.

\PCsubsection{Theorem V.9 --- Split-boundary extension domain}

The bounded variables extend continuously to

\[
\widehat{\mathcal D} = \bigsqcup_{k\in\mathbb Z}\overline Q_k.
\]


Status: Proved theorem.

\PCsubsection{Negative result V.10 --- No canonical seam dynamics}

The interior identities and one-sided boundary limits do not uniquely determine a boundary-crossing rule.

Status: Proved underdetermination theorem.


\PCsection{10. Minimality of the Primitive Presentation}

The word minimality can refer to several different questions. They must not be conflated.

\PCsubsection{Definition 10.1 --- Display minimality}

A displayed family is minimal when no displayed member may be removed without losing a desired symmetry, naming pattern, or visual role.

\PCsubsection{Definition 10.2 --- Algebraic generator minimality}

A set of functions is algebraically minimal when no smaller set generates the same expression field under the declared operations.

\PCsubsection{Definition 10.3 --- Axiomatic minimality}

An axiom set is minimal when no axiom can be removed without losing at least one theorem.

\PCsubsection{Definition 10.4 --- Interpretive minimality}

An interpretive vocabulary is minimal when each label contributes meaning not supplied by another label.

Interpretive minimality is outside the scope of Book 0 mathematics.

\PCsubsection{Theorem 10.5 --- The four displayed primitives are not algebraically independent}

The four displayed operators contain two exact reciprocal dependencies:

\[
\operatorname{sxp} = \frac1{\operatorname{srx}},
\]

and

\[
\operatorname{crx} = \frac1{\operatorname{cxp}}.
\]

Therefore, the four displayed primitives are not algebraically independent.

\PCsubsection{Corollary 10.6 --- Two primary envelopes suffice before quadrant reduction}

If reciprocation is allowed, the displayed primitive family is generated by

\[
\operatorname{srx} \qquad\text{and}\qquad \operatorname{cxp}.
\]

Thus two generators suffice without using the Flatwave branch relation.

\PCsubsection{Theorem 10.7 --- One generator suffices quadrant-locally}

On a fixed open quadrant, every named Book 0 operator is rational in

\[
z=\operatorname{srx}.
\]

Therefore, one nonconstant generator is sufficient quadrant-locally.

\begin{proof}
S0.II supplies explicit rational expressions for every named operator in terms of \(z\) and the constant branch sign \(\varepsilon\).

Once a quadrant is fixed, \(\varepsilon\) is a real constant equal to \(+1\) or \(-1\).

Therefore, every named expression belongs to \(\mathbb R(z)\).

\end{proof}

\PCsubsection{Theorem 10.8 --- Zero nonconstant generators are insufficient}

A field generated over \(\mathbb R\) by no nonconstant function contains only constant functions.

Since

\[
z'=-\frac{1+z^2}{2}<0,
\]

the function \(z\) is nonconstant on every open quadrant.

Therefore, zero nonconstant generators cannot produce the Book 0 calculus.

\PCsubsection{Corollary 10.9 --- Quadrant-local generator-count minimality}

Exactly one nonconstant generator is necessary and sufficient on each open quadrant.

\[
\boxed{ \text{Minimal quadrant-local generator count}=1. }
\]

\PCsubsection{Clarification 10.10 --- The generator is not uniquely determined}

The theorem proves the number of required nonconstant generators.

It does not prove that

\[
z=\operatorname{srx}
\]

is the only possible generator.

For example, any invertible rational transform of \(z\) may generate the same field where the inverse transform is defined.

Possible alternative generators include other primitive operators on appropriate quadrants.

Therefore:

\[
\boxed{ \text{generator-count minimality does not imply generator uniqueness.} }
\]

\PCsubsection{Clarification 10.11 --- Global generation requires branch data}

The formula

\[
\operatorname{cxp} = \frac{z+\varepsilon}{\varepsilon z-1}
\]

contains the quadrant sign \(\varepsilon\).

Although

\[
\varepsilon=\operatorname{sgn}(z-1),
\]

the sign operation is not rational.

Thus the global calculus is not a single ordinary rational field in \(z\) without branch qualification.

A global presentation requires at least one of:

\begin{itemize}
\item a quadrant index;
\item the branch sign \(\varepsilon\);
\item a piecewise expression rule;
\item or an enlarged function class containing sign functions.
\end{itemize}

\PCsection{11. Correct Meaning of Closure}

\PCsubsection{Definition 11.1 --- Finite named catalogue}

Let

\[
\mathcal N
\]

denote the finite list consisting of the four primitive operators, the two UNA operators, the two saw responses, and Flatwave.

\PCsubsection{Proposition 11.2 --- The finite named catalogue is not closed under unrestricted rational operations}

For example, the expression

\[
\operatorname{srx}+\operatorname{cxp}
\]

is generally not identical to any single member of \(\mathcal N\).

Likewise,

\[
\operatorname{srx}^2, \qquad \frac{\operatorname{srx}+2} {\operatorname{cxp}-3},
\]

and indefinitely many other rational expressions need not reduce to one of the nine named functions.

Therefore, \(\mathcal N\) is not a finite algebra closed under unrestricted rational operations.

\PCsubsection{Definition 11.3 --- Quadrant-local rational field}

On a fixed open quadrant, define

\[
\boxed{ \mathcal F_k=\mathbb R(z). }
\]

Elements of \(\mathcal F_k\) are rational functions

\[
F(z)=\frac{P(z)}{Q(z)}
\]

with real polynomial numerator and denominator, evaluated where \(Q(z)\neq0\).

\PCsubsection{Definition 11.4 --- Phase derivation}

Define

\[
\boxed{ \mathcal D_z = -\frac{1+z^2}{2}\frac{d}{dz}. }
\]

Then

\[
\frac{d}{dx}F(z(x)) = \mathcal D_zF.
\]

\PCsubsection{Theorem 11.5 --- Differential-field closure}

For every fixed open quadrant,

\[
\boxed{ \mathcal D_z(\mathcal F_k) \subseteq \mathcal F_k. }
\]

Therefore,

\[
\boxed{ (\mathbb R(z),\mathcal D_z) }
\]

is the correct quadrant-local differential closure object.

\PCsubsection{Definition 11.6 --- Global piecewise-rational class}

Define

\[
\mathscr F = \left\{ F:\mathcal D\to\mathbb R: F|_{Q_k}\in\mathbb R(z) \text{ for every }k \right\},
\]


subject to the requirement that every denominator remain nonzero on its declared evaluation set.

\PCsubsection{Theorem 11.7 --- Component-wise global closure}

The class \(\mathscr F\) is closed under:

\begin{itemize}
\item addition;
\item subtraction;
\item multiplication;
\item division where the divisor is nonzero;
\item and differentiation on each open quadrant.
\end{itemize}
\begin{proof}
Each operation is performed separately on every connected component \(Q_k\).

The rational functions are closed under the four field operations where denominators are nonzero.

S0.II proves closure under differentiation.

No equality across an excluded boundary is required.

Therefore, \(\mathscr F\) is component-wise closed.

\end{proof}

\PCsubsection{Final closure statement}

The defensible closure theorem is:

\begin{center}
\fbox{\parbox{0.92\linewidth}{\centering The Book 0 calculus is a piecewise rational differential calculus, not a finite closed list of named functions.}}
\end{center}


\PCsection{12. Uniqueness Audit}

\PCsubsection{12.1 Reciprocal-conjugate uniqueness}

Given the definitions

\[
\operatorname{cxp}=M+\tan x
\]

and

\[
M^2-\tan^2x=1,
\]

the reciprocal is uniquely determined:

\[
\frac1{\operatorname{cxp}} = M-\tan x.
\]

Likewise,

\[
\frac1{\operatorname{srx}} = N-\cot x.
\]


This is ordinary functional uniqueness of a reciprocal.

\PCsubsection{12.2 Flatwave value}

Once Flatwave is defined as

\[
\frac1{\operatorname{urx}} + \frac1{\operatorname{uxp}},
\]

its value is uniquely determined:

\[
\operatorname{Flatwave} = \operatorname{sgn}(\sin 2x).
\]

This is a theorem about the defined expression.

\PCsubsection{12.3 Flatwave-form uniqueness is not proved}

Book 0 does not prove that Flatwave is the only possible composite satisfying some broad list of desired properties.

There may exist many functions that:

\begin{itemize}
\item equal \(+1\) on Quadrants I and III;
\item equal \(-1\) on Quadrants II and IV;
\item are invariant under half-turn;
\item reverse under quarter-turn.
\end{itemize}
Indeed,

\[
\operatorname{sgn}(\sin 2x)
\]

itself may be represented in multiple algebraically equivalent forms.

A nontrivial uniqueness theorem would require a restricted candidate class, such as:

\begin{itemize}
\item symmetric rational expressions of bounded degree;
\item expressions built from specified generators;
\item fixed pole orders;
\item fixed normalization;
\item fixed behavior under reciprocal exchange;
\item fixed boundary asymptotics.
\end{itemize}
No such complete classification is proved in Book 0.

\PCsubsection{12.4 Primitive-family uniqueness is not proved}

Book 0 does not prove that the four named operators are the only positive reciprocal trigonometric family with related properties.

Such a theorem would require a formally declared search space and equivalence relation.

\PCsubsection{Final uniqueness statement}

\[
\boxed{ \text{Book 0 proves exact evaluation, not universal classification uniqueness.} }
\]


\PCsection{13. The Status of “Non-Additivity”}

Addition occurs throughout the core definitions and proofs:

\[
\operatorname{srx}=N+\cot x,
\]

\[
\operatorname{cxp}=M+\tan x,
\]

\[
\operatorname{Flatwave} = \operatorname{saw}_{r}+\operatorname{saw}_{x},
\]

and

\[
\operatorname{urx}+\operatorname{uxp} = \frac4{\sin 2x}.
\]

Therefore, a general axiom forbidding addition would contradict the formal construction.

A narrower modeling rule may later prohibit arbitrary superposition in a particular projection or physical interpretation. Such a rule would belong to a later book.

The mathematical conclusion is:

\[
\boxed{ \text{Book 0 contains no general non-additivity axiom.} }
\]


\PCsection{14. The Status of Growth and Decay Language}

The following derivative signs are exact:

\[
\operatorname{srx}'<0, \qquad \operatorname{crx}'<0,
\]

\[
\operatorname{sxp}'>0, \qquad \operatorname{cxp}'>0.
\]

Therefore, with respect to increasing phase \(x\):

\begin{itemize}
\item the inverex pair decreases;
\item the expon pair increases.
\end{itemize}
This justifies the restricted mathematical vocabulary:

\[
\text{phase-increasing} \qquad\text{and}\qquad \text{phase-decreasing}.
\]

It does not yet justify:

\begin{itemize}
\item temporal growth;
\item temporal decay;
\item causal production;
\item dissipation;
\item biological growth;
\item economic expansion;
\item cosmological development.
\end{itemize}
Those require a map

\[
x=x(\tau)
\]

or another declared projection relating phase to the external variable of interest.

The Book 0 status is:

\[
\boxed{ \text{growth and decay are phase-direction labels only.} }
\]


\PCsection{15. The Status of Conservation Language}

Book 0 proves:

\[
\operatorname{saw}_{r}+\operatorname{saw}_{x}=\varepsilon,
\]

and

\[
|\operatorname{saw}_{r}|+|\operatorname{saw}_{x}|=1.
\]


It also proves:

\[
\operatorname{saw}_{r}' + \operatorname{saw}_{x}' = 0
\]


inside every open quadrant.

These are exact normalized partition identities.

They justify the phrase:

\[
\boxed{ \text{signed partition invariant}. }
\]

They do not by themselves prove conservation of:

\begin{itemize}
\item mass;
\item energy;
\item momentum;
\item electric charge;
\item probability;
\item thermodynamic information;
\item algorithmic information.
\end{itemize}
To establish a physical conservation law, a later theory must provide:

\begin{enumerate}
\item a state space;
\item a physical quantity \(C\);
\item a projection identifying \(C\) with a function of the saw variables;
\item an evolution parameter;
\item a proof that \(dC/d\tau=0\);
\item dimensional consistency;
\item empirical validation.
\end{enumerate}
Therefore:

\begin{center}
\fbox{\parbox{0.92\linewidth}{\centering Flatwave is mathematically invariant within a quadrant, but no physical conservation law is proved in Book 0.}}
\end{center}


\PCsection{16. The Status of Boundaries and Phase Change}

Book 0 proves exact one-sided limits.

It does not assign values to the original operators at

\[
x=\frac{k\pi}{2}.
\]

It proves that:

\begin{itemize}
\item primitive zero–pole exchanges occur;
\item the saw responses have finite one-sided limits;
\item Flatwave reverses sign;
\item the transfer coordinate runs from \(0\) to \(1\);
\item adjacent boundary states are incompatible under ordinary identification.
\end{itemize}
The term phase reversal may therefore denote the exact sign change

\[
\varepsilon\longrightarrow-\varepsilon.
\]

It does not establish:

\begin{itemize}
\item a singularity in physical spacetime;
\item a Big Bang;
\item a heat death;
\item a bounce;
\item annihilation;
\item creation;
\item a mirror universe;
\item reversal of physical time.
\end{itemize}
Any rule mapping one exit state to another entry state is a seam rule introduced after Book 0.

Therefore:

\[
\boxed{ \text{Book 0 proves boundary approach, not boundary passage.} }
\]


\PCsection{17. Domain Ledger}

Every theorem in Book 0 must be read with its domain.

\PCsubsection{17.1 Global admissible-domain theorems}

The following hold for every \(x\in\mathcal D\):

\begin{itemize}
\item reciprocal-conjugate identities;
\item primitive positivity;
\item reciprocal products;
\item UNA sum and product;
\item common UNA sign;
\item exact Flatwave;
\item signed and unsigned saw partitions.
\end{itemize}
\PCsubsection{17.2 Quadrant-local theorems}

The following require a fixed open quadrant:

\begin{itemize}
\item constant \(\varepsilon\);
\item one-generator rational reduction;
\item rational differential-field closure;
\item ordinary differentiation without crossing a boundary;
\item transfer-coordinate monotonicity;
\item saw-transfer rates;
\item curvature classification.
\end{itemize}
\PCsubsection{17.3 Octant-local theorems}

The following are naturally stated on open octants:

\begin{itemize}
\item unique primitive dominance;
\item fixed saw-curvature sign.
\end{itemize}
\PCsubsection{17.4 Boundary theorems}

The following are one-sided statements:

\begin{itemize}
\item primitive zero–pole limits;
\item UNA asymptotics;
\item saw endpoint limits;
\item Flatwave sign jumps;
\item componentwise quadrant compactification.
\end{itemize}
\PCsubsection{17.5 Split-domain theorems}

Continuous boundary extensions of the bounded variables hold on

\[
\widehat{\mathcal D} = \bigsqcup_k\overline Q_k,
\]

not on the ordinary identified phase circle.


\PCsection{18. Book 0 Negative Theorems}

Negative theorems are important because they prevent later interpretation from being mistaken for deduction.

\PCsubsection{Negative Theorem 18.1 --- No finite named algebra}

The nine named expressions do not form a finite algebra under unrestricted rational operations.

\PCsubsection{Negative Theorem 18.2 --- No global rational single branch}

The complete global calculus is not represented by one branch-independent rational formula in \(z\) alone without a sign, quadrant, or piecewise rule.

\PCsubsection{Negative Theorem 18.3 --- No universal uniqueness theorem}

Book 0 does not classify all possible reciprocal operator families or all possible Flatwave-like expressions.

\PCsubsection{Negative Theorem 18.4 --- No continuously varying Flatwave interior}

Inside every open quadrant,

\[
\operatorname{Flatwave}'=0.
\]

Therefore, Flatwave itself does not encode a continuously varying interior magnitude.

\PCsubsection{Negative Theorem 18.5 --- No continuous Flatwave seam}

Flatwave cannot be continuously extended across an identified quadrant wall.

\PCsubsection{Negative Theorem 18.6 --- No canonical seam map}

The core calculus does not determine how a state passes from one quadrant to another.

\PCsubsection{Negative Theorem 18.7 --- No physical ontology}

Book 0 does not establish physical meanings for its mathematical operators.

\PCsubsection{Negative Theorem 18.8 --- No temporal dynamics}

Differentiation with respect to \(x\) is not, without further assumptions, differentiation with respect to physical time.


\PCsection{19. Conjecture Ledger}

Book 0 leaves the following possible research questions open.

\PCsubsection{Conjecture C.1 --- Restricted Flatwave uniqueness}

Within a precisely declared class of symmetric rational expressions of bounded degree, Flatwave may be unique up to algebraic equivalence and normalization.

No proof is currently supplied.

\PCsubsection{Conjecture C.2 --- Classification of positive reciprocal families}

The R-family may belong to a broader class of positive reciprocal half-angle constructions admitting comparable partition identities.

A complete classification is not supplied.

\PCsubsection{Conjecture C.3 --- Canonical generator classification}

The generators of the quadrant-local field \(\mathbb R(z)\) may be classifiable under Möbius equivalence.

Book 0 proves only the sufficiency of \(z=\operatorname{srx}\).

\PCsubsection{Conjecture C.4 --- Natural seam selection}

A later dynamical or topological extension may select a seam map satisfying additional symmetry, continuity, reversibility, or optimization conditions.

No seam is selected by Book 0 alone.

\PCsubsection{Conjecture C.5 --- External projection utility}

The transfer coordinate, saw partition, or primitive rates may provide useful projections into external systems.

This is a research hypothesis requiring the projection discipline of Book 1.


\PCsection{20. Book 0 Formal Summary}

\PCsubsection{20.1 Mathematical carrier}

\[
\boxed{ \mathcal D = \mathbb R \setminus \left\{ \frac{k\pi}{2}:k\in\mathbb Z \right\}. }
\]

\PCsubsection{20.2 Primitive reciprocal family}

\[
\boxed{\operatorname{srx}=|\csc x|+\cot x}
\]

\[
\boxed{\operatorname{sxp}=|\csc x|-\cot x=\frac1{\operatorname{srx}}}
\]

\[
\boxed{\operatorname{cxp}=|\sec x|+\tan x}
\]

\[
\boxed{\operatorname{crx}=|\sec x|-\tan x=\frac1{\operatorname{cxp}}}
\]


\PCsubsection{20.3 UNA closure}

\[
\boxed{ \operatorname{urx} + \operatorname{uxp} = \frac4{\sin 2x} }
\]

and

\[
\boxed{ \operatorname{urx}\operatorname{uxp} = \frac4{|\sin 2x|}. }
\]

\PCsubsection{20.4 Flatwave}

\[
\boxed{ \operatorname{Flatwave} = \operatorname{sgn}(\sin 2x). }
\]

\PCsubsection{20.5 Saw partition}

\[
\boxed{ \operatorname{saw}_{r}=\varepsilon\lambda }
\]

\[
\boxed{ \operatorname{saw}_{x}=\varepsilon(1-\lambda) }
\]

and

\[
\boxed{ |\operatorname{saw}_{r}|+|\operatorname{saw}_{x}|=1 }
\]


\PCsubsection{20.6 Generator reduction}

On a fixed open quadrant,

\[
\boxed{ \mathcal N\subset\mathbb R(z), \qquad z=\operatorname{srx}. }
\]

\PCsubsection{20.7 Differential law}

\[
\boxed{ z' = -\frac{1+z^2}{2}. }
\]

Therefore,

\[
\boxed{ \mathcal D_z = -\frac{1+z^2}{2}\frac{d}{dz} }
\]

preserves \(\mathbb R(z)\).

\PCsubsection{20.8 Transfer law}

\[
\boxed{ p^2+q^2=2, \qquad p'=-q, \qquad q'=p. }
\]

\PCsubsection{20.9 Symmetry}

The phase transformations generate \(D_4\), while the induced primitive operator action is

\[
\boxed{ V_4. }
\]

\PCsubsection{20.10 Boundary structure}

Within each quadrant,

\[
\boxed{ \lambda:0\to1 }
\]

and

\[
\boxed{ (\operatorname{saw}_{r},\operatorname{saw}_{x}) : (0,\varepsilon)\to(\varepsilon,0). }
\]

Across each identified wall,

\[
\boxed{ \operatorname{Flatwave} \text{ reverses sign}. }
\]


\PCsection{21. Book 0 Certification}

Book 0 establishes a rigorous mathematical calculus with the following certified properties.

\PCsubsubsection{Certified C0.1 --- Exactness}

The principal identities are algebraic equalities, not numerical approximations.

\PCsubsubsection{Certified C0.2 --- Domain control}

Every theorem has an explicit admissible domain or one-sided boundary domain.

\PCsubsubsection{Certified C0.3 --- Minimal generator count}

One nonconstant generator is necessary and sufficient quadrant-locally.

\PCsubsubsection{Certified C0.4 --- Differential closure}

The quadrant-local rational field is closed under phase differentiation.

\PCsubsubsection{Certified C0.5 --- Complementary transfer}

The saw responses form an exact signed and unsigned partition.

\PCsubsubsection{Certified C0.6 --- Discrete symmetry}

The primitive family is a symmetry orbit under the induced Klein action, while octant dominance follows a four-label cycle.

\PCsubsubsection{Certified C0.7 --- Boundary classification}

Primitive poles, reciprocal zeros, saw limits, Flatwave jumps, and split-boundary extensions are explicitly classified.

\PCsubsubsection{Certified C0.8 --- Interpretive restraint}

No physical law is included among the certified Book 0 results.


\PCsection{22. Terminal Theorem of Book 0}

\PCsubsection{Theorem 22.1 --- Foundational closure of Book 0}

Relative to ordinary real trigonometry and elementary real analysis, the bounded phase operator calculus of Book 0 is completely specified by its definitions.

On every open quadrant:

\begin{enumerate}
\item all named operators reduce to rational functions of one nonconstant generator;
\item the resulting rational field is closed under phase differentiation;
\item the saw responses form a complementary signed partition;
\item the primitive and transfer structures obey exact discrete symmetries;
\item every approach to an excluded boundary has a classified one-sided limit.
\end{enumerate}
The calculus does not determine a physical interpretation, dimensional calibration, temporal evolution, or seam-crossing law.

\begin{proof}
Statements 1–5 are the principal results of S0.I–S0.V and are recorded in the theorem ledger above.

The absence of physical interpretation, calibration, temporal evolution, and seam dynamics follows from the fact that no such structures are included among the Book 0 definitions or inherited mathematical background.

Therefore, Book 0 is mathematically closed at the level of its declared structural purpose, while remaining intentionally open to later projection and model layers.


\end{proof}

\PCsection{23. Closing Declaration}

Book 0 began with the Pythagorean identity

\[
\sin^2x+\cos^2x=1.
\]

From that starting point, it derived:

\[
\boxed{ \text{positive reciprocal operators} }
\]

\[
\boxed{ \text{signed UNA differences} }
\]

\[
\boxed{ \text{a double-angle closure} }
\]

\[
\boxed{ \text{a unit Flatwave} }
\]

\[
\boxed{ \text{complementary saw transfer} }
\]

\[
\boxed{ \text{one-generator rational reduction} }
\]

\[
\boxed{ \text{differential-field closure} }
\]

\[
\boxed{ \text{phase and operator symmetry} }
\]

and

\[
\boxed{ \text{exact boundary classification}. }
\]

Its strongest justified conclusion is:

\[
\boxed{ \begin{gathered} \text{The bounded phase operators form a piecewise rational}\\ \text{differential calculus generated quadrant-locally by one}\\ \text{nonconstant function, with exact reciprocal, transfer,}\\ \text{symmetry, and boundary structures.} \end{gathered} }
\]

Its strongest necessary restraint is:

\[
\boxed{ \begin{gathered} \text{No physical conservation law, evolution law,}\\ \text{measurement rule, or ontology follows from}\\ \text{the structural calculus alone.} \end{gathered} }
\]

This completes Book 0.

\[
\boxed{ \text{END OF BOOK 0} }
\]

% ===== END S0_VI =====
\part*{Book 1: Projection and Measurement}
\addcontentsline{toc}{part}{Book 1: Projection and Measurement}
% ===== BEGIN S1_I =====
% Source: S1.i.docx
\PCchapter{S1.I}{Structural Mathematics and Projection Contracts}{Book 1 - Projection and Measurement}

\PCsubsection{Abstract}

Book 0 established a dimensionless bounded phase operator calculus on the admissible domain

\[
\mathcal D = \mathbb R \setminus \left\{ \frac{k\pi}{2}:k\in\mathbb Z \right\}.
\]

Its theorems concern exact relations among phase functions. They do not, by themselves, identify external objects, measurable quantities, physical units, or empirical observations.

This essay introduces the projection layer.

A projection is defined as an explicit composition

\[
\Omega_{\mathcal P} \xrightarrow{\ \Phi\ } \mathcal D \xrightarrow{\ \mathbf F\ } \mathbb R^m \xrightarrow{\ \Gamma_\theta\ } \mathcal Y,
\]

where:

\begin{itemize}
\item \(\Omega_{\mathcal P}\) is the declared valid subset of an external domain;
\item \(\Phi\) assigns an admissible phase to each valid external instance;
\item \(\mathbf F\) evaluates one or more Book 0 structural expressions;
\item \(\Gamma_\theta\) converts the resulting dimensionless structural profile into a declared output space;
\item \(\theta\) contains every calibration parameter, scale, unit, or convention required by the readout.
\end{itemize}
The complete projection contract also declares uncertainty and failure behavior.

It is proved that exact structural identities remain exact after lawful composition. The converse is weaker: agreement of projected outputs implies structural equality only when the readout is injective and the inferred phases cover the structural domain under examination.

It is further proved that all dimensional content enters through the external domain, the calibration parameters, or the readout map. Book 0 alone supplies no dimensional scale.

A boundary guard is established: an interior Book 0 expression may not be evaluated when the inferred phase lies in the forbidden set. Boundary-aware projections must instead use the split-boundary domain and the bounded extensions proved in S0.V.

Finally, an admissible reparameterization theorem shows that different phase assignments and structural profiles can generate identical external predictions. Projection data therefore do not uniquely determine the internal phase description without additional identification conditions.

The projection contract is the formal boundary between structural mathematics and external interpretation.


\PCsection{IDENTIFICATION DIVISION}

\PCsubsection{PROGRAM-ID}

S1-I-PROJECTION-CONTRACTS

\PCsubsection{PURPOSE}

To define the minimal formal machinery required to apply the dimensionless structural results of Book 0 to an external domain.

\PCsubsection{PRINCIPAL QUESTION}

Given an external object or observation \(\omega\), under what declared conditions may a Book 0 expression be used to produce an external output?

\PCsubsection{NON-PURPOSE}

This essay does not:

\begin{itemize}
\item identify a particular physical system;
\item define time, distance, mass, energy, charge, or information;
\item select a phase-inference rule for any empirical domain;
\item calibrate any physical constant;
\item prove that an external system is governed by the R-operators;
\item introduce a boundary-crossing dynamics;
\item establish that a projection is unique;
\item replace empirical validation with structural algebra.
\end{itemize}

\PCsection{ENVIRONMENT DIVISION}

\PCsubsection{1. Inherited Structural Environment}

All definitions and theorems of Book 0 are inherited.

\PCsubsection{Definition 1.1 --- Interior structural domain}

\[
\boxed{ \mathcal D = \mathbb R \setminus \left\{ \frac{k\pi}{2}:k\in\mathbb Z \right\}. }
\]

\PCsubsection{Definition 1.2 --- Forbidden boundary set}

\[
\boxed{ \mathbb B = \left\{ \frac{k\pi}{2}:k\in\mathbb Z \right\}. }
\]

\PCsubsection{Definition 1.3 --- Split-boundary domain}

From S0.V, define

\[
\boxed{ \widehat{\mathcal D} = \bigsqcup_{k\in\mathbb Z}\overline Q_k, }
\]

where each entering and exiting copy of a numerical boundary remains distinct.

\PCsubsection{Definition 1.4 --- Structural expression class}

Let

\[
\mathscr F
\]

denote the global piecewise-rational differential class of Book 0.

Thus, for every open quadrant \(Q_k\),

\[
F|_{Q_k}\in\mathbb R(z)
\]

for each

\[
F\in\mathscr F,
\]

where

\[
z=\operatorname{srx}.
\]

\PCsubsection{Definition 1.5 --- Structural profile}

For a positive integer \(m\), a structural profile is an ordered tuple

\[
\boxed{ \mathbf F = (F_1,\ldots,F_m) }
\]

with

\[
F_i\in\mathscr F.
\]

The profile defines a map

\[
\mathbf F:\mathcal D_{\mathbf F}\to\mathbb R^m,
\]

where

\[
\mathcal D_{\mathbf F} \subseteq \mathcal D
\]

is the common domain on which every component \(F_i\) is defined.

The common domain may be smaller than \(\mathcal D\) if a composite expression introduces additional denominator restrictions.


\PCsection{AXIOM DIVISION}

Book 1 introduces one methodological axiom.

\PCsubsection{Axiom P.1 --- Explicit Projection Mediation}

Every claim assigning Book 0 structure to an external object or observation must be mediated by an explicitly declared projection contract.

Equivalently, no external interpretation is admitted merely because a Book 0 expression resembles, correlates with, or can be verbally associated with an external phenomenon.

The external claim must declare:

\begin{enumerate}
\item what is being projected;
\item how phase is inferred;
\item which structural expressions are evaluated;
\item how structural values become external outputs;
\item where the projection is valid;
\item what uncertainty applies;
\item what happens when the contract fails.
\end{enumerate}
This is a methodological axiom governing the S Essay series. It is not a theorem of trigonometry.


\PCsection{DATA DIVISION}

\PCsection{2. External Objects and Outputs}

\PCsubsection{Definition 2.1 --- External instance space}

Let

\[
\Omega
\]

be a nonempty set whose elements represent external instances.

An element

\[
\omega\in\Omega
\]

may represent an observation, state, object, event, record, configuration, or other externally specified item.

No topology, metric, probability distribution, temporal order, or physical interpretation is assumed merely by declaring \(\Omega\).

\PCsubsection{Definition 2.2 --- Output space}

Let

\[
\mathcal Y
\]

be the declared output space of a projection.

Depending on the contract, \(\mathcal Y\) may be:

\begin{itemize}
\item a subset of \(\mathbb R\);
\item a vector space;
\item a set of labels;
\item an interval;
\item a probability space;
\item or another explicitly defined set.
\end{itemize}
The structure of \(\mathcal Y\) must be declared before it is used.

\PCsubsection{Definition 2.3 --- Validity domain}

A projection contract declares a subset

\[
\boxed{ \Omega_{\mathcal P}\subseteq\Omega }
\]

called its declared validity domain.

This is a scope declaration, not a claim that the domain is maximal. S1.II distinguishes the declared domain from the maximal operational admissibility envelope determined by the non-domain guards.

The contract makes no prediction outside \(\Omega_{\mathcal P}\), except through an explicitly declared failure rule.


\PCsection{3. Phase Inference}

\PCsubsection{Definition 3.1 --- Interior phase-inference map}

An interior phase-inference map is a function

\[
\boxed{ \Phi: \Omega_{\mathcal P}\to\mathcal D. }
\]

For each valid external instance \(\omega\),

\[
\Phi(\omega)
\]

is the phase at which the selected Book 0 structural profile is evaluated.

\PCsubsection{Definition 3.2 --- Boundary-aware phase-inference map}

A boundary-aware phase-inference map is a function

\[
\boxed{ \widehat\Phi: \Omega_{\mathcal P}\to\widehat{\mathcal D}. }
\]

Such a map distinguishes entering and exiting copies of a boundary.

A boundary-aware contract may use only those Book 0 quantities for which the required extensions to \(\widehat{\mathcal D}\) have been proved.

\PCsubsection{Definition 3.3 --- Phase-inference rule}

The phase-inference rule consists of:

\begin{enumerate}
\item the map \(\Phi\) or \(\widehat\Phi\);
\item the data from \(\omega\) used to determine phase;
\item the method of determination;
\item any branch-selection rule;
\item any parameters used;
\item the conditions under which phase inference fails.
\end{enumerate}
A numerical phase assignment without this declaration is not a complete projection rule.


\PCsection{4. Structural Evaluation}

\PCsubsection{Definition 4.1 --- Structural evaluation map}

Given a structural profile

\[
\mathbf F=(F_1,\ldots,F_m),
\]

define

\[
\boxed{ E_{\mathbf F}(x) = \bigl( F_1(x),\ldots,F_m(x) \bigr). }
\]

Thus,

\[
E_{\mathbf F}: \mathcal D_{\mathbf F}\to\mathbb R^m.
\]

\PCsubsection{Definition 4.2 --- Evaluated structural state}

For

\[
\omega\in\Omega_{\mathcal P},
\]

the evaluated structural state is

\[
\boxed{ \mathbf s_{\mathcal P}(\omega) = E_{\mathbf F}(\Phi(\omega)). }
\]

This value is defined only when

\[
\Phi(\omega)\in\mathcal D_{\mathbf F}.
\]


\PCsection{5. Readout and Calibration}

\PCsubsection{Definition 5.1 --- Calibration parameter}

Let

\[
\theta\in\Theta
\]

denote the complete collection of parameters used by the readout map.

The parameter object \(\theta\) may contain:

\begin{itemize}
\item scale factors;
\item offsets;
\item unit definitions;
\item reference values;
\item empirical coefficients;
\item thresholds;
\item orientation conventions;
\item fitted parameters;
\item conversion constants.
\end{itemize}
Every such element must be declared.

\PCsubsection{Definition 5.2 --- Readout map}

A readout map is a function

\[
\boxed{ \Gamma_\theta: \mathcal S_{\mathcal P}\to\mathcal Y, }
\]

where

\[
\mathcal S_{\mathcal P} \subseteq\mathbb R^m
\]

contains the structural states produced by the contract.

The readout map converts dimensionless structural values into the declared external output.

\PCsubsection{Definition 5.3 --- Calibration rule}

A calibration rule specifies:

\begin{enumerate}
\item the parameter space \(\Theta\);
\item the selected parameter \(\theta\);
\item the origin or method of obtaining \(\theta\);
\item the units carried by \(\theta\), if any;
\item whether \(\theta\) is fixed, measured, estimated, or fitted;
\item the uncertainty attached to \(\theta\).
\end{enumerate}
Book 0 supplies no calibration parameter.


\PCsection{6. Uncertainty and Failure}

\PCsubsection{Definition 6.1 --- Uncertainty space}

Let

\[
\mathfrak U(\mathcal Y)
\]

be a declared class of uncertainty descriptions over \(\mathcal Y\).

Examples may include:

\begin{itemize}
\item nonnegative error bounds;
\item intervals;
\item covariance matrices;
\item probability distributions;
\item confidence sets;
\item qualitative uncertainty labels.
\end{itemize}
No particular uncertainty formalism is assumed in S1.I.

\PCsubsection{Definition 6.2 --- Uncertainty assignment}

An uncertainty assignment is a map

\[
\boxed{ U_{\mathcal P}: \Omega_{\mathcal P} \to \mathfrak U(\mathcal Y). }
\]

\PCsubsection{Definition 6.3 --- Failure-witness space}

Let

\[
\mathfrak W
\]

be a declared set of failure witnesses.

A failure witness records why a projection could not be lawfully evaluated.

\PCsubsection{Definition 6.4 --- Failure map}

A failure map is

\[
\boxed{ W_{\mathcal P}: \Omega\setminus\Omega_{\mathcal P} \to \mathfrak W. }
\]

Possible failure witnesses include:

\begin{itemize}
\item PHASE-NOT-INFERABLE;
\item PHASE-AT-FORBIDDEN-BOUNDARY;
\item BRANCH-AMBIGUOUS;
\item STRUCTURAL-DENOMINATOR-ZERO;
\item OUTSIDE-CALIBRATION-RANGE;
\item READOUT-NOT-DEFINED;
\item UNCERTAINTY-NOT-AVAILABLE;
\item MODEL-DOMAIN-EXCEEDED.
\end{itemize}
The presence of a failure witness is not a defect in the theory. It is part of a complete contract.


\PCsection{7. Projection Contract}

\PCsubsection{Definition 7.1 --- Projection contract}

A projection contract is the tuple

\[
\boxed{ \mathcal P = \left( \Omega, \Omega_{\mathcal P}, \Phi, \mathbf F, \Gamma_\theta, U_{\mathcal P}, W_{\mathcal P} \right). }
\]

For a boundary-aware contract, \(\Phi\) is replaced by \(\widehat\Phi\).

\PCsubsection{Definition 7.2 --- Prediction map}

For

\[
\omega\in\Omega_{\mathcal P},
\]

define

\[
\boxed{ \widehat y_{\mathcal P}(\omega) = \Gamma_\theta \left( E_{\mathbf F}(\Phi(\omega)) \right). }
\]

Equivalently,

\[
\boxed{ \widehat y_{\mathcal P} = \Gamma_\theta \circ E_{\mathbf F} \circ \Phi. }
\]

\PCsubsection{Definition 7.3 --- Complete contract}

A projection contract is complete when all of the following are declared:

\begin{enumerate}
\item the external instance space \(\Omega\);
\item the valid subset \(\Omega_{\mathcal P}\);
\item the phase-inference map;
\item the structural profile;
\item the common structural domain;
\item the readout map;
\item every calibration parameter;
\item the output space;
\item the uncertainty assignment;
\item the failure map.
\end{enumerate}
\PCsubsection{Definition 7.4 --- Lawful contract}

A complete projection contract is lawful when:

\begin{enumerate}
\item \(\Phi(\Omega_{\mathcal P})\subseteq\mathcal D_{\mathbf F}\);
\item \(E_{\mathbf F}\) is defined on every inferred phase;
\item \(\Gamma_\theta\) is defined on every produced structural state;
\item all dimensional scales and parameters are declared;
\item boundary states are rejected or treated by a declared boundary-aware contract;
\item no theorem is asserted outside its structural or projection domain.
\end{enumerate}

\PCsection{PROCEDURE DIVISION}

\PCsection{8. Well-Defined Composition}

\PCsubsection{Theorem 8.1 --- Projection composition theorem}

Let

\[
\mathcal P = \left( \Omega, \Omega_{\mathcal P}, \Phi, \mathbf F, \Gamma_\theta, U_{\mathcal P}, W_{\mathcal P} \right)
\]

be a lawful interior projection contract.

Then the prediction map

\[
\widehat y_{\mathcal P} = \Gamma_\theta \circ E_{\mathbf F} \circ \Phi
\]

is well-defined on \(\Omega_{\mathcal P}\).

\begin{proof}
Take any

\[
\omega\in\Omega_{\mathcal P}.
\]

Because the contract is lawful,

\[
\Phi(\omega)\in\mathcal D_{\mathbf F}.
\]

Therefore,

\[
E_{\mathbf F}(\Phi(\omega))
\]

is defined and belongs to the declared structural-state domain

\[
\mathcal S_{\mathcal P}.
\]

Because

\[
\Gamma_\theta: \mathcal S_{\mathcal P}\to\mathcal Y,
\]

the value

\[
\Gamma_\theta(E_{\mathbf F}(\Phi(\omega)))
\]

is defined and belongs to \(\mathcal Y\).

This holds for every

\[
\omega\in\Omega_{\mathcal P}.
\]

Therefore, the composed prediction map is well-defined on the validity domain.

\end{proof}

\PCsubsection{Corollary 8.2 --- No lawful prediction outside the contract}

For

\[
\omega\notin\Omega_{\mathcal P},
\]

the contract produces a failure witness rather than a prediction:

\[
W_{\mathcal P}(\omega)\in\mathfrak W.
\]

Extending the prediction beyond \(\Omega_{\mathcal P}\) requires a new or enlarged contract.


\PCsection{9. Preservation of Structural Identities}

\PCsubsection{Theorem 9.1 --- Identity preservation under phase inference}

Let

\[
F,G\in\mathscr F
\]

satisfy

\[
F(x)=G(x)
\]

for every

\[
x\in S\subseteq\mathcal D.
\]

Let

\[
\Phi: \Omega_{\mathcal P}\to S.
\]

Then

\[
\boxed{ F(\Phi(\omega)) = G(\Phi(\omega)) }
\]

for every

\[
\omega\in\Omega_{\mathcal P}.
\]

\begin{proof}
For each valid \(\omega\),

\[
\Phi(\omega)\in S.
\]

The structural identity holds at every point of \(S\). Therefore, it holds at the particular point \(\Phi(\omega)\):

\[
F(\Phi(\omega)) = G(\Phi(\omega)).
\]

\end{proof}

\PCsubsection{Theorem 9.2 --- Identity preservation under readout}

Under the assumptions of Theorem 9.1, let

\[
\Gamma_\theta
\]

be any function defined on the common values of \(F\) and \(G\).

Then

\[
\boxed{ \Gamma_\theta(F(\Phi(\omega))) = \Gamma_\theta(G(\Phi(\omega))) }
\]

for every valid \(\omega\).

\begin{proof}
Theorem 9.1 gives equal inputs to \(\Gamma_\theta\).

A function assigns the same output to equal inputs.

Therefore, the readout values are equal.

\end{proof}

\PCsubsection{Corollary 9.3 --- Projected UNA identity}

For every lawful phase inference into \(\mathcal D\),

\[
\operatorname{urx}(\Phi(\omega)) + \operatorname{uxp}(\Phi(\omega)) = \frac4{\sin(2\Phi(\omega))}.
\]

Any common readout applied to these equal quantities preserves their equality.

\PCsubsection{Corollary 9.4 --- Projected Flatwave identity}

For every lawful phase inference,

\[
\operatorname{Flatwave}(\Phi(\omega)) = \operatorname{sgn} \left( \sin(2\Phi(\omega)) \right).
\]

The projection does not weaken the structural theorem.


\PCsection{10. Limits of the Converse}

Projected agreement does not generally prove structural equality.

\PCsubsection{Theorem 10.1 --- Injective converse on the inferred image}

Let

\[
F,G\in\mathscr F,
\]

and suppose

\[
\Gamma_\theta
\]

is injective on the union of the values

\[
F(\Phi(\Omega_{\mathcal P})) \cup G(\Phi(\Omega_{\mathcal P})).
\]

If

\[
\Gamma_\theta(F(\Phi(\omega))) = \Gamma_\theta(G(\Phi(\omega)))
\]

for every valid \(\omega\), then

\[
\boxed{ F(x)=G(x) }
\]

for every

\[
x\in\Phi(\Omega_{\mathcal P}).
\]

\begin{proof}
For every valid \(\omega\),

\[
\Gamma_\theta(F(\Phi(\omega))) = \Gamma_\theta(G(\Phi(\omega))).
\]

Because \(\Gamma_\theta\) is injective on the relevant values,

\[
F(\Phi(\omega)) = G(\Phi(\omega)).
\]

Therefore, \(F\) and \(G\) agree at every point in the image

\[
\Phi(\Omega_{\mathcal P}).
\]

\end{proof}

\PCsubsection{Corollary 10.2 --- Structural equality requires phase coverage}

If, in addition,

\[
\Phi(\Omega_{\mathcal P})=S,
\]

then

\[
F=G
\]

throughout \(S\).

If the inferred image is only a proper subset of \(S\), equality outside the sampled image does not follow.

\PCsubsection{Negative Theorem 10.3 --- Noninjective readout cannot establish structural equality}

If \(\Gamma_\theta\) is noninjective, projected equality may occur even when the structural values differ.

\begin{proof}
Since \(\Gamma_\theta\) is noninjective, there exist distinct values

\[
a\neq b
\]

such that

\[
\Gamma_\theta(a)=\Gamma_\theta(b).
\]

Choose structural expressions and an inferred phase for which

\[
F(\Phi(\omega))=a
\]

and

\[
G(\Phi(\omega))=b.
\]

Then the projected outputs agree even though the structural values do not.

Therefore, projected agreement alone does not prove structural equality.

\end{proof}

\PCsubsection{Example 10.4 --- Sign-only readout}

Suppose

\[
\Gamma(u)=\operatorname{sgn}(u).
\]

Then every positive structural value has the same projected output \(+1\), regardless of its magnitude.

A sign-only readout cannot distinguish

\[
1,\quad 2,\quad 100.
\]


\PCsection{11. Dimensional Isolation}

\PCsubsection{Definition 11.1 --- Structural dimensionlessness}

The phase variable \(x\), the Book 0 primitive operators, UNA operators, saw responses, Flatwave, \(p\), \(q\), and \(\lambda\) are dimensionless.

\PCsubsection{Definition 11.2 --- Dimension-bearing element}

A dimension-bearing element is any external datum, parameter, or output assigned a physical or conventional unit.

Examples include units of:

\begin{itemize}
\item seconds;
\item metres;
\item kilograms;
\item kelvin;
\item volts;
\item dollars.
\end{itemize}
\PCsubsection{Theorem 11.3 --- Dimensional isolation theorem}

In a lawful projection contract, every dimension carried by the output must enter through at least one of:

\begin{enumerate}
\item the external instance data;
\item the calibration parameter \(\theta\);
\item the readout map \(\Gamma_\theta\);
\item the declared structure of the output space \(\mathcal Y\).
\end{enumerate}
No dimensional unit is generated by the Book 0 structural profile alone.

\begin{proof}
The evaluated structural state

\[
E_{\mathbf F}(\Phi(\omega))
\]

is a tuple of dimensionless real values.

A composition of dimensionless structural expressions remains dimensionless until an operation involving a dimension-bearing element is applied.

By definition, the only post-structural operation in the contract is the readout

\[
\Gamma_\theta.
\]

Any external data or calibration parameters used by this map must be explicitly declared.

Therefore, every output dimension must enter through the external or readout layer rather than arising from the Book 0 expressions alone.

\end{proof}

\PCsubsection{Corollary 11.4 --- Identity readout remains dimensionless}

If

\[
\Gamma(u)=u,
\]

then the projected output remains dimensionless.

\PCsubsection{Corollary 11.5 --- A unit cannot be derived from a pure number alone}

A dimensionless structural value cannot, without a declared scale or unit assignment, determine a dimensional quantity.

For example, the equation

\[
\lambda=\frac12
\]

does not determine:

\begin{itemize}
\item one-half second;
\item one-half metre;
\item one-half joule;
\item one-half dollar.
\end{itemize}
The unit and scale must enter through the contract.


\PCsection{12. Boundary Guard}

\PCsubsection{Theorem 12.1 --- Interior boundary guard}

Let \(\mathcal P\) be an interior projection contract using

\[
\Phi:\Omega_{\mathcal P}\to\mathcal D.
\]

If an external instance \(\omega\) would require

\[
\Phi(\omega)\in\mathbb B,
\]

then \(\omega\) is not admissible under that contract.

\begin{proof}
By definition,

\[
\mathcal D\cap\mathbb B=\varnothing.
\]

An interior phase-inference map has codomain \(\mathcal D\). It therefore cannot lawfully assign a forbidden boundary phase.

Moreover, the primitive and UNA expressions are not defined at those points.

Hence the contract must reject the instance or replace the interior map with a boundary-aware contract.

\end{proof}

\PCsubsection{Corollary 12.2 --- Required boundary failure witness}

An interior contract encountering a boundary state must return a witness such as

\[
\texttt{PHASE-AT-FORBIDDEN-BOUNDARY}.
\]

It may not silently evaluate the interior formulas by substituting a boundary value.

\PCsubsection{Theorem 12.3 --- Boundary-aware admissibility}

A boundary-aware contract may use

\[
\widehat\Phi: \Omega_{\mathcal P}\to\widehat{\mathcal D}
\]

provided that:

\begin{enumerate}
\item the entering or exiting boundary copy is declared;
\item every selected structural quantity has a proved extension to that copy;
\item no divergent primitive or UNA quantity is assigned a finite value without an additional rule;
\item no seam transition is inferred merely from the boundary state.
\end{enumerate}
\begin{proof}
S0.V proves continuous split-boundary extensions for the bounded transfer coordinate, saw responses, and Flatwave.

It does not provide finite extensions for all primitive and UNA expressions, nor does it provide a seam map.

Therefore, a boundary-aware contract is lawful only for the quantities and side labels actually supplied by S0.V or by an explicitly added extension rule.


\end{proof}

\PCsection{13. Reparameterization and Non-Identification}

Different internal descriptions may produce the same external prediction.

\PCsubsection{Definition 13.1 --- Admissible phase reparameterization}

Let

\[
S,T\subseteq\mathcal D.
\]

An admissible phase reparameterization is a bijection

\[
\boxed{ h:S\to T }
\]

such that, for the selected structural profile,

\[
\widetilde{\mathbf F} = \mathbf F\circ h^{-1}
\]

is defined on \(T\).

When a contract requires \(\widetilde{\mathbf F}\) to remain inside a restricted structural class, that requirement must also be satisfied.

\PCsubsection{Theorem 13.2 --- Reparameterization invariance}

Let

\[
\Phi: \Omega_{\mathcal P}\to S
\]

and define

\[
\widetilde\Phi=h\circ\Phi.
\]

Define

\[
\widetilde{\mathbf F} = \mathbf F\circ h^{-1}.
\]

Then

\[
\boxed{ \widetilde{\mathbf F} \circ \widetilde\Phi = \mathbf F\circ\Phi. }
\]

Consequently, for the same readout map,

\[
\boxed{ \Gamma_\theta \circ \widetilde{\mathbf F} \circ \widetilde\Phi = \Gamma_\theta \circ \mathbf F \circ \Phi. }
\]

\begin{proof}
For any valid \(\omega\),

\[
\widetilde{\mathbf F} (\widetilde\Phi(\omega)) = \mathbf F \left( h^{-1}(h(\Phi(\omega))) \right).
\]

Since \(h\) is bijective,

\[
h^{-1}(h(\Phi(\omega))) = \Phi(\omega).
\]

Therefore,

\[
\widetilde{\mathbf F} (\widetilde\Phi(\omega)) = \mathbf F(\Phi(\omega)).
\]

Applying the same readout map preserves equality.

\end{proof}

\PCsubsection{Corollary 13.3 --- Prediction does not uniquely identify phase}

External prediction data alone do not necessarily determine a unique phase-inference map.

A different phase map may be compensated by a corresponding transformed structural profile.

\PCsubsection{Corollary 13.4 --- Identification conditions are additional}

To identify a unique phase description, a contract must impose additional conditions such as:

\begin{itemize}
\item fixing the structural profile in advance;
\item fixing a phase origin;
\item fixing orientation;
\item fixing a branch convention;
\item requiring a particular symmetry;
\item requiring a particular calibration form;
\item requiring identifiability from data.
\end{itemize}
These conditions do not follow from Book 0.


\PCsection{14. Structural Truth and Empirical Adequacy}

\PCsubsection{Definition 14.1 --- Structural truth}

A statement is structurally true when it is proved for the Book 0 functions on its declared mathematical domain.

Example:

\[
\operatorname{Flatwave}(x) = \operatorname{sgn}(\sin 2x)
\]

for every

\[
x\in\mathcal D.
\]

\PCsubsection{Definition 14.2 --- Contract validity}

A projection contract is valid for an external domain when its assumptions, inference rule, readout, and calibration are appropriate on its declared validity set.

\PCsubsection{Definition 14.3 --- Empirical adequacy}

A projection is empirically adequate when its outputs agree with relevant observations to the standard declared by the contract.

\PCsubsection{Theorem 14.4 --- Structural truth does not imply empirical adequacy}

The correctness of a Book 0 identity does not, by itself, establish that a proposed external projection is empirically adequate.

\begin{proof}
A projected prediction depends on three distinct components:

\[
\Phi, \qquad \mathbf F, \qquad \Gamma_\theta.
\]

A structural theorem constrains \(\mathbf F\).

It does not prove that:

\begin{itemize}
\item \(\Phi\) assigns the correct phase to the external instance;
\item the selected structural profile is relevant to the external phenomenon;
\item \(\Gamma_\theta\) is the correct readout;
\item \(\theta\) is correctly calibrated;
\item the uncertainty model is adequate.
\end{itemize}
Therefore, a true structural identity may participate in an inaccurate projection.

\end{proof}

\PCsubsection{Theorem 14.5 --- Empirical agreement does not prove structural necessity}

A projection may agree with observations without proving that its internal structural representation is uniquely necessary.

\begin{proof}
The reparameterization theorem shows that distinct internal descriptions may generate identical predictions.

Noninjective readouts may also erase structural distinctions.

Finite observations may fail to cover the full structural domain.

Therefore, empirical agreement establishes compatibility with the tested observations, not unique structural necessity.

\end{proof}

\PCsubsection{Corollary 14.6 --- Three separate burdens}

Every external application carries three separate burdens:

\begin{enumerate}
\item mathematical correctness of the structural expressions;
\item contract validity of the inference and readout;
\item empirical adequacy of the resulting predictions.
\end{enumerate}
Passing one burden does not automatically pass the others.


\PCsection{15. Minimal Projection Contract}

\PCsubsection{Theorem 15.1 --- Minimal field count}

A projection cannot be fully specified without at least the following four functional elements:

\begin{enumerate}
\item a validity domain;
\item a phase-inference map;
\item a structural profile;
\item a readout map.
\end{enumerate}
\begin{proof}
Without a validity domain, the contract does not state where it applies.

Without phase inference, no external instance is assigned a structural phase.

Without a structural profile, no Book 0 quantity is selected.

Without a readout map, no structural value becomes an external output.

Therefore, removing any one of the four elements leaves the external prediction undefined or its domain unspecified.

\end{proof}

\PCsubsection{Corollary 15.2 --- Operational completeness requires more than mathematical composition}

For responsible empirical use, a contract must additionally declare:

\begin{itemize}
\item calibration parameters;
\item uncertainty;
\item failure behavior.
\end{itemize}
These do not alter the formal composition but are required to make its use auditable.


\PCsection{16. Abstract Worked Projection}

This example demonstrates the contract form without assigning a physical interpretation.

\PCsubsection{Definition 16.1 --- External phase-labelled instances}

Let

\[
\Omega
\]

be a set of instances for which a phase-inference rule

\[
\Phi:\Omega_{\mathcal P}\to Q_k
\]

has been declared.

\PCsubsection{Definition 16.2 --- Structural profile}

Choose the single bounded structural function

\[
\mathbf F=(\lambda).
\]

Thus,

\[
E_{\mathbf F}(x)=\lambda(x).
\]

From Book 0,

\[
0<\lambda(x)<1
\]

on the open quadrant.

\PCsubsection{Definition 16.3 --- Affine readout}

Let

\[
a,b\in\mathbb R
\]

with

\[
a<b.
\]

Define

\[
\Gamma_{a,b}(u) = a+(b-a)u.
\]

The parameters \(a\) and \(b\) may carry units if the output space is dimensional.

\PCsubsection{Definition 16.4 --- Projected output}

\[
\widehat y(\omega) = a+(b-a)\lambda(\Phi(\omega)).
\]

\PCsubsection{Theorem 16.5 --- Bounded affine projection}

For every valid \(\omega\),

\[
\boxed{ a<\widehat y(\omega)<b. }
\]

\begin{proof}
Since

\[
0<\lambda(\Phi(\omega))<1
\]

and

\[
b-a>0,
\]

we have

\[
0< (b-a)\lambda(\Phi(\omega)) < b-a.
\]

Adding \(a\) gives

\[
a< a+(b-a)\lambda(\Phi(\omega)) < b.
\]

Therefore,

\[
a<\widehat y(\omega)<b.
\]

\end{proof}

\PCsubsection{Corollary 16.6 --- Dimensional source of the bounds}

If \(a\) and \(b\) carry units, then \(\widehat y\) carries those units.

The units originate in the calibration parameters, not in \(\lambda\).

\PCsubsection{Clarification 16.7 --- What the example does not prove}

The example does not prove that any actual phenomenon should be modeled by \(\lambda\).

It proves only that, once such a projection contract is declared, its output is bounded between its declared anchors.


\PCsection{17. Projection Contract Ledger}

Every later projection essay or applied model should begin with the following ledger.

\PCsubsection{17.1 External declaration}

\[
\Omega \text{declared external instance space}.
\]

\PCsubsection{17.2 Validity declaration}

\[
\Omega_{\mathcal P} \text{declared valid subset}.
\]

\PCsubsection{17.3 Phase declaration}

\[
\Phi \text{declared phase-inference map}.
\]

\PCsubsection{17.4 Structural declaration}

\[
\mathbf F \text{declared Book 0 profile}.
\]

\PCsubsection{17.5 Readout declaration}

\[
\Gamma_\theta \text{declared conversion to external output}.
\]

\PCsubsection{17.6 Calibration declaration}

\[
\theta \text{declared external scales and parameters}.
\]

\PCsubsection{17.7 Uncertainty declaration}

\[
U_{\mathcal P} \text{declared uncertainty assignment}.
\]

\PCsubsection{17.8 Failure declaration}

\[
W_{\mathcal P} \text{declared failure behavior}.
\]

\PCsubsection{17.9 Prediction declaration}

\[
\boxed{ \widehat y_{\mathcal P} = \Gamma_\theta \circ E_{\mathbf F} \circ \Phi. }
\]

No applied claim is complete until these fields are supplied.


\PCsection{18. Negative Theorems of S1.I}

\PCsubsection{Negative Theorem 18.1 --- No external meaning by resemblance}

Similarity between a structural graph and an observed graph does not, by itself, define a projection contract.

\PCsubsection{Negative Theorem 18.2 --- No hidden units}

A dimensional unit may not be attributed to a Book 0 quantity unless the unit enters through a declared external or readout element.

\PCsubsection{Negative Theorem 18.3 --- No silent boundary substitution}

An interior structural expression may not be evaluated at an excluded boundary by inserting an arbitrary finite value.

\PCsubsection{Negative Theorem 18.4 --- No inference from output equality alone}

Equal projected outputs do not prove equal structural states unless the readout is injective on the relevant image.

\PCsubsection{Negative Theorem 18.5 --- No unique phase from fit alone}

A successful external fit does not, by itself, establish a unique phase-inference map.

\PCsubsection{Negative Theorem 18.6 --- No physical law from structural identity alone}

A Book 0 theorem becomes part of a physical or empirical model only through a declared and tested projection contract.


\PCsection{19. Terminal Theorem}

\PCsubsection{Theorem 19.1 --- Projection Separation Theorem}

Let Book 0 provide a structural expression class \(\mathscr F\) on \(\mathcal D\).

Every lawful external prediction using that structure factors through three logically distinct stages:

\[
\boxed{ \text{external instance} \longrightarrow \text{inferred phase} \longrightarrow \text{structural state} \longrightarrow \text{external output}. }
\]

More formally,

\[
\boxed{ \widehat y_{\mathcal P} = \Gamma_\theta \circ E_{\mathbf F} \circ \Phi. }
\]

The structural theorems constrain \(E_{\mathbf F}\). They do not independently determine \(\Phi\), \(\Gamma_\theta\), \(\theta\), the validity domain, uncertainty, or failure behavior.

\begin{proof}
The factorization is the defining form required by Axiom P.1 and Definition 7.1.

Book 0 defines and constrains the structural expressions comprising \(E_{\mathbf F}\).

The external instance space, phase inference, readout, calibration, uncertainty, and failure behavior are introduced only in the projection contract.

They are therefore logically distinct from the Book 0 structural calculus.


\end{proof}

\PCsection{20. Closing Declaration}

Book 0 answered:

\[
\boxed{ \text{What does the bounded phase calculus prove?} }
\]

Book 1 begins by asking:

\[
\boxed{ \text{How may that structure be lawfully applied?} }
\]

The answer established in S1.I is the projection contract:

\[
\boxed{ \mathcal P = \left( \Omega, \Omega_{\mathcal P}, \Phi, \mathbf F, \Gamma_\theta, U_{\mathcal P}, W_{\mathcal P} \right). }
\]

Its prediction map is

\[
\boxed{ \widehat y_{\mathcal P} = \Gamma_\theta \circ E_{\mathbf F} \circ \Phi. }
\]

Exact structural identities survive lawful composition.

Dimensional quantities enter only through declared external or readout elements.

Boundary states require explicit side-aware treatment.

Projected agreement does not guarantee structural uniqueness.

Structural truth, contract validity, and empirical adequacy remain separate burdens.

Thus the opening law of Book 1 is:

\[
\boxed{ \begin{gathered} \text{No external interpretation without explicit inference;}\\ \text{no measurement without declared readout;}\\ \text{no unit without declared scale;}\\ \text{no application outside its validity contract.} \end{gathered} }
\]

\[
\boxed{ \text{END OF S1.I} }
\]

% ===== END S1_I =====
% ===== BEGIN S1_II =====
% Source: S1.ii.docx
\PCchapter{S1.II}{Admissibility, Failure Witnesses, and Falsification}{Book 1 - Projection and Measurement}

\PCsubsection{Abstract}

S1.I defined a projection contract as the explicit composition

\[
\Omega_{\mathcal P} \xrightarrow{\ \Phi\ } \mathcal D \xrightarrow{\ E_{\mathbf F}\ } \mathbb R^m \xrightarrow{\ \Gamma_\theta\ } \mathcal Y.
\]

This essay determines how such a contract is audited.

A projection is not admitted merely because its formulas can be evaluated at selected examples. Every instance must pass a declared sequence of guards concerning domain membership, phase inference, branch selection, structural evaluation, boundary status, calibration, readout, and uncertainty.

For each external instance \(\omega\), a failure-witness set

\[
\mathcal W_{\mathcal P}(\omega) \subseteq \mathfrak W
\]

records every failed guard. An instance is admissible exactly when this set is empty.

The ordinary partial prediction map is then extended to a total diagnostic evaluator

\[
\mathcal E_{\mathcal P}: \Omega \longrightarrow \mathcal Y\sqcup\mathcal P^{+}(\mathfrak W),
\]

which returns either a projected output or a nonempty set of failure witnesses.

This construction distinguishes four logically different outcomes:

\begin{enumerate}
\item admissible prediction;
\item undefined evaluation;
\item contract invalidity or insufficiency;
\item empirical contradiction.
\end{enumerate}
Failure to produce a prediction is not empirical falsification. A lawful empirical contradiction can be declared only after the instance has passed the projection guards and the predicted and observed uncertainty sets are disjoint.

A mismatch between an external observation and a projection does not falsify a Book 0 identity. It falsifies, at most, the particular projection contract, calibration, phase-inference rule, or external interpretation under test.

The essay also proves that restriction of a lawful validity domain preserves lawfulness, whereas extension to a larger domain requires a new admissibility proof. It introduces repair locality: a failure witness identifies the contract field that must be revised without silently altering unrelated fields.

The result is a complete operational airlock between structural mathematics and external claims.


\PCsection{IDENTIFICATION DIVISION}

\PCsubsection{PROGRAM-ID}

S1-II-ADMISSIBILITY-AND-FAILURE

\PCsubsection{PURPOSE}

To define the procedures by which a proposed projection is:

\begin{itemize}
\item admitted;
\item rejected;
\item diagnosed;
\item repaired;
\item empirically tested;
\item or declared underdetermined.
\end{itemize}
\PCsubsection{PRINCIPAL QUESTION}

Given a projection contract \(\mathcal P\) and an external instance \(\omega\), what result is logically justified?

\PCsubsection{NON-PURPOSE}

This essay does not:

\begin{itemize}
\item choose a particular physical application;
\item define a universal statistical test;
\item guarantee that every failure can be repaired;
\item prove that a successful projection is unique;
\item convert an undefined expression into an empirical refutation;
\item permit silent extrapolation beyond a declared domain;
\item alter any theorem of Book 0.
\end{itemize}

\PCsection{ENVIRONMENT DIVISION}

\PCsubsection{1. Inherited Projection Contract}

From S1.I, a projection contract is

\[
\boxed{ \mathcal P = \left( \Omega, \Omega_{\mathcal P}, \Phi, \mathbf F, \Gamma_\theta, U_{\mathcal P}, W_{\mathcal P} \right). }
\]

Here:

\begin{itemize}
\item \(\Omega\) is the external instance space;
\item \(\Omega_{\mathcal P}\subseteq\Omega\) is the declared validity domain;
\item \(\Phi\) is the phase-inference map;
\item \(\mathbf F\) is the selected Book 0 structural profile;
\item \(\Gamma_\theta\) is the readout map;
\item \(\theta\) contains all calibration parameters;
\item \(U_{\mathcal P}\) assigns projection uncertainty;
\item \(W_{\mathcal P}\) reports failure.
\end{itemize}
For an admissible instance,

\[
\widehat y_{\mathcal P}(\omega) = \Gamma_\theta \left( E_{\mathbf F}(\Phi(\omega)) \right).
\]

\PCsubsection{1.1 Interior and boundary-aware forms}

An interior contract uses

\[
\Phi: \Omega_{\mathcal P}\to\mathcal D.
\]

A boundary-aware contract uses

\[
\widehat\Phi: \Omega_{\mathcal P}\to\widehat{\mathcal D}.
\]

The two forms may not be silently interchanged.

\PCsubsection{1.2 No new structural axiom}

S1.II introduces no new statement about the Book 0 operators.

Its definitions govern only the logical and operational status of external projections.


\PCsection{DATA DIVISION}

\PCsection{2. Evaluation Outcomes}

\PCsubsection{Definition 2.1 --- Admitted instance}

An external instance

\[
\omega\in\Omega
\]

is admitted by \(\mathcal P\) when every required projection guard succeeds.

\PCsubsection{Definition 2.2 --- Rejected instance}

An instance is rejected by \(\mathcal P\) when at least one required guard fails.

Rejection is contract-relative. The same instance may be admissible under another contract.

\PCsubsection{Definition 2.3 --- Undefined evaluation}

An evaluation is undefined when one of the required mathematical functions is not defined at the value supplied to it.

Examples include:

\begin{itemize}
\item an inferred phase at a forbidden wall;
\item a zero structural denominator;
\item a readout outside its declared input domain.
\end{itemize}
Undefinedness is a mathematical status, not an empirical discrepancy.

\PCsubsection{Definition 2.4 --- Contract insufficiency}

A contract is insufficient for an instance when it lacks information required to decide or evaluate that instance.

Examples include:

\begin{itemize}
\item no phase can be inferred;
\item two branches remain equally admissible;
\item calibration is missing;
\item uncertainty is required but unavailable.
\end{itemize}
\PCsubsection{Definition 2.5 --- Empirical contradiction}

An empirical contradiction occurs when:

\begin{enumerate}
\item the instance is admitted;
\item a lawful prediction exists;
\item a lawful observation record exists;
\item the predicted and observed admissible-output sets are disjoint.
\end{enumerate}
This definition is formalized in Section 11.

\PCsubsection{Definition 2.6 --- Underdetermined outcome}

An outcome is underdetermined when available information permits more than one contract-consistent conclusion and the contract supplies no selection rule.

Underdetermination is neither successful prediction nor empirical contradiction.


\PCsection{3. Guard System}

\PCsubsection{Definition 3.1 --- Guard}

A guard is a Boolean condition

\[
G_j(\mathcal P,\omega) \in \{\operatorname{true},\operatorname{false}\}
\]

that must be satisfied before a prediction is admitted.

\PCsubsection{Definition 3.2 --- Contract-completeness guard}

\[
G_0(\mathcal P)
\]

is true exactly when all required contract fields are present and internally typed.

At minimum, the contract must declare:

\begin{itemize}
\item \(\Omega\);
\item \(\Omega_{\mathcal P}\);
\item \(\Phi\) or \(\widehat\Phi\);
\item \(\mathbf F\);
\item \(\Gamma_\theta\);
\item \(\theta\);
\item \(U_{\mathcal P}\);
\item \(W_{\mathcal P}\).
\end{itemize}
\PCsubsection{Definition 3.3 --- Instance-domain guard}

\[
G_1(\mathcal P,\omega)
\]

is true exactly when

\[
\omega\in\Omega_{\mathcal P}.
\]

\PCsubsection{Definition 3.4 --- Phase-inference guard}

\[
G_2(\mathcal P,\omega)
\]

is true exactly when the contract's declared procedure returns a phase value for \(\omega\).

\PCsubsection{Definition 3.5 --- Phase-uniqueness guard}

\[
G_3(\mathcal P,\omega)
\]

is true exactly when the phase-inference procedure returns one contract-admissible phase state, or when the contract explicitly permits a set-valued phase result and supplies a corresponding readout rule.

If two or more unresolved phase states remain and no set-valued rule is declared, then \(G_3\) is false.

\PCsubsection{Definition 3.6 --- Structural-domain guard}

For an interior contract,

\[
G_4(\mathcal P,\omega)
\]

is true exactly when

\[
\Phi(\omega)\in\mathcal D_{\mathbf F}.
\]

For a boundary-aware contract, it is true exactly when:

\begin{enumerate}
\item the inferred split-boundary point is declared;
\item every selected structural quantity has a proved or explicitly assumed extension to that point.
\end{enumerate}
\PCsubsection{Definition 3.7 --- Structural-evaluation guard}

\[
G_5(\mathcal P,\omega)
\]

is true exactly when every component of

\[
E_{\mathbf F}(\Phi(\omega))
\]

can be evaluated without division by zero, undefined branching, or an undeclared singular convention.

\PCsubsection{Definition 3.8 --- Calibration guard}

\[
G_6(\mathcal P,\omega)
\]

is true exactly when all calibration parameters required for the instance are:

\begin{itemize}
\item present;
\item assigned;
\item dimensionally compatible;
\item and within their declared validity conditions.
\end{itemize}
\PCsubsection{Definition 3.9 --- Readout-domain guard}

\[
G_7(\mathcal P,\omega)
\]

is true exactly when the evaluated structural state belongs to the declared domain of \(\Gamma_\theta\).

\PCsubsection{Definition 3.10 --- Uncertainty guard}

\[
G_8(\mathcal P,\omega)
\]

is true exactly when the uncertainty description required by the contract can be produced.

A purely exact mathematical illustration may declare that no empirical uncertainty is being asserted. An empirical contract may not silently omit required uncertainty.

\PCsubsection{Definition 3.11 --- Observation guard}

When empirical comparison is attempted,

\[
G_9(\mathcal P,\omega,\mathcal O)
\]

is true exactly when the observation record \(\mathcal O\) is valid under the declared observation protocol.


\PCsection{4. Admissibility Predicate}

\PCsubsection{Definition 4.1 --- Instance admissibility}

Define

\[
\boxed{ A_{\mathcal P}(\omega) = \bigwedge_{j=0}^{8} G_j(\mathcal P,\omega). }
\]

Thus,

\[
A_{\mathcal P}(\omega)=\operatorname{true}
\]

exactly when all non-observational projection guards succeed.

\PCsubsection{Definition 4.2 --- Empirical-test admissibility}

For an observation record \(\mathcal O\), define

\[
\boxed{ A_{\mathcal P}^{\mathrm{test}}(\omega,\mathcal O) = A_{\mathcal P}(\omega) \wedge G_9(\mathcal P,\omega,\mathcal O). }
\]

\PCsubsection{Definition 4.3 --- Maximal operational admissibility predicate}

The instance-domain guard \(G_1\) tests membership in the already declared set \(\Omega_{\mathcal P}\). It must therefore not be used to define the maximal domain against which that declaration is assessed.

Define the non-domain operational predicate

\[
\boxed{ A_{\mathcal P}^{\max}(\omega) = G_0(\mathcal P) \wedge \bigwedge_{j=2}^{8}G_j(\mathcal P,\omega). }
\]

The corresponding maximal operational admissibility envelope is

\[
\boxed{ \Omega_{\mathcal P}^{\max} = \left\{ \omega\in\Omega: A_{\mathcal P}^{\max}(\omega)=\operatorname{true} \right\}. }
\]

The ordinary admissibility predicate factors as

\[
\boxed{ A_{\mathcal P}(\omega) = G_1(\mathcal P,\omega) \wedge A_{\mathcal P}^{\max}(\omega). }
\]

\PCsubsection{Theorem 4.4 --- Declared-domain lawfulness and maximality}

A declared validity domain is operationally lawful when

\[
\boxed{ \Omega_{\mathcal P} \subseteq \Omega_{\mathcal P}^{\max}. }
\]

It is maximal relative to the declared contract fields and guards when

\[
\boxed{ \Omega_{\mathcal P} = \Omega_{\mathcal P}^{\max}. }
\]

\begin{proof}
If \(\omega\in\Omega_{\mathcal P}\), then \(G_1(\mathcal P,\omega)\) is true by definition. Lawful prediction additionally requires every non-domain guard to succeed, which is exactly the condition \(\omega\in\Omega_{\mathcal P}^{\max}\). Hence lawfulness requires \(\Omega_{\mathcal P}\subseteq\Omega_{\mathcal P}^{\max}\).

If equality holds, every instance admitted by the non-domain guards is included in the declared scope, so the declaration is maximal relative to those guards. If the inclusion is strict, the contract is conservatively scoped but remains lawful.
\end{proof}

\PCsubsection{Clarification 4.5 --- Conservative domains}

A contract may intentionally choose a conservative validity domain

\[
\Omega_{\mathcal P} \subsetneq \Omega_{\mathcal P}^{\max}.
\]

Such a restriction is lawful when explicitly stated. The contract must not claim that the restricted domain is maximal.


\PCsection{5. Failure-Witness Sets}

\PCsubsection{Definition 5.1 --- Witness assignment by guard}

To each guard \(G_j\), assign a witness code

\[
w_j\in\mathfrak W.
\]

A recommended base vocabulary is:

\[
\begin{array}{c|l} j & w_j\\ \hline 0 & \texttt{CONTRACT-INCOMPLETE}\\ 1 & \texttt{INSTANCE-OUTSIDE-VALIDITY-DOMAIN}\\ 2 & \texttt{PHASE-NOT-INFERABLE}\\ 3 & \texttt{PHASE-OR-BRANCH-AMBIGUOUS}\\ 4 & \texttt{PHASE-OUTSIDE-STRUCTURAL-DOMAIN}\\ 5 & \texttt{STRUCTURAL-EVALUATION-UNDEFINED}\\ 6 & \texttt{CALIBRATION-MISSING-OR-INVALID}\\ 7 & \texttt{READOUT-OUTSIDE-DOMAIN}\\ 8 & \texttt{UNCERTAINTY-NOT-AVAILABLE}\\ 9 & \texttt{OBSERVATION-RECORD-INVALID} \end{array}
\]

\PCsubsection{Definition 5.2 --- Failure-witness set}

For an external instance \(\omega\), define

\[
\boxed{ \mathcal W_{\mathcal P}(\omega) = \left\{ w_j: G_j(\mathcal P,\omega)=\operatorname{false}, \ 0\leq j\leq8 \right\}. }
\]

For an empirical test record,

\[
\boxed{ \mathcal W_{\mathcal P}^{\mathrm{test}} (\omega,\mathcal O) = \mathcal W_{\mathcal P}(\omega) \cup \left\{ w_9: G_9(\mathcal P,\omega,\mathcal O)=\operatorname{false} \right\}. }
\]

\PCsubsection{Theorem 5.3 --- Empty-witness criterion}

\[
\boxed{ A_{\mathcal P}(\omega)=\operatorname{true} \iff \mathcal W_{\mathcal P}(\omega)=\varnothing. }
\]

\begin{proof}
By Definition 4.1, \(A_{\mathcal P}(\omega)\) is true exactly when every guard \(G_0,\ldots,G_8\) is true.

By Definition 5.2, a witness is included exactly when its corresponding guard is false.

Therefore, every guard is true exactly when no witness is included.

\end{proof}

\PCsubsection{Corollary 5.4 --- Rejection criterion}

\[
\boxed{ A_{\mathcal P}(\omega)=\operatorname{false} \iff \mathcal W_{\mathcal P}(\omega)\neq\varnothing. }
\]

\PCsubsection{Definition 5.5 --- Primary witness}

A contract may declare a precedence relation

\[
\prec
\]

on \(\mathfrak W\).

When

\[
\mathcal W_{\mathcal P}(\omega)\neq\varnothing,
\]

a primary witness is a \(\prec\)-minimal element of the witness set.

\PCsubsection{Clarification 5.6 --- Primary witness need not be unique}

If the precedence relation is partial rather than total, more than one minimal witness may exist.

The complete witness set should be preserved even when a primary witness is displayed.


\PCsection{6. Totalized Diagnostic Evaluation}

\PCsubsection{Definition 6.1 --- Disjoint outcome space}

Let

\[
\mathcal Y\sqcup\mathcal P^{+}(\mathfrak W)
\]

denote the disjoint union of:

\begin{itemize}
\item the output space \(\mathcal Y\);
\item the nonempty subsets of the witness space.
\end{itemize}
The disjoint union ensures that an output value cannot be confused with a failure code.

\PCsubsection{Definition 6.2 --- Totalized evaluator}

Define

\[
\boxed{ \mathcal E_{\mathcal P}: \Omega \to \mathcal Y\sqcup\mathcal P^{+}(\mathfrak W) }
\]

by

\[
\mathcal E_{\mathcal P}(\omega) = \begin{cases} \widehat y_{\mathcal P}(\omega), & \mathcal W_{\mathcal P}(\omega)=\varnothing, \\[6pt] \mathcal W_{\mathcal P}(\omega), & \mathcal W_{\mathcal P}(\omega)\neq\varnothing. \end{cases}
\]


\PCsubsection{Theorem 6.3 --- Totality}

For every

\[
\omega\in\Omega,
\]

the evaluator \(\mathcal E_{\mathcal P}(\omega)\) returns exactly one of:

\begin{enumerate}
\item an output in \(\mathcal Y\);
\item a nonempty failure-witness set.
\end{enumerate}
\begin{proof}
By classical logic, either

\[
\mathcal W_{\mathcal P}(\omega)=\varnothing
\]

or

\[
\mathcal W_{\mathcal P}(\omega)\neq\varnothing.
\]

The two conditions are mutually exclusive.

Definition 6.2 assigns one result in each case.

Therefore, the evaluator is total and single-valued.

\end{proof}

\PCsubsection{Corollary 6.4 --- No silent failure}

A complete implementation of the contract cannot return neither an output nor a witness.

\PCsubsection{Corollary 6.5 --- No output with unresolved failure}

A complete implementation cannot lawfully return a projected output while retaining a nonempty unresolved witness set, unless the contract explicitly classifies that witness as a warning rather than a failed guard.


\PCsection{7. Failure Taxonomy}

\PCsubsection{Definition 7.1 --- Contract-definition failure}

A contract-definition failure concerns the specification itself rather than one particular external instance.

Examples include:

\begin{itemize}
\item missing readout map;
\item undeclared units;
\item incomplete parameter list;
\item no failure behavior;
\item contradictory domain declarations.
\end{itemize}
\PCsubsection{Definition 7.2 --- Instance-admission failure}

An instance-admission failure occurs when the contract is complete but the particular instance lies outside its lawful applicability.

Example:

\[
\omega\notin\Omega_{\mathcal P}.
\]

\PCsubsection{Definition 7.3 --- Inference failure}

An inference failure occurs when phase or branch cannot be determined under the contract.

\PCsubsection{Definition 7.4 --- Structural failure}

A structural failure occurs when the inferred state cannot be lawfully evaluated by the selected Book 0 expressions.

\PCsubsection{Definition 7.5 --- Readout failure}

A readout failure occurs when the structural state exists but cannot be converted to an output under the declared calibration and readout map.

\PCsubsection{Definition 7.6 --- Uncertainty failure}

An uncertainty failure occurs when the contract requires an uncertainty statement but none can be lawfully produced.

\PCsubsection{Definition 7.7 --- Observation failure}

An observation failure occurs when the empirical record is missing, malformed, outside protocol, or otherwise unsuitable for comparison.

\PCsubsection{Definition 7.8 --- Empirical contradiction witness}

An empirical contradiction witness records a valid mismatch after every admission and observation guard has passed.

Recommended code:

\[
\texttt{PREDICTION-OBSERVATION-DISJOINT}.
\]

\PCsubsection{Theorem 7.9 --- Failure-class separation}

Contract-definition failure, undefined evaluation, and empirical contradiction are logically distinct.

\begin{proof}
A contract-definition failure concerns whether a projection has been completely specified.

Undefined evaluation concerns whether declared mathematical functions possess a value at the supplied arguments.

Empirical contradiction concerns disagreement between a lawfully produced prediction and a lawfully obtained observation.

Each may occur without either of the others.

For example:

\begin{itemize}
\item a contract may be incomplete before any instance is evaluated;
\item a complete contract may infer a forbidden phase, producing undefined evaluation;
\item a complete and evaluable contract may produce a prediction contradicted by observation.
\end{itemize}
Therefore, the classes are distinct.


\end{proof}

\PCsection{8. Undefinedness Is Not Falsification}

\PCsubsection{Theorem 8.1 --- Non-evaluation theorem}

If

\[
\mathcal W_{\mathcal P}(\omega)\neq\varnothing,
\]

then the instance does not produce a lawful projected output under \(\mathcal P\).

Consequently, that instance cannot by itself empirically contradict the contract's prediction.

\begin{proof}
A nonempty witness set implies, by Corollary 5.4, that the instance is not admitted.

By Definition 6.2, the evaluator returns witnesses rather than a projected output.

Without a projected output, there is no predicted value or predicted admissible-output set to compare with an observation.

Therefore, the instance demonstrates non-applicability or insufficiency, not empirical contradiction.

\end{proof}

\PCsubsection{Corollary 8.2 --- Boundary failure is not a failed prediction}

If an interior contract returns

\[
\texttt{PHASE-OUTSIDE-STRUCTURAL-DOMAIN}
\]

because the inferred phase lies at

\[
x=\frac{k\pi}{2},
\]

the result is a domain failure.

It is not evidence that a numerical prediction was wrong.

\PCsubsection{Corollary 8.3 --- Ambiguity is not contradiction}

If two unresolved phases remain admissible and the contract has no branch-selection rule, the result is underdetermination.

It is not a contradiction between the model and the observation.


\PCsection{9. Failure Soundness and Relative Completeness}

\PCsubsection{Definition 9.1 --- Sound witness}

A witness is sound when its presence implies that the associated guard has genuinely failed.

\PCsubsection{Definition 9.2 --- Guard-complete witness system}

A witness system is guard-complete when every failed guard produces at least one witness.

\PCsubsection{Theorem 9.3 --- Soundness of the canonical witness system}

The witness system of Section 5 is sound.

\begin{proof}
By Definition 5.2, \(w_j\) is included only when

\[
G_j=\operatorname{false}.
\]

Therefore, the presence of \(w_j\) directly certifies failure of its corresponding guard.

\end{proof}

\PCsubsection{Theorem 9.4 --- Relative completeness}

The witness system is complete relative to the declared guard set \(G_0,\ldots,G_8\).

\begin{proof}
If a declared guard fails, its corresponding witness is included by Definition 5.2.

Therefore, no failure among the declared guards is left unwitnessed.

\end{proof}

\PCsubsection{Clarification 9.5 --- Relative completeness is not universal completeness}

A contract may omit a guard that later proves necessary.

The witness system cannot diagnose a condition the contract has not represented.

Thus guard completeness is always relative to the declared audit structure.


\PCsection{10. Prediction and Observation Sets}

A single predicted number is often insufficient to express calibration and uncertainty. We therefore compare sets of admissible outputs.

\PCsubsection{Definition 10.1 --- Predicted output set}

For an admitted instance \(\omega\), let

\[
\boxed{ P_{\mathcal P}(\omega) \subseteq \mathcal Y }
\]

be the set of outputs permitted by the projection prediction and its declared uncertainty.

For an exact deterministic prediction,

\[
P_{\mathcal P}(\omega) = \left\{ \widehat y_{\mathcal P}(\omega) \right\}.
\]

\PCsubsection{Definition 10.2 --- Observation record}

An observation record is a tuple

\[
\boxed{ \mathcal O = (\omega,\mathcal M,y_{\mathrm{obs}},U_{\mathrm{obs}}) }
\]

where:

\begin{itemize}
\item \(\omega\) identifies the observed instance;
\item \(\mathcal M\) is the declared measurement procedure;
\item \(y_{\mathrm{obs}}\) is the reported observation;
\item \(U_{\mathrm{obs}}\) is its uncertainty description.
\end{itemize}
\PCsubsection{Definition 10.3 --- Observed output set}

A lawful observation record determines a set

\[
\boxed{ O_{\mathcal O} \subseteq \mathcal Y }
\]

of outputs compatible with the observation and its declared uncertainty.

For an exact observation,

\[
O_{\mathcal O} = \left\{ y_{\mathrm{obs}}\right\}.
\]

\PCsubsection{Definition 10.4 --- Empirical compatibility}

The prediction and observation are compatible when

\[
\boxed{ P_{\mathcal P}(\omega) \cap O_{\mathcal O} \neq \varnothing. }
\]

\PCsubsection{Definition 10.5 --- Record-level empirical contradiction}

The prediction and observation contradict one another when

\[
\boxed{ P_{\mathcal P}(\omega) \cap O_{\mathcal O} = \varnothing. }
\]


\PCsection{11. Falsification Conditions}

\PCsubsection{Theorem 11.1 --- Record-level contradiction theorem}

A record-level empirical contradiction is lawfully established exactly when:

\begin{enumerate}
\item \(A_{\mathcal P}(\omega)=\operatorname{true}\);
\item \(G_9(\mathcal P,\omega,\mathcal O)=\operatorname{true}\);
\item both \(P_{\mathcal P}(\omega)\) and \(O_{\mathcal O}\) are defined;
\item their intersection is empty.
\end{enumerate}
\begin{proof}
Conditions 1–3 establish that both the prediction and observation are lawful objects in the common output space.

Condition 4 states that no output is compatible with both.

Therefore, the prediction and observation cannot simultaneously be correct under the declared contract and observation protocol.

Conversely, if any of Conditions 1–3 fails, the comparison is not lawfully formed. If Condition 4 fails, at least one common output remains possible.

Thus all four conditions are necessary and sufficient.

\end{proof}

\PCsubsection{Corollary 11.2 --- Exact deterministic case}

If both prediction and observation are exact singletons, contradiction occurs exactly when

\[
\widehat y_{\mathcal P}(\omega) \neq y_{\mathrm{obs}}.
\]

\PCsubsection{Definition 11.3 --- Contract-level rejection rule}

A contract-level empirical decision rule is a declared function

\[
\boxed{ \Delta: \mathcal R \to \left\{ \operatorname{retain}, \operatorname{revise}, \operatorname{reject}, \operatorname{undetermined} \right\} }.
\]


where \(\mathcal R\) is a permitted collection of observation records and contradiction results.

\PCsubsection{Clarification 11.4 --- One contradiction need not reject every contract}

Whether one record-level contradiction rejects an entire contract depends on the declared decision rule.

The contradiction may instead indicate:

\begin{itemize}
\item a bad observation;
\item a calibration failure;
\item a phase-inference error;
\item an incorrect validity domain;
\item an omitted variable;
\item or failure of the complete projection hypothesis.
\end{itemize}
The decision rule must state how these possibilities are treated.


\PCsection{12. What Empirical Contradiction Can Falsify}

\PCsubsection{Theorem 12.1 --- Projection-specific falsification}

Suppose a lawful observation contradicts the output of \(\mathcal P\).

Then at least one empirical or projection-layer claim used by \(\mathcal P\) is false, inadequate, or incomplete.

The contradiction does not, by itself, falsify an exact Book 0 identity.

\begin{proof}
The prediction has the form

\[
\widehat y_{\mathcal P} = \Gamma_\theta \circ E_{\mathbf F} \circ \Phi.
\]

A contradiction can arise from failure of:

\begin{itemize}
\item the phase-inference map \(\Phi\);
\item relevance of the selected structural profile \(\mathbf F\);
\item calibration \(\theta\);
\item readout \(\Gamma_\theta\);
\item validity-domain declaration;
\item uncertainty model;
\item observation protocol;
\item or an external assumption coupling these components.
\end{itemize}
The Book 0 identity constrains the structural expressions internally on their mathematical domain.

An external mismatch does not exhibit a counterexample to that mathematical identity unless it supplies an admissible \(x\in\mathcal D\) for which the two sides of the identity are unequal.

Therefore, the contradiction targets the projection or external model, not automatically the structural theorem.

\end{proof}

\PCsubsection{Corollary 12.2 --- Structural identity versus relevance hypothesis}

A structural identity may remain true while the hypothesis that it models a particular external system is false.

\PCsubsection{Corollary 12.3 --- Failed physical interpretation}

If a projection interpreting Flatwave as a particular conserved physical quantity is contradicted, the contradiction may falsify that interpretation without affecting

\[
\operatorname{Flatwave} = \operatorname{sgn}(\sin 2x).
\]



\PCsection{13. Successful Prediction and Its Limits}

\PCsubsection{Theorem 13.1 --- Compatibility is weaker than confirmation of uniqueness}

If

\[
P_{\mathcal P}(\omega) \cap O_{\mathcal O} \neq \varnothing,
\]

then the projection is compatible with that observation record.

It does not follow that:

\begin{itemize}
\item the projection is uniquely correct;
\item the phase inference is unique;
\item the calibration is uniquely determined;
\item competing projections are false;
\item the structural profile is necessary.
\end{itemize}
\begin{proof}
Compatibility means only that at least one output is permitted by both prediction and observation.

S1.I proved that reparameterized internal descriptions may produce the same output.

A noninjective readout may erase structural differences.

Multiple projections may have overlapping predicted sets.

Therefore, compatibility does not establish uniqueness or necessity.

\end{proof}

\PCsubsection{Corollary 13.2 --- Repeated success strengthens evidence but not logical necessity}

A larger collection of successful tests may reduce the set of viable alternatives.

It does not transform empirical adequacy into a mathematical proof of unique ontology.


\PCsection{14. Domain Restriction}

\PCsubsection{Definition 14.1 --- Restricted contract}

Let

\[
\Omega_{\mathcal P'}\subseteq\Omega_{\mathcal P}.
\]

Define the restricted contract \(\mathcal P'\) by retaining the same:

\begin{itemize}
\item phase-inference rule;
\item structural profile;
\item readout;
\item calibration;
\item uncertainty formalism;
\end{itemize}
but limiting evaluation to \(\Omega_{\mathcal P'}\).

\PCsubsection{Theorem 14.2 --- Restriction preserves instance lawfulness}

If \(\mathcal P\) is lawful on \(\Omega_{\mathcal P}\), then its restriction \(\mathcal P'\) is lawful on every retained instance in \(\Omega_{\mathcal P'}\).

\begin{proof}
Every retained instance already belonged to the lawful validity domain of \(\mathcal P\).

All maps and guards are unchanged on those instances.

Therefore, every guard that succeeded under \(\mathcal P\) continues to succeed under \(\mathcal P'\).

\end{proof}

\PCsubsection{Corollary 14.3 --- Restriction reduces coverage}

A restricted contract may become more conservative, but it makes fewer predictions.

Improved reliability due solely to removing difficult cases must not be reported as increased explanatory range.


\PCsection{15. Domain Extension}

\PCsubsection{Definition 15.1 --- Extended contract}

An extension proposes a validity domain

\[
\Omega_{\mathcal P^{+}} \supsetneq \Omega_{\mathcal P}.
\]

\PCsubsection{Theorem 15.2 --- Extension burden}

Lawfulness on \(\Omega_{\mathcal P}\) does not imply lawfulness on the added region

\[
\Omega_{\mathcal P^{+}} \setminus \Omega_{\mathcal P}.
\]

Every added instance must pass the complete guard sequence.

\begin{proof}
The original lawfulness proof concerns only instances in \(\Omega_{\mathcal P}\).

No conclusion has been established for instances outside that set.

The added instances may require:

\begin{itemize}
\item new phase branches;
\item extrapolated calibration;
\item new structural domains;
\item boundary handling;
\item different uncertainty;
\item or a new readout.
\end{itemize}
Therefore, the old proof does not cover the new region.

\end{proof}

\PCsubsection{Corollary 15.3 --- No silent extrapolation}

A contract may not extend its prediction by merely evaluating the same formula outside its declared validity domain.

Formula evaluability is not sufficient evidence of projection validity.

\PCsubsection{Corollary 15.4 --- Extension may create a new contract}

When an extension changes inference, calibration, branch handling, or uncertainty, it should be assigned a new contract identifier.


\PCsection{16. Calibration Failure and Drift}

\PCsubsection{Definition 16.1 --- Calibration provenance}

Calibration provenance records how each component of \(\theta\) was obtained.

Possible statuses include:

\begin{itemize}
\item fixed by convention;
\item directly measured;
\item estimated from a training set;
\item inferred jointly with the model;
\item imported from another experiment.
\end{itemize}
\PCsubsection{Definition 16.2 --- Calibration domain}

Let

\[
\Omega_\theta\subseteq\Omega
\]

be the set on which the calibration is declared valid.

\PCsubsection{Theorem 16.3 --- Calibration-domain guard}

A contract using calibration \(\theta\) is not lawful for

\[
\omega\notin\Omega_\theta
\]

unless an extrapolation rule is explicitly declared and tested.

\begin{proof}
The readout depends on \(\theta\).

The validity of that parameter has been established only on \(\Omega_\theta\).

Outside this domain, the required calibration guard lacks support.

Therefore, the instance must be rejected or treated under a declared extrapolation contract.

\end{proof}

\PCsubsection{Definition 16.4 --- Calibration drift}

Calibration drift occurs when a parameter treated as fixed changes with external conditions not represented by the contract.

\PCsubsection{Corollary 16.5 --- Drift may mimic structural failure}

An empirical contradiction produced by calibration drift does not prove that the selected Book 0 profile is structurally unsuitable.

The calibration and structural hypotheses must be tested separately where possible.


\PCsection{17. Branch Ambiguity}

\PCsubsection{Definition 17.1 --- Candidate phase set}

For an external instance \(\omega\), let

\[
C_\Phi(\omega) \subseteq \mathcal D
\]

be the set of phase values compatible with the inference data before branch selection.

\PCsubsection{Definition 17.2 --- Unique inference}

Phase inference is unique when

\[
|C_\Phi(\omega)|=1.
\]

\PCsubsection{Definition 17.3 --- Branch ambiguity}

Branch ambiguity occurs when

\[
|C_\Phi(\omega)|>1
\]

and no declared rule selects one candidate or combines them.

\PCsubsection{Theorem 17.4 --- Ambiguous phase produces ambiguous structural state}

If there exist

\[
x_1,x_2\in C_\Phi(\omega)
\]

such that

\[
E_{\mathbf F}(x_1) \neq E_{\mathbf F}(x_2),
\]

then the structural state is not uniquely determined by the instance.

\begin{proof}
Both \(x_1\) and \(x_2\) satisfy the inference data.

They produce distinct structural states.

Without a selection or aggregation rule, the contract does not determine which state to use.

\end{proof}

\PCsubsection{Corollary 17.5 --- Flatwave branch ambiguity}

If two candidate phases lie in quadrants of opposite sign, then

\[
\operatorname{Flatwave}(x_1) = -\operatorname{Flatwave}(x_2).
\]

A contract cannot assign a unique Flatwave sign until the quadrant ambiguity is resolved.


\PCsection{18. Failure Locality and Contract Repair}

\PCsubsection{Definition 18.1 --- Contract field map}

Associate each witness with the contract field or fields it directly concerns.

For example:

\[
\begin{array}{l|l} \text{Witness} & \text{Primary field}\\ \hline \texttt{PHASE-NOT-INFERABLE} & \Phi \\ \texttt{PHASE-OR-BRANCH-AMBIGUOUS} & \Phi \\ \texttt{STRUCTURAL-EVALUATION-UNDEFINED} & \mathbf F,\mathcal D_{\mathbf F} \\ \texttt{CALIBRATION-MISSING-OR-INVALID} & \theta \\ \texttt{READOUT-OUTSIDE-DOMAIN} & \Gamma_\theta \\ \texttt{UNCERTAINTY-NOT-AVAILABLE} & U_{\mathcal P} \end{array}
\]

\PCsubsection{Definition 18.2 --- Local repair}

A repair is local when it modifies only the contract fields implicated by the witness, together with any fields that depend on them.

\PCsubsection{Theorem 18.3 --- Repair noninterference condition}

Suppose a repaired contract \(\mathcal P'\) changes a field \(C\).

Every theorem or validation result depending on \(C\) must be rechecked.

Results independent of \(C\) remain unchanged.

\begin{proof}
A derived result can change only if one of its premises or definitions changes.

If the result depends on \(C\), changing \(C\) may alter the result and requires reevaluation.

If its dependency graph excludes \(C\), the repair does not alter any of its premises.

Therefore, the result remains valid.

\end{proof}

\PCsubsection{Corollary 18.4 --- Calibration repair does not rewrite Book 0}

Changing \(\theta\) requires rechecking projected outputs.

It does not require reproving the structural identities of Book 0.

\PCsubsection{Corollary 18.5 --- Phase-rule repair requires broad retesting}

Changing \(\Phi\) may alter every downstream structural state and projected output.

All instance-level results depending on phase inference must therefore be recomputed.


\PCsection{19. Audit Record}

\PCsubsection{Definition 19.1 --- Evaluation record}

For each attempted instance, define an audit record

\[
\boxed{ \mathcal A_{\mathcal P}(\omega) = \left( \operatorname{id}_{\mathcal P}, \omega, \mathbf G, \Phi(\omega), E_{\mathbf F}(\Phi(\omega)), \theta, \widehat y_{\mathcal P}(\omega), U_{\mathcal P}(\omega), \mathcal W_{\mathcal P}(\omega) \right), }
\]

with undefined entries explicitly marked when a prior guard fails.

Here

\[
\mathbf G = (G_0,\ldots,G_8)
\]

is the guard vector.

\PCsubsection{Definition 19.2 --- Contract identifier}

Every materially distinct contract should have a distinct identifier or version.

A material change includes:

\begin{itemize}
\item validity-domain change;
\item phase-inference change;
\item structural-profile change;
\item calibration change;
\item readout change;
\item uncertainty-model change;
\item boundary-rule change.
\end{itemize}
\PCsubsection{Theorem 19.3 --- Reproducibility requirement}

A projected result is reproducible from its audit record when the record contains every non-derived input required by the contract.

\begin{proof}
The prediction is a deterministic composition of declared maps once all non-derived inputs and parameters are supplied.

If the audit record contains these values and the contract version, another evaluator can repeat the same composition.

If any required external input or parameter is absent, exact reproduction may be impossible.


\end{proof}

\PCsection{PROCEDURE DIVISION}

\PCsection{20. Canonical Evaluation Procedure}

The following procedure expresses the projection airlock in a COBOL-like form.

EVALUATE-PROJECTION.


PERFORM CHECK-CONTRACT-COMPLETENESS.

IF CONTRACT-INCOMPLETE

RETURN FAILURE-WITNESS-SET

END-IF.


PERFORM CHECK-INSTANCE-DOMAIN.

IF INSTANCE-OUTSIDE-VALIDITY-DOMAIN

RETURN FAILURE-WITNESS-SET

END-IF.


PERFORM INFER-PHASE.

IF PHASE-NOT-INFERABLE

RETURN FAILURE-WITNESS-SET

END-IF.


PERFORM RESOLVE-BRANCH.

IF PHASE-OR-BRANCH-AMBIGUOUS

RETURN FAILURE-WITNESS-SET

END-IF.


PERFORM CHECK-STRUCTURAL-DOMAIN.

IF PHASE-OUTSIDE-STRUCTURAL-DOMAIN

RETURN FAILURE-WITNESS-SET

END-IF.


PERFORM EVALUATE-STRUCTURAL-PROFILE.

IF STRUCTURAL-EVALUATION-UNDEFINED

RETURN FAILURE-WITNESS-SET

END-IF.


PERFORM CHECK-CALIBRATION.

IF CALIBRATION-MISSING-OR-INVALID

RETURN FAILURE-WITNESS-SET

END-IF.


PERFORM CHECK-READOUT-DOMAIN.

IF READOUT-OUTSIDE-DOMAIN

RETURN FAILURE-WITNESS-SET

END-IF.


PERFORM APPLY-READOUT.

PERFORM ASSIGN-UNCERTAINTY.


IF UNCERTAINTY-NOT-AVAILABLE

RETURN FAILURE-WITNESS-SET

END-IF.


RETURN PROJECTED-OUTPUT

PROJECTED-UNCERTAINTY

COMPLETE-AUDIT-RECORD.

\PCsubsection{20.1 Empirical comparison procedure}

COMPARE-WITH-OBSERVATION.


PERFORM EVALUATE-PROJECTION.


IF FAILURE-WITNESS-SET IS NONEMPTY

RETURN NOT-TESTABLE-UNDER-THIS-CONTRACT

END-IF.


PERFORM VALIDATE-OBSERVATION-RECORD.


IF OBSERVATION-RECORD-INVALID

RETURN OBSERVATION-FAILURE

END-IF.


FORM PREDICTED-OUTPUT-SET.

FORM OBSERVED-OUTPUT-SET.


IF INTERSECTION IS EMPTY

RETURN PREDICTION-OBSERVATION-DISJOINT

ELSE

RETURN EMPIRICALLY-COMPATIBLE

END-IF.


\PCsection{21. Theorem Ledger}

\PCsubsection{Theorem II.1 --- Empty-witness criterion}

\[
A_{\mathcal P}(\omega) \iff \mathcal W_{\mathcal P}(\omega)=\varnothing.
\]

\PCsubsection{Theorem II.2 --- Totalized evaluation}

Every instance returns either:

\begin{itemize}
\item an output;
\item or a nonempty witness set.
\end{itemize}
\PCsubsection{Theorem II.3 --- Non-evaluation theorem}

An inadmissible instance cannot produce an empirical contradiction because it produces no lawful prediction.

\PCsubsection{Theorem II.4 --- Witness soundness}

Every returned canonical witness corresponds to a genuinely failed declared guard.

\PCsubsection{Theorem II.5 --- Relative witness completeness}

Every failed declared guard produces a witness.

\PCsubsection{Theorem II.6 --- Record-level contradiction}

A lawful empirical contradiction occurs exactly when the admitted prediction set and valid observation set are disjoint.

\PCsubsection{Theorem II.7 --- Projection-specific falsification}

External contradiction targets the projection contract or its interpretation, not automatically the Book 0 identity.

\PCsubsection{Theorem II.8 --- Restriction preservation}

Restricting a lawful contract to a smaller validity domain preserves lawfulness on retained instances.

\PCsubsection{Theorem II.9 --- Extension burden}

Extending the validity domain requires a new guard proof on every added region.

\PCsubsection{Theorem II.10 --- Ambiguous-phase theorem}

Distinct admissible phase candidates that produce distinct structural states prevent unique projection unless a selection or aggregation rule is declared.

\PCsubsection{Theorem II.11 --- Repair noninterference}

Only results depending on a modified contract field require revalidation.


\PCsection{22. Negative Theorems}

\PCsubsection{Negative Theorem 22.1 --- No prediction from undefinedness}

An undefined structural expression is not a numerical prediction.

\PCsubsection{Negative Theorem 22.2 --- No falsification without admission}

An instance outside the contract's validity domain cannot falsify a prediction the contract did not make.

\PCsubsection{Negative Theorem 22.3 --- No silent branch choice}

A projection may not select one of several admissible phases without a declared selection rule.

\PCsubsection{Negative Theorem 22.4 --- No extrapolation by formula alone}

The ability to evaluate a formula outside its calibration or validity domain does not establish projection validity there.

\PCsubsection{Negative Theorem 22.5 --- No structural refutation by external mismatch alone}

An external mismatch does not refute a Book 0 identity unless it supplies an actual mathematical counterexample on the theorem's domain.

\PCsubsection{Negative Theorem 22.6 --- No unique explanation from compatibility}

Empirical compatibility does not prove that the tested projection is the only compatible projection.

\PCsubsection{Negative Theorem 22.7 --- No witness suppression}

A failed guard may not be hidden merely because a downstream numerical value can still be calculated.


\PCsection{23. Terminal Theorem}

\PCsubsection{Theorem 23.1 --- Projection Airlock Theorem}

For every projection contract \(\mathcal P\) and external instance \(\omega\), exactly one of the following operational results is justified:

\begin{enumerate}
\item a lawful projected output with declared uncertainty;
\item a nonempty failure-witness set;
\item after a lawful observation comparison, empirical compatibility;
\item after a lawful observation comparison, record-level empirical contradiction.
\end{enumerate}
These outcomes are mutually distinguishable.

Undefined evaluation, insufficient contract specification, and empirical contradiction may not be substituted for one another.

\begin{proof}
The totalized evaluator separates lawful output from guard failure.

The observation guard separates evaluable instances from valid empirical tests.

The intersection rule separates empirical compatibility from empirical contradiction.

The outcome categories are defined by mutually exclusive conditions:

\begin{itemize}
\item empty versus nonempty projection witness sets;
\item valid versus invalid observation records;
\item nonempty versus empty prediction–observation intersections.
\end{itemize}
Therefore, each evaluated case has a determinate logical status.


\end{proof}

\PCsection{24. Closing Declaration}

S1.I established the form of a projection contract:

\[
\boxed{ \widehat y_{\mathcal P} = \Gamma_\theta \circ E_{\mathbf F} \circ \Phi. }
\]

S1.II establishes the admission procedure governing its use.

A projection is admitted only when every declared guard succeeds:

\[
\boxed{ A_{\mathcal P}(\omega) = \bigwedge_{j=0}^{8}G_j(\mathcal P,\omega). }
\]

The declared scope must satisfy

\[
\boxed{ \Omega_{\mathcal P}\subseteq\Omega_{\mathcal P}^{\max}, \qquad \Omega_{\mathcal P}^{\max}=\{\omega\in\Omega:A_{\mathcal P}^{\max}(\omega)=\operatorname{true}\}. }
\]

Failure is reported explicitly:

\[
\boxed{ \mathcal W_{\mathcal P}(\omega) = \left\{ w_j:G_j(\mathcal P,\omega)=\operatorname{false} \right\}. }
\]

The evaluator returns either a prediction or witnesses:

\[
\boxed{ \mathcal E_{\mathcal P}(\omega) \in \mathcal Y \sqcup \mathcal P^{+}(\mathfrak W). }
\]

Empirical contradiction requires lawful comparison:

\[
\boxed{ P_{\mathcal P}(\omega) \cap O_{\mathcal O} = \varnothing. }
\]

The resulting discipline is:

\[
\boxed{ \begin{gathered} \text{No output before admission;}\\ \text{no admission without passed guards;}\\ \text{no rejection without a witness;}\\ \text{no falsification without a lawful prediction and observation;}\\ \text{no extrapolation without a new domain proof.} \end{gathered} }
\]

\[
\boxed{ \text{END OF S1.II} }
\]

% ===== END S1_II =====
% ===== BEGIN S1_III =====
% Source: S1.iii.docx
\PCchapter{S1.III}{Measurement Protocols, Error, and Uncertainty}{Book 1 - Projection and Measurement}

\PCsubsection{Abstract}

S1.I defined the projection contract

\[
\widehat y_{\mathcal P} = \Gamma_\theta \circ E_{\mathbf F} \circ \Phi,
\]

and S1.II established the guards, failure witnesses, and empirical comparison conditions governing its use.

This essay defines the observation side of that comparison.

A measurement is not merely a reported number. It is a protocol specifying:

\begin{itemize}
\item the measurand;
\item the external instances to which the measurement applies;
\item sampling and preparation;
\item instrument response;
\item data reduction;
\item units;
\item quality guards;
\item uncertainty;
\item and failure behavior.
\end{itemize}
A measurement protocol is represented by

\[
\mathcal M = \left( \Omega_{\mathcal M}, \mu, S, I, A, \mathcal U_{\mathcal M}, \mathcal Q_{\mathcal M}, W_{\mathcal M} \right).
\]

The realized observation is produced by

\[
y_{\mathrm{obs}} = (A\circ I\circ S)(\omega),
\]

when all measurement guards succeed.

The essay distinguishes measurement error from model residual. If \(\mu(\omega)\) is the measurand value and \(\widehat y_{\mathcal P}(\omega)\) is the projected prediction, then

\[
y_{\mathrm{obs}} - \widehat y_{\mathcal P} = \underbrace{ \left( y_{\mathrm{obs}}-\mu(\omega) \right) }_{\text{measurement error}} + \underbrace{ \left( \mu(\omega)-\widehat y_{\mathcal P} \right) }_{\text{model discrepancy}}.
\]

A residual therefore cannot be assigned wholly to measurement error or wholly to the model without additional evidence.

Uncertainty is represented generally as a set, distribution, covariance structure, or other declared description of values compatible with the measurement. Exact propagation through the projection layer is defined by image sets and pushforward measures.

Within a single open quadrant, continuous structural functions propagate bounded phase uncertainty continuously. At Flatwave boundaries, however, arbitrarily small unresolved phase uncertainty can produce the two-valued result

\[
\{-1,+1\}.
\]

Near primitive poles, small phase uncertainty can also produce arbitrarily large output uncertainty. Linearized uncertainty propagation is therefore forbidden across a discontinuity or uncontrolled singularity.

Repeated measurement may reduce random dispersion under declared independence assumptions. It does not remove systematic bias. Correlation between repeated measurements limits the benefit of averaging.

The final result is a formal separation among:

\begin{itemize}
\item the quantity intended to be measured;
\item the realized observation;
\item measurement error;
\item uncertainty;
\item model discrepancy;
\item and prediction residual.
\end{itemize}

\PCsection{IDENTIFICATION DIVISION}

\PCsubsection{PROGRAM-ID}

S1-III-MEASUREMENT-ERROR-UNCERTAINTY

\PCsubsection{PURPOSE}

To define the observation protocol required for lawful comparison between projected predictions and external measurements.

\PCsubsection{PRINCIPAL QUESTION}

What must be declared before a reported value may be used as evidence for or against a projection contract?

\PCsubsection{NON-PURPOSE}

This essay does not:

\begin{itemize}
\item prescribe one universal instrument;
\item require one universal statistical philosophy;
\item assume that all uncertainty is probabilistic;
\item equate uncertainty with error;
\item equate precision with accuracy;
\item infer a true value merely from repeated measurements;
\item select the calibration model of S1.IV;
\item determine whether one projection is superior to another;
\item alter any theorem of Book 0.
\end{itemize}

\PCsection{ENVIRONMENT DIVISION}

\PCsubsection{1. Inherited Projection Environment}

From S1.I, a projection contract is

\[
\mathcal P = \left( \Omega, \Omega_{\mathcal P}, \Phi, \mathbf F, \Gamma_\theta, U_{\mathcal P}, W_{\mathcal P} \right).
\]

Its prediction map is

\[
\widehat y_{\mathcal P}(\omega) = \Gamma_\theta \left( E_{\mathbf F}(\Phi(\omega)) \right).
\]

From S1.II, empirical comparison is lawful only when:

\begin{enumerate}
\item the instance passes every projection guard;
\item the observation record passes every measurement guard;
\item the prediction and observation belong to a common output space;
\item their uncertainty descriptions are declared.
\end{enumerate}
\PCsubsection{1.1 Projection and measurement are distinct}

The projection contract answers:

\[
\text{What does the model predict?}
\]

The measurement protocol answers:

\[
\text{What was observed, and under what conditions?}
\]

Neither answer may be substituted for the other.


\PCsection{METHODOLOGICAL AXIOM DIVISION}

\PCsubsection{Axiom M.1 --- Observation Declaration}

Every empirical claim used to test a projection must be supported by an explicitly declared measurement protocol.

A reported number without its measurand, units, method, applicability conditions, uncertainty, and failure status is not a complete observation record.

\PCsubsection{Axiom M.2 --- Prediction–Observation Separation}

The observation record must declare every way in which prediction, model choice, or expected outcome influenced:

\begin{itemize}
\item sampling;
\item exclusion;
\item preprocessing;
\item instrument settings;
\item data reduction;
\item uncertainty estimation;
\item or reporting.
\end{itemize}
A prediction-dependent procedure is not automatically invalid, but the dependence may not remain hidden.


\PCsection{DATA DIVISION}

\PCsection{2. Measurands}

\PCsubsection{Definition 2.1 --- Measurement instance space}

Let

\[
\Omega_{\mathcal M} \subseteq \Omega
\]

be the external instances for which a measurement protocol is declared.

\PCsubsection{Definition 2.2 --- Measurand space}

Let

\[
\mathcal Y
\]

be the declared space in which the target quantity takes values.

The measurand space must include its relevant structure, such as:

\begin{itemize}
\item scalar or vector form;
\item units;
\item coordinate convention;
\item reference frame;
\item category definitions;
\item permissible range.
\end{itemize}
\PCsubsection{Definition 2.3 --- Measurand}

A measurand is a declared target map

\[
\boxed{ \mu: \Omega_{\mathcal M} \to \mathcal Y. }
\]

The value

\[
\mu(\omega)
\]

is the quantity the protocol intends to determine for instance \(\omega\).

\PCsubsection{Clarification 2.4 --- Measurand value need not be known}

Declaring the measurand does not imply that its value is exactly known.

The measurand identifies the target of measurement.

The measurement result is an estimate, report, or constraint concerning that target.

\PCsubsection{Example 2.5 --- Incomplete measurand}

The phrase “measure the temperature” is incomplete when the following remain unspecified:

\begin{itemize}
\item which object or region;
\item at what time;
\item under which equilibrium assumptions;
\item on which temperature scale;
\item with what spatial or temporal averaging.
\end{itemize}
Different specifications define different measurands.


\PCsection{3. Measurement Protocol}

\PCsubsection{Definition 3.1 --- Sampling map}

A sampling map is

\[
\boxed{ S: \Omega_{\mathcal M} \to \mathcal X, }
\]

where \(\mathcal X\) is the space of samples, signals, specimens, or prepared inputs supplied to the instrument.

The sampling map includes any declared:

\begin{itemize}
\item selection;
\item timing;
\item spatial window;
\item averaging;
\item preparation;
\item filtering;
\item or transformation occurring before instrument response.
\end{itemize}
\PCsubsection{Definition 3.2 --- Instrument-response map}

An instrument-response map is

\[
\boxed{ I: \mathcal X \to \mathcal R, }
\]

where \(\mathcal R\) is the raw-record space.

The instrument map may be deterministic, stochastic, or partially unknown.

\PCsubsection{Definition 3.3 --- Analysis map}

An analysis map is

\[
\boxed{ A: \mathcal R \to \mathcal Y_{\mathrm{rep}}, }
\]

where \(\mathcal Y_{\mathrm{rep}}\) is the reported-value space.

The analysis map includes:

\begin{itemize}
\item digitization;
\item preprocessing;
\item baseline subtraction;
\item transformation;
\item fitting;
\item aggregation;
\item classification;
\item rounding;
\item and reporting conventions.
\end{itemize}
\PCsubsection{Definition 3.4 --- Realized observation}

For a valid measurement instance,

\[
\boxed{ y_{\mathrm{obs}}(\omega) = A(I(S(\omega))). }
\]

This value is the reported observation, not necessarily the measurand value.

\PCsubsection{Definition 3.5 --- Measurement uncertainty assignment}

The measurement uncertainty assignment is

\[
\boxed{ \mathcal U_{\mathcal M}: \Omega_{\mathcal M} \to \mathfrak U(\mathcal Y). }
\]

\PCsubsection{Definition 3.6 --- Measurement quality guards}

Let

\[
\mathcal Q_{\mathcal M} = (Q_0,\ldots,Q_n)
\]

be the declared collection of guards required for a valid observation.

\PCsubsection{Definition 3.7 --- Measurement failure map}

\[
\boxed{ W_{\mathcal M}: \Omega\setminus\Omega_{\mathcal M}^{\ast} \to \mathfrak W_{\mathcal M}, }
\]

where \(\Omega_{\mathcal M}^{\ast}\) is the set of instances passing every measurement guard.

\PCsubsection{Definition 3.8 --- Measurement protocol}

A measurement protocol is the tuple

\[
\boxed{ \mathcal M = \left( \Omega_{\mathcal M}, \mu, S, I, A, \mathcal U_{\mathcal M}, \mathcal Q_{\mathcal M}, W_{\mathcal M} \right). }
\]


\PCsection{4. Measurement Guards}

A protocol should declare guards appropriate to its instrument and measurand.

The following form a recommended minimum.

\PCsubsection{Definition 4.1 --- Protocol-completeness guard}

The measurand, sampling rule, instrument, analysis, units, uncertainty method, and failure behavior must be declared.

Recommended witness:

\[
\texttt{MEASUREMENT-PROTOCOL-INCOMPLETE}.
\]

\PCsubsection{Definition 4.2 --- Instance-applicability guard}

The instance must lie within the protocol's declared measurement domain.

Recommended witness:

\[
\texttt{INSTANCE-OUTSIDE-MEASUREMENT-DOMAIN}.
\]

\PCsubsection{Definition 4.3 --- Sampling-validity guard}

The obtained sample or signal must be suitable for the declared measurand.

Recommended witness:

\[
\texttt{SAMPLE-NOT-REPRESENTATIVE}.
\]

\PCsubsection{Definition 4.4 --- Instrument-range guard}

The instrument response must lie inside its declared operating range.

Recommended witnesses include:

\begin{itemize}
\item INSTRUMENT-SATURATED;
\item SIGNAL-BELOW-DETECTION-LIMIT;
\item INSTRUMENT-OUTSIDE-RATED-RANGE.
\end{itemize}
\PCsubsection{Definition 4.5 --- Resolution guard}

The instrument resolution must be sufficient for the reported number of digits or classification.

Recommended witness:

\[
\texttt{REPORTED-PRECISION-EXCEEDS-RESOLUTION}.
\]

\PCsubsection{Definition 4.6 --- Calibration-status guard}

The instrument must possess a valid calibration status for the measurement conditions.

Recommended witness:

\[
\texttt{INSTRUMENT-CALIBRATION-INVALID}.
\]

The full theory of calibration is deferred to S1.IV.

\PCsubsection{Definition 4.7 --- Analysis-validity guard}

The analysis procedure must be defined for the raw record and must not conceal failed preprocessing assumptions.

Recommended witness:

\[
\texttt{ANALYSIS-PROCEDURE-INVALID}.
\]

\PCsubsection{Definition 4.8 --- Unit-consistency guard}

The reported value, uncertainty, and projection output must use compatible units or a declared conversion.

Recommended witness:

\[
\texttt{MEASUREMENT-UNIT-MISMATCH}.
\]

\PCsubsection{Definition 4.9 --- Uncertainty-completeness guard}

An empirical result must include the uncertainty representation required by its protocol.

Recommended witness:

\[
\texttt{MEASUREMENT-UNCERTAINTY-MISSING}.
\]

\PCsubsection{Definition 4.10 --- Independence-disclosure guard}

Any prediction-dependent sampling, filtering, exclusion, or analysis must be disclosed.

Recommended witness:

\[
\texttt{PREDICTION-DEPENDENCE-UNDISCLOSED}.
\]


\PCsection{5. Observation Record}

\PCsubsection{Definition 5.1 --- Complete observation record}

A complete observation record is

\[
\boxed{ \mathcal O = \left( \operatorname{id}_{\mathcal M}, \omega, \mu, y_{\mathrm{obs}}, \mathcal U_{\mathcal M}(\omega), \mathbf Q, \mathbf m \right), }
\]

where:

\begin{itemize}
\item \(\operatorname{id}_{\mathcal M}\) identifies the protocol version;
\item \(\omega\) identifies the measured instance;
\item \(\mu\) identifies the measurand;
\item \(y_{\mathrm{obs}}\) is the reported observation;
\item \(\mathcal U_{\mathcal M}(\omega)\) is its uncertainty description;
\item \(\mathbf Q\) is the measurement-guard vector;
\item \(\mathbf m\) contains required metadata.
\end{itemize}
\PCsubsection{Definition 5.2 --- Measurement witness set}

\[
\boxed{ \mathcal W_{\mathcal M}(\omega) = \left\{ w_j: Q_j(\mathcal M,\omega)=\operatorname{false} \right\} }.
\]


\PCsubsection{Theorem 5.3 --- Valid-observation criterion}

An observation record is valid exactly when

\[
\boxed{ \mathcal W_{\mathcal M}(\omega)=\varnothing. }
\]

\begin{proof}
The record is valid precisely when every declared measurement guard succeeds.

A witness is included precisely when its guard fails.

Therefore, every guard succeeds exactly when the witness set is empty.

\end{proof}

\PCsubsection{Corollary 5.4 --- A failed measurement is not an observation against the model}

If

\[
\mathcal W_{\mathcal M}(\omega)\neq\varnothing,
\]

the record cannot be used as a lawful contradiction of the projection until the failure is resolved or the protocol is replaced.


\PCsection{6. Error}

\PCsubsection{Definition 6.1 --- Measurement error}

When subtraction is defined in \(\mathcal Y\), the measurement error is

\[
\boxed{ e_{\mathcal M}(\omega) = y_{\mathrm{obs}}(\omega)-\mu(\omega). }
\]

\PCsubsection{Clarification 6.2 --- Error is usually not directly known}

The exact measurement error requires knowledge of the measurand value.

If \(\mu(\omega)\) is unknown, the exact error is also unknown.

An uncertainty estimate is not the same thing as the unknown error.

\PCsubsection{Definition 6.3 --- Signed error}

For a scalar measurand,

\[
e_{\mathcal M}>0
\]

indicates an observation above the measurand, while

\[
e_{\mathcal M}<0
\]

indicates an observation below it.

\PCsubsection{Definition 6.4 --- Absolute error}

\[
\boxed{ |e_{\mathcal M}| = |y_{\mathrm{obs}}-\mu|. }
\]

\PCsubsection{Definition 6.5 --- Relative error}

Where

\[
\mu(\omega)\neq0,
\]

define

\[
\boxed{ e_{\mathrm{rel}} = \frac{y_{\mathrm{obs}}-\mu}{\mu}. }
\]

Relative error is undefined when the reference measurand is zero unless another normalization is declared.


\PCsection{7. Bias and Random Error}

The following definitions require a declared probability model for repeated outcomes of the measurement process.

\PCsubsection{Definition 7.1 --- Expected observation}

Suppose the realized observation is represented by a random variable \(Y\) with finite expectation.

Then

\[
\mathbb E[Y]
\]

is the expected reported value under the declared repetition model.

\PCsubsection{Definition 7.2 --- Measurement bias}

\[
\boxed{ b = \mathbb E[Y]-\mu. }
\]

\PCsubsection{Definition 7.3 --- Random error component}

\[
\boxed{ \eta = Y-\mathbb E[Y]. }
\]

Then

\[
\mathbb E[\eta]=0.
\]

\PCsubsection{Theorem 7.4 --- Error decomposition}

\[
\boxed{ Y-\mu = b+\eta. }
\]

\begin{proof}
Add and subtract \(\mathbb E[Y]\):

\[
Y-\mu = Y-\mathbb E[Y] + \mathbb E[Y]-\mu.
\]

By definition,

\[
Y-\mathbb E[Y]=\eta
\]

and

\[
\mathbb E[Y]-\mu=b.
\]

Therefore,

\[
Y-\mu=b+\eta.
\]

\end{proof}

\PCsubsection{Clarification 7.5 --- Systematic and random are model-relative}

An effect may appear random under one description and systematic under a richer description that includes an omitted condition.

The classification depends on the declared repetition and conditioning model.


\PCsection{8. Accuracy and Precision}

\PCsubsection{Definition 8.1 --- Accuracy}

Accuracy concerns closeness of a measurement result, or expected measurement result, to the measurand.

A scalar measure of inaccuracy may be based on:

\[
|y_{\mathrm{obs}}-\mu|
\]

or

\[
|\mathbb E[Y]-\mu|.
\]

\PCsubsection{Definition 8.2 --- Precision}

Precision concerns the dispersion of repeated measurement results under declared repeatability or reproducibility conditions.

Possible precision measures include:

\begin{itemize}
\item variance;
\item standard deviation;
\item interquartile range;
\item repeatability interval;
\item resolution.
\end{itemize}
\PCsubsection{Theorem 8.3 --- Precision does not imply accuracy}

Arbitrarily precise measurements may remain inaccurate.

\begin{proof}
Let every repeated measurement return the same value

\[
Y=\mu+c
\]

for a fixed nonzero constant \(c\).

Then

\[
\operatorname{Var}(Y)=0,
\]

so the measurement is perfectly precise under the repetition model.

However,

\[
Y-\mu=c\neq0.
\]

Therefore, perfect precision does not imply accuracy.

\end{proof}

\PCsubsection{Theorem 8.4 --- Accuracy does not imply high precision}

A measurement process may be unbiased but highly dispersed.

\begin{proof}
Let \(Y\) satisfy

\[
\mathbb E[Y]=\mu
\]

while

\[
\operatorname{Var}(Y)
\]

is large.

Then the bias is zero, but individual results may be far from the measurand.

Therefore, zero bias does not imply high precision.


\end{proof}

\PCsection{9. Resolution, Sensitivity, and Detection}

\PCsubsection{Definition 9.1 --- Resolution}

Resolution is the smallest declared change in input or measurand that the instrument or reporting system can distinguish under specified conditions.

\PCsubsection{Definition 9.2 --- Sensitivity}

For a differentiable instrument-response function \(r(u)\), local sensitivity is

\[
\boxed{ \frac{dr}{du}. }
\]

Sensitivity measures response change per input change.

It is not identical to resolution or accuracy.

\PCsubsection{Definition 9.3 --- Detection limit}

A detection limit is a declared threshold below which the protocol does not reliably distinguish signal from background under its stated decision rule.

\PCsubsection{Definition 9.4 --- Quantization}

If an instrument reports values only on a discrete grid of spacing \(\Delta\), then the reporting process introduces quantization.

Under nearest-grid rounding, the quantization error satisfies

\[
\boxed{ |e_q| \leq \frac{\Delta}{2}. }
\]

This bound does not include other sources of error.


\PCsection{10. Measurement Uncertainty}

\PCsubsection{Definition 10.1 --- Uncertainty description}

A measurement uncertainty description is a declared representation of values, distributions, or parameter states compatible with the measurement information.

It may be:

\begin{itemize}
\item a bounded set;
\item an interval;
\item a covariance matrix;
\item a probability distribution;
\item a confidence region;
\item a credible region;
\item an ensemble;
\item a qualitative classification.
\end{itemize}
\PCsubsection{Definition 10.2 --- Observed admissible-value set}

For a set-valued uncertainty representation, define

\[
\boxed{ O_{\mathcal O} \subseteq \mathcal Y }
\]

as the set of measurand values compatible with the observation under the declared protocol.

\PCsubsection{Definition 10.3 --- Interval uncertainty}

For a scalar observation, an interval representation has the form

\[
\boxed{ O_{\mathcal O} = [y_{\mathrm{obs}}-u_{-}, y_{\mathrm{obs}}+u_{+}]. }
\]

The symmetric case is

\[
[y_{\mathrm{obs}}-u, y_{\mathrm{obs}}+u].
\]

\PCsubsection{Definition 10.4 --- Probabilistic uncertainty}

A probabilistic uncertainty description assigns a probability measure

\[
\boxed{ \nu_{\mathcal O} }
\]

to the measurand or to the measurement process under declared assumptions.

The interpretation of \(\nu_{\mathcal O}\) must be stated.

\PCsubsection{Clarification 10.5 --- Uncertainty is not automatically probability}

Incomplete information may be represented by a set without assigning relative probabilities to the values inside it.

\PCsubsection{Clarification 10.6 --- Confidence and credibility are different declarations}

A confidence interval is defined by the long-run coverage properties of a procedure under a sampling model.

A credible interval is defined by probability under a declared posterior distribution.

Neither term should be used without its supporting interpretation.


\PCsection{11. Measurement Error and Model Residual}

\PCsubsection{Definition 11.1 --- Projected prediction}

For an admitted instance,

\[
\widehat y_{\mathcal P}(\omega)
\]

is the output of the projection contract.

\PCsubsection{Definition 11.2 --- Residual}

Where subtraction is defined,

\[
\boxed{ r(\omega) = y_{\mathrm{obs}}(\omega) - \widehat y_{\mathcal P}(\omega). }
\]

\PCsubsection{Definition 11.3 --- Model discrepancy}

\[
\boxed{ d_{\mathcal P}(\omega) = \mu(\omega) - \widehat y_{\mathcal P}(\omega). }
\]

\PCsubsection{Theorem 11.4 --- Residual decomposition}

\[
\boxed{ r = e_{\mathcal M} + d_{\mathcal P}. }
\]

\begin{proof}
By definition,

\[
r = y_{\mathrm{obs}} - \widehat y_{\mathcal P}.
\]

Add and subtract \(\mu\):

\[
r = y_{\mathrm{obs}} - \mu + \mu - \widehat y_{\mathcal P}.
\]

The first difference is measurement error:

\[
e_{\mathcal M} = y_{\mathrm{obs}}-\mu.
\]

The second is model discrepancy:

\[
d_{\mathcal P} = \mu-\widehat y_{\mathcal P}.
\]

Therefore,

\[
r=e_{\mathcal M}+d_{\mathcal P}.
\]

\end{proof}

\PCsubsection{Corollary 11.5 --- A residual is not measurement error}

Unless

\[
d_{\mathcal P}=0
\]

has been independently established,

\[
r\neq e_{\mathcal M}
\]

in general.

\PCsubsection{Corollary 11.6 --- A residual is not model discrepancy}

Unless

\[
e_{\mathcal M}=0
\]

has been independently established,

\[
r\neq d_{\mathcal P}
\]

in general.

\PCsubsection{Negative Theorem 11.7 --- Residual attribution is underdetermined}

A nonzero residual alone does not determine how much arises from measurement error and how much arises from projection discrepancy.


\PCsection{12. Repeated Measurement}

Suppose repeated observations satisfy

\[
Y_i = \mu+b+\eta_i,
\]

where

\[
\mathbb E[\eta_i]=0.
\]

\PCsubsection{Definition 12.1 --- Sample mean}

\[
\boxed{ \overline Y_n = \frac1n \sum_{i=1}^{n}Y_i. }
\]

\PCsubsection{Theorem 12.2 --- Averaging does not remove bias}

\[
\boxed{ \mathbb E[\overline Y_n] = \mu+b. }
\]

\begin{proof}
By linearity of expectation,

\[
\mathbb E[\overline Y_n] = \frac1n \sum_{i=1}^{n} \mathbb E[Y_i].
\]

Each expectation is

\[
\mu+b.
\]

Therefore,

\[
\mathbb E[\overline Y_n] = \mu+b.
\]

\end{proof}

\PCsubsection{Corollary 12.3 --- Infinite repetition does not correct fixed bias}

Increasing \(n\) does not cause the expected mean to approach \(\mu\) unless \(b=0\).

\PCsubsection{Theorem 12.4 --- Independent equal-variance averaging}

If the random errors are independent and each has variance \(\sigma^2\), then

\[
\boxed{ \operatorname{Var}(\overline Y_n) = \frac{\sigma^2}{n}. }
\]

\begin{proof}
\[
\operatorname{Var}(\overline Y_n) = \frac1{n^2} \operatorname{Var} \left( \sum_{i=1}^{n}\eta_i \right).
\]

Under independence,

\[
\operatorname{Var} \left( \sum_{i=1}^{n}\eta_i \right) = n\sigma^2.
\]

Therefore,

\[
\operatorname{Var}(\overline Y_n) = \frac{n\sigma^2}{n^2} = \frac{\sigma^2}{n}.
\]

\end{proof}

\PCsubsection{Theorem 12.5 --- Equal-correlation averaging}

Suppose each random error has variance \(\sigma^2\), and every distinct pair has correlation \(\rho\). Then

\[
\boxed{ \operatorname{Var}(\overline Y_n) = \frac{\sigma^2}{n} \left( 1+(n-1)\rho \right). }
\]

\begin{proof}
The variance of the sum is

\[
n\sigma^2 + 2\binom n2\rho\sigma^2.
\]

Thus,

\[
\operatorname{Var} \left( \sum_{i=1}^{n}\eta_i \right) = n\sigma^2 + n(n-1)\rho\sigma^2.
\]

Divide by \(n^2\):

\[
\operatorname{Var}(\overline Y_n) = \frac{\sigma^2}{n} \left( 1+(n-1)\rho \right).
\]

\end{proof}

\PCsubsection{Corollary 12.6 --- Correlation limits averaging benefit}

When \(\rho>0\), repeated measurements contain less independent information than the formula \(\sigma^2/n\) would imply.

When \(\rho=1\),

\[
\operatorname{Var}(\overline Y_n)=\sigma^2.
\]

Averaging provides no dispersion reduction.


\PCsection{13. Set-Valued Phase Inference}

Measurement uncertainty may affect the phase-inference map.

\PCsubsection{Definition 13.1 --- Candidate external-state set}

Let

\[
X_{\mathcal O} \subseteq \Omega
\]

be the external states compatible with an observation record.

\PCsubsection{Definition 13.2 --- Candidate phase set}

\[
\boxed{ C_\Phi(\mathcal O) = \left\{ \Phi(\omega): \omega\in X_{\mathcal O} \right\}. }
\]

When phase inference itself is set-valued, \(C_\Phi\) includes every contract-admissible inferred phase.

\PCsubsection{Definition 13.3 --- Candidate structural set}

For a structural profile \(\mathbf F\),

\[
\boxed{ C_{\mathbf F}(\mathcal O) = E_{\mathbf F} \left( C_\Phi(\mathcal O) \right). }
\]

\PCsubsection{Definition 13.4 --- Candidate projected-output set}

\[
\boxed{ P_{\mathcal P}(\mathcal O) = \Gamma_\theta \left( C_{\mathbf F}(\mathcal O) \right), }
\]

with calibration uncertainty included where required.

\PCsubsection{Theorem 13.5 --- Exact set propagation}

Every output generated by a candidate phase in \(C_\Phi\) belongs to \(P_{\mathcal P}\), and every member of \(P_{\mathcal P}\) arises from at least one candidate phase.

\begin{proof}
The set

\[
C_{\mathbf F} = E_{\mathbf F}(C_\Phi)
\]

is the image of the candidate phase set under structural evaluation.

The set

\[
P_{\mathcal P} = \Gamma_\theta(C_{\mathbf F})
\]

is its image under the readout.

This is precisely the definition of an image set.


\end{proof}

\PCsection{14. Continuous Uncertainty Propagation}

\PCsubsection{Theorem 14.1 --- Continuous-image compactness}

Let

\[
X\subset Q_k
\]

be compact, and let

\[
F:X\to\mathbb R
\]

be continuous.

Then

\[
F(X)
\]

is compact and therefore bounded.

\begin{proof}
The continuous image of a compact set is compact.

Every compact subset of \(\mathbb R\) is bounded.

\end{proof}

\PCsubsection{Theorem 14.2 --- Monotone interval propagation}

Let

\[
X=[a,b]\subset Q_k
\]

and suppose \(F\) is monotone increasing on \(X\). Then

\[
\boxed{ F(X) = [F(a),F(b)]. }
\]

If \(F\) is monotone decreasing, then

\[
\boxed{ F(X) = [F(b),F(a)]. }
\]

\begin{proof}
A continuous monotone function on an interval attains its extrema at the endpoints and assumes every intermediate value.

\end{proof}

\PCsubsection{Theorem 14.3 --- Lipschitz uncertainty bound}

Suppose

\[
|F(x)-F(y)| \leq L|x-y|
\]

for all \(x,y\in X\).

Then

\[
\boxed{ \operatorname{diam}(F(X)) \leq L\operatorname{diam}(X). }
\]

\begin{proof}
For any \(x,y\in X\),

\[
|F(x)-F(y)| \leq L|x-y| \leq L\operatorname{diam}(X).
\]

Taking the supremum over \(x,y\in X\) gives the result.

\end{proof}

\PCsubsection{Corollary 14.4 --- Derivative-based bound}

If \(F\) is differentiable on an interval \(X\) and

\[
|F'(x)|\leq L
\]

throughout \(X\), then

\[
\operatorname{diam}(F(X)) \leq L\operatorname{diam}(X).
\]

This follows from the mean-value theorem.


\PCsection{15. Probabilistic Propagation}

\PCsubsection{Definition 15.1 --- Phase probability measure}

Let

\[
\nu_\Phi
\]

be a declared probability measure over candidate phases.

\PCsubsection{Definition 15.2 --- Pushforward measure}

For a measurable structural function \(F\), define

\[
\boxed{ F_{\#}\nu_\Phi(B) = \nu_\Phi(F^{-1}(B)) }
\]

for measurable output sets \(B\).

\PCsubsection{Theorem 15.3 --- Exact probabilistic propagation}

The probability distribution of

\[
Y=F(X)
\]

for \(X\sim\nu_\Phi\) is

\[
\boxed{ \nu_Y = F_{\#}\nu_\Phi. }
\]

This is exact and does not require linearization.

\PCsubsection{Definition 15.4 --- Readout pushforward}

With readout \(\Gamma_\theta\),

\[
\boxed{ \nu_{\mathrm{out}} = (\Gamma_\theta\circ F)_{\#}\nu_\Phi. }
\]


\PCsection{16. Linearized Uncertainty Propagation}

\PCsubsection{Definition 16.1 --- First-order propagation}

For a differentiable scalar function \(F\) near a nominal phase \(x_0\),

\[
F(x) \approx F(x_0) + F'(x_0)(x-x_0).
\]

If the phase uncertainty has variance \(\sigma_x^2\), the first-order propagated variance is

\[
\boxed{ \sigma_F^2 \approx F'(x_0)^2\sigma_x^2. }
\]

\PCsubsection{Definition 16.2 --- Multivariate first-order propagation}

For a differentiable map

\[
g:\mathbb R^n\to\mathbb R^m
\]

with Jacobian \(J\), the first-order covariance approximation is

\[
\boxed{ \Sigma_y \approx J\Sigma_xJ^{\mathsf T}. }
\]

\PCsubsection{Theorem 16.3 --- Exactness for affine maps}

If

\[
g(x)=Ax+b,
\]

then

\[
\boxed{ \Sigma_y = A\Sigma_xA^{\mathsf T} }
\]

exactly.

\PCsubsection{Clarification 16.4 --- Nonlinear propagation is approximate}

For nonlinear \(g\), the Jacobian covariance formula neglects higher-order terms.

Its adequacy must be checked against:

\begin{itemize}
\item uncertainty width;
\item curvature;
\item singularity distance;
\item branch structure;
\item discontinuities.
\end{itemize}

\PCsection{17. Flatwave Boundary Uncertainty}

\PCsubsection{Theorem 17.1 --- Quadrant-stable Flatwave}

If the candidate phase set satisfies

\[
C_\Phi\subset Q_k,
\]

then

\[
\boxed{ \operatorname{Flatwave}(C_\Phi) = \{\varepsilon_k\}. }
\]

\begin{proof}
Flatwave is constant on \(Q_k\):

\[
\operatorname{Flatwave}(x)=\varepsilon_k.
\]

Therefore, every candidate phase produces the same value.

\end{proof}

\PCsubsection{Theorem 17.2 --- Boundary-straddling Flatwave uncertainty}

Suppose \(C_\Phi\) contains at least one phase in \(Q_{k-1}\) and at least one phase in \(Q_k\).

Then

\[
\boxed{ \{-1,+1\} \subseteq \operatorname{Flatwave}(C_\Phi). }
\]

If no other Flatwave values are permitted, then

\[
\boxed{ \operatorname{Flatwave}(C_\Phi) = \{-1,+1\}. }
\]

\begin{proof}
Adjacent quadrants have opposite signs:

\[
\varepsilon_{k-1}=-\varepsilon_k.
\]

One candidate phase therefore produces \(+1\), and the other produces \(-1\).

Flatwave assumes only those two values on its interior domain.

\end{proof}

\PCsubsection{Corollary 17.3 --- Arbitrarily narrow phase uncertainty may produce maximal sign uncertainty}

For every \(\delta>0\), a candidate phase interval of width less than \(\delta\) can straddle an excluded boundary and produce both Flatwave signs.

Thus small phase width does not imply small Flatwave uncertainty near a boundary.

\PCsubsection{Negative Theorem 17.4 --- No linearized Flatwave propagation across a wall}

Because Flatwave is discontinuous across an identified boundary, derivative-based uncertainty propagation cannot represent boundary-straddling uncertainty.

Inside a quadrant,

\[
\operatorname{Flatwave}'=0,
\]

but across the wall the output changes by \(2\).


\PCsection{18. Primitive Singular Amplification}

\PCsubsection{Theorem 18.1 --- No uniform primitive bound near a pole}

Let \(b\) be a boundary at which a primitive \(F\) has asymptotic form

\[
F(b+h)\sim\frac2h
\]

as \(h\to0^+\).

Then \(F\) is unbounded on every interval

\[
(b,b+\delta)
\]

with \(\delta>0\).

\begin{proof}
Since

\[
F(b+h)\sim\frac2h
\]

and \(2/h\to+\infty\), the function exceeds every finite bound for sufficiently small \(h\).

\end{proof}

\PCsubsection{Corollary 18.2 --- Small phase uncertainty may produce arbitrarily large primitive uncertainty}

A candidate phase set approaching a primitive pole may have arbitrarily small width while its structural image has arbitrarily large diameter.

\PCsubsection{Theorem 18.3 --- No uniform Lipschitz constant at a primitive pole}

No finite \(L\) provides a Lipschitz bound for the divergent primitive on every sufficiently small interval approaching the pole.

\begin{proof}
A Lipschitz function on a bounded interval has bounded variation and finite values throughout that interval.

The primitive diverges at the boundary approach.

Therefore, no finite uniform Lipschitz constant exists on intervals whose closure reaches the pole.

\end{proof}

\PCsubsection{Corollary 18.4 --- Linear propagation requires a singularity-distance guard}

A projection using primitive values must declare a minimum distance from the relevant excluded boundary or use an exact nonlinear propagation procedure.


\PCsection{19. Saw Uncertainty}

\PCsubsection{Theorem 19.1 --- Saw boundedness}

For every admissible phase,

\[
\boxed{ -1< \operatorname{saw}_{r}(x) < 1 }
\]

and

\[
\boxed{ -1< \operatorname{saw}_{x}(x) < 1. }
\]

Their quadrant-relative endpoint extensions lie in \([-1,1]\).

\begin{proof}
From Book 0,

\[
\operatorname{saw}_{r} = \varepsilon\lambda,
\]

and

\[
\operatorname{saw}_{x} = \varepsilon(1-\lambda),
\]

with

\[
0<\lambda<1.
\]

Therefore, both magnitudes are strictly less than \(1\) in the interior.

\end{proof}

\PCsubsection{Corollary 19.2 --- Saw uncertainty cannot diverge}

Unlike primitive uncertainty near a pole, saw-value uncertainty remains bounded.

\PCsubsection{Clarification 19.3 --- Bounded does not mean continuous across identified seams}

Each individual saw has unequal one-sided limits at a quadrant boundary.

A phase set straddling the boundary may therefore produce separated bounded values.


\PCsection{20. Uncertainty Budgets}

\PCsubsection{Definition 20.1 --- Uncertainty source}

An uncertainty source is a declared contribution arising from one stage of measurement or projection.

Possible sources include:

\begin{itemize}
\item sampling uncertainty;
\item instrument noise;
\item instrument bias uncertainty;
\item preprocessing uncertainty;
\item phase-inference uncertainty;
\item branch ambiguity;
\item calibration uncertainty;
\item readout approximation;
\item numerical approximation;
\item model discrepancy allowance.
\end{itemize}
\PCsubsection{Definition 20.2 --- Uncertainty budget}

An uncertainty budget is a ledger

\[
\boxed{ \mathcal B = \left\{ (B_i,\text{magnitude}_i,\text{dependence}_i,\text{method}_i) \right\} }.
\]


recording each included source and how it enters the final uncertainty description.

\PCsubsection{Definition 20.3 --- Independence declaration}

Sources may be combined by independent-error formulas only when their independence or negligible dependence is justified.

\PCsubsection{Theorem 20.4 --- Variance addition under zero covariance}

For random variables \(X_1,\ldots,X_n\),

\[
\operatorname{Var} \left( \sum_iX_i \right) = \sum_i\operatorname{Var}(X_i)
\]

exactly when the total covariance contribution is zero:

\[
\sum_{i\neq j} \operatorname{Cov}(X_i,X_j)=0.
\]

\begin{proof}
The general variance identity is

\[
\operatorname{Var} \left( \sum_iX_i \right) = \sum_i\operatorname{Var}(X_i) + \sum_{i\neq j} \operatorname{Cov}(X_i,X_j).
\]

The reduced formula holds exactly when the covariance sum vanishes.

\end{proof}

\PCsubsection{Negative Theorem 20.5 --- No root-sum-square without dependence assumptions}

The formula

\[
u_{\mathrm{total}} = \sqrt{ u_1^2+\cdots+u_n^2 }
\]

is not universally valid.

It requires an interpretation under which the corresponding contributions are uncorrelated standard uncertainties or satisfy an equivalent declared condition.

\PCsubsection{Negative Theorem 20.6 --- No double counting}

The same uncertainty contribution may not be included once in the measurement interval and again as an independent projection term unless the duplication is explicitly intended and justified.


\PCsection{21. Missing, Censored, and Selected Data}

\PCsubsection{Definition 21.1 --- Missing observation}

A value is missing when no report is available for an instance expected under the protocol.

\PCsubsection{Definition 21.2 --- Censored observation}

A value is censored when only a constraint is reported, such as

\[
y<L
\]

or

\[
y>U.
\]

\PCsubsection{Definition 21.3 --- Truncated protocol}

A protocol is truncated when instances outside a selection range are excluded from the recorded sample.

\PCsubsection{Theorem 21.4 --- A censored value is a set-valued observation}

The report

\[
y<L
\]

corresponds to an observed admissible set of the form

\[
O_{\mathcal O} \subseteq (-\infty,L),
\]

subject to the declared protocol and uncertainty.

It may not be replaced by the single value \(L\) without an added convention.

\PCsubsection{Negative Theorem 21.5 --- Missingness is not zero}

An absent observation may not be assigned the numerical value \(0\) unless zero is the explicit coding convention and is kept distinct from an actual measured zero.

\PCsubsection{Negative Theorem 21.6 --- Selection may alter empirical distributions}

A sample selected according to the observed or predicted value need not represent the distribution of the full instance population.

The selection rule must be included in the measurement protocol.


\PCsection{22. Compatibility with Projection Testing}

\PCsubsection{Definition 22.1 --- Predicted admissible-output set}

From S1.II,

\[
P_{\mathcal P}(\omega) \subseteq \mathcal Y
\]

is the output set compatible with the projection and its uncertainty.

\PCsubsection{Definition 22.2 --- Observed admissible-output set}

From the measurement protocol,

\[
O_{\mathcal O} \subseteq \mathcal Y
\]

is the output set compatible with the observation and its uncertainty.

\PCsubsection{Theorem 22.3 --- Compatibility condition}

The prediction and observation are compatible when

\[
\boxed{ P_{\mathcal P}(\omega) \cap O_{\mathcal O} \neq \varnothing. }
\]

\PCsubsection{Theorem 22.4 --- Contradiction condition}

They are contradictory at the record level when

\[
\boxed{ P_{\mathcal P}(\omega) \cap O_{\mathcal O} = \varnothing. }
\]

\PCsubsection{Clarification 22.5 --- Larger uncertainty weakens discrimination}

Enlarging either admissible set cannot convert a nonempty intersection into an empty one.

Therefore, wider uncertainty generally makes contradiction harder to establish.

\PCsubsection{Clarification 22.6 --- Artificially narrow uncertainty inflates apparent contradiction}

Understating uncertainty may create a disjoint-set result that disappears under a complete uncertainty budget.


\PCsection{PROCEDURE DIVISION}

\PCsection{23. Canonical Measurement Procedure}

PERFORM-MEASUREMENT.


PERFORM CHECK-MEASUREMENT-PROTOCOL.


IF MEASUREMENT-PROTOCOL-INCOMPLETE

RETURN MEASUREMENT-WITNESS-SET

END-IF.


PERFORM CHECK-INSTANCE-APPLICABILITY.


IF INSTANCE-OUTSIDE-MEASUREMENT-DOMAIN

RETURN MEASUREMENT-WITNESS-SET

END-IF.


PERFORM OBTAIN-SAMPLE.


IF SAMPLE-NOT-REPRESENTATIVE

RETURN MEASUREMENT-WITNESS-SET

END-IF.


PERFORM ACQUIRE-INSTRUMENT-RECORD.


IF INSTRUMENT-SATURATED

RETURN MEASUREMENT-WITNESS-SET

END-IF.


IF SIGNAL-BELOW-DETECTION-LIMIT

RETURN CENSORED-OBSERVATION

WITH DECLARED LIMIT

END-IF.


PERFORM APPLY-ANALYSIS-MAP.


IF ANALYSIS-PROCEDURE-INVALID

RETURN MEASUREMENT-WITNESS-SET

END-IF.


PERFORM CHECK-UNITS.

PERFORM CHECK-RESOLUTION.

PERFORM BUILD-UNCERTAINTY-DESCRIPTION.

PERFORM DISCLOSE-PREDICTION-DEPENDENCE.


IF ANY REQUIRED GUARD FAILED

RETURN MEASUREMENT-WITNESS-SET

ELSE

RETURN COMPLETE-OBSERVATION-RECORD

END-IF.

\PCsubsection{23.1 Uncertainty propagation procedure}

PROPAGATE-UNCERTAINTY.


FORM CANDIDATE-EXTERNAL-STATE-SET.

FORM CANDIDATE-PHASE-SET.


IF CANDIDATE-PHASE-SET CROSSES FORBIDDEN BOUNDARY

DO NOT USE SINGLE-BRANCH LINEARIZATION

END-IF.


IF STRUCTURAL-FUNCTION HAS UNCONTROLLED POLE

USE EXACT SET OR DISTRIBUTIONAL PROPAGATION

OR RETURN SINGULARITY-PROXIMITY-WITNESS

END-IF.


APPLY STRUCTURAL-IMAGE MAP.

APPLY CALIBRATION AND READOUT MAP.

RECORD CORRELATIONS AND ASSUMPTIONS.

RETURN PREDICTED-OUTPUT-SET

OR OUTPUT-DISTRIBUTION.


\PCsection{24. Theorem Ledger}

\PCsubsection{Theorem III.1 --- Valid-observation criterion}

\[
\mathcal W_{\mathcal M}(\omega)=\varnothing
\]

if and only if every declared measurement guard succeeds.

\PCsubsection{Theorem III.2 --- Error decomposition}

\[
Y-\mu=b+\eta.
\]

\PCsubsection{Theorem III.3 --- Precision does not imply accuracy}

Zero dispersion is compatible with nonzero bias.

\PCsubsection{Theorem III.4 --- Residual decomposition}

\[
\boxed{ y_{\mathrm{obs}} - \widehat y_{\mathcal P} = \left( y_{\mathrm{obs}}-\mu \right) + \left( \mu-\widehat y_{\mathcal P} \right). }
\]

\PCsubsection{Theorem III.5 --- Averaging does not remove bias}

\[
\mathbb E[\overline Y_n]=\mu+b.
\]

\PCsubsection{Theorem III.6 --- Independent averaging}

\[
\operatorname{Var}(\overline Y_n) = \frac{\sigma^2}{n}
\]

under independence and equal variance.

\PCsubsection{Theorem III.7 --- Correlated averaging}

\[
\operatorname{Var}(\overline Y_n) = \frac{\sigma^2}{n} \left( 1+(n-1)\rho \right)
\]

under equal pairwise correlation.

\PCsubsection{Theorem III.8 --- Exact set propagation}

\[
P_{\mathcal P} = \Gamma_\theta \left( E_{\mathbf F}(C_\Phi) \right).
\]

\PCsubsection{Theorem III.9 --- Lipschitz propagation bound}

\[
\operatorname{diam}(F(X)) \leq L\operatorname{diam}(X).
\]

\PCsubsection{Theorem III.10 --- Pushforward propagation}

\[
\nu_{\mathrm{out}} = (\Gamma_\theta\circ E_{\mathbf F})_{\#}\nu_\Phi.
\]

\PCsubsection{Theorem III.11 --- Flatwave quadrant stability}

If

\[
C_\Phi\subset Q_k,
\]

then

\[
\operatorname{Flatwave}(C_\Phi)=\{\varepsilon_k\}.
\]

\PCsubsection{Theorem III.12 --- Flatwave boundary ambiguity}

If candidate phases occupy adjacent quadrants, then

\[
\operatorname{Flatwave}(C_\Phi)=\{-1,+1\}.
\]

\PCsubsection{Theorem III.13 --- Primitive singular amplification}

Arbitrarily small phase uncertainty near a primitive pole may produce arbitrarily large primitive-value uncertainty.

\PCsubsection{Theorem III.14 --- Saw boundedness}

Saw-value uncertainty remains inside \([-1,1]\).

\PCsubsection{Theorem III.15 --- Record-level compatibility}

\[
P_{\mathcal P}(\omega) \cap O_{\mathcal O} \neq\varnothing
\]

is the general compatibility condition.


\PCsection{25. Negative Theorems}

\PCsubsection{Negative Theorem 25.1 --- No number without a measurand}

A numerical report without a declared measurand is not a complete measurement result.

\PCsubsection{Negative Theorem 25.2 --- No error without a reference value}

Exact measurement error cannot be calculated without a measurand value or accepted reference value.

\PCsubsection{Negative Theorem 25.3 --- No uncertainty–error identity}

An uncertainty interval is not the unknown realized error.

\PCsubsection{Negative Theorem 25.4 --- No precision–accuracy identity}

Small dispersion does not prove closeness to the measurand.

\PCsubsection{Negative Theorem 25.5 --- No residual attribution without further evidence}

A residual does not uniquely identify measurement error or model discrepancy.

\PCsubsection{Negative Theorem 25.6 --- No independent-error formula without dependence analysis}

Covariance may not be discarded without justification.

\PCsubsection{Negative Theorem 25.7 --- No linearization across a Flatwave wall}

A zero interior derivative does not describe the two-valued uncertainty created by a boundary-straddling phase set.

\PCsubsection{Negative Theorem 25.8 --- No linearization through an uncontrolled pole}

Derivative-based propagation is invalid when the required derivative bound does not exist on the uncertainty domain.

\PCsubsection{Negative Theorem 25.9 --- No failed measurement as model falsification}

A record with unresolved measurement witnesses cannot serve as a lawful contradiction.

\PCsubsection{Negative Theorem 25.10 --- No hidden prediction-dependent filtering}

Data selection influenced by model expectations must be disclosed.


\PCsection{26. Terminal Theorem}

\PCsubsection{Theorem 26.1 --- Measurement–Projection Separation Theorem}

A lawful empirical comparison requires two independently auditable constructions:

\[
\boxed{ \text{projection contract} \quad \mathcal P }
\]

and

\[
\boxed{ \text{measurement protocol} \quad \mathcal M. }
\]

The projection contract produces

\[
P_{\mathcal P}(\omega),
\]

the set or distribution of outputs permitted by the model.

The measurement protocol produces

\[
O_{\mathcal O},
\]

the set or distribution of measurand values permitted by the observation.

The comparison is meaningful only when:

\begin{enumerate}
\item both constructions pass their guards;
\item both refer to the same measurand;
\item both use compatible units and reference conventions;
\item uncertainty and dependence have been declared.
\end{enumerate}
A residual between central values does not, by itself, identify either measurement error or model discrepancy.

\begin{proof}
The projection result depends on phase inference, structural evaluation, calibration, and readout.

The observation result depends on sampling, instrument response, analysis, and measurement uncertainty.

These are logically distinct processes.

The residual decomposition proves that their difference contains contributions from both measurement error and model discrepancy.

Therefore, lawful comparison requires separate audit and compatible definitions before intersection, residual, or falsification statements are made.


\end{proof}

\PCsection{27. Closing Declaration}

S1.I established the projection contract.

S1.II established admissibility, witnesses, and falsification conditions.

S1.III establishes the observation protocol.

A measurement protocol is

\[
\boxed{ \mathcal M = \left( \Omega_{\mathcal M}, \mu, S, I, A, \mathcal U_{\mathcal M}, \mathcal Q_{\mathcal M}, W_{\mathcal M} \right). }
\]

Its realized report is

\[
\boxed{ y_{\mathrm{obs}}(\omega) = (A\circ I\circ S)(\omega) }.
\]


Measurement error is

\[
\boxed{ e_{\mathcal M} = y_{\mathrm{obs}}-\mu. }
\]

Model discrepancy is

\[
\boxed{ d_{\mathcal P} = \mu-\widehat y_{\mathcal P}. }
\]

The residual is their sum:

\[
\boxed{ r = e_{\mathcal M} + d_{\mathcal P}. }
\]

Uncertainty propagates by exact image sets or pushforward measures:

\[
\boxed{ P_{\mathcal P} = \Gamma_\theta \left( E_{\mathbf F}(C_\Phi) \right). }
\]

Inside one quadrant, continuous propagation is available.

Across a Flatwave wall,

\[
\boxed{ \operatorname{Flatwave}(C_\Phi) = \{-1,+1\}. }
\]

Near a primitive pole, arbitrarily small phase uncertainty may be amplified without finite bound.

The resulting measurement discipline is:

\[
\boxed{ \begin{gathered} \text{No observation without a measurand;}\\ \text{no error without a reference;}\\ \text{no uncertainty without a declared representation;}\\ \text{no precision claim without repetition conditions;}\\ \text{no residual attribution without separating}\\ \text{measurement error from model discrepancy;}\\ \text{no linear propagation across a discontinuity or pole.} \end{gathered} }
\]

\[
\boxed{ \text{END OF S1.III} }
\]

% ===== END S1_III =====
% ===== BEGIN S1_IV =====
% Source: S1.iv.docx
\PCchapter{S1.IV}{Calibration and Dimensional Readout}{Book 1 - Projection and Measurement}

\PCsubsection{Abstract}

Book 0 established a dimensionless structural calculus. S1.I defined the projection contract

\[
\widehat y_{\mathcal P} = \Gamma_\theta \circ E_{\mathbf F} \circ \Phi,
\]

S1.II established admissibility and failure witnesses, and S1.III separated measurement error, uncertainty, and model discrepancy.

This essay formalizes the readout and calibration layer.

Calibration is the process by which external reference information is used to select the parameters

\[
\theta\in\Theta
\]

of a readout map

\[
\Gamma_\theta: \mathcal S_{\mathcal P}\to\mathcal Y.
\]

Because the Book 0 structural variables are dimensionless, every dimensional scale, offset, unit, and reference value appearing in the projected output must enter through the readout, the calibration data, or another explicitly declared external quantity.

A calibration is represented by

\[
\widehat\theta = K(\mathcal C),
\]

where \(\mathcal C\) is the calibration record and \(K\) is the declared calibration procedure.

For a scalar affine readout,

\[
y=a+bs,
\]

two distinct exact structural anchors determine a unique pair \((a,b)\). If the structural anchors coincide, the calibration is either inconsistent or underdetermined.

For a multivariable linear readout,

\[
y=b+\boldsymbol\beta^{\mathsf T}\mathbf s,
\]

parameter identifiability requires full column rank of the design matrix. A numerically optimized parameter is not automatically identifiable; multiple parameter values may produce the same outputs.

Calibration uncertainty is propagated jointly with phase, measurement, and readout uncertainty:

\[
P_{\mathcal P}(\omega) = \left\{ \Gamma_\theta(\mathbf s): \theta\in C_\theta,\; \mathbf s\in C_{\mathbf F}(\omega) \right\}.
\]

The essay distinguishes fitting data from validation data. Reusing the same information to select parameters and to claim independent confirmation does not provide an independent test.

Calibration is valid only on its declared domain. Evaluation outside the calibrated structural or external region is extrapolation and requires an explicit guard. Parameter drift may invalidate a formerly lawful calibration without altering the Book 0 identities.

The resulting principle is:

\[
\boxed{ \text{Book 0 supplies form; calibration supplies scale.} }
\]


\PCsection{IDENTIFICATION DIVISION}

\PCsubsection{PROGRAM-ID}

S1-IV-CALIBRATION-AND-DIMENSIONAL-READOUT

\PCsubsection{PURPOSE}

To define:

\begin{itemize}
\item dimensional readout;
\item calibration anchors;
\item parameter estimation;
\item parameter identifiability;
\item calibration uncertainty;
\item unit transformation;
\item interpolation and extrapolation;
\item calibration drift;
\item and calibration failure witnesses.
\end{itemize}
\PCsubsection{PRINCIPAL QUESTION}

How may dimensionless structural quantities be converted into externally meaningful, dimensionally declared outputs without hiding scale choices or empirical assumptions?

\PCsubsection{NON-PURPOSE}

This essay does not:

\begin{itemize}
\item select a particular physical constant;
\item calibrate the R-operators to hydrogen, gravity, cosmology, or another specific system;
\item prove that a selected external anchor is fundamental;
\item infer units from dimensionless identities alone;
\item guarantee that a fitted model is correct;
\item establish model superiority;
\item replace independent validation with parameter fitting;
\item remove Book 0 boundary restrictions.
\end{itemize}

\PCsection{ENVIRONMENT DIVISION}

\PCsubsection{1. Inherited Projection Structure}

A lawful projection contract has the form

\[
\mathcal P = \left( \Omega, \Omega_{\mathcal P}, \Phi, \mathbf F, \Gamma_\theta, U_{\mathcal P}, W_{\mathcal P} \right).
\]

For an admitted instance \(\omega\),

\[
\mathbf s(\omega) = E_{\mathbf F}(\Phi(\omega))
\]

is its dimensionless structural state.

The predicted output is

\[
\boxed{ \widehat y_{\mathcal P}(\omega) = \Gamma_\theta(\mathbf s(\omega)). }
\]

\PCsubsection{1.1 Structural state space}

Let

\[
\mathcal S_{\mathcal P} \subseteq \mathbb R^m
\]

be the set of structural states admitted by the contract.

Every coordinate of \(\mathbf s\in\mathcal S_{\mathcal P}\) is dimensionless unless a later contract explicitly introduces a dimension-bearing external coordinate.

\PCsubsection{1.2 Measurement environment}

A valid calibration record must be built from observation records satisfying the measurement requirements of S1.III.

A failed, censored, ambiguous, or dimensionally incompatible observation may not be used as an exact calibration anchor without an explicit treatment rule.


\PCsection{METHODOLOGICAL AXIOM DIVISION}

\PCsubsection{Axiom C.1 --- Explicit Scale Provenance}

Every dimensional scale, offset, unit, conversion constant, and reference value used by a readout must have declared provenance.

The provenance must state whether the quantity is:

\begin{itemize}
\item fixed by unit convention;
\item measured externally;
\item selected as a reference;
\item estimated from calibration data;
\item fitted jointly with other parameters;
\item imported from another model;
\item or assumed.
\end{itemize}
No dimensional constant may be presented as though it were generated by the dimensionless structural calculus.

\PCsubsection{Axiom C.2 --- Calibration–Validation Separation}

Every datum used to select, constrain, or tune the calibration parameters must be identified as calibration information.

A claim of independent validation must use information not already exhausted in selecting those parameters, or must explicitly account for the dependence created by reuse.


\PCsection{DATA DIVISION}

\PCsection{2. Dimensions and Units}

\PCsubsection{Definition 2.1 --- Dimension set}

Let

\[
\mathfrak D
\]

be the set of declared quantity dimensions.

Examples may include:

\[
\mathsf L,\quad \mathsf T,\quad \mathsf M,\quad \mathsf Q,
\]

representing length, time, mass, charge, or other base dimensions.

The dimensionless element is denoted

\[
\mathbf 1.
\]

\PCsubsection{Definition 2.2 --- Dimension assignment}

For a quantity \(q\), write

\[
[q]\in\mathfrak D
\]

for its dimension.

For every Book 0 structural variable \(s\),

\[
\boxed{ [s]=\mathbf 1. }
\]

\PCsubsection{Definition 2.3 --- Unit representation}

A unit representation assigns a numerical coordinate to a dimensional quantity relative to a declared unit.

Changing units changes the numerical representation but not the underlying quantity.

\PCsubsection{Definition 2.4 --- Dimensional compatibility}

Two quantities may be added or compared numerically only when they have compatible dimensions and reference conventions.

Thus,

\[
q_1+q_2
\]

requires

\[
[q_1]=[q_2].
\]

\PCsubsection{Theorem 2.5 --- Dimensional isolation}

Let

\[
\mathbf s\in\mathcal S_{\mathcal P}
\]

be dimensionless.

If

\[
y=\Gamma_\theta(\mathbf s)
\]

has nontrivial dimension \([y]\neq\mathbf1\), then at least one component of \(\theta\), the readout definition, or an external input to the readout must carry that dimension.

\begin{proof}
The structural state consists only of dimensionless numbers.

Algebraic combinations of dimensionless numbers remain dimensionless unless combined with a dimension-bearing quantity or assigned a dimension through an external convention.

Therefore, any nontrivial output dimension must enter through the readout layer or another declared external element.

\end{proof}

\PCsubsection{Corollary 2.6 --- No unit from Flatwave alone}

Since

\[
\operatorname{Flatwave}\in\{-1,+1\},
\]

Flatwave does not by itself determine seconds, metres, joules, hertz, or another unit.

\PCsubsection{Corollary 2.7 --- No dimensional scale from \(\lambda\)}

Since

\[
0\leq\lambda\leq1
\]

on its closed-quadrant extension, \(\lambda\) provides a normalized coordinate but not a dimensional magnitude.


\PCsection{3. Readout Maps}

\PCsubsection{Definition 3.1 --- Parametric readout}

A parametric readout is a family of maps

\[
\boxed{ \Gamma: \Theta\times\mathcal S_{\mathcal P}\to\mathcal Y }
\]

written

\[
\Gamma_\theta(\mathbf s) = \Gamma(\theta,\mathbf s).
\]

\PCsubsection{Definition 3.2 --- Readout parameter}

The parameter

\[
\theta\in\Theta
\]

contains every free or externally supplied quantity required to convert structural values into the output.

\PCsubsection{Definition 3.3 --- Fixed parameter}

A parameter is fixed when its value is specified before the calibration data under consideration are analyzed.

\PCsubsection{Definition 3.4 --- Estimated parameter}

A parameter is estimated when calibration information is used to determine or constrain its value.

\PCsubsection{Definition 3.5 --- Fitted parameter}

A parameter is fitted when it is selected by optimizing a declared objective function over calibration data.

Every fitted parameter is estimated, but not every estimate must arise from optimization.

\PCsubsection{Definition 3.6 --- Readout domain}

The readout domain is the set

\[
\mathcal S_\theta \subseteq \mathcal S_{\mathcal P}
\]

on which \(\Gamma_\theta\) is declared valid.

It may be smaller than the complete structural state space.


\PCsection{4. Calibration Records}

\PCsubsection{Definition 4.1 --- Calibration instance set}

Let

\[
\Omega_{\mathcal C} \subseteq \Omega_{\mathcal P}\cap\Omega_{\mathcal M}
\]

be the external instances used for calibration.

\PCsubsection{Definition 4.2 --- Structural calibration state}

For each calibration instance \(\omega_i\),

\[
\boxed{ \mathbf s_i = E_{\mathbf F}(\Phi(\omega_i)). }
\]

\PCsubsection{Definition 4.3 --- Calibration observation}

Let

\[
\mathcal O_i
\]

be a valid observation record for the same measurand the readout is intended to predict.

Its admissible observed-value set is

\[
O_i\subseteq\mathcal Y.
\]

\PCsubsection{Definition 4.4 --- Exact calibration anchor}

An exact calibration anchor is a pair

\[
\boxed{ (\mathbf s_i,y_i) }
\]

for which the structural state and output value are both treated as exact under the declared calibration procedure.

Exact anchors are idealizations unless their uncertainty is demonstrably negligible for the intended purpose.

\PCsubsection{Definition 4.5 --- Uncertain calibration anchor}

An uncertain anchor is a pair

\[
\boxed{ (S_i,O_i), }
\]

where

\[
S_i\subseteq\mathcal S_{\mathcal P}
\]

and

\[
O_i\subseteq\mathcal Y
\]

represent admissible structural and observed values.

\PCsubsection{Definition 4.6 --- Calibration record}

The calibration record is

\[
\boxed{ \mathcal C = \left\{ (\omega_i,S_i,O_i,\mathcal O_i) \right\}_{i=1}^{n}. }
\]

\PCsubsection{Definition 4.7 --- Calibration procedure}

A calibration procedure is a map or algorithm

\[
\boxed{ K: \mathfrak C\to\mathfrak T, }
\]

where \(\mathfrak C\) is the class of admissible calibration records and \(\mathfrak T\) is a class of:

\begin{itemize}
\item parameter values;
\item parameter sets;
\item or parameter distributions.
\end{itemize}
\PCsubsection{Definition 4.8 --- Calibrated parameter result}

The calibrated result is

\[
\boxed{ \widehat\Theta = K(\mathcal C). }
\]

It may be:

\begin{itemize}
\item a single estimate \(\widehat\theta\);
\item a confidence region;
\item a credible distribution;
\item an admissible parameter set;
\item or a nonidentifiability report.
\end{itemize}

\PCsection{5. Calibration Lawfulness}

\PCsubsection{Definition 5.1 --- Calibration completeness}

A calibration is complete when it declares:

\begin{enumerate}
\item the readout family \(\Gamma_\theta\);
\item the parameter space \(\Theta\);
\item the calibration instances;
\item the structural states or structural uncertainty sets;
\item the observation records;
\item the objective, constraint, or anchor equations;
\item the treatment of uncertainty and dependence;
\item the calibration domain;
\item the parameter result;
\item failure behavior.
\end{enumerate}
\PCsubsection{Definition 5.2 --- Calibration admissibility}

A calibration record is admissible when:

\begin{itemize}
\item all observation records are valid;
\item all structural evaluations are lawful;
\item units and measurands are compatible;
\item prediction-dependent selection is disclosed;
\item the calibration procedure is defined for the record;
\item parameter constraints are declared.
\end{itemize}
\PCsubsection{Definition 5.3 --- Calibration witness set}

Recommended witnesses include:

\begin{itemize}
\item CALIBRATION-SPECIFICATION-INCOMPLETE;
\item CALIBRATION-DATA-INVALID;
\item CALIBRATION-MEASURAND-MISMATCH;
\item CALIBRATION-UNIT-MISMATCH;
\item CALIBRATION-ANCHORS-DEGENERATE;
\item PARAMETERS-NOT-IDENTIFIABLE;
\item CALIBRATION-OBJECTIVE-UNDEFINED;
\item CALIBRATION-UNCERTAINTY-MISSING;
\item CALIBRATION-DOMAIN-UNDECLARED;
\item CALIBRATION-VALIDATION-LEAKAGE;
\item CALIBRATION-DRIFT-DETECTED;
\item EXTRAPOLATION-UNDECLARED.
\end{itemize}

\PCsection{6. Scalar Affine Calibration}

\PCsubsection{Definition 6.1 --- Scalar affine readout}

Let

\[
s\in\mathbb R
\]

be dimensionless.

Define

\[
\boxed{ \Gamma_{a,b}(s)=a+bs. }
\]

For dimensional consistency,

\[
[a]=[b]=[y].
\]

\PCsubsection{Theorem 6.2 --- Unique two-anchor affine calibration}

Let

\[
(s_1,y_1)
\]

and

\[
(s_2,y_2)
\]

be exact anchors with

\[
s_1\neq s_2.
\]

Then there exists exactly one affine readout

\[
y=a+bs
\]

passing through both anchors.

Its parameters are

\[
\boxed{ b = \frac{y_2-y_1}{s_2-s_1} }
\]

and

\[
\boxed{ a = \frac{s_2y_1-s_1y_2}{s_2-s_1}. }
\]

\begin{proof}
The anchor equations are

\[
y_1=a+bs_1
\]

and

\[
y_2=a+bs_2.
\]

Subtracting,

\[
y_2-y_1=b(s_2-s_1).
\]

Since

\[
s_2-s_1\neq0,
\]

the slope is uniquely determined:

\[
b=\frac{y_2-y_1}{s_2-s_1}.
\]

Substituting into either anchor equation uniquely determines \(a\).

\end{proof}

\PCsubsection{Corollary 6.3 --- Offset-free calibration}

If the readout is constrained to pass through zero,

\[
y=bs,
\]

then one nonzero exact anchor determines

\[
\boxed{ b=\frac ys. }
\]

\PCsubsection{Theorem 6.4 --- Degenerate equal-state anchors}

Suppose

\[
s_1=s_2.
\]

Then:

\begin{enumerate}
\item if \(y_1\neq y_2\), no single-valued affine readout can satisfy both anchors;
\item if \(y_1=y_2\), infinitely many affine readouts satisfy both anchors.
\end{enumerate}
\begin{proof}
Both anchor equations have the same structural input.

A function must assign one output to that input.

If the outputs differ, the requirements are inconsistent.

If the outputs agree, the single equation

\[
y_1=a+bs_1
\]

leaves one free parameter.

\end{proof}

\PCsubsection{Corollary 6.5 --- Distinct outputs require structurally distinguishable anchors}

Calibration cannot assign different deterministic outputs to two calibration instances that the selected structural profile and readout treat as the same state.


\PCsection{7. Normalized Transfer Readout}

\PCsubsection{Definition 7.1 --- Two-endpoint transfer readout}

For a declared lower anchor \(y_0\) and upper anchor \(y_1\), define

\[
\boxed{ \Gamma_{y_0,y_1}(\lambda) = y_0+(y_1-y_0)\lambda. }
\]

\PCsubsection{Theorem 7.2 --- Endpoint correspondence}

On the closed-quadrant extension,

\[
\lambda=0 \implies \Gamma(\lambda)=y_0,
\]

and

\[
\lambda=1 \implies \Gamma(\lambda)=y_1.
\]

\PCsubsection{Corollary 7.3 --- Interior boundedness}

If

\[
y_0<y_1
\]

and

\[
0<\lambda<1,
\]

then

\[
\boxed{ y_0<\Gamma(\lambda)<y_1. }
\]

\PCsubsection{Clarification 7.4 --- Endpoint anchors are external choices}

The structural calculus supplies the normalized endpoints \(0\) and \(1\).

It does not supply the dimensional values \(y_0\) and \(y_1\).

\PCsubsection{Clarification 7.5 --- Boundary-aware requirement}

If \(y_0\) and \(y_1\) are interpreted as values at the exact quadrant boundaries, the contract must use the split-boundary extensions of S0.V.

The primitive interior formulas may not be evaluated directly at those boundaries.


\PCsection{8. Polynomial and Nonlinear Readouts}

\PCsubsection{Definition 8.1 --- Polynomial readout}

For a dimensionless scalar \(s\),

\[
\boxed{ y = \sum_{k=0}^{d}a_ks^k. }
\]

\PCsubsection{Theorem 8.2 --- Polynomial coefficient dimensions}

Every nonzero coefficient must satisfy

\[
\boxed{ [a_k]=[y]. }
\]

\begin{proof}
Since \(s\) is dimensionless,

\[
[s^k]=\mathbf1.
\]

Therefore,

\[
[a_ks^k]=[a_k].
\]

All terms being added must have the same dimension as \(y\).

\end{proof}

\PCsubsection{Definition 8.3 --- Nonlinear readout}

A nonlinear readout is any declared map

\[
\Gamma_\theta(\mathbf s)
\]

that is not affine in \(\mathbf s\).

Examples include:

\begin{itemize}
\item rational maps;
\item exponentials;
\item thresholds;
\item logistic functions;
\item piecewise maps;
\item saturation functions.
\end{itemize}
The use of such a readout is a projection choice. It is not a new theorem of Book 0.

\PCsubsection{Negative Theorem 8.4 --- A flexible readout can conceal structural weakness}

If the readout class is sufficiently flexible, it may reproduce many external datasets independently of the distinctive structure of \(\mathbf F\).

Therefore, goodness of fit must be interpreted relative to:

\begin{itemize}
\item parameter count;
\item readout flexibility;
\item competing models;
\item and out-of-sample performance.
\end{itemize}

\PCsection{9. Linear Multivariable Calibration}

\PCsubsection{Definition 9.1 --- Linear structural readout}

Let

\[
\mathbf s\in\mathbb R^m.
\]

Define the scalar readout

\[
\boxed{ y = b+\boldsymbol\beta^{\mathsf T}\mathbf s. }
\]

The parameter vector is

\[
\theta = (b,\beta_1,\ldots,\beta_m)^{\mathsf T}.
\]

\PCsubsection{Definition 9.2 --- Design matrix}

For calibration states \(\mathbf s_1,\ldots,\mathbf s_n\), define

\[
\boxed{ X = \begin{pmatrix} 1 & \mathbf s_1^{\mathsf T}\\ \vdots & \vdots\\ 1 & \mathbf s_n^{\mathsf T} \end{pmatrix}. }
\]

The observed-output vector is

\[
\mathbf y = (y_1,\ldots,y_n)^{\mathsf T}.
\]

The calibration equations are

\[
X\theta=\mathbf y.
\]

\PCsubsection{Theorem 9.3 --- Exact linear identifiability}

The parameter vector \(\theta\) is uniquely determined by exact calibration equations only if

\[
\boxed{ \operatorname{rank}(X)=m+1. }
\]

When the system is consistent and \(X\) has full column rank, the solution is unique.

\begin{proof}
If \(X\) does not have full column rank, there exists a nonzero vector \(v\) such that

\[
Xv=0.
\]

If \(\theta\) is one solution, then

\[
X(\theta+v) = X\theta+Xv = \mathbf y.
\]

Thus the solution is not unique.

If \(X\) has full column rank, its null space is trivial. Hence no two distinct parameter vectors can produce the same exact calibration outputs.

\end{proof}

\PCsubsection{Corollary 9.4 --- Duplicate structural information does not add rank}

Repeating calibration instances with identical structural states may improve precision under a measurement model, but it does not increase the algebraic rank of \(X\).

\PCsubsection{Corollary 9.5 --- Collinear structural coordinates cause parameter degeneracy}

If one selected structural coordinate is a linear combination of the others throughout the calibration record, its coefficient cannot be independently identified by the linear readout.


\PCsection{10. Weighted Linear Calibration}

\PCsubsection{Definition 10.1 --- Weighted quadratic objective}

Let \(W\) be a symmetric positive-semidefinite matrix.

Define

\[
\boxed{ J(\theta) = (\mathbf y-X\theta)^{\mathsf T} W (\mathbf y-X\theta). }
\]

\PCsubsection{Theorem 10.2 --- Unique weighted least-squares estimate}

If

\[
X^{\mathsf T}WX
\]

is positive definite, then \(J\) has the unique minimizer

\[
\boxed{ \widehat\theta = (X^{\mathsf T}WX)^{-1} X^{\mathsf T}W\mathbf y. }
\]

\begin{proof}
Differentiating,

\[
\nabla J(\theta) = -2X^{\mathsf T}W(\mathbf y-X\theta).
\]

Setting the gradient to zero gives

\[
X^{\mathsf T}WX\theta = X^{\mathsf T}W\mathbf y.
\]

Positive definiteness makes \(X^{\mathsf T}WX\) invertible and the quadratic objective strictly convex.

Therefore, the stationary point is unique and is the global minimizer.

\end{proof}

\PCsubsection{Clarification 10.3 --- Weight interpretation must be declared}

Weights may represent:

\begin{itemize}
\item inverse variances;
\item relative importance;
\item sampling corrections;
\item robust-loss approximations;
\item or another convention.
\end{itemize}
Different weight interpretations produce different calibrations.

\PCsubsection{Negative Theorem 10.4 --- Numerical uniqueness is not physical necessity}

A unique minimizer of a chosen objective proves only that the optimization problem has one solution.

It does not prove that:

\begin{itemize}
\item the readout family is correct;
\item the weights are correct;
\item the structural profile is externally relevant;
\item or the parameter has fundamental physical status.
\end{itemize}

\PCsection{11. Nonlinear Parameter Identifiability}

\PCsubsection{Definition 11.1 --- Calibration prediction vector}

For calibration instances \(\omega_1,\ldots,\omega_n\), define

\[
\boxed{ \mathbf g(\theta) = \begin{pmatrix} \Gamma_\theta(\mathbf s_1)\\ \vdots\\ \Gamma_\theta(\mathbf s_n) \end{pmatrix}. }
\]

\PCsubsection{Definition 11.2 --- Global identifiability}

The parameter is globally identifiable on \(\Theta\) when

\[
\boxed{ \mathbf g(\theta_1)=\mathbf g(\theta_2) \implies \theta_1=\theta_2. }
\]

\PCsubsection{Definition 11.3 --- Observationally equivalent parameters}

Parameters \(\theta_1\) and \(\theta_2\) are observationally equivalent on the calibration record when

\[
\mathbf g(\theta_1)=\mathbf g(\theta_2).
\]

\PCsubsection{Definition 11.4 --- Sensitivity matrix}

When differentiable, define

\[
\boxed{ J_\theta = \frac{\partial\mathbf g}{\partial\theta}. }
\]

\PCsubsection{Theorem 11.5 --- Rank-deficiency obstruction}

If \(J_\theta\) has a nontrivial null direction at \(\theta^\ast\), then there exists a first-order parameter perturbation that produces no first-order change in the calibration predictions.

\begin{proof}
If

\[
J_\theta v=0
\]

for nonzero \(v\), then

\[
\mathbf g(\theta^\ast+tv) = \mathbf g(\theta^\ast) + tJ_\theta v + o(t).
\]

Thus,

\[
\mathbf g(\theta^\ast+tv) = \mathbf g(\theta^\ast)+o(t).
\]

The perturbation is invisible to first order.

\end{proof}

\PCsubsection{Corollary 11.6 --- Full rank is a local diagnostic}

Full column rank supports local parameter distinguishability.

It does not, by itself, prove global uniqueness across the complete parameter space.

\PCsubsection{Negative Theorem 11.7 --- Parameter estimation does not imply parameter identification}

An optimization routine may return one parameter value even when many values produce indistinguishable predictions.


\PCsection{12. Calibration Uncertainty}

\PCsubsection{Definition 12.1 --- Parameter admissibility set}

Let

\[
\boxed{ C_\theta \subseteq \Theta }
\]

be the set of parameter values compatible with the calibration record and the declared calibration procedure.

\PCsubsection{Definition 12.2 --- Parameter distribution}

Under a probabilistic calibration, let

\[
\nu_\theta
\]

be the declared distribution over parameter values.

\PCsubsection{Definition 12.3 --- Joint structural–calibration output set}

Let

\[
C_{\mathbf F}(\omega) \subseteq \mathcal S_{\mathcal P}
\]

be the structural-state uncertainty set.

Define

\[
\boxed{ P_{\mathcal P}(\omega) = \left\{ \Gamma_\theta(\mathbf s): \theta\in C_\theta,\; \mathbf s\in C_{\mathbf F}(\omega) \right\}. }
\]

\PCsubsection{Theorem 12.4 --- Exact joint set propagation}

The complete predicted-output set is the image of

\[
C_\theta\times C_{\mathbf F}(\omega)
\]

under the map

\[
(\theta,\mathbf s) \mapsto \Gamma_\theta(\mathbf s).
\]

\begin{proof}
The displayed set contains exactly the outputs generated by every admissible parameter and structural-state pair.

This is the definition of the image of the Cartesian product under the readout map.

\end{proof}

\PCsubsection{Definition 12.5 --- Joint probabilistic propagation}

Let

\[
\nu_{\theta,\mathbf s}
\]

be the declared joint distribution of calibration parameters and structural state.

Then

\[
\boxed{ \nu_y = \Gamma_{\#}\nu_{\theta,\mathbf s}. }
\]

\PCsubsection{Clarification 12.6 --- Independence may not be assumed silently}

If the same observations influence both phase inference and calibration, then \(\theta\) and \(\mathbf s\) may be statistically dependent.

Replacing the joint distribution by

\[
\nu_\theta\otimes\nu_{\mathbf s}
\]

requires an independence justification.


\PCsection{13. Calibration Contribution to Residuals}

\PCsubsection{Definition 13.1 --- Reference calibration}

Suppose \(\theta^\ast\) denotes the parameter value that would produce the correct readout within the selected readout family.

This is a conceptual reference and need not be known.

\PCsubsection{Definition 13.2 --- Calibration discrepancy}

For structural state \(\mathbf s\), define

\[
\boxed{ c_\theta(\mathbf s) = \Gamma_{\theta^\ast}(\mathbf s) - \Gamma_{\widehat\theta}(\mathbf s). }
\]

\PCsubsection{Definition 13.3 --- Structural-model discrepancy}

Define

\[
\boxed{ d_{\mathrm{str}} = \mu - \Gamma_{\theta^\ast}(\mathbf s). }
\]

\PCsubsection{Theorem 13.4 --- Three-part residual decomposition}

\[
\boxed{ y_{\mathrm{obs}} - \Gamma_{\widehat\theta}(\mathbf s) = e_{\mathcal M} + d_{\mathrm{str}} + c_\theta. }
\]

\begin{proof}
Add and subtract \(\mu\) and \(\Gamma_{\theta^\ast}(\mathbf s)\):

\[
y_{\mathrm{obs}} - \Gamma_{\widehat\theta}(\mathbf s) = \left( y_{\mathrm{obs}}-\mu \right) + \left( \mu-\Gamma_{\theta^\ast}(\mathbf s) \right) + \left( \Gamma_{\theta^\ast}(\mathbf s) - \Gamma_{\widehat\theta}(\mathbf s) \right).
\]

The three terms are measurement error, structural-model discrepancy, and calibration discrepancy.

\end{proof}

\PCsubsection{Corollary 13.5 --- Calibration error may mimic structural failure}

A nonzero residual does not by itself distinguish incorrect structural modeling from incorrect calibration.


\PCsection{14. Calibration Domain}

\PCsubsection{Definition 14.1 --- External calibration domain}

Let

\[
\boxed{ \Omega_\theta \subseteq \Omega_{\mathcal P} }
\]

be the external region on which the calibration is declared valid.

\PCsubsection{Definition 14.2 --- Structural calibration domain}

Let

\[
\boxed{ \mathcal S_\theta \subseteq \mathcal S_{\mathcal P} }
\]

be the corresponding region of structural states covered by the calibration declaration.

\PCsubsection{Definition 14.3 --- Calibration range}

For a scalar structural coordinate, the calibration range is the interval or union of intervals represented by the calibration states.

For a multivariable state, it is the declared region containing the calibration states.

\PCsubsection{Definition 14.4 --- Interpolation}

In a Euclidean structural state space, a prediction is interpolative when its structural state lies inside the declared interpolation region, often the convex hull of the calibration states.

\PCsubsection{Definition 14.5 --- Extrapolation}

A prediction is extrapolative when its structural state lies outside the declared calibration region.

\PCsubsection{Theorem 14.6 --- Calibration lawfulness does not automatically extend}

A calibration validated on \(\mathcal S_\theta\) is not thereby validated on

\[
\mathcal S_{\mathcal P}\setminus\mathcal S_\theta.
\]

\begin{proof}
The calibration record constrains the readout only on the represented calibration region.

Different readout functions may agree on that region and diverge outside it.

Therefore, agreement within the calibration domain does not determine or validate behavior outside it.

\end{proof}

\PCsubsection{Corollary 14.7 --- Formula evaluation is not calibration validation}

The numerical ability to evaluate \(\Gamma_{\widehat\theta}(\mathbf s)\) outside \(\mathcal S_\theta\) does not establish that the output is reliable there.

\PCsubsection{Required witness}

Undeclared extrapolation must return

\[
\boxed{ \texttt{EXTRAPOLATION-UNDECLARED}. }
\]


\PCsection{15. Singular and Boundary Calibration Guards}

\PCsubsection{Theorem 15.1 --- Primitive-pole calibration instability}

Suppose a readout uses a primitive operator that diverges at a boundary.

Calibration based on states approaching that boundary may be arbitrarily sensitive to small phase error.

\begin{proof}
S1.III proved that primitive values are unbounded on every neighborhood of the relevant pole.

Therefore, arbitrarily small phase perturbations may produce arbitrarily large changes in the structural input to the readout.

Unless the readout suppresses that divergence in a controlled and declared way, calibration sensitivity is unbounded.

\end{proof}

\PCsubsection{Corollary 15.2 --- Minimum boundary distance}

A primitive-based calibration should declare a minimum admissible distance from excluded boundaries or use an exact singular uncertainty treatment.

\PCsubsection{Theorem 15.3 --- Flatwave calibration cannot resolve interior phase}

Within a fixed open quadrant,

\[
\operatorname{Flatwave}(x)=\varepsilon
\]

is constant.

Therefore, a readout using Flatwave alone cannot distinguish two phases within the same quadrant.

\begin{proof}
Equal structural inputs produce equal deterministic readout outputs.

Since every phase in the quadrant has the same Flatwave value, no Flatwave-only readout can distinguish them.

\end{proof}

\PCsubsection{Corollary 15.4 --- Flatwave supplies a sign channel, not a continuous calibration coordinate}

A continuously varying external quantity cannot be calibrated uniquely from Flatwave alone within one quadrant.

\PCsubsection{Theorem 15.5 --- Saw-based calibration remains bounded}

A readout based continuously on the saw coordinates receives structural inputs bounded in \([-1,1]\).

This avoids primitive divergence but does not remove seam discontinuity or branch ambiguity.


\PCsection{16. Unit Transformations}

\PCsubsection{Definition 16.1 --- Affine unit conversion}

Let two numerical unit systems be related by

\[
\boxed{ u(y)=cy+d, }
\]

with \(c\neq0\).

Pure scale conversions have \(d=0\).

Offset units may have \(d\neq0\).

\PCsubsection{Definition 16.2 --- Converted readout}

Define

\[
\boxed{ \Gamma_\theta^{(u)} = u\circ\Gamma_\theta. }
\]

\PCsubsection{Theorem 16.3 --- Unit-covariant prediction}

If

\[
\widehat y = \Gamma_\theta(\mathbf s),
\]

then the numerical prediction in the converted unit is

\[
\boxed{ \widehat y^{(u)} = u(\widehat y). }
\]

\begin{proof}
By definition,

\[
\widehat y^{(u)} = \Gamma_\theta^{(u)}(\mathbf s) = u(\Gamma_\theta(\mathbf s)) = u(\widehat y).
\]

\end{proof}

\PCsubsection{Corollary 16.4 --- Affine parameter transformation}

For

\[
\Gamma_{a,b}(s)=a+bs,
\]

the converted readout is

\[
u(\Gamma_{a,b}(s)) = (ca+d)+(cb)s.
\]

Thus,

\[
\boxed{ a^{(u)}=ca+d, \qquad b^{(u)}=cb. }
\]

\PCsubsection{Clarification 16.5 --- Numerical parameter values depend on units}

A dimensional calibration parameter may change numerically under unit conversion while representing the same underlying quantity.

\PCsubsection{Negative Theorem 16.6 --- Unit-dependent numerical equality is not physical identity}

Two parameters with the same numerical value in different unit systems need not represent the same dimensional quantity.


\PCsection{17. Calibration Drift}

\PCsubsection{Definition 17.1 --- Context-dependent parameter}

Let \(c\) denote external conditions such as:

\begin{itemize}
\item time;
\item temperature;
\item instrument state;
\item population;
\item environment;
\item manufacturing batch;
\item reference frame.
\end{itemize}
A drifting parameter has form

\[
\theta=\theta(c).
\]

\PCsubsection{Definition 17.2 --- Fixed-calibration assumption}

A contract using one constant \(\widehat\theta\) assumes

\[
\theta(c)=\widehat\theta
\]

throughout its declared calibration domain.

\PCsubsection{Definition 17.3 --- Calibration drift}

Calibration drift occurs when the fixed-calibration assumption fails materially.

\PCsubsection{Theorem 17.4 --- Drift invalidates fixed readout without altering structural identities}

If the true readout parameter changes from \(\theta_1\) to \(\theta_2\) while the contract continues to use \(\theta_1\), projected outputs may become inaccurate even though every Book 0 structural identity remains exact.

\begin{proof}
The structural state \(\mathbf s\) and its identities do not depend on the external calibration parameter.

The output does:

\[
\widehat y=\Gamma_{\theta_1}(\mathbf s).
\]

If the appropriate readout is instead

\[
\Gamma_{\theta_2}(\mathbf s),
\]

the output discrepancy arises in the readout layer.

The Book 0 identities are unchanged.

\end{proof}

\PCsubsection{Definition 17.5 --- Drift detection rule}

A calibration contract should declare:

\begin{itemize}
\item monitoring variables;
\item acceptable parameter variation;
\item recalibration threshold;
\item and failure action.
\end{itemize}
\PCsubsection{Required witness}

Detected material drift should return

\[
\boxed{ \texttt{CALIBRATION-DRIFT-DETECTED}. }
\]


\PCsection{18. Calibration Transfer}

\PCsubsection{Definition 18.1 --- Source calibration domain}

Let a calibration be established on domain

\[
\Omega_A.
\]

\PCsubsection{Definition 18.2 --- Target domain}

Let

\[
\Omega_B
\]

be a distinct external domain to which the same parameters are proposed to apply.

\PCsubsection{Definition 18.3 --- Calibration transfer assumption}

Calibration transfer assumes that the same readout parameters remain valid on \(\Omega_B\).

\PCsubsection{Theorem 18.4 --- Transfer is an additional hypothesis}

Calibration validity on \(\Omega_A\) does not imply validity on \(\Omega_B\).

\begin{proof}
The calibration record contains evidence from \(\Omega_A\), not necessarily from \(\Omega_B\).

Conditions affecting the readout may differ between the domains.

Therefore, parameter invariance across the two domains is an additional claim requiring justification or testing.

\end{proof}

\PCsubsection{Corollary 18.5 --- Imported constants require provenance}

A parameter imported from another experiment or field must state:

\begin{itemize}
\item its source;
\item original calibration domain;
\item uncertainty;
\item unit convention;
\item and transfer justification.
\end{itemize}

\PCsection{19. Calibration and Validation Data}

\PCsubsection{Definition 19.1 --- Calibration set}

Let

\[
\mathcal C_{\mathrm{fit}}
\]

contain the records used to select \(\widehat\theta\).

\PCsubsection{Definition 19.2 --- Validation set}

Let

\[
\mathcal C_{\mathrm{val}}
\]

contain records used to evaluate the calibrated projection after parameter selection.

\PCsubsection{Definition 19.3 --- Independent validation}

Validation is independent when its information did not influence:

\begin{itemize}
\item parameter choice;
\item model form;
\item structural-profile selection;
\item preprocessing choices;
\item stopping rules;
\item or reporting thresholds.
\end{itemize}
\PCsubsection{Theorem 19.4 --- Training fit is not independent validation}

Agreement with

\[
\mathcal C_{\mathrm{fit}}
\]

does not constitute an independent test of the calibration selected from that same information.

\begin{proof}
The parameter was chosen, at least in part, to improve agreement with the fitting data.

The observed agreement is therefore statistically and procedurally dependent on the parameter-selection process.

It may demonstrate optimization success, but not independent predictive performance.

\end{proof}

\PCsubsection{Corollary 19.5 --- Data reuse must be disclosed}

If the same data are used for both calibration and evaluation, the dependence must be represented in the uncertainty or validation procedure.

\PCsubsection{Definition 19.6 --- Validation leakage}

Validation leakage occurs when information from the nominal validation set influences model or parameter selection without being accounted for.

\PCsubsection{Required witness}

\[
\boxed{ \texttt{CALIBRATION-VALIDATION-LEAKAGE}. }
\]


\PCsection{20. Overfitting and Readout Complexity}

\PCsubsection{Definition 20.1 --- Calibration fit error}

Let

\[
L_{\mathrm{fit}}(\widehat\theta)
\]

be the declared loss on the calibration data.

\PCsubsection{Definition 20.2 --- Validation error}

Let

\[
L_{\mathrm{val}}(\widehat\theta)
\]

be the corresponding loss on independent validation data.

\PCsubsection{Definition 20.3 --- Overfitting}

Overfitting occurs when a readout captures features specific to the calibration record that do not generalize to the declared target domain.

\PCsubsection{Negative Theorem 20.4 --- Perfect fit does not prove correct projection}

A sufficiently flexible readout may interpolate a finite calibration dataset exactly while making poor predictions elsewhere.

\begin{proof}
For finitely many distinct scalar structural states, a polynomial of sufficiently high degree can interpolate arbitrary output values.

Exact agreement on those points therefore does not uniquely determine behavior between or beyond them.

\end{proof}

\PCsubsection{Corollary 20.5 --- Parameter count must be reported}

The number of fitted parameters and the flexibility of the readout class are part of the calibration record.

\PCsubsection{Corollary 20.6 --- Simpler readouts require fewer calibration constraints}

A low-parameter readout is generally easier to identify and audit, though simplicity alone does not establish correctness.


\PCsection{21. Calibration Failure Locality}

\PCsubsection{Definition 21.1 --- Calibration-field dependency}

The readout output depends on:

\[
\Gamma,\quad \theta,\quad \mathcal C,\quad K,\quad \mathcal S_\theta.
\]

\PCsubsection{Theorem 21.2 --- Calibration repair locality}

Changing calibration parameters requires recalculating projected outputs and their calibration uncertainty.

It does not require reproving Book 0 identities.

\begin{proof}
Book 0 identities depend on the structural definitions and phase domain.

Calibration parameters appear only after structural evaluation.

Therefore, changing them affects the readout layer but not the structural theorem layer.

\end{proof}

\PCsubsection{Corollary 21.3 --- Readout-form revision requires broader reevaluation}

Changing the readout family \(\Gamma\) may alter:

\begin{itemize}
\item parameter identifiability;
\item dimensional consistency;
\item calibration domain;
\item extrapolation behavior;
\item uncertainty propagation;
\item and validation results.
\end{itemize}
All dependent claims must be rechecked.


\PCsection{PROCEDURE DIVISION}

\PCsection{22. Canonical Calibration Procedure}

CALIBRATE-READOUT.


PERFORM CHECK-CALIBRATION-SPECIFICATION.


IF CALIBRATION-SPECIFICATION-INCOMPLETE

RETURN CALIBRATION-WITNESS-SET

END-IF.


PERFORM VALIDATE-CALIBRATION-OBSERVATIONS.


IF CALIBRATION-DATA-INVALID

RETURN CALIBRATION-WITNESS-SET

END-IF.


PERFORM VERIFY-MEASURAND-AND-UNIT-COMPATIBILITY.


IF CALIBRATION-MEASURAND-MISMATCH

RETURN CALIBRATION-WITNESS-SET

END-IF.


IF CALIBRATION-UNIT-MISMATCH

RETURN CALIBRATION-WITNESS-SET

END-IF.


PERFORM EVALUATE-CALIBRATION-STRUCTURAL-STATES.


IF STRUCTURAL-EVALUATION-UNDEFINED

RETURN CALIBRATION-WITNESS-SET

END-IF.


PERFORM CHECK-ANCHOR-DEGENERACY.

PERFORM CHECK-PARAMETER-IDENTIFIABILITY.


IF CALIBRATION-ANCHORS-DEGENERATE

RETURN CALIBRATION-WITNESS-SET

END-IF.


IF PARAMETERS-NOT-IDENTIFIABLE

RETURN PARAMETER-SET

OR IDENTIFIABILITY-WITNESS

END-IF.


PERFORM APPLY-CALIBRATION-PROCEDURE.

PERFORM ESTIMATE-CALIBRATION-UNCERTAINTY.

PERFORM DECLARE-CALIBRATION-DOMAIN.

PERFORM CHECK-CALIBRATION-VALIDATION-SEPARATION.


IF CALIBRATION-UNCERTAINTY-MISSING

RETURN CALIBRATION-WITNESS-SET

END-IF.


IF CALIBRATION-DOMAIN-UNDECLARED

RETURN CALIBRATION-WITNESS-SET

END-IF.


RETURN CALIBRATED-PARAMETER-RESULT

CALIBRATION-UNCERTAINTY

CALIBRATION-DOMAIN

COMPLETE-AUDIT-RECORD.

\PCsubsection{22.1 Readout procedure}

APPLY-DIMENSIONAL-READOUT.


PERFORM CHECK-INSTANCE-STRUCTURAL-STATE.

PERFORM CHECK-CALIBRATION-STATUS.

PERFORM CHECK-CALIBRATION-DOMAIN.

PERFORM CHECK-UNIT-CONVENTION.

PERFORM CHECK-DRIFT-STATUS.


IF STRUCTURAL-STATE OUTSIDE CALIBRATION-DOMAIN

IF EXTRAPOLATION-RULE IS NOT DECLARED

RETURN EXTRAPOLATION-UNDECLARED

END-IF

END-IF.


IF CALIBRATION-DRIFT-DETECTED

RETURN CALIBRATION-DRIFT-DETECTED

END-IF.


APPLY GAMMA-THETA TO STRUCTURAL-STATE.

PROPAGATE JOINT STRUCTURAL

AND CALIBRATION UNCERTAINTY.


RETURN DIMENSIONAL-PREDICTION

OUTPUT-UNIT

PREDICTED-OUTPUT-SET

CALIBRATION-PROVENANCE.


\PCsection{23. Theorem Ledger}

\PCsubsection{Theorem IV.1 --- Dimensional isolation}

Nontrivial output dimensions must enter through the calibration or readout layer.

\PCsubsection{Theorem IV.2 --- Unique two-anchor affine calibration}

Two exact anchors with distinct structural coordinates determine one affine scalar readout.

\PCsubsection{Theorem IV.3 --- Degenerate-anchor classification}

Equal structural anchors with unequal outputs are inconsistent; equal outputs are underdetermining.

\PCsubsection{Theorem IV.4 --- Polynomial coefficient dimensions}

For dimensionless structural input, every polynomial coefficient carries the output dimension.

\PCsubsection{Theorem IV.5 --- Linear rank identifiability}

A linear multivariable readout requires full column rank for unique parameter determination.

\PCsubsection{Theorem IV.6 --- Weighted least-squares uniqueness}

If

\[
X^{\mathsf T}WX
\]

is positive definite, the weighted quadratic calibration has a unique minimizer.

\PCsubsection{Theorem IV.7 --- Nonlinear rank obstruction}

A null direction of the parameter sensitivity matrix is invisible to first order.

\PCsubsection{Theorem IV.8 --- Exact joint uncertainty propagation}

\[
P_{\mathcal P}(\omega) = \left\{ \Gamma_\theta(\mathbf s): \theta\in C_\theta,\; \mathbf s\in C_{\mathbf F}(\omega) \right\}.
\]

\PCsubsection{Theorem IV.9 --- Three-part residual decomposition}

\[
y_{\mathrm{obs}} - \Gamma_{\widehat\theta}(\mathbf s) = e_{\mathcal M} + d_{\mathrm{str}} + c_\theta.
\]

\PCsubsection{Theorem IV.10 --- Calibration-domain restriction}

Calibration validity does not automatically extend beyond its declared structural or external domain.

\PCsubsection{Theorem IV.11 --- Primitive-pole instability}

Primitive-based calibration may become arbitrarily sensitive near excluded boundaries.

\PCsubsection{Theorem IV.12 --- Flatwave nonresolution}

Flatwave alone cannot resolve phase location within one open quadrant.

\PCsubsection{Theorem IV.13 --- Unit covariance}

Changing units transforms the numerical readout and parameters while preserving the underlying projected quantity.

\PCsubsection{Theorem IV.14 --- Drift separation}

Calibration drift may invalidate projected outputs without affecting Book 0 identities.

\PCsubsection{Theorem IV.15 --- Calibration-transfer burden}

Parameter validity in one domain does not imply validity in another.

\PCsubsection{Theorem IV.16 --- Fit–validation separation}

Agreement with data used to select parameters is not an independent validation result.


\PCsection{24. Negative Theorems}

\PCsubsection{Negative Theorem 24.1 --- No dimensional scale from structural form alone}

A dimensionless identity cannot select a physical unit or numerical scale without external input.

\PCsubsection{Negative Theorem 24.2 --- No unique calibration from degenerate anchors}

Identical structural anchors do not uniquely determine a nontrivial readout.

\PCsubsection{Negative Theorem 24.3 --- No identifiability from optimizer output alone}

A numerical routine returning one estimate does not prove that alternative parameter values are distinguishable.

\PCsubsection{Negative Theorem 24.4 --- No physical necessity from best fit alone}

The parameter minimizing a declared loss is not thereby a fundamental constant.

\PCsubsection{Negative Theorem 24.5 --- No independent confirmation from reused data}

Calibration data cannot simultaneously function as unqualified independent validation data.

\PCsubsection{Negative Theorem 24.6 --- No extrapolation by analytic continuation alone}

A readout formula remaining finite outside its calibration domain does not establish external validity there.

\PCsubsection{Negative Theorem 24.7 --- No fixed calibration under undeclared drift}

A parameter known to vary with external conditions may not be treated as constant without restricting or modeling those conditions.

\PCsubsection{Negative Theorem 24.8 --- No unit-free comparison of dimensional parameters}

Numerical parameter values in incompatible unit systems may not be compared without conversion.

\PCsubsection{Negative Theorem 24.9 --- No residual attribution to structure alone}

A residual may contain measurement error, structural discrepancy, and calibration discrepancy.

\PCsubsection{Negative Theorem 24.10 --- No continuous phase information from Flatwave alone}

A two-valued quadrant sign cannot uniquely encode a continuous within-quadrant coordinate.


\PCsection{25. Terminal Theorem}

\PCsubsection{Theorem 25.1 --- Dimensional Readout Separation Theorem}

Let

\[
\mathbf s(\omega) = E_{\mathbf F}(\Phi(\omega))
\]

be a dimensionless structural state.

Every lawful dimensional prediction has the form

\[
\boxed{ \widehat y_{\mathcal P}(\omega) = \Gamma_{\widehat\theta} \left( \mathbf s(\omega) \right), }
\]

where the calibrated parameter result

\[
\widehat\Theta = K(\mathcal C)
\]

is obtained from explicitly declared external information.

The Book 0 structure determines neither:

\begin{itemize}
\item the output dimension;
\item the unit;
\item the scale;
\item the offset;
\item the calibration anchors;
\item the parameter values;
\item nor the calibration domain.
\end{itemize}
These belong to the readout and calibration contract.

\begin{proof}
The structural state is dimensionless by Book 0 and S1.I.

The output dimension and numerical scale enter only when the readout map and calibration parameters are applied.

The parameter result is selected from external calibration information by \(K\).

None of those external choices is contained in the structural definitions.

Therefore, structural form and dimensional calibration are logically distinct.


\end{proof}

\PCsection{26. Closing Declaration}

S1.I established the projection contract.

S1.II established admission and falsification.

S1.III established measurement and uncertainty.

S1.IV establishes calibration and dimensional readout.

The calibrated parameter result is

\[
\boxed{ \widehat\Theta = K(\mathcal C). }
\]

The projected output is

\[
\boxed{ \widehat y_{\mathcal P} = \Gamma_{\widehat\theta} \circ E_{\mathbf F} \circ \Phi. }
\]

For an affine scalar readout,

\[
\boxed{ y=a+bs, }
\]

two distinct exact anchors determine

\[
\boxed{ b = \frac{y_2-y_1}{s_2-s_1} }
\]

and

\[
\boxed{ a = \frac{s_2y_1-s_1y_2}{s_2-s_1}. }
\]

For a linear multivariable readout, unique parameter determination requires

\[
\boxed{ \operatorname{rank}(X)=m+1. }
\]

Calibration uncertainty combines with structural uncertainty through

\[
\boxed{ P_{\mathcal P}(\omega) = \left\{ \Gamma_\theta(\mathbf s): \theta\in C_\theta,\; \mathbf s\in C_{\mathbf F}(\omega) \right\}. }
\]

The complete residual may contain

\[
\boxed{ \text{measurement error} + \text{structural discrepancy} + \text{calibration discrepancy}. }
\]

The resulting calibration discipline is:

\[
\boxed{ \begin{gathered} \text{No unit without declared provenance;}\\ \text{no scale without an external anchor;}\\ \text{no fitted parameter without an objective;}\\ \text{no unique parameter without identifiability;}\\ \text{no extrapolation without a domain declaration;}\\ \text{no independent validation from calibration data alone;}\\ \text{no fixed readout after material drift.} \end{gathered} }
\]

\[
\boxed{ \text{Book 0 supplies form; calibration supplies scale.} }
\]

\[
\boxed{ \text{END OF S1.IV} }
\]

% ===== END S1_IV =====
% ===== BEGIN S1_V =====
% Source: S1.v.docx
\PCchapter{S1.V}{Projection Equivalence and Model Comparison}{Book 1 - Projection and Measurement}

\PCsubsection{Abstract}

S1.I defined projection contracts. S1.II established admissibility and failure witnesses. S1.III formalized measurement, error, and uncertainty. S1.IV established calibration and dimensional readout.

This essay determines when two projections are genuinely different, when they are merely different descriptions of the same prediction, and what is required to prefer one model over another.

Several levels of equivalence are distinguished:

\begin{enumerate}
\item representational equivalence, in which two contracts differ only by invertible reparameterization;
\item structural equivalence, in which their internal structural states correspond through an invertible map;
\item predictive equivalence, in which they produce the same output sets or distributions;
\item observational equivalence, in which a declared measurement protocol cannot distinguish their outputs;
\item finite-record indistinguishability, in which they agree only on the observations presently available.
\end{enumerate}
These relations are not interchangeable.

An invertible change of phase coordinate or structural variable does not create a new empirical model when the readout is transformed correspondingly. Conversely, two structurally different projections may remain observationally equivalent because calibration, readout, uncertainty, or instrument resolution erases their differences.

Within a fixed open quadrant, the transfer coordinate

\[
\lambda
\]

is strictly increasing and therefore identifies phase uniquely. Globally, however,

\[
\lambda\left(x+\frac{\pi}{2}\right)=\lambda(x),
\]

so \(\lambda\) alone cannot identify the quadrant. Flatwave supplies quadrant sign but cannot identify position within the quadrant. The pair

\[
(\operatorname{Flatwave},\lambda)
\]

distinguishes positive from negative quadrants and within-quadrant position, but still does not distinguish phases separated by \(\pi\). Additional orientation or quadrant information is required for global phase identification.

A comparison contract is introduced:

\[
\mathcal K = \left( \{\mathcal P_j\}_{j=1}^{r}, \Omega_{\mathcal K}, \mathcal M, \mathcal D_{\mathrm{fit}}, \mathcal D_{\mathrm{val}}, L, C, \Delta \right),
\]

where \(L\) is a declared predictive loss or score, \(C\) is a declared complexity measure, and \(\Delta\) is a declared decision rule.

For nested model classes, the more flexible class cannot have worse minimum fitting loss on the same calibration data. This does not imply better validation performance. A perfect fit may result from excessive readout flexibility, parameter nonidentifiability, or data reuse.

No universal comparison metric is imposed. Deterministic error, interval compatibility, likelihood, proper scoring rules, residual diagnostics, and other methods may be used only when their assumptions are declared.

The final result is:

\[
\boxed{ \text{equivalent predictions cannot be distinguished by prediction-only evidence.} }
\]

A model may be preferred only relative to a common domain, observation protocol, uncertainty model, validation record, scoring rule, and complexity criterion. Even then, comparative preference does not prove unique ontology.


\PCsection{IDENTIFICATION DIVISION}

\PCsubsection{PROGRAM-ID}

S1-V-PROJECTION-EQUIVALENCE-AND-MODEL-COMPARISON

\PCsubsection{PURPOSE}

To formalize:

\begin{itemize}
\item equivalence between projection contracts;
\item reparameterization;
\item identifiability;
\item observational indistinguishability;
\item model nesting;
\item predictive scoring;
\item complexity;
\item validation;
\item comparative preference;
\item and the limits of model selection.
\end{itemize}
\PCsubsection{PRINCIPAL QUESTIONS}

\begin{enumerate}
\item When do two apparently different projections make the same claim?
\item When can observations distinguish them?
\item What evidence justifies preferring one over another?
\item What does comparative success fail to prove?
\end{enumerate}
\PCsubsection{NON-PURPOSE}

This essay does not:

\begin{itemize}
\item select one physical projection;
\item prescribe one universal statistical method;
\item claim that the simplest model is always true;
\item infer ontology from predictive success;
\item equate parameter count with complete model complexity;
\item permit comparison across incompatible measurands or units;
\item convert finite-data agreement into universal equivalence;
\item alter the structural theorems of Book 0.
\end{itemize}

\PCsection{ENVIRONMENT DIVISION}

\PCsubsection{1. Inherited Projection Contract}

For \(j=1,\ldots,r\), let

\[
\mathcal P_j = \left( \Omega_j, \Omega_{\mathcal P_j}, \Phi_j, \mathbf F_j, \Gamma_{j,\theta_j}, U_j, W_j \right)
\]

be lawful projection contracts.

Their predictions may be represented as:

\begin{itemize}
\item point outputs
\end{itemize}
\[
\widehat y_j(\omega);
\]

\begin{itemize}
\item admissible output sets
\end{itemize}
\[
P_j(\omega)\subseteq\mathcal Y;
\]

\begin{itemize}
\item or probability measures
\end{itemize}
\[
\nu_j(\,\cdot\mid\omega).
\]

\PCsubsection{Definition 1.1 --- Common external domain}

The common external domain is

\[
\boxed{ \Omega_{\cap} = \bigcap_{j=1}^{r} \Omega_{\mathcal P_j}. }
\]

A comparative domain must satisfy

\[
\Omega_{\mathcal K} \subseteq \Omega_{\cap}.
\]

\PCsubsection{Definition 1.2 --- Common measurand requirement}

Two projections are directly comparable only when their outputs refer to the same declared measurand or are connected by a lawful conversion.

\PCsubsection{Definition 1.3 --- Common output representation}

Before comparison, all predictions must be expressed in a common:

\begin{itemize}
\item output space;
\item unit system;
\item reference frame;
\item coordinate convention;
\item and uncertainty interpretation.
\end{itemize}
A numerical comparison without this normalization is not lawful.


\PCsection{METHODOLOGICAL AXIOM DIVISION}

\PCsubsection{Axiom E.1 --- Comparison Declaration}

Every claim that one projection is equivalent, superior, simpler, better supported, or falsified relative to another must be mediated by an explicit comparison contract.

\PCsubsection{Axiom E.2 --- Equivalence-Level Declaration}

Every equivalence claim must identify its level.

The statement “the models are equivalent” is incomplete unless it specifies whether the equivalence is:

\begin{itemize}
\item representational;
\item structural;
\item predictive;
\item observational;
\item or restricted to a finite record.
\end{itemize}
\PCsubsection{Axiom E.3 --- Common-Evidence Requirement}

Comparative preference must be based on evidence and scoring procedures applied under compatible conditions.

A model may not be credited for one domain while its competitor is judged on a different, more difficult domain without explicit adjustment.


\PCsection{DATA DIVISION}

\PCsection{2. Projection Representations}

\PCsubsection{Definition 2.1 --- Complete prediction map}

For deterministic comparison, define

\[
\boxed{ H_j = \Gamma_{j,\theta_j} \circ E_{\mathbf F_j} \circ \Phi_j. }
\]

Thus,

\[
H_j: \Omega_{\mathcal P_j} \to \mathcal Y.
\]

\PCsubsection{Definition 2.2 --- Structural-state map}

Define

\[
\boxed{ S_j = E_{\mathbf F_j} \circ \Phi_j. }
\]

Thus,

\[
S_j: \Omega_{\mathcal P_j} \to \mathcal S_j.
\]

\PCsubsection{Definition 2.3 --- Prediction family}

When parameters remain free, define

\[
\boxed{ \mathfrak H_j = \left\{ H_{j,\theta}: \theta\in\Theta_j \right\}. }
\]

The projection family, rather than one fitted parameter value, is the relevant object when comparing model flexibility.


\PCsection{3. Representational Equivalence}

\PCsubsection{Definition 3.1 --- External-coordinate reparameterization}

Let

\[
h: \Omega_{\mathcal P_1} \to \Omega_{\mathcal P_2}
\]

be a bijection preserving the external instances under a change of representation.

\PCsubsection{Definition 3.2 --- Phase reparameterization}

Let

\[
r: \Phi_1(\Omega_{\mathcal K}) \to \Phi_2(\Omega_{\mathcal K})
\]

be a bijection.

\PCsubsection{Definition 3.3 --- Structural reparameterization}

Let

\[
g: S_1(\Omega_{\mathcal K}) \to S_2(\Omega_{\mathcal K})
\]

be a bijection.

\PCsubsection{Definition 3.4 --- Readout compensation}

The readouts compensate for \(g\) when

\[
\boxed{ \Gamma_{2,\theta_2}\circ g = \Gamma_{1,\theta_1} }
\]

on the relevant structural image.

\PCsubsection{Definition 3.5 --- Representational equivalence}

Two projections are representationally equivalent on \(\Omega_{\mathcal K}\) when they differ only by invertible coordinate or parameter transformations that leave the complete prediction unchanged.

\PCsubsection{Theorem 3.6 --- Reparameterization invariance}

Suppose

\[
S_2=g\circ S_1
\]

and

\[
\Gamma_{2,\theta_2}\circ g = \Gamma_{1,\theta_1}.
\]

Then

\[
\boxed{ H_2(\omega)=H_1(\omega) }
\]

for every

\[
\omega\in\Omega_{\mathcal K}.
\]

\begin{proof}
For each \(\omega\),

\[
H_2(\omega) = \Gamma_{2,\theta_2}(S_2(\omega)).
\]

Using

\[
S_2=g\circ S_1,
\]

we obtain

\[
H_2(\omega) = \Gamma_{2,\theta_2}(g(S_1(\omega))).
\]

By readout compensation,

\[
\Gamma_{2,\theta_2}\circ g = \Gamma_{1,\theta_1}.
\]

Therefore,

\[
H_2(\omega) = \Gamma_{1,\theta_1}(S_1(\omega)) = H_1(\omega).
\]

\end{proof}

\PCsubsection{Corollary 3.7 --- Coordinate change does not create new predictive content}

An invertible change of internal coordinates, together with the corresponding transformed readout, does not produce a new empirical prediction.

\PCsubsection{Corollary 3.8 --- Parameter names have no independent empirical force}

Renaming, rescaling, or bijectively transforming parameters does not create a distinct empirical model when the prediction family remains unchanged.


\PCsection{4. Structural Equivalence}

\PCsubsection{Definition 4.1 --- Structural equivalence}

Two projections are structurally equivalent on \(\Omega_{\mathcal K}\) when there exists a bijection

\[
g: S_1(\Omega_{\mathcal K}) \to S_2(\Omega_{\mathcal K})
\]

such that

\[
\boxed{ S_2=g\circ S_1. }
\]

Structural equivalence concerns the internal state descriptions before readout.

\PCsubsection{Definition 4.2 --- Exact structural identity}

Two projections are exactly structurally identical when

\[
S_1(\omega)=S_2(\omega)
\]

for every \(\omega\) in the comparison domain.

\PCsubsection{Clarification 4.3 --- Structural equivalence need not imply identical notation}

One model may describe the same state with:

\begin{itemize}
\item a primitive operator;
\item its reciprocal;
\item a half-angle coordinate;
\item a rational generator;
\item or another invertible state coordinate.
\end{itemize}
\PCsubsection{Clarification 4.4 --- Structural equivalence alone does not imply predictive equivalence}

If the readouts do not compensate for the structural transformation, equivalent structural states may be assigned different outputs.


\PCsection{5. Predictive Equivalence}

\PCsubsection{Definition 5.1 --- Point-predictive equivalence}

Two deterministic projections are point-predictively equivalent on \(\Omega_{\mathcal K}\) when

\[
\boxed{ H_1(\omega)=H_2(\omega) }
\]

for every

\[
\omega\in\Omega_{\mathcal K}.
\]

\PCsubsection{Definition 5.2 --- Set-predictive equivalence}

Two uncertain projections are set-predictively equivalent when

\[
\boxed{ P_1(\omega)=P_2(\omega) }
\]

for every \(\omega\) in the comparison domain.

\PCsubsection{Definition 5.3 --- Distributional predictive equivalence}

Two probabilistic projections are distributionally equivalent when

\[
\boxed{ \nu_1(\,\cdot\mid\omega) = \nu_2(\,\cdot\mid\omega) }
\]

for every \(\omega\in\Omega_{\mathcal K}\).

\PCsubsection{Theorem 5.4 --- Predictive equivalence is an equivalence relation}

On a fixed comparison domain and output representation, predictive equivalence is:

\begin{enumerate}
\item reflexive;
\item symmetric;
\item transitive.
\end{enumerate}
\begin{proof}
For point predictions, equality of functions has all three properties.

The same is true for equality of output sets and equality of probability measures.

\end{proof}

\PCsubsection{Corollary 5.5 --- Predictive equivalence classes}

Projection contracts may be partitioned into equivalence classes whose members produce identical predictions on the declared comparison domain.

\PCsubsection{Negative Theorem 5.6 --- Distinct derivations need not imply distinct predictions}

Two projections may use different structural profiles, calibrations, or narratives and still belong to the same predictive equivalence class.


\PCsection{6. Observational Equivalence}

The measurement protocol may erase differences present in the projected outputs.

\PCsubsection{Definition 6.1 --- Measurement channel}

Let

\[
\mathcal C_{\mathcal M}
\]

be the declared measurement channel mapping an output state \(y\) to a distribution or admissible set of possible observation records.

For a probabilistic protocol,

\[
\mathcal C_{\mathcal M}(\,\cdot\mid y)
\]

is the observation distribution conditional on output \(y\).

\PCsubsection{Definition 6.2 --- Induced observation distribution}

For a deterministic prediction \(H_j(\omega)\), define

\[
\boxed{ \zeta_j(\,\cdot\mid\omega) = \mathcal C_{\mathcal M} \left( \,\cdot\mid H_j(\omega) \right). }
\]

For an uncertain prediction, integrate or propagate through the measurement channel according to the declared uncertainty formalism.

\PCsubsection{Definition 6.3 --- Observational equivalence}

Two projections are observationally equivalent under \(\mathcal M\) when

\[
\boxed{ \zeta_1(\,\cdot\mid\omega) = \zeta_2(\,\cdot\mid\omega) }
\]

for every

\[
\omega\in\Omega_{\mathcal K}.
\]

\PCsubsection{Theorem 6.4 --- Predictive equivalence implies observational equivalence}

If two projections are predictively equivalent and are passed through the same measurement channel, then they are observationally equivalent.

\begin{proof}
Equal predictions provide equal inputs to the same measurement channel.

A function or conditional distribution assigns the same observation law to equal inputs.

\end{proof}

\PCsubsection{Negative Theorem 6.5 --- Observational equivalence does not imply predictive equivalence}

Two distinct outputs may induce the same observation distribution under a noninjective or resolution-limited measurement channel.

\begin{proof}
Suppose there exist

\[
y_1\neq y_2
\]

such that

\[
\mathcal C_{\mathcal M}(\,\cdot\mid y_1) = \mathcal C_{\mathcal M}(\,\cdot\mid y_2).
\]

A projection predicting \(y_1\) and another predicting \(y_2\) are observationally equivalent under that channel despite differing predictions.

\end{proof}

\PCsubsection{Example 6.6 --- Resolution-limited equivalence}

If an instrument rounds every value to the nearest integer, predictions

\[
2.1 \qquad\text{and}\qquad 2.4
\]

may produce the same reported result.

The projections are observationally indistinguishable under that instrument even though their exact predictions differ.


\PCsection{7. Finite-Record Indistinguishability}

\PCsubsection{Definition 7.1 --- Evaluation record}

Let

\[
\mathcal R = \left\{ (\omega_i,\mathcal O_i) \right\}_{i=1}^{n}.
\]


be a finite valid observation record.

\PCsubsection{Definition 7.2 --- Finite-record indistinguishability}

Two projections are indistinguishable on \(\mathcal R\) when the declared comparison procedure returns equal or unresolved comparative status for every record in \(\mathcal R\).

For exact deterministic observations, a simple form is

\[
H_1(\omega_i)=H_2(\omega_i)
\]

for every \(i\).

\PCsubsection{Theorem 7.3 --- Finite agreement does not imply global predictive equivalence}

Two functions may agree at every point in a finite comparison record and differ elsewhere.

\begin{proof}
Let the sampled external coordinates be

\[
t_1,\ldots,t_n.
\]

Define

\[
H_1(t)=0
\]

and

\[
H_2(t) = \prod_{i=1}^{n}(t-t_i).
\]

Then

\[
H_2(t_i)=0=H_1(t_i)
\]

at every sampled point, but \(H_2\) is not identically zero.

Therefore, finite agreement does not imply global equality.

\end{proof}

\PCsubsection{Corollary 7.4 --- Finite-data equivalence is record-relative}

A new observation outside the existing record may distinguish previously indistinguishable projections.


\PCsection{8. Identification of Book 0 Phase}

\PCsubsection{Definition 8.1 --- Structural phase identifiability}

A structural coordinate \(F(x)\) identifies phase on a domain \(S\) when

\[
F|_S
\]

is injective.

\PCsubsection{Theorem 8.2 --- Transfer-coordinate identifiability within a quadrant}

On every open quadrant \(Q_k\),

\[
\lambda: Q_k\to(0,1)
\]

is a bijection.

\begin{proof}
Book 0 proves that \(\lambda\) is continuous and strictly increasing on every open quadrant.

It also proves the endpoint limits

\[
\lambda\to0
\]

at the entering wall and

\[
\lambda\to1
\]

at the exiting wall.

A continuous strictly increasing function with those limits maps the quadrant bijectively onto \((0,1)\).

\end{proof}

\PCsubsection{Corollary 8.3 --- Within-quadrant phase recovery}

If the quadrant \(Q_k\) is known and \(\lambda\) is known exactly, then the phase \(x\in Q_k\) is uniquely determined.

\PCsubsection{Theorem 8.4 --- Global nonidentifiability of \(\lambda\)}

\[
\boxed{ \lambda\left(x+\frac{\pi}{2}\right)=\lambda(x). }
\]

Therefore, \(\lambda\) alone does not identify global phase modulo \(2\pi\).

\begin{proof}
S0.IV proves quarter-turn invariance of \(\lambda\).

Thus four distinct phase positions separated by quarter-turns may possess the same \(\lambda\) value.

\end{proof}

\PCsubsection{Theorem 8.5 --- Flatwave cannot identify within-quadrant position}

For every \(x,y\in Q_k\),

\[
\operatorname{Flatwave}(x) = \operatorname{Flatwave}(y) = \varepsilon_k.
\]

Therefore, Flatwave is noninjective on every open quadrant.

\PCsubsection{Theorem 8.6 --- Combined \((\operatorname{Flatwave},\lambda)\) limitation}

The pair

\[
\left( \operatorname{Flatwave}(x), \lambda(x) \right)
\]

distinguishes:

\begin{itemize}
\item positive-sign quadrants from negative-sign quadrants;
\item and position within a quadrant.
\end{itemize}
However,

\[
\boxed{ \left( \operatorname{Flatwave}(x+\pi), \lambda(x+\pi) \right) = \left( \operatorname{Flatwave}(x), \lambda(x) \right). }
\]

Therefore, the pair does not distinguish phases separated by a half-turn.

\begin{proof}
Flatwave is \(\pi\)-periodic, and \(\lambda\) is already \(\pi/2\)-periodic.

Thus both coordinates repeat after \(\pi\).

\end{proof}

\PCsubsection{Corollary 8.7 --- Global phase identification requires additional information}

A complete phase identification modulo \(2\pi\) requires at least one additional orientation, quadrant, or half-turn label beyond \((\operatorname{Flatwave},\lambda)\).


\PCsection{9. Parameter Identifiability}

\PCsubsection{Definition 9.1 --- Parameter prediction map}

For projection family \(\mathfrak H\), define

\[
\boxed{ \Psi: \Theta \to \mathcal Y^{\Omega_{\mathcal K}} }
\]

by

\[
\Psi(\theta) = H_\theta|_{\Omega_{\mathcal K}}.
\]

\PCsubsection{Definition 9.2 --- Global parameter identifiability}

Parameters are globally identifiable when

\[
\boxed{ \Psi(\theta_1)=\Psi(\theta_2) \implies \theta_1=\theta_2. }
\]

\PCsubsection{Definition 9.3 --- Parameter equivalence class}

Define

\[
\theta_1\sim\theta_2
\]

when

\[
\Psi(\theta_1)=\Psi(\theta_2).
\]

The equivalence class is

\[
[\theta] = \left\{ \vartheta\in\Theta: \Psi(\vartheta)=\Psi(\theta) \right\}.
\]


\PCsubsection{Theorem 9.4 --- Nonidentifiable parameters cannot be distinguished by output evidence}

If

\[
\theta_1\sim\theta_2,
\]

then no observation depending only on the projected outputs over \(\Omega_{\mathcal K}\) can distinguish \(\theta_1\) from \(\theta_2\).

\begin{proof}
The parameters produce identical predictions at every comparison instance.

Passing identical predictions through the same observation protocol produces identical observation laws.

Therefore, output evidence on the comparison domain cannot distinguish them.

\end{proof}

\PCsubsection{Corollary 9.5 --- One fitted value need not be uniquely supported}

An optimizer returning \(\widehat\theta\) does not prove that

\[
[\widehat\theta] = \left\{ \widehat\theta\right\}.
\]


\PCsection{10. Model Classes and Nesting}

\PCsubsection{Definition 10.1 --- Model class}

A model class is a set of prediction maps

\[
\boxed{ \mathfrak H = \left\{ H_\theta:\theta\in\Theta \right\} }.
\]



\PCsubsection{Definition 10.2 --- Nested model class}

Model class \(\mathfrak H_0\) is nested in \(\mathfrak H_1\) when

\[
\boxed{ \mathfrak H_0\subseteq\mathfrak H_1. }
\]

\PCsubsection{Definition 10.3 --- Strict nesting}

The nesting is strict when

\[
\mathfrak H_0 \subsetneq \mathfrak H_1.
\]

\PCsubsection{Theorem 10.4 --- Training-loss monotonicity under nesting}

Let \(L_{\mathrm{fit}}(H)\) be the same fitting loss applied to both model classes.

If

\[
\mathfrak H_0\subseteq\mathfrak H_1,
\]

then

\[
\boxed{ \inf_{H\in\mathfrak H_1} L_{\mathrm{fit}}(H) \leq \inf_{H\in\mathfrak H_0} L_{\mathrm{fit}}(H). }
\]

\begin{proof}
The optimization domain for \(\mathfrak H_1\) contains every candidate in \(\mathfrak H_0\), possibly with additional candidates.

The infimum over a larger set cannot exceed the infimum over its subset.

\end{proof}

\PCsubsection{Corollary 10.5 --- Better fitting loss may be purely a flexibility effect}

A more flexible nested class is guaranteed the opportunity to match or improve the best fitting loss, even when its additional structure has no external validity.

\PCsubsection{Negative Theorem 10.6 --- Training improvement does not imply validation improvement}

The inequality of Theorem 10.4 applies to the fitting record.

It imposes no corresponding inequality on independent validation loss.


\PCsection{11. Comparison Contract}

\PCsubsection{Definition 11.1 --- Comparison contract}

A model-comparison contract is

\[
\boxed{ \mathcal K = \left( \{\mathcal P_j\}_{j=1}^{r}, \Omega_{\mathcal K}, \mathcal M, \mathcal D_{\mathrm{fit}}, \mathcal D_{\mathrm{val}}, L, C, \Delta, W_{\mathcal K} \right), }
\]

where:

\begin{itemize}
\item \(\{\mathcal P_j\}\) is the candidate projection set;
\item \(\Omega_{\mathcal K}\) is the common comparison domain;
\item \(\mathcal M\) is the common observation protocol or lawful protocol conversion;
\item \(\mathcal D_{\mathrm{fit}}\) is the calibration or fitting record;
\item \(\mathcal D_{\mathrm{val}}\) is the validation record;
\item \(L\) is the declared loss or score;
\item \(C\) is the declared complexity measure;
\item \(\Delta\) is the decision rule;
\item \(W_{\mathcal K}\) is the comparison failure map.
\end{itemize}
\PCsubsection{Definition 11.2 --- Complete comparison}

A comparison is complete when it declares:

\begin{enumerate}
\item all candidate contracts and versions;
\item a common validity domain;
\item a common measurand;
\item common units and conventions;
\item the fitting and validation records;
\item uncertainty treatment;
\item scoring or loss;
\item complexity treatment;
\item model-search scope;
\item the decision rule;
\item failure behavior.
\end{enumerate}
\PCsubsection{Recommended failure witnesses}

\begin{itemize}
\item COMPARISON-CONTRACT-INCOMPLETE;
\item COMMON-VALIDITY-DOMAIN-EMPTY;
\item MEASURAND-MISMATCH;
\item OUTPUT-UNIT-MISMATCH;
\item OBSERVATION-PROTOCOL-INCOMPATIBLE;
\item UNCERTAINTY-INTERPRETATION-INCOMPATIBLE;
\item COMPARISON-SCORE-UNDECLARED;
\item COMPLEXITY-MEASURE-UNDECLARED;
\item VALIDATION-RECORD-NOT-INDEPENDENT;
\item MODEL-SEARCH-SCOPE-UNDISCLOSED;
\item CANDIDATE-MODELS-OBSERVATIONALLY-EQUIVALENT;
\item PARAMETER-NONIDENTIFIABILITY-UNRESOLVED;
\item COMPARISON-EXTRAPOLATION-MISMATCH;
\item COMPARISON-INCONCLUSIVE.
\end{itemize}

\PCsection{12. Deterministic Comparison}

\PCsubsection{Definition 12.1 --- Pointwise residual}

For a scalar observation \(y_i\),

\[
\boxed{ r_{j,i} = y_i-H_j(\omega_i). }
\]

\PCsubsection{Definition 12.2 --- Absolute-error loss}

\[
\boxed{ L_{\mathrm{AE}}(H_j) = \sum_{i=1}^{n} w_i |r_{j,i}|. }
\]

\PCsubsection{Definition 12.3 --- Squared-error loss}

\[
\boxed{ L_{\mathrm{SE}}(H_j) = \sum_{i=1}^{n} w_i r_{j,i}^{2}. }
\]

The weights and their interpretation must be declared.

\PCsubsection{Definition 12.4 --- Maximum-error loss}

\[
\boxed{ L_{\infty}(H_j) = \max_i|r_{j,i}|. }
\]

\PCsubsection{Clarification 12.5 --- Loss choice changes preference}

A projection minimizing squared error need not minimize absolute error or maximum error.

Therefore, “best fit” is incomplete without the loss definition.

\PCsubsection{Theorem 12.6 --- Equivalent predictions receive equal prediction-only loss}

If

\[
H_1(\omega_i)=H_2(\omega_i)
\]

for every record \(i\), then every loss depending only on those predictions and observations assigns equal values:

\[
\boxed{ L(H_1)=L(H_2). }
\]

\begin{proof}
The complete list of loss inputs is identical for the two projections.

Therefore, the loss outputs are identical.


\end{proof}

\PCsection{13. Set-Valued Comparison}

\PCsubsection{Definition 13.1 --- Compatibility indicator}

For prediction set \(P_j(\omega_i)\) and observed set \(O_i\), define

\[
\boxed{ c_{j,i} = \begin{cases} 1, & P_j(\omega_i)\cap O_i\neq\varnothing, \\ 0, & P_j(\omega_i)\cap O_i=\varnothing. \end{cases} }
\]

\PCsubsection{Definition 13.2 --- Contradiction count}

\[
\boxed{ N_{\mathrm{contra}}(j) = \sum_{i=1}^{n}(1-c_{j,i}). }
\]

\PCsubsection{Clarification 13.3 --- Compatibility count ignores sharpness}

A projection making extremely wide predictions may achieve compatibility by admitting nearly every result.

Compatibility alone therefore does not reward informative narrowness.

\PCsubsection{Definition 13.4 --- Set width or volume}

When meaningful, define a width or volume measure

\[
V(P_j(\omega)).
\]

A comparison may jointly evaluate:

\begin{itemize}
\item empirical coverage;
\item and predictive sharpness.
\end{itemize}
\PCsubsection{Negative Theorem 13.5 --- Universal prediction is not strong prediction}

A projection returning the entire output space

\[
P_j(\omega)=\mathcal Y
\]

is compatible with every observation but excludes nothing.

It provides no discriminatory prediction.


\PCsection{14. Probabilistic Comparison}

\PCsubsection{Definition 14.1 --- Predictive probability distribution}

Let

\[
\nu_j(\,\cdot\mid\omega)
\]

be the declared predictive distribution for projection \(j\).

\PCsubsection{Definition 14.2 --- Scoring rule}

A scoring rule assigns a numerical score

\[
\boxed{ S(\nu_j,y_{\mathrm{obs}}) }
\]

to a predictive distribution and realized observation.

The direction of preference—larger or smaller—must be declared.

\PCsubsection{Definition 14.3 --- Aggregate predictive score}

\[
\boxed{ L_j = \sum_{i=1}^{n} S \left( \nu_j(\,\cdot\mid\omega_i), y_i \right). }
\]

\PCsubsection{Clarification 14.4 --- Probability models require interpretation}

A probability distribution may represent:

\begin{itemize}
\item aleatory variation;
\item measurement noise;
\item uncertainty concerning parameters;
\item uncertainty concerning states;
\item or a mixture.
\end{itemize}
The interpretation must be stated.

\PCsubsection{Theorem 14.5 --- Distributional equivalence implies equal probabilistic score}

If

\[
\nu_1(\,\cdot\mid\omega_i) = \nu_2(\,\cdot\mid\omega_i)
\]

for every record \(i\), then any prediction-only scoring rule assigns equal aggregate score.

\PCsubsection{Negative Theorem 14.6 --- Likelihood alone does not determine ontology}

A greater likelihood under one projection shows stronger support relative to the declared probabilistic models and data.

It does not prove that the internal structural description is uniquely real.


\PCsection{15. Residual Diagnostics}

\PCsubsection{Definition 15.1 --- Residual sequence}

For projection \(j\),

\[
\mathbf r_j = (r_{j,1},\ldots,r_{j,n}).
\]

\PCsubsection{Definition 15.2 --- Residual structure}

Residual structure includes systematic dependence on:

\begin{itemize}
\item phase;
\item external covariates;
\item time;
\item magnitude;
\item location;
\item instrument;
\item calibration batch;
\item or previous residuals.
\end{itemize}
\PCsubsection{Theorem 15.3 --- Zero mean residual does not imply adequate modeling}

A residual sequence may have mean zero while retaining systematic structure.

\begin{proof}
Consider residuals alternating between \(+1\) and \(-1\).

Their mean may be zero, but they display perfect serial structure.

Therefore, zero average alone does not establish adequacy.

\end{proof}

\PCsubsection{Corollary 15.4 --- Residual diagnostics must exceed one summary statistic}

A complete comparison should examine patterns relevant to the declared model assumptions.

\PCsubsection{Clarification 15.5 --- Residual structure may arise from several layers}

Patterns may indicate:

\begin{itemize}
\item phase-inference failure;
\item readout misspecification;
\item calibration drift;
\item omitted external variables;
\item measurement artifacts;
\item dependence;
\item or structural inadequacy.
\end{itemize}

\PCsection{16. Complexity}

\PCsubsection{Definition 16.1 --- Nominal parameter count}

\[
\boxed{ k_{\mathrm{nom}} = \dim(\Theta) }
\]

when the parameter space has a declared finite dimension.

\PCsubsection{Definition 16.2 --- Identifiable parameter count}

The identifiable parameter count is the dimension of the predictionally distinguishable parameter space after quotienting by parameter equivalence.

Symbolically,

\[
\boxed{ k_{\mathrm{id}} = \dim(\Theta/\!\sim) }
\]

where defined.

\PCsubsection{Definition 16.3 --- Functional complexity}

Functional complexity concerns the richness of the prediction family

\[
\mathfrak H
\]

rather than merely the number of symbols used to parameterize it.

\PCsubsection{Definition 16.4 --- Procedural complexity}

Procedural complexity includes:

\begin{itemize}
\item branch rules;
\item preprocessing;
\item calibration stages;
\item exception handling;
\item selected data transformations;
\item and post-hoc adjustments.
\end{itemize}
\PCsubsection{Definition 16.5 --- Declared complexity measure}

A comparison contract may use a complexity function

\[
\boxed{ C(\mathcal P)\geq0. }
\]

The measure and its rationale must be stated.

\PCsubsection{Theorem 16.6 --- Parameter count is representation-sensitive under redundancy}

The same prediction family may be represented with additional redundant parameters.

Therefore, nominal parameter count alone is not an invariant measure of empirical complexity.

\begin{proof}
Let a model depend on parameter \(a\) through prediction \(H_a\).

Introduce two parameters \(b,c\) with

\[
a=b+c.
\]

The prediction family is unchanged, but the nominal parameter count increases from one to two.

Because many pairs \((b,c)\) produce the same \(a\), the added parameterization is redundant.

\end{proof}

\PCsubsection{Corollary 16.7 --- Identifiability must accompany parameter count}

Reporting only \(k_{\mathrm{nom}}\) may overstate the number of empirically distinguishable degrees of freedom.


\PCsection{17. Parsimony}

\PCsubsection{Definition 17.1 --- Empirical equivalence tolerance}

Let

\[
\epsilon\geq0
\]

be a declared tolerance for considering validation performance practically equivalent.

\PCsubsection{Definition 17.2 --- Conditional parsimony rule}

A comparison may prefer \(\mathcal P_1\) over \(\mathcal P_2\) by parsimony when:

\begin{enumerate}
\item their validation performance differs by no more than \(\epsilon\);
\item both satisfy required adequacy guards;
\item the same complexity measure is applied;
\item and
\end{enumerate}
\[
C(\mathcal P_1)<C(\mathcal P_2).
\]

\PCsubsection{Clarification 17.3 --- Parsimony is a decision rule}

Parsimony is not a theorem that the simpler model is ontologically true.

It is a rule for choosing among models with sufficiently similar empirical performance.

\PCsubsection{Negative Theorem 17.4 --- Simplicity cannot rescue empirical failure}

A simpler projection contradicted by valid observations is not preferable merely because it has fewer parameters.


\PCsection{18. Calibration and Validation Separation}

\PCsubsection{Definition 18.1 --- Fitting record}

\[
\mathcal D_{\mathrm{fit}}
\]

is the information used to select:

\begin{itemize}
\item model form;
\item structural profile;
\item readout;
\item parameters;
\item thresholds;
\item preprocessing;
\item or comparison criteria.
\end{itemize}
\PCsubsection{Definition 18.2 --- Validation record}

\[
\mathcal D_{\mathrm{val}}
\]

is the information reserved for evaluating decisions made using the fitting record.

\PCsubsection{Theorem 18.3 --- Selection dependence}

Performance measured on data used to select among candidate models is statistically and procedurally dependent on the selection process.

\begin{proof}
The selected model is chosen in response to the fitting data.

Its measured fitting performance is therefore not the performance of a model specified independently of that data.

\end{proof}

\PCsubsection{Definition 18.4 --- Model-search scope}

The model-search scope is the complete set of:

\begin{itemize}
\item projection families considered;
\item transformations tried;
\item variables screened;
\item thresholds tested;
\item and stopping rules used.
\end{itemize}
\PCsubsection{Negative Theorem 18.5 --- Undisclosed search exaggerates evidence}

Reporting only the best-performing model without disclosing the broader search omits information required to evaluate selection effects.

\PCsubsection{Required witness}

\[
\boxed{ \texttt{MODEL-SEARCH-SCOPE-UNDISCLOSED}. }
\]


\PCsection{19. Cross-Validation and Repeated Assessment}

\PCsubsection{Definition 19.1 --- Data partition}

A data partition divides the available record into disjoint or partially overlapping fitting and validation subsets.

\PCsubsection{Definition 19.2 --- Cross-validation procedure}

A cross-validation procedure repeatedly fits on one subset and evaluates on another according to a declared partition rule.

\PCsubsection{Clarification 19.3 --- Cross-validation estimates procedural performance}

Cross-validation evaluates the complete fitting procedure, including parameter selection, when the procedure is repeated correctly within each training partition.

\PCsubsection{Negative Theorem 19.4 --- Preprocessing leakage invalidates nominal separation}

If preprocessing, structural-profile selection, or parameter tuning uses the complete dataset before partitioned evaluation, the validation subsets are not fully independent of the selection process.


\PCsection{20. Model Dominance}

\PCsubsection{Definition 20.1 --- Record-wise predictive dominance}

Projection \(\mathcal P_1\) dominates \(\mathcal P_2\) on a validation record under loss \(L\) when

\[
L_i(\mathcal P_1) \leq L_i(\mathcal P_2)
\]

for every record \(i\), with strict inequality for at least one.

\PCsubsection{Definition 20.2 --- Aggregate dominance}

Projection \(\mathcal P_1\) has lower aggregate loss when

\[
L(\mathcal P_1) < L(\mathcal P_2).
\]

\PCsubsection{Clarification 20.3 --- Aggregate dominance may hide subgroup failure}

A model may have lower total loss while performing worse on a scientifically important subset.

Comparison domains and subgroup requirements must be declared.

\PCsubsection{Definition 20.4 --- Constraint dominance}

Projection \(\mathcal P_1\) constraint-dominates \(\mathcal P_2\) when:

\begin{enumerate}
\item every output permitted by \(\mathcal P_1\) is permitted by \(\mathcal P_2\);
\item \(\mathcal P_1\) excludes at least one output that \(\mathcal P_2\) permits;
\item both remain compatible with the validation evidence.
\end{enumerate}
Symbolically,

\[
P_1(\omega)\subseteq P_2(\omega)
\]

throughout the domain, with strict inclusion somewhere.

\PCsubsection{Clarification 20.5 --- Narrower is better only when calibrated}

An unjustifiably narrow prediction set may appear more informative while possessing incorrect coverage.

Constraint dominance requires valid uncertainty construction.


\PCsection{21. Equivalence and Comparison Theorems}

\PCsubsection{Theorem 21.1 --- Prediction-only indistinguishability}

If two projections are predictively equivalent on \(\Omega_{\mathcal K}\), then no comparison method using only their predictions and observations from that domain can prefer one on predictive grounds.

\begin{proof}
Predictive equivalence gives identical predictions for every comparison instance.

Every prediction-only score, residual, compatibility result, and induced observation distribution is therefore identical.

No predictive difference remains for the comparison method to use.

\end{proof}

\PCsubsection{Corollary 21.2 --- Preference between predictively equivalent projections requires nonpredictive criteria}

Possible criteria include:

\begin{itemize}
\item computational cost;
\item interpretability;
\item portability;
\item calibration stability;
\item procedural simplicity;
\item or compatibility with separately established theory.
\end{itemize}
These criteria must be declared.

\PCsubsection{Theorem 21.3 --- Observational equivalence blocks protocol-specific discrimination}

If two projections are observationally equivalent under protocol \(\mathcal M\), no amount of repetition using only that unchanged protocol can distinguish them in expectation.

\begin{proof}
Every repeated observation has the same declared probability law under both projections.

Therefore, the joint distribution of any finite or countable sequence of repeated observations is the same under both, assuming the same repetition model.

\end{proof}

\PCsubsection{Corollary 21.4 --- Distinguishing observationally equivalent models requires a new intervention or protocol}

The new protocol must access a difference that the former measurement channel erased.

\PCsubsection{Theorem 21.5 --- Better score is comparison-relative}

The statement

\[
L(\mathcal P_1)<L(\mathcal P_2)
\]

establishes preference only relative to:

\begin{itemize}
\item the chosen data;
\item the chosen loss;
\item the chosen weighting;
\item the chosen uncertainty representation;
\item the chosen domain;
\item and the chosen candidate set.
\end{itemize}
\begin{proof}
Changing any of these elements may change the numerical losses or their interpretation.

Therefore, the preference is conditional on the comparison contract.


\end{proof}

\PCsection{22. Comparison Decision Rules}

\PCsubsection{Definition 22.1 --- Decision set}

Let

\[
\mathfrak D_{\mathcal K} = \left\{ \operatorname{equivalent}, \operatorname{prefer}\mathcal P_j, \operatorname{reject}\mathcal P_j, \operatorname{retain\_multiple}, \operatorname{inconclusive} \right\}.
\]


\PCsubsection{Definition 22.2 --- Comparison decision rule}

\[
\boxed{ \Delta: \mathcal Z_{\mathcal K} \to \mathfrak D_{\mathcal K}, }
\]

where \(\mathcal Z_{\mathcal K}\) contains the declared scores, diagnostics, uncertainty results, complexity values, and guard statuses.

\PCsubsection{Definition 22.3 --- Predeclared decision rule}

A decision rule is predeclared when it is fixed before examining the validation outcomes to which it will be applied.

\PCsubsection{Clarification 22.4 --- Exploratory comparison}

A post-hoc rule may be used for exploration, but its result must be labeled exploratory rather than confirmatory.

\PCsubsection{Theorem 22.5 --- Inconclusive comparison}

A comparison is inconclusive when:

\begin{itemize}
\item candidate predictions overlap within uncertainty;
\item observational equivalence has not been broken;
\item validation evidence is insufficient;
\item scoring rules disagree without a declared priority;
\item or required guards fail.
\end{itemize}
\PCsubsection{Required witness}

\[
\boxed{ \texttt{COMPARISON-INCONCLUSIVE}. }
\]


\PCsection{23. Comparison Failure Locality}

\PCsubsection{Definition 23.1 --- Equivalence failure}

An equivalence claim fails when the required equality, bijection, or induced observation identity does not hold on the declared domain.

\PCsubsection{Definition 23.2 --- Comparability failure}

Comparability fails when the projections do not share:

\begin{itemize}
\item a common measurand;
\item compatible units;
\item a common domain;
\item or a lawful observation protocol.
\end{itemize}
\PCsubsection{Definition 23.3 --- Selection failure}

Selection fails when no declared decision rule resolves the available evidence.

\PCsubsection{Theorem 23.4 --- Comparison failure does not invalidate individual contracts}

Two individually lawful projections may be impossible to compare directly because their measurands, domains, or protocols differ.

\begin{proof}
Lawfulness is assessed relative to each projection's own contract.

Direct comparison imposes additional common-domain and common-output requirements.

Failure of those additional requirements does not imply either projection is internally unlawful.


\end{proof}

\PCsection{PROCEDURE DIVISION}

\PCsection{24. Canonical Equivalence Procedure}

CHECK-PROJECTION-EQUIVALENCE.


DECLARE COMPARISON-DOMAIN.

VERIFY COMMON MEASURAND.

VERIFY COMMON UNITS.

VERIFY COMMON OBSERVATION PROTOCOL.


IF COMMON-VALIDITY-DOMAIN-EMPTY

RETURN COMPARISON-WITNESS-SET

END-IF.


CHECK FOR INVERTIBLE PHASE

OR STRUCTURAL REPARAMETERIZATION.


IF REPARAMETERIZATION EXISTS

AND READOUTS COMPENSATE

RECORD REPRESENTATIONAL-EQUIVALENCE

END-IF.


COMPARE STRUCTURAL-STATE MAPS.

COMPARE PREDICTED OUTPUTS.

COMPARE PREDICTED OUTPUT SETS

OR DISTRIBUTIONS.

PROPAGATE THROUGH MEASUREMENT CHANNEL.


CLASSIFY AS:

STRUCTURALLY-EQUIVALENT,

PREDICTIVELY-EQUIVALENT,

OBSERVATIONALLY-EQUIVALENT,

FINITE-RECORD-INDISTINGUISHABLE,

OR DISTINCT.


RETURN EQUIVALENCE-CLASSIFICATION

DOMAIN

ASSUMPTIONS

WITNESS-SET.

\PCsection{25. Canonical Model-Comparison Procedure}

COMPARE-PROJECTION-MODELS.


PERFORM CHECK-COMPARISON-CONTRACT.


IF COMPARISON-CONTRACT-INCOMPLETE

RETURN COMPARISON-WITNESS-SET

END-IF.


PERFORM NORMALIZE-MEASURAND-UNITS

AND OBSERVATION-PROTOCOL.


IF MEASURAND-MISMATCH

OR OUTPUT-UNIT-MISMATCH

OR OBSERVATION-PROTOCOL-INCOMPATIBLE

RETURN COMPARISON-WITNESS-SET

END-IF.


PERFORM CHECK-PROJECTION-EQUIVALENCE.


IF CANDIDATE-MODELS-OBSERVATIONALLY-EQUIVALENT

RETURN NO-PROTOCOL-SPECIFIC-DISCRIMINATION

END-IF.


DECLARE FITTING-RECORD.

DECLARE VALIDATION-RECORD.

DECLARE MODEL-SEARCH-SCOPE.

DECLARE LOSS-OR-SCORE.

DECLARE COMPLEXITY-MEASURE.

DECLARE DECISION-RULE.


FIT EACH CANDIDATE

USING FITTING-RECORD ONLY.


EVALUATE EACH FITTED CANDIDATE

ON VALIDATION-RECORD.


COMPUTE:

PREDICTIVE SCORES,

UNCERTAINTY COMPATIBILITY,

RESIDUAL DIAGNOSTICS,

CALIBRATION STABILITY,

COMPLEXITY VALUES,

AND FAILURE WITNESSES.


APPLY PREDECLARED DECISION-RULE.


RETURN:

EQUIVALENCE STATUS,

COMPARATIVE SCORES,

VALIDATION RESULTS,

COMPLEXITY RESULTS,

DECISION,

AND COMPLETE AUDIT RECORD.


\PCsection{26. Theorem Ledger}

\PCsubsection{Theorem V.1 --- Reparameterization invariance}

Invertible structural reparameterization with compensating readout preserves prediction.

\PCsubsection{Theorem V.2 --- Predictive equivalence relation}

Predictive equivalence is reflexive, symmetric, and transitive.

\PCsubsection{Theorem V.3 --- Predictive-to-observational implication}

Predictive equivalence under a common measurement channel implies observational equivalence.

\PCsubsection{Theorem V.4 --- Failure of the converse}

Observational equivalence may occur despite distinct predictions.

\PCsubsection{Theorem V.5 --- Finite-record limitation}

Agreement on a finite record does not prove global predictive equivalence.

\PCsubsection{Theorem V.6 --- Quadrant-local \(\lambda\) identifiability}

\[
\lambda:Q_k\to(0,1)
\]

is bijective.

\PCsubsection{Theorem V.7 --- Global \(\lambda\) nonidentifiability}

\[
\lambda\left(x+\frac{\pi}{2}\right)=\lambda(x).
\]

\PCsubsection{Theorem V.8 --- Flatwave nonidentifiability}

Flatwave cannot identify within-quadrant phase.

\PCsubsection{Theorem V.9 --- Half-turn ambiguity of \((\operatorname{Flatwave},\lambda)\)}

The combined pair repeats after \(\pi\).

\PCsubsection{Theorem V.10 --- Parameter-equivalence limitation}

Predictionally equivalent parameters cannot be distinguished by output evidence on the comparison domain.

\PCsubsection{Theorem V.11 --- Nested fitting monotonicity}

A larger nested model class cannot have worse minimum fitting loss on the same record.

\PCsubsection{Theorem V.12 --- Equivalent-prediction score equality}

Prediction-only losses assign equal scores to equal predictions.

\PCsubsection{Theorem V.13 --- Training improvement limitation}

Improved fitting loss does not imply improved validation loss.

\PCsubsection{Theorem V.14 --- Parameter-count noninvariance under redundancy}

Nominal parameter count may change while the prediction family remains unchanged.

\PCsubsection{Theorem V.15 --- Prediction-only indistinguishability}

Predictively equivalent models cannot be distinguished on predictive evidence from the equivalence domain.

\PCsubsection{Theorem V.16 --- Observational repetition limitation}

Repeated use of the same nondiscriminating measurement protocol does not distinguish observationally equivalent projections.

\PCsubsection{Theorem V.17 --- Comparison-relative preference}

A better score establishes preference only relative to the complete comparison contract.


\PCsection{27. Negative Theorems}

\PCsubsection{Negative Theorem 27.1 --- No universal meaning of “equivalent”}

An equivalence claim without a declared level is incomplete.

\PCsubsection{Negative Theorem 27.2 --- No new empirical model from notation alone}

Changing symbols or invertibly reparameterizing internal states does not create distinct predictions.

\PCsubsection{Negative Theorem 27.3 --- No global equivalence from finite agreement}

Agreement on finitely many records does not prove equality outside those records.

\PCsubsection{Negative Theorem 27.4 --- No structural equality from observational equality alone}

A measurement protocol may erase genuine structural or predictive differences.

\PCsubsection{Negative Theorem 27.5 --- No global phase from \(\lambda\) alone}

Quarter-turn-related phases share the same transfer coordinate.

\PCsubsection{Negative Theorem 27.6 --- No continuous phase from Flatwave alone}

Flatwave is constant within each quadrant.

\PCsubsection{Negative Theorem 27.7 --- No unique parameters without identifiability}

One optimizer output does not prove one empirically distinguishable parameter state.

\PCsubsection{Negative Theorem 27.8 --- No validation conclusion from fitting loss alone}

Nested model flexibility guarantees opportunities for improved training fit.

\PCsubsection{Negative Theorem 27.9 --- No universal best loss function}

Different valid comparison goals may require different declared scoring procedures.

\PCsubsection{Negative Theorem 27.10 --- No strong prediction from universal compatibility}

A prediction set containing every possible output excludes nothing.

\PCsubsection{Negative Theorem 27.11 --- No universal truth from parsimony}

Simplicity is a conditional decision criterion, not proof of ontology.

\PCsubsection{Negative Theorem 27.12 --- No independent validation after hidden data reuse}

Validation evidence is compromised when it influenced model selection without adjustment.

\PCsubsection{Negative Theorem 27.13 --- No fair comparison across unequal domains}

A model tested only on favorable cases may not be declared superior to one evaluated on a broader domain without domain adjustment.

\PCsubsection{Negative Theorem 27.14 --- No ontology from comparative victory}

Outperforming declared competitors does not prove that the winning model is the uniquely real underlying structure.


\PCsection{28. Terminal Theorem}

\PCsubsection{Theorem 28.1 --- Projection Comparison Separation Theorem}

Let \(\mathcal P_1,\ldots,\mathcal P_r\) be projection contracts.

A lawful comparative conclusion requires:

\begin{enumerate}
\item a common comparison domain;
\item a common measurand and output representation;
\item compatible observation protocols;
\item declared uncertainty treatment;
\item declared calibration and validation records;
\item a declared score or loss;
\item a declared complexity measure when complexity affects the decision;
\item a declared model-search scope;
\item and a declared decision rule.
\end{enumerate}
If two projections are predictively equivalent on the comparison domain, predictive evidence from that domain cannot distinguish them.

If they are observationally equivalent under the declared measurement protocol, repetition of that same protocol cannot distinguish them.

If one projection receives a better comparative score, the preference is conditional on the comparison contract and does not establish unique structural or physical ontology.

\begin{proof}
The first requirements follow from the definitions of lawful comparability and the inherited measurement and calibration disciplines.

Predictive equivalence supplies identical prediction inputs to every prediction-only comparison.

Observational equivalence supplies identical observation laws under the declared protocol.

A comparative score depends on the chosen domain, records, uncertainty representation, score, weighting, calibration, and candidate class.

Therefore, equivalence blocks discrimination at its declared level, while comparative preference remains conditional on the complete comparison design.


\end{proof}

\PCsection{29. Closing Declaration}

Book 1 began by separating structural mathematics from external interpretation.

S1.I established projection contracts:

\[
\boxed{ \widehat y_{\mathcal P} = \Gamma_\theta \circ E_{\mathbf F} \circ \Phi. }
\]

S1.II established admission, failure witnesses, and falsification.

S1.III established measurement protocols, error, and uncertainty.

S1.IV established calibration and dimensional readout.

S1.V establishes equivalence and comparison.

The principal equivalence ladder is:

\[
\boxed{ \begin{gathered} \text{representational equivalence}\\ \Downarrow\\ \text{structural correspondence}\\ \Downarrow\\ \text{predictive equivalence}\\ \Downarrow\\ \text{observational equivalence}\\ \Downarrow\\ \text{finite-record indistinguishability}. \end{gathered} }
\]

The downward implications are not all reversible.

Within a known quadrant,

\[
\boxed{ \lambda \text{ identifies phase uniquely}. }
\]

Globally,

\[
\boxed{ \lambda\left(x+\frac{\pi}{2}\right)=\lambda(x). }
\]

Flatwave supplies only quadrant sign:

\[
\boxed{ \operatorname{Flatwave} = \operatorname{sgn}(\sin 2x). }
\]

The pair

\[
(\operatorname{Flatwave},\lambda)
\]

still retains half-turn ambiguity.

For model classes,

\[
\boxed{ \mathfrak H_0\subseteq\mathfrak H_1 \implies \inf_{\mathfrak H_1}L_{\mathrm{fit}} \leq \inf_{\mathfrak H_0}L_{\mathrm{fit}}. }
\]

This is a fitting theorem, not a validation theorem.

The resulting comparison discipline is:

\[
\boxed{ \begin{gathered} \text{No equivalence without a declared level;}\\ \text{no comparison without a common domain;}\\ \text{no preference without a declared score;}\\ \text{no simplicity claim without a complexity measure;}\\ \text{no validation after hidden data reuse;}\\ \text{no discrimination between equivalent predictions;}\\ \text{no ontology from predictive victory alone.} \end{gathered} }
\]

This completes the projection and measurement framework.

\[
\boxed{ \text{END OF S1.V} }
\]

\[
\boxed{ \text{END OF BOOK 1} }
\]

% ===== END S1_V =====
\end{document}
